

\documentclass[11pt]{book}

\usepackage{amsmath,amssymb,amsfonts,amsthm}
\usepackage{mathtools}
\usepackage{geometry}
\geometry{a4paper,margin=1in}
\usepackage{enumitem}
\usepackage{bm}
\usepackage{graphicx}
\usepackage{hyperref} \usepackage{tikz-cd}

\usepackage{fancyhdr}
\pagestyle{fancy}
\fancyhf{} 

\fancyhead[LE,RO]{\footnotesize \thepage\quad CHAPTER~\thechapter
}



\renewcommand{\sectionmark}[1]{}
\renewcommand{\chaptermark}[1]{}

\setlength{\headheight}{13.6pt} 

\theoremstyle{plain}
\newtheorem{theorem}{Theorem}[chapter]      \newtheorem{proposition}[theorem]{Proposition} \newtheorem{lemma}[theorem]{Lemma}          \newtheorem{corollary}[theorem]{Corollary}  

\theoremstyle{definition}
\newtheorem{definition}[theorem]{Definition} \theoremstyle{remark}
\newtheorem{remark}[theorem]{Remark}         



\newcommand{\PLPIMMTCN}{\PLPIMMTEN} \newcommand{\PLPIMMTEN}{\text{Meta-Mathematical Theory based on Pan-Logic Analysis and Pan-Iterative Analysis}}
\newcommand{\PLPIMMT}{\text{PL-PI MMT}}  

\newcommand{\GFramework}{\text{G-Framework}} \newcommand{\GAlgebra}{\text{G-Algebra}}     

\newcommand{\GZLStxt}{\textsf{GZ-LS}}   \newcommand{\GZLTtxt}{\textsf{GZ-LT}}   \newcommand{\GZLOCtxt}{\textsf{GZ-LOC}} \newcommand{\GZLCurvtxt}{\textsf{GZ-LCurv}} 

\newcommand{\GZLS}{\ensuremath{\mathfrak{L}_{\mathrm{GZ}}}}       \newcommand{\GZLT}{\ensuremath{M_{!w}}}                           \newcommand{\GZLOC}{\ensuremath{A_M}}                             \newcommand{\GZLCurv}{\ensuremath{\mathcal{F}_{\mathrm{law}}}}    

\newcommand{\GZTLCtxt}{\textsf{GZ-TLC}}                           \newcommand{\GZTLC}{\ensuremath{\mathcal{A}}}                     \newcommand{\GZGMS}{\textsf{GZ-GMS}}                              

\newcommand{\LawSpace}{\ensuremath{\mathcal{L}}}                  \newcommand{\Base}{\ensuremath{M}}                                \newcommand{\Bundle}{\ensuremath{P}}                              \newcommand{\pr}{\operatorname{pr}}                               \newcommand{\Group}{\ensuremath{G}}                               \newcommand{\Conn}{\ensuremath{\omega}}                           \newcommand{\Curv}{\ensuremath{\Omega}}                           \newcommand{\Homotopy}{\ensuremath{\mathcal{H}}}                  \newcommand{\Jac}{\ensuremath{\mathsf{Jac}}}                      \newcommand{\Layer}{\ensuremath{\ell}}                            



\title{\textbf{Meta-Mathematical Theory based on Pan-Logic Analysis and Pan-Iterative Analysis\\[0.3em]
(GaoZheng G-Framework, \GFramework; also referred to as the PL-PI Meta-Mathematical Theory, \PLPIMMT)\\[0.8em]
and the Principal-Bundle-Based Generalized Noncommutative Lie Algebra\\[0.3em]
(GaoZheng G-Algebra, \GAlgebra; PFB-GNLA) -- An Integrated Construction\\[0.8em]
\large
Research Monograph on Law-Connections, Homotopy, and CH-Layering\\[0.3em]
Based on GZ-Nomenclature
}}
\author{GaoZheng}
\date{Version 1.1, 2025}

\begin{document}

\maketitle



\frontmatter

\chapter*{License and Permissions}


Copyright \textcopyright{} 2025 GaoZheng.

This arXiv version of the monograph is made available under the
Creative Commons Attribution 4.0 International license (CC BY 4.0).
In particular, any reuse or adaptation of this text must provide
appropriate credit to the author, include a reference to this
monograph, and indicate if changes were made, but is not otherwise
restricted by non-commercial or no-derivative clauses.

Earlier project-level source materials in the public repository: \\
\url{https://github.com/CTaiDeng/open_meta_mathematical_theory} \\
may be released under more restrictive licenses (such as
CC BY-NC-ND 4.0 for markdown source documents or GPL-3.0-only for
code). Those repository licenses govern public use of the
corresponding source files. The present arXiv monograph is an
author-prepared, curated derivative work of that project and is
separately licensed under CC BY 4.0 as stated above.


\chapter*{Abstract}


We develop a unified law-space and principal-bundle framework for
principal-bundle-based generalised noncommutative Lie algebras
(PFB-GNLA) and for the GaoZheng G-Algebra \\(\GAlgebra). At the
conceptual level, the construction is rooted in a meta-mathematical
theory based on Pan-Logic Analysis and Pan-Iterative Analysis
(PL-PI MMT), which treats “laws” as dynamically generated and iterated
procedures rather than fixed axioms. Geometrically, this provenance is
realised in the GaoZheng G-Framework (\GFramework) through a law-space
(\GZLS), triple-layer connections (\GZTLC), and a family of
law-curvatures (\GZLCurv) controlling the admissible flows of geometric,
parametric, and law-level data.



Within this setting we construct four integrated versions of PFB-GNLA.
The strict version encodes a principal bundle with connection and
curvature data into a near-Lie algebra of operators, in which the
Jacobi identity holds exactly and curvature quantifies controlled
deviations from flatness. The homotopy version relaxes Jacobi by
introducing a closed 3-form $H$ and a ternary operation $l_3$, yielding
a 2-term $L_\infty$-algebra that captures structured failures of ideal
behaviour. The \emph{variable functional--operator} version upgrades
parameter directions to operator-valued law-functionals
(e.g.\ $A_M$), so that “laws of evolution” become first-class
geometric objects, parallel to the way classifying maps
$\Base\to BG$ encode topological types; it also provides the
natural home for classical differential dynamics and for the
generalized Feynman path integral on law-space
($\mathcal{Z}_{\mathrm{law}}$, the GaoZheng law-space path
integral, GZ-LSPI). The meta-mathematical original version relocates
these structures to PL-PI law-space,
encoding redefined connections and higher corrections directly at the
level of rules of construction.



In parallel, we introduce a CH-layering scheme and the GZ-CBT
(Continuum Breakdown and Rewriting) machinery, which together recast
continuum questions as geometric and homotopical phenomena in law-space.
Classical continuum hypotheses (such as CH and its variants) are
assigned to a binary-law layer $l_2$, where size is measured by
cardinalities, while higher layers $l_{\ge3}$ carry homotopy, curvature,
and semantic metrics. Law-curvature diagnostics, Jacobiator-gap
invariants, and CH-breakdown indicators are combined into
certificate-type quantities that constrain admissible paths in
law-space and provide quantitative “gap statements” between strict and
homotopy regimes.



Finally, we sketch how the resulting law-space and PFB-GNLA structures
interface with mainstream theoretical physics. Gauge theories,
gravitation, quantum field theories, and renormalisation flows are
reinterpreted as projections of law-space geometry: principal bundles
and classical connections appear as slices of \GZLS, renormalisation
group trajectories as \GZTLC-horizontal flows, and anomalies or
obstructions as components of \GZLCurv. In this way the
GaoZheng G-Framework provides a coherent bridge between PL-PI
meta-mathematical provenance, classical principal bundle geometry,
and higher homotopy structures in mathematical physics.



\chapter*{Structure of the Research Monograph}


This research monograph is organised in five parts. Each part can be read on its
own, but they form a cumulative story: from PL-PI meta-mathematical
provenance, through principal-bundle–based algebraic structures, to
CH-layering and applications in physics. The aim of this summary is to
serve as a structural guide rather than a second abstract.



\paragraph*{Part I: PL-PI meta-paradigm and law-level structures.}


Part~I introduces the PL-PI Meta-Mathematical Theory (\PLPIMMT) and the
GaoZheng G-Framework (\GFramework). Chapter~1 formulates the
pan-iteration problem (PIP) and the pan-logic problem (PLP) in a purely
structural setting of heterogeneous iterative systems, and motivates a
law-space point of view as a unified solution framework. Chapter~2
(Triple-Layer Connection \GZTLCtxt) then introduces the
law-space \GZLS, triple-layer connections \GZTLC, and the GZ Law Suite
(\GZLS, \GZLT, \GZLOC, \GZLCurv). These two chapters provide the
conceptual and geometric background for all later PFB-GNLA constructions.



\paragraph*{Part II: Principal bundles and four PFB-GNLA constructions.}


Part~II moves to classical principal bundle geometry and constructs the
strict, homotopy, and variable functional–operator versions of
PFB-GNLA. Chapter~3 recalls principal bundles, connections, curvature,
and generalised noncommutative Lie algebras (GNLA). Chapters~4–6 then
build, respectively, the strict PFB-GNLA, the homotopy PFB-GNLA (with
closed 3-forms $H$ and ternary operations $l_3$), and the variable
functional–operator PFB-GNLA (with law-operator connections $A_M$ and
mixed curvatures). The dependence structure is:



\begin{itemize}
  \item Chapter~3 is the geometric and algebraic prerequisite;
  \item Chapter~4 assumes only Chapter~3;
  \item Chapter~5 requires Chapter~4 (Jacobiator-gap and $H$-twisting);
  \item Chapter~6 builds on Chapters~4–5 (adding operator-level law
        dynamics).
\end{itemize}



\paragraph*{Part III: Meta-mathematical original version and four-way comparison.}


Part~III returns to the PL-PI provenance and reconstructs the original
law-level \GAlgebra inside PL-PI law-space, then compares all four
versions. Chapter~7 starts from PL (Pan-Logic) and PI (Pan-Iterative)
analysis, builds law-objects for logics and procedures, and realises an
“original” PFB-GNLA directly at the law level. Chapter~8 then aligns
the strict, homotopy, variable functional–operator, and original
versions via a common set of structural data and gap-type invariants
(such as $\mathsf{JacGap}$). Conceptually, Part~III closes the loop:



\begin{itemize}
  \item Parts~I–II: “from geometry to law-objects”;
  \item Part~III: “from PL-PI law-objects back to $\GAlgebra$ and PFB-GNLA”.
\end{itemize}



\paragraph*{Part IV: CH-layering and continuum semantics.}


Part~IV applies the law-space machinery to the Continuum Hypothesis and
related set-theoretic phenomena. Chapter~9 reviews CH, GCH, and
independence results in ZFC, and introduces the notion of CH-layering:
CH is seen as a statement at a binary-law layer $l_2$, while higher
layers $l_{\ge3}$ carry homotopy and law-curvature data. Chapter~10
develops GZ-CBT (Continuum Breakdown and Rewriting) and presents
scenario-based analyses of how continuum breakdowns and phase transitions
can be organised as flows in law-space, monitored by certificate-type
invariants.



Readers primarily interested in PFB-GNLA may treat Part~IV as an
application module; readers focused on foundations may read Part~IV
directly after Part~I, using Parts~II–III as a structural backdrop.



\paragraph*{Part V: Outlook towards physics.}


Part~V sketches how the \GFramework\ and PFB-GNLA structures interface
with mainstream theoretical physics. Chapter~11 discusses mappings from
principal bundles, gauge fields, gravitation, and quantum field theory
to law-space objects and flows, and shows how renormalisation group
trajectories can be seen as \GZTLC-horizontal or curvature-driven paths
in \GZLS, and how standard path-integral formulations can be
reinterpreted as projections of the generalized Feynman path integral
on law-space (GZ-LSPI).



\medskip

For a first reading, a natural path is:
Part~I (to understand the meta-paradigm and law-space),
then Chapter~3 and Chapter~4 (strict PFB-GNLA),
followed by Chapter~8 (four-way comparison), and finally selected
sections of Parts~IV–V according to interest.







\tableofcontents

\mainmatter



\part{PL-PI Meta-Mathematical Theory (GaoZheng G-Framework) and Law-Level Structures}


\chapter[Pan-Iterative Systems and Pan-Logical Metrics]{From Heterogeneous Iterative Systems to the GaoZheng G-Framework:\\
         A Meta-Mathematical Theory of Pan-Logic and Pan-Iteration (PL-PI MMT)}


\section{Pan-Iterative Systems and Pan-Logical Metrics: Abstract Motivation}
\label{sec:pan-iter-motivation}


In this chapter we do not start from any particular empirical
discipline. Instead, we pose a purely structural problem about
heterogeneous iterative systems and the logical metrics that compare
them.

\begin{definition}[Heterogeneous iterative systems]
A heterogeneous iterative system is a triple
\[
  S_i = (\mathcal{X}_i,\mathcal{L}_i,\Phi_i),
\]
where:
\begin{itemize}
  \item $\mathcal{X}_i$ is a state space (not assumed to be of any
    fixed type: it may be a manifold, a discrete set, a function
    space, \emph{etc.});
  \item $\mathcal{L}_i$ is a collection of law fragments or local
    update rules acting on $\mathcal{X}_i$;
  \item $\Phi_i$ is an iterative mechanism (for example, a flow,
    semigroup, or discrete-time evolution) generated by
    $\mathcal{L}_i$.
\end{itemize}
We write
\[
  \mathfrak{S} = \{S_i\}_{i\in I}
\]
for a family of such systems, possibly with very different state
spaces, law sets, and iteration mechanisms.
\end{definition}

The \emph{pan-iteration problem} asks whether there exists a single
meta-structure in which all systems $S_i$ can be realised as instances
of one unified iterative dynamics.

\begin{definition}[Pan-iteration problem]
Given a family $\mathfrak{S} = \{S_i\}_{i\in I}$ of heterogeneous
iterative systems, the pan-iteration problem (PIP) is to construct a
law-space
\[
  \mathcal{L}_{\mathrm{tot}}
\]
and a pan-iterative mechanism
\[
  \Pi:\mathcal{L}_{\mathrm{tot}}\times\mathsf{State}_{\mathrm{tot}}
  \longrightarrow \mathsf{State}_{\mathrm{tot}}
\]
such that each $S_i$ appears as a suitable restriction, quotient, or
projection of $(\mathcal{L}_{\mathrm{tot}},\Pi)$.
\end{definition}

Once one tries to solve PIP, a second design problem immediately
emerges. In order to compare and transform different $S_i$ inside a
single law-space, we need abstract \emph{logical placeholders} and
\emph{logicality metrics} that are not tied to any particular
underlying state space.

\begin{definition}[Pan-logic problem]
The pan-logic problem (PLP) is to define, on the same law-space
$\mathcal{L}_{\mathrm{tot}}$,
\begin{itemize}
  \item logical predicates (truth, admissibility, equivalence) that
    apply uniformly to law fragments coming from different $S_i$;
  \item quantitative metrics of logicality that measure, for example,
    consistency, explanatory power, or structural proximity between
    laws.
\end{itemize}
\end{definition}

From the perspective of this monograph, the GaoZheng G-Framework or
Pan-Logic and Pan-Iteration Meta-Mathematical Theory (PL-PI MMT) is
proposed precisely as a solution to PIP and PLP: it provides a concrete
law-space with pan-iterative dynamics and pan-logical metrics in which
heterogeneous systems of very different origins can be embedded,
analysed, and transformed.



\section{Curved Spacetime and Hilbert Spaces as a Heterogeneous Evolution Model}


As a first logical model instance of the abstract pan-iteration and
pan-logic setting introduced in Section~\ref{sec:pan-iter-motivation},
we extract a classical configuration from theoretical physics and
recast it in purely mathematical terms as a heterogeneous evolution
model.

\begin{definition}[Geometric component]
Let $(M,g)$ be a four-dimensional smooth manifold equipped with a
Lorentzian (pseudo-Riemannian) metric $g$ of signature $(-,+,+,+)$.
Let $\nabla$ denote the Levi-Civita connection of $g$ on $TM$. Timelike
worldlines are smooth curves
\[
  \gamma:\mathbb{R}\to M,\qquad
  \langle \dot\gamma(t),\dot\gamma(t)\rangle_g < 0,
\]
and their geometric evolution may be described by an equation of the
form
\[
  \nabla_{\dot\gamma(t)}\dot\gamma(t)
  \;=\; F\big(\gamma(t),\dot\gamma(t)\big),
\]
where $F$ is a prescribed smooth vector field on $TM$ encoding external
fields or effective interactions.
\end{definition}



\begin{definition}[Hilbert-space component]
Let $(\mathcal{H},\langle\cdot,\cdot\rangle)$ be a complex separable
Hilbert space. We regard pure states as rays in the projective Hilbert
space $\mathbb{P}(\mathcal{H})$. For each spacetime point $x\in M$, we
assign a self-adjoint operator
\[
  H(x):\mathcal{D}(H(x))\subset\mathcal{H}\to\mathcal{H},
\]
interpreted as a Hamiltonian depending on the geometric location $x$.

Along a given worldline $\gamma$, a state vector $\psi(t)\in
\mathcal{H}$ evolves according to a Schr\"odinger-type equation
\[
  \frac{\mathrm{d}}{\mathrm{d}t}\psi(t)
  \;=\;
  -\mathrm{i}\,H\big(\gamma(t)\big)\,\psi(t),
  \qquad \psi(t)\in\mathcal{H}.
\]
\end{definition}



\begin{definition}[Coupled heterogeneous evolution]
A coupled heterogeneous evolution model is specified by the data
\[
  \mathfrak{E}
  \;=\;
  \big(
    M,g,\nabla,F;
    \mathcal{H},\langle\cdot,\cdot\rangle,H
  \big)
\]
together with a family of curves
\[
  t\longmapsto\big(\gamma(t),\psi(t)\big)\in M\times\mathcal{H},
\]
satisfying the system
\begin{align*}
  \nabla_{\dot\gamma(t)}\dot\gamma(t)
  &= F\big(\gamma(t),\dot\gamma(t),\psi(t)\big), \\
  \frac{\mathrm{d}}{\mathrm{d}t}\psi(t)
  &= -\mathrm{i}\,H\big(\gamma(t),\dot\gamma(t)\big)\,\psi(t).
\end{align*}
Here the right-hand sides are allowed to depend on both the geometric
variables $(\gamma(t),\dot\gamma(t))$ and the Hilbert state $\psi(t)$.

We call such a system a \emph{curved-spacetime--Hilbert-space
heterogeneous evolution model}.
\end{definition}



\begin{remark}[Status of the model]
\label{rem:GR-QM-fit}
The construction above is deliberately presented in purely mathematical
terms. It should \emph{not} be read as a proposed physical solution to
the problem of reconciling general relativity and quantum mechanics.

Instead, $\mathfrak{E}$ is used here only as a \emph{fitting scheme}:
it provides a well-defined heterogeneous evolution model in which a
four-dimensional pseudo-Riemannian manifold and a complex Hilbert space
are coupled through common law-objects. From the viewpoint of the
GaoZheng G-Framework and the Pan-Logic--Pan-Iteration
meta-mathematical theory, $\mathfrak{E}$ is simply one instance of a
law-space in which
\[
  \big(\text{geometric iteration on } (M,g)\big)
  \quad\text{and}\quad
  \big(\text{Hilbert-space iteration on } \mathcal{H}\big)
\]
are jointly encoded and compared.

Any potential physical interpretation---for example, as a candidate
approximation to a reconciliation of general relativity and quantum
mechanics---lies strictly beyond the scope of the present mathematical
discussion.
\end{remark}



\section[From Abstract PL-PI to the GZ Law Suite]{From Abstract PL-PI to the GZ Law Suite (\GZLS, \GZLT, \GZLOC, \GZLCurv)}


Sections~\ref{sec:pan-iter-motivation} and~1.2 have presented the
Pan-Logic and Pan-Iteration viewpoint at two levels: an abstract family
of heterogeneous iterative systems and a concrete curved-spacetime–
Hilbert-space model. We now indicate how these ingredients are realised
in the GaoZheng G-Framework by the GZ Law Suite.

At a high level, the law-space \GZLS\ is designed to host law-objects
that play the role of the elements of $\mathcal{L}_{\mathrm{tot}}$ in
the pan-iteration and pan-logic problems. Each law-object carries both
iterative content (how states are updated) and logical content (how
truth, admissibility, and equivalence are evaluated). The triple-layer
connection \GZTLC\ organises how these law-objects vary over suitable
parameter spaces, while \GZLOC\ and \GZLCurv\ encode, respectively,
law-level transformations and law-level curvature data.

Informally, one may think of the GZ Law Suite as providing:

\begin{itemize}
  \item a concrete realisation of the total law-space
    $\mathcal{L}_{\mathrm{tot}}$ as a structured collection of
    law-objects (\GZLS);
  \item a connection-like mechanism (\GZLT) that transports and compares
    these law-objects across different regions of law-space;
  \item operators (\GZLOC) describing admissible law-level
    transformations, capturing the ``morphisms of laws'';
  \item curvature-like quantities (\GZLCurv) that measure the
    path-dependence and incompatibility phenomena in law-level
    evolution.
\end{itemize}

The purpose of Chapter~2 and the subsequent development is to turn this
informal picture into a precise geometric and algebraic theory. From
the present chapter's perspective, the key point is that the PL-PI
meta-mathematical problems (PIP and PLP) can be instantiated inside the
GZ Law Suite: heterogeneous systems such as the model in Section~1.2
are embedded into law-space as trajectories governed by \GZLS,
\GZLT, \GZLOC, and \GZLCurv.



\section[Rigid Type Architectures and Heterogeneous Flow-Based Evolution]{Rigid Type Architectures and Their Mismatch with Heterogeneous Flow-Based Evolution}


In the abstract setting of Pan-Logic and Pan-Iteration introduced in
Section~\ref{sec:pan-iter-motivation}, it is natural to ask which
existing mathematical and applied formalisms can be regarded as
special cases of a pan-iterative law-space. A particularly important
class is formed by what we shall call \emph{rigid type architectures}.

Intuitively, a rigid type architecture is one in which the underlying
universe is partitioned into mutually exclusive and hierarchically
organised types, and the allowed evolutions are constrained to stay
within this static type scaffold. Such architectures work well when
all relevant dynamics can be modelled as small variations inside fixed
types, but they become inadequate as soon as heterogeneous,
cross-type, or law-changing flows are taken seriously.

\begin{definition}[Rigid type architecture]
A rigid type architecture is a tuple
\[
  \mathsf{RigidArch}
  = (X,\mathcal{T},\leq_{\mathcal{T}},\tau,\mathrm{cl}),
\]
where:
\begin{itemize}
  \item $X$ is a set of objects or states;
  \item $(\mathcal{T},\leq_{\mathcal{T}})$ is a partially ordered set
    of types;
  \item $\tau:X\to\mathcal{T}$ assigns to each $x\in X$ a unique type
    $\tau(x)$, with $\tau$ required to be stable under the intended
    dynamics on $X$;
  \item $\mathrm{cl}:\mathcal{P}(X)\to\mathcal{P}(X)$ is a closure
    operator (for example, generated by local rules or admissibility
    constraints) that respects the type partition, in the sense that
    $\mathrm{cl}(A)$ decomposes along the fibres of $\tau$.
\end{itemize}
The intended evolutions in a rigid type architecture are flows or
iterative procedures
\[
  \Phi:\mathbb{T}\times X\to X
\]
(for some time parameter set $\mathbb{T}$) such that $\tau$ is
invariant along the trajectories:
\[
  \tau\big(\Phi(t,x)\big)=\tau(x)
  \quad\text{for all } t\in\mathbb{T},\;x\in X.
\]
\end{definition}



From the PL-PI viewpoint, such a rigid architecture can be embedded
into a pan-iterative law-space only as a highly constrained corner:
the law-space trajectories are forced to remain inside type fibres,
and law-level changes that would move objects across types or deform
the type scaffold are ruled out by construction.

In contrast, the heterogeneous systems considered in
Section~\ref{sec:pan-iter-motivation} and the curved-spacetime–
Hilbert-space model of Section~1.2 intrinsically involve cross-type
and law-changing evolutions. In those examples, what counts as a
``type'' may itself depend on the current law-object, and meaningful
dynamics take place precisely at the level where types split, merge,
or are reconfigured.

\begin{remark}[From rigid type architectures to flow-based unification]
Rigid type architectures reappear in the PL-PI framework as special
cases in which:
\begin{itemize}
  \item the total law-space $\mathcal{L}_{\mathrm{tot}}$ admits a
    decomposition compatible with a fixed type partition on the
    underlying state space;
  \item pan-iterative flows $\Pi$ are restricted to trajectories that
    preserve this partition and do not access law-level directions
    that would deform it.
\end{itemize}
In this sense, rigid architectures can be viewed as degenerate
pan-iterative systems: they capture only the intra-type component of
heterogeneous flow-based evolution, while suppressing the inter-type
and law-changing components that are central to the GaoZheng
G-Framework and to the law-space constructions developed in later
chapters.
\end{remark}



These rigid architectures have been extraordinarily successful. They
support clear definitions, sharp theorems, and powerful classification
results inside each discipline. However, viewed from the meta-paradigm
of the PL-PI Meta-Mathematical Theory and the GaoZheng G-Framework
(\GFramework), they expose three structural limitations that become
decisive once one tries to model heterogeneous, law-level evolution.



\subsubsection*{(1) Fragmentation across disciplines}

Each major discipline carries its own rigid type architecture
$\mathsf{RigidArch}_i=(X_i,\mathcal{C}_i,\leq_i,\mathrm{cl}_i)$, with
its own vocabulary, scales, and admissible compositions. The “world”
seen by the sciences is then a family of disconnected charts:
\[
  (X_1,\mathcal{C}_1)\;\sqcup\;
  (X_2,\mathcal{C}_2)\;\sqcup\;\dots
\]
together with ad hoc correspondences between some of the $\mathcal{C}_i$
and $\mathcal{C}_j$.

	In practice, moving from one rigid type architecture to another is handled
	by informal “translation tables” or empirical correlations
	\[
	   T_{ij}\colon \mathcal{C}_i \dashrightarrow \mathcal{C}_j,
	\]
	which are context-dependent and rarely composable. There is no
	\emph{law-level} mechanism that tells us how a change of law in one
	domain induces a compatible change of law in another; the “bridges”
	are themselves extra-scientific constructions.
	


\subsubsection*{(2) Externalised dynamics and heterogeneous evolution}

In a rigid type architecture, dynamics is usually imposed
\emph{on top of} a fixed $(X,\mathcal{C})$. One studies
time-indexed states $x_t\in X$ and writes evolution equations such as
\[
   x_{t+1} = F(x_t),\qquad
   \mathrm{cl}(x_{t+1}) = \mathrm{cl}(x_t)\ \text{or}\ 
   \mathrm{cl}(x_{t+1})\ \text{changes along a fixed rule}.
\]
The classification lattice itself does not evolve; the allowed types,
their partial order, and their compositional rules are considered
background data.

	However, the central phenomena that motivate the PL-PI MMT and the
	GaoZheng G-Algebra (\GAlgebra) are precisely those in which
	\emph{the laws and the admissible types themselves} must be allowed to
	change. A schematic heterogeneous evolution path
	\[
	  x_0 \xrightarrow{\Phi_1} x_1 \xrightarrow{\Phi_2} \cdots
	      \xrightarrow{\Phi_n} x_n
\]
	may pass through regimes where the very meaning of “state” and “type”
	changes (for example, from a biological description to a socio-economic
description, or from a continuum model to a discrete agent model). In a
rigid type-architecture view, such a path is broken into disjoint segments,
	each studied in its own language, and the global trajectory is not
	captured as a single, law-governed object.
	


\subsubsection*{(3) Metrics of “how many” versus metrics of “how generated”}

From the perspective of CH-layering, a rigid type-architecture viewpoint corresponds
to staying in a purely extensional, binary-law layer $l_2$, where the
primary metric is “how many elements” or “how large a subset”. At this
level, one can meaningfully state continuum questions (such as the
classical Continuum Hypothesis) in terms of cardinalities and their
comparisons.

Once one moves to layers $l_{\ge3}$, where one must keep track of higher
homotopy data, law-connections, and law-curvatures, the natural metrics
are no longer purely extensional. They are dominated by generative and
semantic quantities such as
\[
\begin{aligned}
  \sigma_{\ge3}(X)
  &\sim
  \text{(generative complexity, homotopy depth, curvature spectrum)},\\
  \vartheta_n(X)
  &\sim
  \text{(degree of semantic continuity)}.
\end{aligned}
\]
Such rigid type-based frameworks, by construction, do not see this change
of metric: they try to continue asking “how many” even in regimes where
the meaningful question has become “how generated, how coupled, and at
what cost”.



\medskip

Summarising, traditional rigid type-based frameworks provide indispensable rigid
cross-sections of reality: they tell us \emph{what} kinds of objects
there are and \emph{how} they can be composed, within a fixed type
universe. What they do not provide is a unified, law-level geometry of
how heterogeneous classification schemes interact, deform, and evolve
together.

The meta-paradigm of \PLPIMMT\ and of the
GaoZheng G-Framework will be constructed precisely to complement this
rigidity. In later parts of this monograph, the strict, homotopy, and
variable functional–operator versions of the GaoZheng G-Algebra
(PFB-GNLA) will realize rigid type architectures as special, “static”
layers inside a larger law-space, and will equip that law-space with
connections and curvatures that make heterogeneous evolution a
first-class, computable object rather than an afterthought.



\section[Flow-Based Unification and PL-PI MMT]{Flow-Based Unification and the PL-PI Meta-Mathematical Theory (\PLPIMMT, \GFramework)}


The meta-paradigm of the PL-PI Meta-Mathematical Theory
(\PLPIMMT) and the GaoZheng G-Framework (\GFramework) can be
summarised, in contrast to such rigid type-based frameworks, as a shift
from “what kinds of objects exist inside a fixed hierarchy” to
“how objects and their governing laws flow, deform, and recombine
across heterogeneous regimes”.



\subsubsection*{(A) Flow-based unification: from static hierarchies to law-space flows}
	
Recall a schematic rigid type architecture
$\mathsf{RigidArch}=(X,\mathcal{C},\leq_{\mathcal{C}},\mathrm{cl})$.
In the PL-PI perspective, what matters is not only that $x\in X$ is
assigned a type $\mathrm{cl}(x)\in\mathcal{C}$ at each instant, but
that there exists a \emph{family of flows}
\[
  \gamma\colon I\to X,\qquad t\mapsto x_t,
\]
together with a law-level description of how the admissible types,
their relations, and the applicable evolution rules \emph{change along
the flow}. The target of this description is the law-space \GZLS,
introduced schematically in Section~1.1.

At an abstract level, a flow-based unification scheme consists of
\begin{itemize}
  \item a law-space $\GZLS=(\mathfrak{L};J,\mathrm{Loc},\Sigma)$,
        whose elements are “law-objects” equipped with composition,
        metric, and threshold data;
	  \item a family of embeddings
	        \[
	          \Psi_i\colon \mathsf{RigidArch}_i
	          \longrightarrow \GZLS,
	        \]
	        sending each discipline’s rigid type architecture
	        to a region or layer of \GZLS;
  \item a collection of curves
        \[
          \Gamma\subseteq \{\Gamma\colon I\to\GZLS\}
        \]
        representing allowed \emph{law-level flows}, i.e. admissible
        ways in which laws, types, and admissible evolutions may
	        continuously deform, bifurcate, or couple across different
	        $\mathsf{RigidArch}_i$.
\end{itemize}

The unification is “flow-based” in the sense that it is not achieved
by collapsing all classification vocabularies into a single static
ontology, but by organising law-objects into a geometry where
\emph{paths} and \emph{homotopies of paths} become first-class
mathematical objects. Interdisciplinary translation then takes the
form of following a law-level curve in $\GZLS$, rather than switching
between disconnected tables of correspondences.



\subsubsection*{(B) PL: Pan-Logic Analysis as law-level organisation of inference}

The “PL” in \PLPIMMT stands for Pan-Logic Analysis. At a schematic
level, it starts from the observation that different domains employ
different logics (deductive calculi, probabilistic updates, causal
inference rules, constraint propagation, etc.), each of which can be
seen as a small category or a sequent-style system
\[
   \mathsf{Logic}_\alpha=(\mathsf{Form}_\alpha,\;\vdash_\alpha).
\]
Pan-Logic Analysis assigns to each such local logic a law-object
$L_\alpha\in\GZLS$ and to each admissible change of logic (e.g.,
adding a modality, relaxing an axiom, embedding one logic into
another) a law-transform
\[
   \GZLT\colon L_\alpha \longrightarrow L_\beta.
\]



\begin{definition}[Schematic pan-logic layer]
  A pan-logic layer inside \GZLS\ is a collection
  \[
     \mathfrak{L}_{\mathrm{PL}}
     = \{L_\alpha\}_{\alpha\in A}
  \]
  of law-objects representing local logics, together with a family of
  morphisms
  \[
     \GZLT^{\alpha\beta}\colon L_\alpha \longrightarrow L_\beta,
  \]
  closed under composition and equipped with a metric/curvature
  structure $(J,\mathrm{Loc})$ that measures the “distance” and
  “direction” between different logics.
\end{definition}



In this picture, logical reasoning and changes of reasoning scheme are
not external to the mathematical model: they become trajectories and
deformations inside the law-space. This is the first pillar of the
meta-paradigm: the organisation of all admissible inference schemes as
a structured region of \GZLS, governed by law-connections and
law-curvatures that will be made explicit later via the GaoZheng
G-Algebra (\GAlgebra) and the triple-layer connection \GZTLC.



\subsubsection*{(C) PI: Pan-Iterative Analysis as law-level organisation of procedures}

The “PI” in \PLPIMMT stands for Pan-Iterative Analysis. While Pan-Logic
Analysis organises logics as law-objects, Pan-Iterative Analysis does
the same for iterative procedures: numerical schemes, learning
algorithms, control loops, renormalisation flows, and so on.



At a schematic level, an iterative scheme is given by
\[
   x_{n+1} = F_\theta(x_n),\qquad n=0,1,2,\dots,
\]
with parameters $\theta$ and possibly stochastic or path-dependent
features. Pan-Iterative Analysis assigns to each such scheme a
law-object $L^{\mathrm{iter}}_\beta\in\GZLS$ and to each admissible
transformation of schemes (e.g., changing step sizes, changing loss
functions, upgrading from deterministic to stochastic) a law-transform
inside the same law-space.



\begin{definition}[Schematic pan-iterative layer]
  A pan-iterative layer inside \GZLS\ is a collection
  \[
    \mathfrak{L}_{\mathrm{PI}}
    = \{L^{\mathrm{iter}}_\beta\}_{\beta\in B},
  \]
  together with
  \begin{itemize}
    \item a family of evaluation functionals
          $E_\beta$ capturing convergence, stability, cost, or other
          performance measures of the associated iterative schemes;
    \item a family of law-transforms
          $\GZLT^{\beta\gamma}\colon
          L^{\mathrm{iter}}_\beta\to L^{\mathrm{iter}}_\gamma$
          encoding admissible design moves in the space of procedures.
  \end{itemize}
\end{definition}



The key point is that these law-objects and transforms live in the
\emph{same} law-space \GZLS as the pan-logic layer. This allows one to
speak of \emph{flows} that simultaneously modify the underlying logic
and the iterative procedures, keeping track of a global “semantic
curvature” induced by their interaction.



\subsubsection*{(D) The PL-PI Meta-Mathematical Theory as an integrated law-level structure}

Putting the pan-logic and pan-iterative layers together, the PL-PI
Meta-Mathematical Theory can be viewed, at a schematic level, as an
integrated law-level structure

\begin{multline*}
   \PLPIMMTCN ;\\(\PLPIMMT)
   \;\sim\;
   (\GZLS,\GZLT,\GZLOC,\GZLCurv;\\
    \mathfrak{L}_{\mathrm{PL}},\mathfrak{L}_{\mathrm{PI}},
    \text{flow/curvature constraints}),
\end{multline*}

where:

\begin{itemize}
  \item $\GZLS$ is the law-space carrying both pan-logic and
        pan-iterative layers as distinguished regions or strata;
  \item $\GZLT$ is the family of law-transforms connecting law-objects
        within and across these layers;
  \item $\GZLOC$ is a law-connection on \GZLS, specifying how to
        transport laws along flows in underlying object-level
        geometries (such as principal bundles);
  \item $\GZLCurv$ is the associated family of law-curvatures,
        measuring the obstruction to trivialising laws along closed
        loops (in both state space and law-space);
  \item flow/curvature constraints are conditions ensuring that flows
        in $\GZLS$ are compatible with object-level dynamics and with
        higher homotopy data (to be expressed later via homotopy
        versions of \GAlgebra\ and CH-layering).
\end{itemize}



\begin{remark}[From rigid type architectures to flow-based unification]
  Rigid type architectures reappear in this picture as special
  cases where $\GZLOC$ is flat along the relevant region of law-space
  and where the law-curvature $\GZLCurv$ vanishes on the flows under
  consideration. In that regime, law-level evolution can be neglected,
  and one is effectively working inside a single, static
  $\mathsf{RigidArch}$ with its own fixed logic and procedure class.
\end{remark}

From the perspective of this monograph, the PL-PI MMT and the
\GFramework provide the meta-level scaffolding in which the four
integrated PFB-GNLA constructions will be placed:

\begin{itemize}
  \item the \emph{strict version} realises object-level principal bundle
        geometries as law-objects in $\GZLS$ with flat or controlled
        law-curvature;
  \item the \emph{homotopy version} incorporates non-trivial Jacobiators
        and higher homotopies as law-curvature effects, to be measured
        by certificate-type invariants (such as Jacobiator-gap
        certificates introduced later);
  \item The \emph{variable functional–operator version} extends
        $\mathsf{Obj}$ to path and functional spaces, aligns with the
        pan-iterative layer in $\mathsf{Law}$, and uses \GZTLC\ to
        implement law-connections directly on operator paths; at this
        level, classical differential dynamics and law-space path
        integrals (in particular the generalized Feynman path integral
        on law-space, $\mathcal{Z}_{\mathrm{law}}$, a.k.a.\ the
        GaoZheng law-space path integral, GZ-LSPI) are encoded as
        operator paths.
\item the \emph{meta-mathematical original version} reconstructs the
        whole structure from the PL-PI provenance, showing that object-
        level geometries are themselves emergent from law-level flows.
\end{itemize}

In this sense, “flow-based unification” is not a slogan but a precise
meta-mathematical demand: all rigid type architectures, all local
logics, and all iterative procedures must be realisable as slices or
layers of a single law-space endowed with connections, curvatures, and
homotopy data, such that heterogeneous evolution paths become
\emph{geodesics, parallel transports, and curvature effects} in a
unified GaoZheng G-Framework.



\section{The GZ Law Suite and Its Abstract Location in the Framework}


\begin{definition}[GaoZheng Law-Space \GZLS\ (schematic)]
We set
  \[
    \GZLS = (\mathfrak{L}; J,\mathrm{Loc},\Sigma),
  \]
  where $\mathfrak{L}$ is the semantic universe for law-objects and their composition rules,
  $J,\mathrm{Loc}$ encode metric and localization data, and $\Sigma$ is a family of thresholds
  that trigger transitions between continuous and discrete regimes.
\end{definition}




\begin{remark}[Law-Space as a Target for Object-Level Theories]
In later parts of this monograph, principal bundles, connections,
  and generalized Lie structures will be treated as \emph{object-level}
  geometries. The role of the law-space $\GZLS$ is to serve as a semantic
  target in which these geometries appear as law-objects with explicit
  composition, metric, and threshold data $(J,\mathrm{Loc},\Sigma)$.
\end{remark}


\begin{proposition}[Object-Level Geometry to Law-Objects in \GZLS\ (schematic)]
There exists a functorial assignment
  \[
    \Phi_{\mathrm{geom}} : \mathsf{GeomObj} \longrightarrow \GZLS,
  \]
  from a suitable category $\mathsf{GeomObj}$ of geometric objects
  (for example, principal bundles with connection and curvature data)
  to law-objects in $\GZLS$, such that:
  \begin{enumerate}[label=\roman*)]
    \item parallel transport, holonomy, and curvature of a geometric object
          are encoded as composition and metric data in the corresponding
          law-object in $\GZLS$;
    \item gauge-equivalent geometric data are mapped to isomorphic law-objects
          in $\GZLS$.
  \end{enumerate}
  The construction of $\Phi_{\mathrm{geom}}$ and the precise definition of
  $\mathsf{GeomObj}$ are developed in Part~II and revisited in Part~III.
  A more systematic treatment of this embedding is outlined in
  Technical Appendix~III (Section~\ref{app:III-P1}).
\end{proposition}


\section[Positioning the GaoZheng Generalized Mathematical Structure]{Positioning the GZ-GMS (\GZGMS)}


The GaoZheng Generalized Mathematical Structure \GZGMS\ is the ambient object in \\ 
which the PL-PI Meta-Mathematical Theory (\PLPIMMT), the GaoZheng G-Framework \\ 
(\GFramework), and the four integrated versions of the GaoZheng G-Algebra (\GAlgebra, PFB-GNLA)
are jointly realised. Informally, \GZGMS is the level at which
“what counts as an object”, “what counts as a law”, and “what counts
as an admissible flow or deformation” are organised into a single,
coherent, law-level geometry.


\subsubsection*{(a) From law-space and flows to a generalized structure}

Sections~1.1 and~1.2 introduced the law-space \GZLS, law-transforms
\GZLT, and the flow-based unification picture. At that level, the
emphasis is on law-objects and their deformations, while the
underlying “object-level geometries” (manifolds, bundles, operator
spaces, etc.) appear only indirectly.

In order to connect this law-level picture with concrete constructions
such as principal bundles with connection, homotopy Lie structures,
and variable functional–operator systems, one needs a generalized
structure that simultaneously carries:



\begin{itemize}
  \item categories (or higher categories) of object-level entities
        (spaces, bundles, operators, path spaces);
  \item the law-space \GZLS\ and its law-transforms \GZLT;
  \item a system of flows and homotopies on both the object side and
        the law side;
  \item a triple-layer connection \GZTLC\ linking these components;
  \item semantic layering data (such as CH-layering) that governs
        which metrics and invariants are visible at each level.
\end{itemize}

The abstract object that packages all of this is what we call the
GaoZheng Generalized Mathematical Structure.

\begin{definition}[GZ-GMS (schematic)]
  A GaoZheng Generalized Mathematical Structure is a tuple
  \[
    \GZGMS =
    \bigl(\mathsf{Obj},\mathsf{Law},\mathsf{Flow},\mathsf{Layer},
          \GZTLC;\,\GZLCurv\bigr),
  \]
  where
  \begin{itemize}
    \item $\mathsf{Obj}$ is a (multi-layer) category of object-level
          geometries, whose typical objects include principal bundles
          with connection and curvature, homotopy-enhanced algebraic
          structures, and operator/path-space systems;
    \item $\mathsf{Law}$ is the law-space \GZLS\ equipped with
          law-transforms \GZLT\ and law-connection \GZLOC, together
          with its law-curvature family \GZLCurv;
    \item $\mathsf{Flow}$ is a class of admissible flows and homotopies
          on $\mathsf{Obj}$ and on $\mathsf{Law}$, together with a
          compatibility rule specifying how object-level flows induce
          law-level flows and conversely;
    \item $\mathsf{Layer}$ encodes semantic and CH-like layering data,
          distinguishing, for example, purely extensional layers
          (cardinal/measure-based) from higher homotopy and law-curvature
          layers;
    \item $\GZTLC$ is a triple-layer connection that couples
          object-level geometry, parameter layers, and law-level
          structures, and that projects to \GZLOC\ on the law-side.
  \end{itemize}
\end{definition}


\subsubsection*{(b) Positioning \GZGMS\ relative to \PLPIMMT\ and \GFramework}

Within the PL-PI Meta-Mathematical Theory, the pan-logic and
pan-iterative layers (Section~1.2) are realised as distinguished
regions of the law-space component $\mathsf{Law}$ of \GZGMS. The
G-Framework (\GFramework) can then be seen as the meta-level scheme
that organises:

\begin{itemize}
  \item which object-level categories are admitted into $\mathsf{Obj}$;
  \item which families of logics and iterative schemes are represented
        as law-objects in $\mathsf{Law}$;
  \item which flows and homotopies are allowed in $\mathsf{Flow}$;
  \item how CH-layering and related semantic stratifications are
        encoded in $\mathsf{Layer}$.
\end{itemize}

Diagrammatically, one may think of \GFramework\ as specifying a
“design envelope”

\[
   \GFramework:
   \Bigl(
     \mathsf{Obj},\mathsf{Law},\mathsf{Flow},\mathsf{Layer}
   \Bigr)
   \text{ equipped with }
   \GZTLC,\;\GZLCurv.
\]

while \PLPIMMT\ provides the meta-mathematical principles that govern
how inference schemes and iterative procedures populate this envelope.
The \GZGMS\ is then the concrete realisation of this envelope for a
given choice of object-level geometries and law-level structures.


\subsubsection*{(c) Positioning \GZGMS\ relative to the four PFB-GNLA constructions}

The four integrated PFB-GNLA constructions, developed later in this
monograph and in the kernel documents, can now be positioned as
\emph{charts} or \emph{slices} inside \GZGMS:

\begin{itemize}
  \item The \emph{strict version} of \GAlgebra\ selects from
        $\mathsf{Obj}$ those principal bundles with parameterised
        connections for which the Jacobi identity holds strictly, and
        realises them as law-objects in $\mathsf{Law}$ with controlled
        or vanishing law-curvature on the relevant flows.
  \item The \emph{homotopy version} enlarges $\mathsf{Obj}$ to include
        homotopy Lie structures and encodes non-trivial Jacobiators as
        components of law-curvature in \GZLCurv, to be measured and
        controlled by certificate-type invariants (such as
        Jacobiator-gap certificates introduced later in Part~III).
  \item The \emph{variable functional–operator homotopy version}
        extends $\mathsf{Obj}$ to path and functional spaces, aligns
        with the pan-iterative layer in $\mathsf{Law}$, and uses
        \GZTLC\ to implement law-connections directly on operator
        paths.
  \item The \emph{meta-mathematical original version} reconstructs the
        whole PFB-GNLA picture from PL-PI provenance data, showing how
        what appears as “object-level geometry” can be reinterpreted
        as emergent from law-level flows and homotopies in \GZLS.
\end{itemize}

In each case, the corresponding PFB-GNLA construction can be seen as
choosing a substructure
\[
   \GZGMS_{\mathrm{strict}},\quad
   \GZGMS_{\mathrm{homo}},\quad
   \GZGMS_{\mathrm{vf-op}},\quad
   \GZGMS_{\mathrm{orig}}
\]
of a common ambient \GZGMS, together with comparison functors between
these substructures. The “four-way comparison” performed in later
chapters is precisely an analysis of how these substructures embed
into and overlap inside a single GaoZheng Generalized Mathematical
Structure.


\subsubsection*{(d) Positioning \GZGMS\ relative to CH-layering and semantic metrics}

Finally, the role of \GZGMS\ with respect to CH-layering can be
summarised as follows. The semantic layering data $\mathsf{Layer}$ in
the definition of \GZGMS\ records, for each region of
$\mathsf{Obj}$ and $\mathsf{Law}$, which kind of metric is meaningful:

\begin{itemize}
  \item at extensional layers (denoted schematically by $l_2$), one
        can use cardinality- or measure-based metrics on object sets
        and law-sets;
  \item at higher layers ($l_{\ge3}$), one must use homotopy- and
        law-curvature-based metrics, such as invariants extracted from
        \GZLCurv\ and from higher brackets in homotopy versions of
        \GAlgebra.
\end{itemize}

Within \GZGMS, CH-layering is thus not an external commentary on
classical set theory, but an internal stratification of the generalized
structure: different layers of \GZGMS\ “see” different aspects of
continuum and discreteness, and the breakdown/reconstruction of
continuum semantics (Part~IV) is realised as a controlled change of
layer inside the same ambient structure.

\paragraph*{Dynamical mechanisms in the GaoZheng G-Framework.}
In the present monograph, the GaoZheng G-Framework adopts a two-layer
dynamical semantics. On the classical side, law-level dynamics are
described by trajectories
\[
  \Gamma:[0,T]\to\GZLS,
\]
together with lifts to the object-level component $\mathsf{Obj}$ that
are horizontal with respect to the triple-layer connection \GZTLC. Such
trajectories are governed schematically by equations of motion of the
form
\[
  \frac{d\Gamma}{dt} = X_{\mathrm{law}}(\Gamma(t)),
\]
where $X_{\mathrm{law}}$ is a vector field determined by the
law-connection \GZLOC\ and the law-curvature family \GZLCurv, encoding
how admissible flows in $\mathsf{Obj}$ and $\mathsf{Law}$ must be
coupled.

On the quantum side, each admissible law-space trajectory $\Gamma$ is
assigned a law-action functional $S_{\mathrm{law}}[\Gamma]$, and the
corresponding law-space path integral is schematically written as
\[
  \mathcal{Z}_{\mathrm{law}}
  = \int_{\mathrm{Hist}(\GZLS)}
    \exp\!\Bigl(\frac{i}{\hbar}S_{\mathrm{law}}[\Gamma]\Bigr)\,
    \mathcal{D}\Gamma,
\]
where $\mathrm{Hist}(\GZLS)$ denotes a suitable space of law-level
histories and $\mathcal{D}\Gamma$ is a formal path measure. Classical
law-level flows then appear as stationary points or semiclassical
approximations of this path integral, so that differential equations and
path-integral quantum descriptions are unified inside the same
law-space.

We shall refer to such path integrals on $\GZLS$ as \emph{generalized
Feynman path integrals on law-space}; in the specific context of the
GaoZheng G-Framework we will also call $\mathcal{Z}_{\mathrm{law}}$ the
\emph{GaoZheng law-space path integral} (GZ-LSPI).


In the semiclassical regime, contributions to $\mathcal{Z}_{\mathrm{law}}$
are concentrated near law-level histories $\Gamma_*$ satisfying
$\delta S_{\mathrm{law}}[\Gamma_*]/\delta\Gamma = 0$, which we may view
as classical or ``optimal'' law-trajectories inside \GZLS.


In this volume, these differential and path-integral mechanisms form the
the primary dynamical content of the GaoZheng G-Framework. In
particular, the generalized Feynman path integral on law-space (the
GaoZheng law-space path integral, GZ-LSPI) plays for \GFramework\ the
role that the usual Feynman path integral plays for quantum
mechanics. Possible engineering applications to
reinforcement-learning-type procedures or certificate-based AI can then
be formulated as special instances of GZ-LSPI, obtained by interpreting
$S_{\mathrm{law}}$ as a reward or performance functional on
law-trajectories; such reinforcement-learning-type specialisations,
however, are left to separate work.


\begin{remark}[Summary: \GZGMS\ as the ambient object of this monograph]
  The remainder of this monograph can be read as a step-by-step
  construction of a concrete \GZGMS\ satisfying the schematic definition
  above, together with a detailed analysis of its four PFB-GNLA charts
  and its CH-layered semantics. In this sense, the GaoZheng Generalized
  Mathematical Structure is the geometric and algebraic “stage” on
  which the PL-PI Meta-Mathematical Theory, the G-Framework, and the
  law-connection/homotopy machinery are jointly enacted.
\end{remark}


\chapter[Triple-Layer Connection \GZTLCtxt]{The Triple-Layer Connection \GZTLCtxt \  and the Unification of Geometry,\\
         Parameters, and Laws}


\section{Formal Structure of the Triple-Layer Connection \GZTLC}


The purpose of this section is to fix a minimal but precise formal model
for the triple-layer connection \GZTLC. We isolate the three constituent
layers (geometry, parameters, laws), package them into a single
principal-bundle–type object, and formulate the compatibility
conditions that make \GZTLC into the basic unifying device for all
subsequent PFB-GNLA constructions.



\subsubsection*{(a) Three-layer base and total configuration space}

We start from the three layers that have appeared implicitly in
Part~I.

\begin{itemize}
  \item \emph{Geometric layer.}
        A principal bundle $(\Bundle,\Base,\Group)$ with classical
        connection and curvature
        \[
           \Conn\in\Omega^1(\Bundle;\mathfrak{g}),\qquad
           \Curv = \mathrm{d}\Conn + \tfrac12[\Conn,\Conn]
           \in\Omega^2(\Bundle;\mathfrak{g}),
        \]
        where $\mathfrak{g}=\mathrm{Lie}(\Group)$.
        This is the standard “object-level” geometry.
  \item \emph{Parameter layer.}
        A smooth parameter manifold $\Lambda_{\mathrm{par}}$ (for
        example, hyperparameters, control variables, model indices).
        In later PFB-GNLA constructions this layer indexes families of
        connections and operators, but at this stage it is just a
        smooth base space.
  \item \emph{Law layer.}
        A law-space
        \[
          \GZLS=(\mathfrak{L};J,\mathrm{Loc},\Sigma),
        \]
        as introduced in Section~1.2, whose points are law-objects and
        whose additional data $(J,\mathrm{Loc},\Sigma)$ encode metric,
        localisation, and threshold information. The law-connection
        \GZLOC\ and law-curvature family \GZLCurv\ will live on bundles
        with typical fiber contained in $\mathfrak{L}$.
\end{itemize}



For the purpose of defining \GZTLC\ we work, locally, over a finite
chart $\mathfrak{L}_0\subset\mathfrak{L}$ and set the \emph{total
configuration base}
\[
   X_{\mathrm{tot}}
   \;:=\;
   \Base \times \Lambda_{\mathrm{par}} \times \mathfrak{L}_0.
\]



\begin{definition}[Triple-layer principal bundle (schematic)]
  A \emph{triple-layer principal bundle} over
  $(\Base,\Lambda_{\mathrm{par}},\GZLS)$ consists of:
  \begin{itemize}
    \item a Lie group decomposition
          \[
             \widehat{G}
             \;=\;
             G_{\mathrm{geom}}\times G_{\mathrm{par}}\times G_{\mathrm{law}},
          \]
          with Lie algebra
          $\widehat{\mathfrak{g}}
           = \mathfrak{g}_{\mathrm{geom}}
             \oplus\mathfrak{g}_{\mathrm{par}}
             \oplus\mathfrak{g}_{\mathrm{law}}$;
          in the basic PFB-GNLA setting one may take
          $G_{\mathrm{geom}}=\Group$ and let
          $G_{\mathrm{par}},G_{\mathrm{law}}$ encode parameter and
          law symmetries;
    \item a principal $\widehat{G}$-bundle
          \[
              \widehat{\pi}:\widehat{\Bundle}\longrightarrow X_{\mathrm{tot}};
          \]
    \item projections
          \[
             \mathrm{pr}_{\mathrm{geom}}:
             \widehat{\Bundle}\to\Bundle,\quad
             \mathrm{pr}_{\mathrm{par}}:
             \widehat{\Bundle}\to\Lambda_{\mathrm{par}},\quad
             \mathrm{pr}_{\mathrm{law}}:
             \widehat{\Bundle}\to\mathfrak{L}_0
          \]
          compatible with $\widehat{\pi}$ in the obvious sense.
  \end{itemize}
\end{definition}

Informally, $\widehat{\Bundle}$ is a principal bundle that remembers, at
each point, a geometric configuration, a parameter choice, and a
law-object, together with their gauge symmetries.



\subsubsection*{(b) Triple-layer connection as a structured principal connection}

On a principal $\widehat{G}$-bundle $\widehat{\Bundle}$, a classical
connection is a $1$-form
\[
  \mathcal{A}
  \;\in\;
  \Omega^1(\widehat{\Bundle};\widehat{\mathfrak{g}})
\]
satisfying the usual $\widehat{G}$-equivariance and reproducing-vertical
conditions. In the present context we require that this connection be
compatible with the three-layer splitting of base and group.

\begin{definition}[Triple-layer connection \GZTLC\ (formal)]
  A \emph{triple-layer connection} on a triple-layer principal bundle
  $\widehat{\Bundle}\to X_{\mathrm{tot}}$ is a principal connection
  \[
     \GZTLC
     \;=\;
     \mathcal{A}
     \;\in\;
     \Omega^1(\widehat{\Bundle};\widehat{\mathfrak{g}})
  \]
  together with fixed projections
  \[
     \pi_{\mathrm{geom}}:
       \widehat{\mathfrak{g}}\to\mathfrak{g}_{\mathrm{geom}},\quad
     \pi_{\mathrm{par}}:
       \widehat{\mathfrak{g}}\to\mathfrak{g}_{\mathrm{par}},\quad
     \pi_{\mathrm{law}}:
       \widehat{\mathfrak{g}}\to\mathfrak{g}_{\mathrm{law}},
  \]
  such that the following conditions hold.
  \begin{enumerate}[label=\roman*)]
    \item (\emph{Geometric component.})
          The projected $1$-form
          \[
             \Conn_{\mathrm{tot}}
             \;:=\;
             \pi_{\mathrm{geom}}\circ\mathcal{A}
             \;\in\;
             \Omega^1(\widehat{\Bundle};
                      \mathfrak{g}_{\mathrm{geom}})
          \]
          restricts, along each slice with fixed parameter and law data,
          to the pull-back of the classical connection $\Conn$ on
          $(\Bundle,\Base,\Group)$. That is, for any local section
          $s:\Base\to\Bundle$ and any fixed
          $(\lambda,\ell)\in
             \Lambda_{\mathrm{par}}\times\mathfrak{L}_0$,
          the induced section
          $\widehat{s}(x)=(s(x),\lambda,\ell)$ satisfies
          \[
             \widehat{s}^*(\Conn_{\mathrm{tot}})=s^*(\Conn).
          \]
    \item (\emph{Parameter component.})
          The projected $1$-form
          \[
             B_{\mathrm{par}}
             \;:=\;
             \pi_{\mathrm{par}}\circ\mathcal{A}
             \;\in\;
             \Omega^1(\widehat{\Bundle};
                      \mathfrak{g}_{\mathrm{par}})
          \]
          is basic with respect to the projection
          $\widehat{\Bundle}\to\Lambda_{\mathrm{par}}$, and thus
          induces a well-defined Ehresmann connection on the parameter
          layer; infinitesimal changes in $\Lambda_{\mathrm{par}}$ are
          lifted horizontally according to $B_{\mathrm{par}}$.
    \item (\emph{Law component and \GZLOC.})
          The projected $1$-form
          \[
             A_{\mathrm{law}}
             \;:=\;
             \pi_{\mathrm{law}}\circ\mathcal{A}
             \;\in\;
             \Omega^1(\widehat{\Bundle};
                      \mathfrak{g}_{\mathrm{law}})
          \]
          restricts, after suitable trivialisation of the law fiber, to
          the law-connection \GZLOC\ on the law-space \GZLS; more
          precisely, there exists a law-bundle
          $\mathcal{P}_{\mathrm{law}}\to\Base\times\Lambda_{\mathrm{par}}$
          with fiber in $\mathfrak{L}$ and a projection
          $\rho:\widehat{\Bundle}\to\mathcal{P}_{\mathrm{law}}$
          such that
          \[
             A_{\mathrm{law}}=\rho^*(\GZLOC).
          \]
  \end{enumerate}
  We refer to the triple
  $(\Conn_{\mathrm{tot}},B_{\mathrm{par}},A_{\mathrm{law}})$ as the
  \emph{layer-wise decomposition} of \GZTLC.
\end{definition}

In words, a triple-layer connection is an ordinary principal connection
on a suitably enlarged principal bundle, whose projections reproduce,
respectively, the object-level geometric connection, a parameter
connection, and the law-connection on \GZLS.



\begin{remark}[Horizontal distribution viewpoint]
  Equivalently, \GZTLC\ determines a horizontal distribution
  $H_{\mathcal{A}}\subset T\widehat{\Bundle}$ complementary to the
  vertical subbundle. Using the splitting
  \[
      T X_{\mathrm{tot}}
      \;\cong\;
      T\Base
      \;\oplus\;
      T\Lambda_{\mathrm{par}}
      \;\oplus\;
      T\mathfrak{L}_0,
  \]
  one may decompose
  \[
      H_{\mathcal{A}}
      \;=\;
      H^{\mathrm{geom}}
      \;\oplus\;
      H^{\mathrm{par}}
      \;\oplus\;
      H^{\mathrm{law}},
  \]
  where $H^{\mathrm{geom}}$ consists of horizontal lifts of vectors in
  $T\Base$, \\ 
  etc. The three components $(H^{\mathrm{geom}},H^{\mathrm{par}},H^{\mathrm{law}})$ encode the
  same data as $(\Conn_{\mathrm{tot}},B_{\mathrm{par}},A_{\mathrm{law}})$.
\end{remark}



\subsubsection*{(c) Curvature decomposition and coupling of layers}

As for any principal connection, \GZTLC\ has a curvature
\[
   \mathcal{F}_{\GZTLC}
   \;:=\;
   \mathrm{d}\mathcal{A}
   + \tfrac12[\mathcal{A},\mathcal{A}]
   \;\in\;
   \Omega^2(\widehat{\Bundle};\widehat{\mathfrak{g}}).
\]
Its projections onto the three Lie algebra components produce
layer-wise curvatures:
\[
\begin{aligned}
   \Curv_{\mathrm{geom}}
   &:= \pi_{\mathrm{geom}}
       \bigl(\mathcal{F}_{\GZTLC}\bigr),\\
   \Curv_{\mathrm{par}}
   &:= \pi_{\mathrm{par}}
       \bigl(\mathcal{F}_{\GZTLC}\bigr),\\
   \GZLCurv
   &:= \pi_{\mathrm{law}}
       \bigl(\mathcal{F}_{\GZTLC}\bigr).
\end{aligned}
\]

In the simplest (decoupled) situation, one would have
\[
   \mathcal{A}
   \;=\;
   \Conn_{\mathrm{tot}}
   \;\oplus\;
   B_{\mathrm{par}}
   \;\oplus\;
   A_{\mathrm{law}},
\]
with $\Curv_{\mathrm{geom}}$, $\Curv_{\mathrm{par}}$, and \GZLCurv\
depending only on their own directions. In the general setting relevant
for PFB-GNLA, however, $\mathcal{A}$ is not block-diagonal: its
horizontal distribution mixes geometric, parameter, and law
directions. This mixing manifests itself as cross-terms in the curvature
decomposition.

\begin{definition}[Coupling terms in the curvature]
  The \emph{coupling curvature components} of \GZTLC\ are the mixed
  pieces of $\mathcal{F}_{\GZTLC}$ that involve at least two distinct
  base directions. Schematically, writing
  \[
    T X_{\mathrm{tot}}
    =
    T\Base
    \oplus T\Lambda_{\mathrm{par}}
    \oplus T\mathfrak{L}_0,
  \]
  one has a decomposition
  \[
    \mathcal{F}_{\GZTLC}
    = \mathcal{F}^{MM}
      + \mathcal{F}^{M\Lambda}
      + \mathcal{F}^{M\mathfrak{L}}
      + \mathcal{F}^{\Lambda\Lambda}
      + \mathcal{F}^{\Lambda\mathfrak{L}}
      + \mathcal{F}^{\mathfrak{L}\mathfrak{L}},
  \]
  where, for example, $\mathcal{F}^{M\Lambda}$ collects components with
  one leg in $T\Base$ and one in $T\Lambda_{\mathrm{par}}$, and so on.
  Their projections
  $\pi_{\mathrm{geom}},\pi_{\mathrm{par}},\pi_{\mathrm{law}}$ define
  geometric–parameter, geometric–law, and parameter–law coupling
  curvatures.
\end{definition}


The law-curvature family \GZLCurv\ defined earlier is obtained by
projecting $\mathcal{F}_{\GZTLC}$ to the law-component and then
restricting to appropriate directions. In particular:

\begin{itemize}
  \item purely law–law components of \GZLCurv\ encode intrinsic
        curvature of law-space flows (e.g. homotopy brackets, Jacobiator
        effects in homotopy versions of \GAlgebra);
  \item mixed geometric–law and parameter–law components measure how
        changes in geometric configuration or parameter values induce
        non-trivial parallel transport and curvature in the law-layer.
\end{itemize}



\begin{proposition}[Compatibility with layer-wise Bianchi identities (schematic)]
  Suppose that each projected curvature
  $\Curv_{\mathrm{geom}}$, $\Curv_{\mathrm{par}}$, \GZLCurv\ satisfies
  the usual Bianchi identity with respect to its projected connection.
  Then the covariant derivative $D_{\mathcal{A}}$ associated with
  \GZTLC\ satisfies a unified Bianchi identity
  \[
      D_{\mathcal{A}}\mathcal{F}_{\GZTLC} = 0,
  \]
  and the coupling curvature components obey consistency conditions
  of the form
  \[
      D_{\mathrm{geom}}\mathcal{F}^{M\Lambda}
      + D_{\mathrm{par}}\mathcal{F}^{\Lambda M}
      + \text{(law terms)} = 0,
  \]
  where $D_{\mathrm{geom}}$ and $D_{\mathrm{par}}$ denote the
  covariant derivatives induced by $\Conn_{\mathrm{tot}}$ and
  $B_{\mathrm{par}}$. These relations express that “failures of
  flatness” in one layer are transported consistently across the other
  layers by \GZTLC.
  Further details and a complete proof are given in Technical Appendix~II
  (Section~\ref{app:II-P2}).
\end{proposition}



We will not need the full detailed form of these Bianchi-type relations
in this monograph; what matters is that \GZTLC\ organises all such
compatibility constraints into a single geometric object.

\subsubsection*{(d) Role of \GZTLC\ in PFB-GNLA and law-level geometry}

With the above formal structure in place, the role of the triple-layer
connection \GZTLC\ in the overall framework can be summarised as
follows.

\begin{itemize}
  \item In the \emph{strict} PFB-GNLA version, \GZTLC\ reduces to a
        situation where $\Curv_{\mathrm{par}}$ and the mixed coupling
        terms are tightly controlled; law-curvature \GZLCurv\ encodes
        essentially the same information as the strict geometric
        curvature $\Curv$, and law-level evolution can be read off from
        object-level geometry.
  \item In the \emph{homotopy} and \emph{variable functional–operator}
        versions, non-trivial law-curvatures and mixed terms appear.
        These encode, for example, controlled failures of the Jacobi
        identity and path-dependent operator flows; \GZTLC\ provides the
        ambient connection in which Jacobiator-gap certificates and
        GZ-OHU constructions will be formulated.
  \item In the \emph{meta-mathematical original} version, the emphasis
        is reversed: law-level flows and \GZLCurv\ become primary, and
        object-level geometry is reconstructed as an emergent slice of
        the triple-layer geometry determined by \GZTLC.
\end{itemize}

From the point of view of the GaoZheng Generalized Mathematical
Structure \GZGMS, the triple-layer connection \GZTLC\ is therefore the
principal device that couples $\mathsf{Obj}$ and $\mathsf{Law}$,
organises admissible flows in $\mathsf{Flow}$, and carries the
curvature information (\GZLCurv) that underlies CH-layering and the
various “Gap” statements developed later in this monograph.



\section[From \GZLStxt\ to \GZTLCtxt: Coupling Law- and Geometric Structures]{From \GZLS\ (\GZLStxt) to \GZTLC\ (\GZTLCtxt):Coupling Law-Level and Geometric Structures}


The triple-layer connection \GZTLC\ is not an arbitrary enlargement of a
classical connection: it is designed to be \emph{fed} by, and
\emph{compatible} with, the GaoZheng Law-Space \GZLS\ and its law-level
connection \GZLOC. This section spells out, at a schematic but precise
level, how law-objects in \GZLS\ and the law-connection \GZLOC\ give
rise to triple-layer connections on principal bundles, and conversely
how object-level geometric data are encoded back into \GZLS\ along
\GZTLC.



\subsubsection*{(a) Law-bundles over geometric backgrounds}

Fix a principal bundle $(\Bundle,\Base,\Group)$ with connection $\Conn$
and curvature $\Curv$ as in Section~2.1, together with a parameter
manifold $\Lambda_{\mathrm{par}}$. The first step is to place the
GaoZheng Law-Space \GZLS\ over this geometric background.

\begin{definition}[Law-bundle over $(\Base,\Lambda_{\mathrm{par}})$ (schematic)]
A \emph{law-bundle} over the product base
  $\Base\times\Lambda_{\mathrm{par}}$ is a fiber bundle
  \[
      \pi_{\mathrm{law}}:\mathcal{P}_{\mathrm{law}}
      \longrightarrow \Base\times\Lambda_{\mathrm{par}}
  \]
  with typical fiber a submanifold $\mathfrak{L}_0\subset\mathfrak{L}$
  of the law-space, together with a structure group
  $G_{\mathrm{law}}$ acting smoothly on $\mathfrak{L}_0$ by law-space
  symmetries. Locally one has
  \[
      \mathcal{P}_{\mathrm{law}}
      \cong
      \bigl(U\times V\bigr)\times \mathfrak{L}_0,
  \]
  for open sets $U\subset\Base$, $V\subset\Lambda_{\mathrm{par}}$.
\end{definition}

A section
\[
   s_{\mathrm{law}}:
      \Base\times\Lambda_{\mathrm{par}}
      \longrightarrow \mathcal{P}_{\mathrm{law}}
\]
assigns, to each geometric point $x\in\Base$ and parameter choice
$\lambda\in\Lambda_{\mathrm{par}}$, a law-object
\[
   L_{x,\lambda}
   \;\in\;
   \mathfrak{L}_0\subset\mathfrak{L},
\]
to be interpreted as the “local law package” governing admissible
dynamics and couplings at $(x,\lambda)$.



\begin{definition}[Law-connection on the law-bundle]
A \emph{law-connection} on $\mathcal{P}_{\mathrm{law}}$ is a principal
  connection
  \[
     \GZLOC\;\in\;\Omega^1(\mathcal{P}_{\mathrm{law}};\mathfrak{g}_{\mathrm{law}})
  \]
  for a chosen principal $G_{\mathrm{law}}$-bundle structure on
  $\mathcal{P}_{\mathrm{law}}$, whose curvature
  \[
     \GZLCurv
     \;=\;
     \mathrm{d}\GZLOC
     + \tfrac12[\GZLOC,\GZLOC]
     \;\in\;
     \Omega^2(\mathcal{P}_{\mathrm{law}};\mathfrak{g}_{\mathrm{law}})
  \]
  realises the law-curvature family associated with \GZLS\ (in the sense
  of Section~1.2).
\end{definition}

Pulling back $\GZLOC$ along the law-section $s_{\mathrm{law}}$ we obtain
a $1$-form on $\Base\times\Lambda_{\mathrm{par}}$:
\[
   A_{M,\Lambda}
   \;:=\;
   s_{\mathrm{law}}^*(\GZLOC)
   \;\in\;
   \Omega^1\bigl(\Base\times\Lambda_{\mathrm{par}};
                 \mathfrak{g}_{\mathrm{law}}\bigr),
\]
which describes how local law-objects $L_{x,\lambda}$ vary when we move
in space $x$ and in parameter directions $\lambda$.



\subsubsection*{(b) Law-driven families of geometric connections}

The second step is to explain how law-objects in \GZLS\ and the
law-connection \GZLOC\ constrain or generate the geometric connection
on $(\Bundle,\Base,\Group)$. For this we introduce a schematic
“law–geometry coupling”.

\begin{definition}[Law–geometry coupling (schematic)]
A \emph{law–geometry coupling} is a rule
  \[
     \mathcal{C}:
       \mathcal{P}_{\mathrm{law}}
       \longrightarrow
       \mathsf{Conn}(\Bundle)
  \]
  which assigns to each law-state $p\in\mathcal{P}_{\mathrm{law}}$ a
  connection
  \[
     \Conn_p\in\Omega^1(\Bundle;\mathfrak{g}),
  \]
  subject to the following schematic conditions.
  \begin{enumerate}[label=\roman*)]
    \item (\emph{Equivariance.})
          For $g_{\mathrm{law}}\in G_{\mathrm{law}}$ acting on
          $\mathcal{P}_{\mathrm{law}}$ and inducing a symmetry on
          $\mathsf{Conn}(\Bundle)$ (e.g.\ via gauge transformations),
          one has
          \[
             \mathcal{C}(p\cdot g_{\mathrm{law}})
             = \mathcal{C}(p)\cdot g_{\mathrm{law}}.
          \]
    \item (\emph{Locality.})
          If two law-states $p,p'$ project to the same
          $(x,\lambda)\in\Base\times\Lambda_{\mathrm{par}}$ and are
          connected by a law-horizontal curve with respect to \GZLOC,
          then $\Conn_p$ and $\Conn_{p'}$ differ only by geometric gauge
          transformations supported in a neighbourhood of $x$.
    \item (\emph{Regularity.})
          The map $\mathcal{C}$ is smooth and compatible with the
          parameter topology on $\Lambda_{\mathrm{par}}$.
  \end{enumerate}
\end{definition}

Composing the law-section $s_{\mathrm{law}}$ with $\mathcal{C}$, we
obtain a family of geometric connections indexed by
$(x,\lambda)$:
\[
   \Conn_{x,\lambda}
   \;:=\;
   \mathcal{C}\bigl(s_{\mathrm{law}}(x,\lambda)\bigr)
   \;\in\;
   \Omega^1(\Bundle;\mathfrak{g}).
\]
In particular, for a fixed parameter $\lambda_0$ one gets a connection
field
\[
   \Conn^{(\lambda_0)}:\Base\longrightarrow\mathsf{Conn}(\Bundle),
\]
whose variation along $\Base$ is constrained by the pull-back law-connection
$A_{M,\Lambda}$.



\begin{remark}[Static vs.\ dynamic viewpoint]
In Part~I, the functor
  \[
     \Phi_{\mathrm{geom}}:\mathsf{GeomObj}\longrightarrow\GZLS
  \]
  encodes how a \emph{given} geometric configuration is read as a
  law-object. The coupling $\mathcal{C}$ plays the complementary role:
  it says how, once law-objects are specified and allowed to flow via
  \GZLOC, one reconstructs or updates geometric connections. The
  triple-layer connection \GZTLC\ will interpolate between these two
  viewpoints.
\end{remark}



\subsubsection*{(c) From \GZLS\ and \GZLOC\ to a triple-layer connection}

We now explain how the data
\[
  \bigl(\GZLS,\GZLOC,\mathcal{P}_{\mathrm{law}},s_{\mathrm{law}},\mathcal{C}\bigr)
\]
give rise to a triple-layer principal bundle
$\widehat{\Bundle}\to X_{\mathrm{tot}}$ equipped with a triple-layer
connection \GZTLC\ as defined in Section~2.1.

Consider the fiber product
\[
   \widehat{\Bundle}
   \;:=\;
   \Bundle
   \times_{\Base}
   \bigl(\mathcal{P}_{\mathrm{law}}\times_{\Base}\Lambda_{\mathrm{par}}\bigr),
\]
whose points can be written schematically as triples
\[
   (u,p,\lambda),
   \qquad
   u\in\Bundle,\;
   p\in\mathcal{P}_{\mathrm{law}},\;
   \lambda\in\Lambda_{\mathrm{par}},
\]
with $\pi(u)=\pi_{\mathrm{law}}(p)=(x,\lambda)$.
Endow $\widehat{\Bundle}$ with the product structure group
\[
   \widehat{G}
   =
   G_{\mathrm{geom}}\times G_{\mathrm{law}}\times G_{\mathrm{par}},
\]
where $G_{\mathrm{geom}}=\Group$ and $G_{\mathrm{par}}$ encodes any
parameter symmetries one wishes to track.

\begin{proposition}[Induced triple-layer connection (schematic)]
There exists a unique triple-layer connection
  \[
     \GZTLC
     \;\in\;
     \Omega^1(\widehat{\Bundle};\widehat{\mathfrak{g}})
  \]
  whose layer-wise components satisfy:
  \begin{enumerate}[label=\roman*)]
    \item (\emph{Geometric component.})
          For fixed $(p,\lambda)$, the restriction of
          $\GZTLC$ to the slice
          $\{(u,p,\lambda)\mid u\in\Bundle\}$ projects to the
          geometric connection $\Conn_p$ produced by the coupling
          $\mathcal{C}$:
          \[
             \pi_{\mathrm{geom}}\bigl(\GZTLC\bigr)
             \big|_{\{(u,p,\lambda)\}}
             = \Conn_p.
          \]
    \item (\emph{Parameter component.})
          The parameter projection
          $\pi_{\mathrm{par}}(\GZTLC)$ is basic with respect to
          $\widehat{\Bundle}\to\Lambda_{\mathrm{par}}$ and coincides
          with a chosen parameter connection $B_{\mathrm{par}}$ on
          the parameter layer.
    \item (\emph{Law component.})
          The law projection $\pi_{\mathrm{law}}(\GZTLC)$ is the pull-back
          of the law-connection \GZLOC\ along the natural projection
          $\widehat{\Bundle}\to\mathcal{P}_{\mathrm{law}}$:
          \[
             \pi_{\mathrm{law}}(\GZTLC)
             =
             \widehat{\rho}^*(\GZLOC),
          \]
          where $\widehat{\rho}(u,p,\lambda)=p$.
  \end{enumerate}
  In particular, \GZTLC\ extends simultaneously the geometric
  connection $\Conn_p$, the parameter connection $B_{\mathrm{par}}$,
  and the law-connection \GZLOC.
  A complete construction and proof are given in Technical Appendix~I
  (Section~\ref{app:I-P3}).
\end{proposition}

\begin{proof}[Proof sketch]
One constructs \GZTLC\ by specifying its horizontal distribution
  $H_{\mathcal{A}}\subset T\widehat{\Bundle}$ as follows.

  \begin{itemize}
    \item Vectors tangent to $\Bundle$ are declared horizontal if and
          only if they are horizontal for the connection $\Conn_p$,
          with $p$ fixed.
    \item Vectors tangent to $\mathcal{P}_{\mathrm{law}}$ are declared
          horizontal if and only if they are horizontal for \GZLOC.
    \item Vectors tangent to $\Lambda_{\mathrm{par}}$ are lifted
          according to $B_{\mathrm{par}}$.
  \end{itemize}

  One then checks that this distribution is $\widehat{G}$-equivariant
  and complementary to vertical spaces, hence defines a unique
  principal connection whose projections are the desired ones. The
  coupling conditions on $\mathcal{C}$ ensure compatibility along
  overlaps and under gauge transformations; details are routine and
  omitted.
\end{proof}



\subsubsection*{(d) Compatibility with the law-functor \texorpdfstring{$\Phi_{\mathrm{geom}}$}{Phi\_geom}}

The construction in Section~1.3 provided a functor
\[
   \Phi_{\mathrm{geom}}:
      \mathsf{GeomObj}\longrightarrow\GZLS,
\]
which associates to geometric objects (bundles with connection and
curvature, homotopy-enhanced versions, etc.) a corresponding
law-object in \GZLS. The triple-layer connection \GZTLC\ couples this
static encoding with the dynamic law-connection \GZLOC\ via the
following schematic compatibility condition.

\begin{definition}[GZ-LS–compatible triple-layer connection]
A triple-layer connection \GZTLC\ on $\widehat{\Bundle}$ is said to be
  \emph{GZ-LS–compatible} if, for any sufficiently regular flow
  \[
     \gamma(t)
       = (x(t),\lambda(t),L(t))
       \in \Base\times\Lambda_{\mathrm{par}}\times\mathfrak{L}_0
  \]
  whose law component $L(t)$ is obtained by \GZLOC-horizontal transport
  of $L_{x(0),\lambda(0)}$ along $(x(t),\lambda(t))$, the following
  holds:
  \begin{enumerate}[label=\roman*)]
    \item the \GZTLC-horizontal lift $\widehat{\gamma}(t)$ projects on
          the geometric side to a family of connections
          $\Conn_{x(t),\lambda(t)}$ that satisfy
          \[
             \Phi_{\mathrm{geom}}\bigl(\Conn_{x(t),\lambda(t)}\bigr)
             = L(t);
          \]
    \item conversely, any \GZTLC-horizontal geometric flow
          $(\Conn_{x(t),\lambda(t)})$ induces a law-flow $L(t)$ that is
          horizontal for \GZLOC\ and satisfies the same relation.
  \end{enumerate}
\end{definition}

In words, moving horizontally in the triple-layer bundle corresponds to
moving simultaneously in geometric, parameter, and law directions in
such a way that the static encoding
$\Phi_{\mathrm{geom}}(\text{geometry}) = \text{law}$ is preserved along
the flow. This is how \GZLS\ and \GZLOC\ “drive” the evolution of
geometric connections via \GZTLC, and how geometric changes are
consistently reinterpreted at the law-level.



\subsubsection*{(e) Summary: from \GZLS\ to \GZTLC\ as a law–geometry interface}

The map “from \GZLS\ (\GZLStxt) to \GZTLC\ (\GZTLCtxt)” can now be
summarised structurally:

\begin{itemize}
  \item Starting from the GaoZheng Law-Space
        $\GZLS=(\mathfrak{L};J,\mathrm{Loc},\Sigma)$ and its
        law-connection \GZLOC, one constructs a law-bundle
        $\mathcal{P}_{\mathrm{law}}$ over $(\Base,\Lambda_{\mathrm{par}})$
        and a law-section $s_{\mathrm{law}}$ that assigns local law
        packages $L_{x,\lambda}$.
  \item A law–geometry coupling $\mathcal{C}$ translates law-states in
        $\mathcal{P}_{\mathrm{law}}$ into geometric connections on
        $(\Bundle,\Base,\Group)$, producing families
        $\Conn_{x,\lambda}$ compatible with law-horizontal motion.
  \item These ingredients assemble into a triple-layer principal bundle
        $\widehat{\Bundle}\to X_{\mathrm{tot}}$ with a triple-layer
        connection \GZTLC\ whose law component is the pull-back of
        \GZLOC\ and whose geometric component is given by $\Conn_{x,\lambda}$.
  \item GZ-LS–compatibility ensures that the static encoding
        $\Phi_{\mathrm{geom}}$ is preserved along horizontal flows,
        thereby turning \GZTLC\ into a genuine interface between
        object-level geometry and law-level structure.
\end{itemize}

From the vantage point of the PFB-GNLA constructions, this interface is
what allows one to regard strict, homotopy, variable functional–operator
and meta-mathematical original versions of \GAlgebra\ as different
\emph{choices} of law–geometry coupling inside a common GZ-LS/\GZTLC\
background. The next section will make this link explicit by introducing
“Gap statements” and certificate-type invariants that are computed from
law-curvature \GZLCurv\ but constrain the existence and shape of
geometric realizations inside the GaoZheng G-Framework.



\section{“Gap” Statements and Certificate-Type Invariants}


The triple-layer connection \GZTLC\ provides a unified geometric object
in which curvature components in the geometric, parameter, and law
layers are organised together. The goal of this section is to explain
how one formulates \emph{“Gap” statements} in this setting, and how
\emph{certificate-type invariants} extract finite, auditable information
from the law-curvature \GZLCurv\ and its couplings, in a way compatible
with the PFB-GNLA constructions and the PL-PI meta-paradigm.



\subsubsection*{(a) Intuitive picture: gaps as protected non-deformability}

At a schematic level, a “Gap statement” in the GaoZheng G-Framework
expresses the fact that certain desirable configurations or behaviours
cannot be reached by a continuous deformation that stays within a
given class, without crossing a non-trivial region of law-curvature or
homotopy obstruction.

\begin{itemize}
  \item In the \emph{strict} PFB-GNLA setting, the target configuration
        might be “a strictly Jacobi-satisfying Lie-type structure” or
        “a flat law-connection along a given family of flows”.
  \item In the \emph{homotopy} PFB-GNLA setting, one allows higher
        brackets and homotopies, but still asks that they remain
        confined below certain quantitative thresholds.
  \item In the \emph{variable functional–operator} setting, the target
        is a class of operator-path dynamics with bounded law-curvature
        along triple-layer flows.
\end{itemize}

A Gap statement then says, schematically:

\begin{quote}
  “If one tries to deform from region $\mathcal{R}_{\mathrm{in}}$
  (e.g.\ a given class of geometric configurations) to region
  $\mathcal{R}_{\mathrm{out}}$ (e.g.\ a strictly flat law-regime),
  along flows that are horizontally admissible for \GZTLC, then
  \emph{necessarily} the law-curvature (or certain of its invariants)
  must cross a positive threshold; in particular, there is no
  deformation that stays below that threshold everywhere.”
\end{quote}

This is analogous to a spectral gap, an isoperimetric inequality, or a
topological obstruction, but formulated at the level of law-space flows
and triple-layer curvature.



\subsubsection*{(b) Schematic notion of a certificate-type invariant}

To make such statements usable—especially for certificate-based AI and for
complex-systems modelling via total-operator principal bundles (developed in
separate applied volumes)—the obstruction must be encoded in a
  \emph{finite} family of quantities that are:

\begin{itemize}
  \item gauge-invariant (or invariant up to a controlled equivalence);
  \item stable under small perturbations and local truncations;
  \item computable, or at least approximable, from local data along
        finite pieces of triple-layer flows.
\end{itemize}

This motivates the following schematic definition.

\begin{definition}[Certificate-type invariant for a triple-layer system]
Let $\mathsf{Sys}_{\mathrm{TLC}}$ be a class of triple-layer systems
  \[
     S =
     (X_{\mathrm{tot}},
      \widehat{\Bundle},\GZTLC;
      \GZLS,\GZLOC,\mathsf{Layer}),
  \]
  where $\GZTLC$ is a GZ-LS–compatible triple-layer connection as in
  Sections~2.1–2.2, and $\mathsf{Layer}$ encodes CH-type semantic
  layering. A \emph{certificate-type invariant} on
  $\mathsf{Sys}_{\mathrm{TLC}}$ is a map
  \[
     \mathsf{Cert}:
       \mathsf{Sys}_{\mathrm{TLC}}
       \longrightarrow
       V,
  \]
  where $V$ is a finite-dimensional normed space, such that:
  \begin{enumerate}[label=\roman*)]
    \item (\emph{Gauge invariance.})
          If $S$ and $S'$ are related by a triple-layer gauge
          transformation (combining geometric, parameter, and law
          symmetries), then $\mathsf{Cert}(S)=\mathsf{Cert}(S')$.
    \item (\emph{Local computability.})
          For any bounded region $U\subset X_{\mathrm{tot}}$,
          $\mathsf{Cert}(S)|_U$ depends only on the restriction of
          $\mathcal{F}_{\GZTLC}$ and CH-layering data to $U$, and
          can be obtained (up to a prescribed precision) from a finite
          sampling of flows and curvature along $U$.
    \item (\emph{Stability under homotopy.})
          If $S_t$ is a triple-layer homotopy
          $t\in[0,1]\mapsto S_t$ with uniformly bounded
          curvature and controlled CH-layer transitions, then
          $t\mapsto\mathsf{Cert}(S_t)$ is continuous and piecewise
          smooth.
  \end{enumerate}
  The components of $\mathsf{Cert}(S)$ are called \emph{certificate
  components} for $S$.
\end{definition}

In later parts of the monograph, concrete examples of such invariants
will be constructed, including Jacobiator-gap certificates for homotopy
versions of \GAlgebra\ and curvature-based certificates for
law-operator dynamics.



\subsubsection*{(c) “Gap” statements in terms of certificates}

Once certificate-type invariants are available, “Gap” statements can be
formulated as inequalities or non-vanishing conditions in the target
space $V$.



\begin{definition}[Gap statement driven by a certificate (schematic)]
Let $\mathsf{Cert}:\mathsf{Sys}_{\mathrm{TLC}}\to V$ be a
  certificate-type invariant, and let $\mathcal{C}_{\mathrm{in}}$,
  $\mathcal{C}_{\mathrm{out}}\subset\mathsf{Sys}_{\mathrm{TLC}}$ be two
  classes of systems (e.g.\ strict vs.\ homotopy PFB-GNLA regimes). A
  \emph{Gap statement driven by $\mathsf{Cert}$} is a theorem of the
  following schematic form:

  \begin{enumerate}[label=\roman*)]
    \item (\emph{Vanishing side.})
          If $S\in\mathcal{C}_{\mathrm{in}}$ and
          $\mathsf{Cert}(S)=0$, then $S$ admits a refinement or
          extension in $\mathcal{C}_{\mathrm{out}}$ with additional
          desirable properties (e.g.\ strict Jacobi identity, flat
          law-curvature along certain flows).
    \item (\emph{Positive side.})
          There exists $\varepsilon>0$ such that if a system $S$ is
          connected by a triple-layer homotopy to a system
          $S'\in\mathcal{C}_{\mathrm{out}}$ while remaining inside
          $\mathcal{C}_{\mathrm{in}}$, then along this homotopy one
          must cross the region
          \[
             \|\mathsf{Cert}(S_t)\|\;\ge\;\varepsilon,
          \]
          for some $t\in[0,1]$. Equivalently, there is no homotopy
          entirely contained in $\{\|\mathsf{Cert}\|<\varepsilon\}$
          connecting $\mathcal{C}_{\mathrm{in}}$ to
          $\mathcal{C}_{\mathrm{out}}$.
  \end{enumerate}
\end{definition}

The parameter $\varepsilon$ plays the role of a quantitative “gap”:
even after allowing for homotopy relaxations and triple-layer flows,
  one cannot bring $\mathsf{Cert}$ arbitrarily close to zero while
  connecting the two classes. When $\varepsilon$ is strictly positive,
  this expresses a genuine obstruction; when the best possible
  $\varepsilon$ is zero but the vanishing side remains non-trivial, one
  obtains a \emph{threshold-type} gap rather than a strict spectral gap.


\subsubsection*{(d) Law-curvature and Jacobiator-gap as basic examples}

Among the certificate-type invariants naturally suggested by \GZTLC,
two families will play a central role in later parts of the monograph.

\paragraph{(1) Law-curvature derived certificates.}

Given the law-curvature
\[
   \GZLCurv
   = \pi_{\mathrm{law}}\bigl(\mathcal{F}_{\GZTLC}\bigr),
\]
  one may construct local functionals of the form
\[
   C_{\mathrm{law}}(U)
   \;:=\;
   \Bigl(
     \int_U \Phi_1(\GZLCurv),\dots,
     \int_U \Phi_k(\GZLCurv)
   \Bigr),
\]
for suitable polynomial or trace-type expressions $\Phi_i$ and regions
$U\subset X_{\mathrm{tot}}$. By assembling these into a finite vector
(after CH-layering and symmetry reduction), one obtains a certificate
  \[
     \mathsf{Cert}_{\mathrm{curv}}(S)\in\mathbb{R}^k,
  \]
  measuring how far the law-layer deviates from flatness along relevant
  flows. Gap statements driven by $\mathsf{Cert}_{\mathrm{curv}}$ express
  that certain law-level unifications or decouplings are impossible
  without crossing a minimal “curvature budget”.


\paragraph{(2) Jacobiator-gap certificates.}

In the homotopy versions of \GAlgebra, the failure of the strict Jacobi
identity is encoded in higher brackets and Jacobiators. As anticipated
in Part~III, one can package these into a Jacobiator-gap certificate
  \[
     \mathsf{JacGap}(L)
     = \bigl(c_1(L),\dots,c_k(L)\bigr),
  \]
  for law-level operator systems $L$, with the properties:

\begin{itemize}
  \item $c_i(L)$ quantify controlled failures of the strict Jacobi
        identity;
  \item $\mathsf{JacGap}(L)=0$ if and only if $L$ can be promoted to a
        strictly Jacobi-satisfying law-object without changing its
        low-order observables;
	  \item $\mathsf{JacGap}$ is stable under \GZTLC-induced deformations
	        in law-space.
  \end{itemize}


When such law-level systems $L$ are realised as law-slices of a
triple-layer system $S$, the assignment
\[
   S \longmapsto \mathsf{JacGap}\bigl(L(S)\bigr)
\]
provides another certificate-type invariant on
$\mathsf{Sys}_{\mathrm{TLC}}$. The associated Gap statements say, in
essence, that one cannot deform certain homotopy PFB-GNLA structures
into strict ones along triple-layer flows without the Jacobiator-gap
crossing a quantitative threshold. This is the geometric content
behind the “GZ-OHU” picture introduced later in Part~III.



\subsubsection*{(e) Interface with applications and certificate-based AI}

From the viewpoint of applications (physics, AI, life and complex systems), the role of
certificate-type invariants is twofold:

\begin{itemize}
  \item \emph{Theoretical side.}
        They make precise, at the level of \GZTLC\ and \GZLS, the
        intuition that “some failures of ideal behaviour” are not
        artefacts of poor modelling but are structurally protected by
        law-curvature and homotopy; the Gap statements identify
        precisely which failures are unavoidable in a given regime.
  \item \emph{Engineering side.}
        They provide finite, auditable quantities that can be monitored
        along learning or control procedures (e.g.\ in law-space
        learning or in complex-systems models based on
        total-operator bundles) to guarantee
        that the system stays in certified-safe regions, or to detect
        when a gap-protected transition is being approached.
\end{itemize}

In this sense, “certificate-type invariants” are the bridge between the
triple-layer curvature geometry of \GZTLC\ and the concrete practice of
certificate-based AI and law-space learning. Subsequent chapters will
instantiate this general scheme in the strict, homotopy, variable
functional–operator, and meta-mathematical original versions of
PFB-GNLA, and will state explicit Gap theorems expressing how
$\mathsf{Cert}_{\mathrm{curv}}$, $\mathsf{JacGap}$ and related
invariants bound what is geometrically and law-level possible inside
the GaoZheng G-Framework.







\part[Principal Bundles and the GaoZheng G-Algebra]{Principal Bundles and the GaoZheng G-Algebra:\\
      Four Integrated PFB-GNLA Constructions}


\chapter[Principal Bundles and GNLA: Preliminaries]{Principal Bundles, Connections, and Generalized Noncommutative\\
         Lie Algebras: Preliminaries}


\section[Principal Bundles and Connections/Curvature]{Principal Bundles \texorpdfstring{$(\Bundle,\Base,\Group)$}{(P,M,G)}and Connections/Curvature}


This section fixes the classical geometric background for the PFB-GNLA
constructions. We recall principal bundles, their connections and
curvatures, and the associated notions of horizontal lift, holonomy,
and covariant derivative. Throughout we use the notation
\[
  (\Bundle,\Base,\Group),\qquad
  \pi:\Bundle\to\Base,\qquad
  \mathfrak{g}=\mathrm{Lie}(\Group),
\]
with \Bundle\ and \Base\ smooth manifolds and \Group\ a Lie group.



\subsubsection*{(a) Principal bundles \texorpdfstring{$(\Bundle,\Base,\Group)$}{(P,M,G)}}

\begin{definition}[Principal \Group-bundle]
A \emph{principal $\Group$-bundle} over a smooth manifold $\Base$ is
  a smooth manifold $\Bundle$ together with
  \begin{itemize}
    \item a smooth surjective submersion
          $\pi:\Bundle\to\Base$;
    \item a smooth right action
          $R:\Bundle\times\Group\to\Bundle$,
          $(u,g)\mapsto u\cdot g$,
          that is free and transitive on the fibers,
  \end{itemize}
  such that for every $x\in\Base$ there exists an open neighbourhood
  $U\subset\Base$ and a $\Group$-equivariant diffeomorphism
  (local trivialisation)
  \[
     \phi_U:\pi^{-1}(U)\longrightarrow U\times\Group
  \]
  making the diagram
  \[
    \begin{array}{ccc}
       \pi^{-1}(U) & \xrightarrow{\ \phi_U\ }
                   & U\times\Group \\
       \pi \downarrow\quad & &
            \quad\downarrow\mathrm{pr}_1 \\
       U & \xrightarrow{\ \mathrm{id}\ }
         & U
    \end{array}
  \]
  commute and intertwining the right action with right multiplication
  on the $\Group$ factor.
\end{definition}

For each $x\in\Base$ the fiber $\pi^{-1}(x)$ is a $\Group$-torsor:
there is no distinguished “origin”, but the group acts simply
transitively on it. The tangent bundle $T\Bundle$ splits at each point
into vertical and horizontal directions once a connection is chosen; for
now we record the vertical part only.



\begin{definition}[Vertical bundle and fundamental vector fields]
The \emph{vertical bundle} $V\Bundle\subset T\Bundle$ is the kernel
  of the differential $\mathrm{d}\pi:T\Bundle\to T\Base$:
  \[
     V_u\Bundle
     :=
     \ker\bigl(\mathrm{d}\pi_u:T_u\Bundle\to T_{\pi(u)}\Base\bigr).
  \]
  For each $\xi\in\mathfrak{g}$, the \emph{fundamental vector field}
  $\xi_\Bundle\in\mathfrak{X}(\Bundle)$ is defined by
  \[
    \xi_\Bundle(u)
    :=
    \left.\frac{\mathrm{d}}{\mathrm{d}t}\right|_{t=0}
    \bigl(u\cdot\exp(t\xi)\bigr).
  \]
  The image of $\mathfrak{g}$ under $\xi\mapsto\xi_\Bundle(u)$ spans
  $V_u\Bundle$.
\end{definition}



\subsubsection*{(b) Ehresmann connection and connection 1-form}

There are two equivalent ways to describe a connection on a principal
bundle: as an Ehresmann horizontal distribution, or as a
$\mathfrak{g}$-valued 1-form satisfying certain properties. We recall
both and fix the notation \Conn\ for the connection 1-form.



\begin{definition}[Ehresmann connection]
An \emph{Ehresmann connection} on $(\Bundle,\Base,\Group)$ is a
  smooth subbundle $H\Bundle\subset T\Bundle$ (the horizontal bundle)
  such that
  \begin{itemize}
    \item $T\Bundle = H\Bundle \oplus V\Bundle$;
    \item $H\Bundle$ is $\Group$-equivariant:
          $(R_g)_*H_u\Bundle = H_{u\cdot g}\Bundle$ for all $u\in\Bundle$,
          $g\in\Group$.
  \end{itemize}
\end{definition}

Given such a horizontal distribution, every tangent vector
$v\in T_x\Base$ admits a unique horizontal lift
$\widetilde{v}\in H_u\Bundle$ at each $u\in\pi^{-1}(x)$ with
$\mathrm{d}\pi_u(\widetilde{v})=v$. In particular, a curve
$x(t)\in\Base$ can be lifted to a horizontal curve $u(t)\in\Bundle$
such that $\pi(u(t))=x(t)$ and $\dot{u}(t)\in H_{u(t)}\Bundle$.



\begin{definition}[Connection 1-form]
A \emph{connection 1-form} on $(\Bundle,\Base,\Group)$ is a
  $\mathfrak{g}$-valued 1-form
  \[
     \Conn\in\Omega^1(\Bundle;\mathfrak{g})
  \]
  such that:
  \begin{enumerate}[label=\roman*)]
    \item (\emph{Reproducing property.})
          For all $\xi\in\mathfrak{g}$,
          \[
             \Conn(\xi_\Bundle)=\xi.
          \]
    \item (\emph{$\Group$-equivariance.})
          For all $g\in\Group$,
          \[
             (R_g)^*\Conn
             = \mathrm{Ad}_{g^{-1}}\circ\Conn.
          \]
  \end{enumerate}
\end{definition}

The horizontal distribution associated to \Conn\ is
\[
   H_u\Bundle
   := \ker\bigl(\Conn_u:T_u\Bundle\to\mathfrak{g}\bigr),
\]
and conversely, \Conn\ is characterised by declaring vertical vectors
to be mapped to their $\mathfrak{g}$-labels and horizontal vectors to
be mapped to $0$. Thus Ehresmann connections and connection 1-forms
are equivalent data; we shall freely pass from one description to the
other.



\subsubsection*{(c) Curvature and the Bianchi identity}

The curvature of a connection measures the failure of horizontal transport to be \\
path-independent. In terms of the connection 1-form
\Conn, it is given by the familiar Cartan formula.



\begin{definition}[Curvature 2-form]
The \emph{curvature} of a connection 1-form $\Conn$ is the
  $\mathfrak{g}$-valued 2-form
  \[
     \Curv
     :=
     \mathrm{d}\Conn
     + \tfrac12[\Conn,\Conn]
     \;\in\;
     \Omega^2(\Bundle;\mathfrak{g}),
  \]
  where $[\Conn,\Conn]$ is the wedge–bracket defined by
  \[
     [\Conn,\Conn](X,Y)
     :=
     [\Conn(X),\Conn(Y)],
     \qquad X,Y\in T_u\Bundle.
  \]
\end{definition}

The curvature \Curv\ is horizontal and $\mathrm{Ad}$-equivariant:
\[
   \Curv(\xi_\Bundle,\cdot)=0,
   \qquad
   (R_g)^*\Curv
   = \mathrm{Ad}_{g^{-1}}\circ\Curv.
\]
In particular, it descends to a well-defined 2-form with values in the
adjoint bundle $\mathrm{Ad}(\Bundle)=\Bundle\times_{\mathrm{Ad}}\mathfrak{g}$
over \Base.



\begin{definition}[Covariant exterior derivative and Bianchi identity]
The \emph{covariant exterior derivative} associated with \Conn\ is
  the operator
  \[
     D_\Conn:
       \Omega^k(\Bundle;\mathfrak{g})
       \longrightarrow
       \Omega^{k+1}(\Bundle;\mathfrak{g}),
  \]
  defined by
  \[
     D_\Conn\eta
     :=
     \mathrm{d}\eta
     + [\Conn,\eta],
  \]
  where $[\Conn,\eta]$ uses the same wedge–bracket convention. The
  curvature 2-form satisfies the \emph{Bianchi identity}
  \[
     D_\Conn\Curv = 0.
  \]
\end{definition}



\subsubsection*{(d) Parallel transport, holonomy, and path-ordered exponentials}

Given a connection, one can define parallel transport along curves in
\Base\ and extract from it holonomy information. This is the first
place where the path-space viewpoint, and ultimately the PFB-GNLA
operator picture, emerge.



Let $x:[0,1]\to\Base$ be a smooth curve and $u_0\in\pi^{-1}(x(0))$ a
starting point in the fiber. A \emph{horizontal lift} of $x$ through
$u_0$ is a curve $u:[0,1]\to\Bundle$ such that
\[
   \pi(u(t)) = x(t),\qquad
   \dot{u}(t)\in H_{u(t)}\Bundle\ \text{for all }t.
\]
Given \Conn, such a lift exists and is unique for any initial condition
$u(0)=u_0$.

\begin{definition}[Parallel transport and holonomy]
The \emph{parallel transport} along $x$ is the map
  \[
     \mathsf{PT}_x:
       \pi^{-1}(x(0))\longrightarrow\pi^{-1}(x(1)),
     \qquad
     u_0\longmapsto u(1),
  \]
  where $u$ is the horizontal lift through $u_0$. When
  $x(0)=x(1)=x_0$, the resulting map
  \[
     \mathsf{Hol}_x:
       \pi^{-1}(x_0)\longrightarrow\pi^{-1}(x_0)
  \]
  is a fiberwise $\Group$-equivariant diffeomorphism and can be
  identified with right multiplication by an element
  $\mathrm{Hol}_\Conn(x)\in\Group$, called the \emph{holonomy} of
  $\Conn$ along $x$.
\end{definition}

At a formal level, if one chooses a local trivialisation and expresses
\Conn\ as a $\mathfrak{g}$-valued 1-form $A$ on \Base, parallel
transport along $x$ corresponds to the path-ordered exponential
\[
   \mathrm{Hol}_\Conn(x)
   =
   \mathcal{P}\exp\left(
     -\int_0^1 A\bigl(\dot{x}(t)\bigr)\,\mathrm{d}t
   \right),
\]
where $\mathcal{P}$ denotes the time-ordering operator. This expression
will later be generalised in the variable functional–operator version
of PFB-GNLA, where path-space operators and law-operator dynamics
replace finite-dimensional matrices.



\subsubsection*{(e) Associated bundles and covariant derivatives}

Given a principal bundle and a representation of its structure group,
one constructs associated vector bundles and defines covariant
derivatives on their sections.



\begin{definition}[Associated vector bundle]
Let $\rho:\Group\to\mathrm{GL}(V)$ be a finite-dimensional
  representation. The \emph{associated vector bundle} is
  \[
     E
     :=
     \Bundle\times_\rho V
     :=
     (\Bundle\times V)/\!\sim,
  \]
  where $(u,v)\sim(u\cdot g,\rho(g^{-1})v)$ for all $g\in\Group$. The
  projection $[u,v]\mapsto\pi(u)$ makes $E$ into a vector bundle over
  \Base.
\end{definition}

Any connection \Conn\ on $(\Bundle,\Base,\Group)$ induces a covariant
derivative on sections of $E$, denoted
\[
   \nabla:\Gamma(E)\longrightarrow\Gamma(T^*\Base\otimes E).
\]
In local trivialisations, \Conn\ is represented by a matrix-valued
1-form $A$ and $\nabla$ acts as
\[
   \nabla = \mathrm{d} + \rho_*(A),
\]
where $\rho_*:\mathfrak{g}\to\mathfrak{gl}(V)$ is the induced Lie
algebra representation. The curvature \Curv\ is encoded at the level
of $E$ by the curvature of $\nabla$,
\[
   R_\nabla = \nabla^2
            = \rho_*(\Curv),
\]
again viewed as an $\mathrm{End}(E)$-valued 2-form on \Base.



\subsubsection*{(f) Parameterised families of connections}

The strict PFB-GNLA construction, as well as its homotopy and variable
functional–operator extensions, work not with a single connection but
with \emph{families} of connections indexed by parameter spaces and/or
law-space coordinates. We record the basic notation.



Let $\Lambda_{\mathrm{par}}$ be a parameter manifold as in Part~II,
and consider a smooth map
\[
   \Lambda_{\mathrm{par}}
   \longrightarrow
   \mathsf{Conn}(\Bundle),
   \qquad
   \lambda\longmapsto\Conn_\lambda,
\]
where each $\Conn_\lambda$ is a connection 1-form on \Bundle. This can
be encoded as a single 1-form
\[
   \widehat{\Conn}
   \;\in\;
   \Omega^1(\Bundle\times\Lambda_{\mathrm{par}};\mathfrak{g}),
\]
whose restriction to each slice $\lambda=\mathrm{const}$ recovers
$\Conn_\lambda$ and whose $\Lambda_{\mathrm{par}}$-components encode
how $\Conn_\lambda$ varies with $\lambda$.

In the language of Section~2, such families are naturally organised by
the parameter component of the triple-layer connection \GZTLC:
\[
   B_{\mathrm{par}}
   \;\in\;
   \Omega^1(\widehat{\Bundle};\mathfrak{g}_{\mathrm{par}}),
\]
and by the pull-back of the law-connection \GZLOC. The strict PFB-GNLA
version corresponds roughly to the case where
$\widehat{\Conn}$ depends on $\lambda$ in a controlled way and the
law-curvature \GZLCurv\ remains in a “small” or “simple” regime; the
homotopy and variable functional–operator versions relax these
conditions.



\subsubsection*{(g) Position of classical geometry in the GaoZheng G-Framework}

From the vantage point of the GaoZheng G-Framework and the PL-PI
Meta-Mathematical Theory, the classical principal bundle geometry
summarised above plays the following role:

\begin{itemize}
  \item It supplies the \emph{object-level} input to the strict
        version of \GAlgebra: the Lie algebra $\mathfrak{g}$, the
        connection \Conn, the curvature \Curv, and associated bundles
        give the raw geometric data to be reinterpreted as law-objects
        in \GZLS\ via the functor $\Phi_{\mathrm{geom}}$.
  \item It provides the geometric substrate on which the triple-layer
        connection \GZTLC\ is built: \Conn\ is the geometric component
        of \GZTLC, while parameter and law components are added as in
        Section~2.
  \item It offers a familiar “strict” baseline against which homotopy
        enhancements, variable functional–operator dynamics, and
        meta-mathematical reconstructions can be measured; departures
        from the classical identities (e.g.\ Jacobi identity, Bianchi
        identity in adjoint form) will be quantified by Gap
        statements and certificate-type invariants.
\end{itemize}

The next section upgrades this classical setting to the algebraic
language that will be used in PFB-GNLA: generalized noncommutative
Lie algebras and their homotopy refinements, positioned explicitly
inside the GaoZheng Law-Space and the GaoZheng Generalized
Mathematical Structure.




\section{Generalized Noncommutative Lie Algebras and Homotopy Lie Structures}


This section upgrades the classical principal-bundle geometry of
Section~3.1 to the algebraic language that underlies the PFB-GNLA
constructions. We first recall ordinary Lie algebras and their geometric
realisation via connections, then introduce \emph{generalized
noncommutative Lie algebras} (GNLA) as a controlled relaxation of the
Jacobi identity, and finally position homotopy Lie structures
($L_\infty$-algebras) as the natural homotopy-complete framework in
which PFB-GNLA will be formulated.



\subsubsection*{(a) Classical Lie algebras and geometric origin}

\begin{definition}[Lie algebra]
A \emph{Lie algebra} over a field $\Bbbk$ is a vector space
  $\mathfrak{g}$ together with a bilinear map
  \[
     [\cdot,\cdot]:
       \mathfrak{g}\times\mathfrak{g}\longrightarrow\mathfrak{g},
  \]
  such that:
  \begin{enumerate}[label=\roman*)]
    \item (\emph{Antisymmetry})
          $[X,Y]=-[Y,X]$ for all $X,Y\in\mathfrak{g}$;
    \item (\emph{Jacobi identity})
          \[
             [X,[Y,Z]] + [Y,[Z,X]] + [Z,[X,Y]] = 0
             \quad
             \text{for all }X,Y,Z\in\mathfrak{g}.
          \]
  \end{enumerate}
\end{definition}

In the context of a principal bundle $(\Bundle,\Base,\Group)$, the Lie
algebra $\mathfrak{g}=\mathrm{Lie}(\Group)$ appears in three closely
related ways:

\begin{itemize}
  \item as the infinitesimal generator of vertical directions in
        $T\Bundle$, via fundamental vector fields $\xi_\Bundle$;
  \item as the target of the connection 1-form
        $\Conn\in\Omega^1(\Bundle;\mathfrak{g})$ and curvature
        $\Curv\in\Omega^2(\Bundle;\mathfrak{g})$;
  \item as the fiber of the adjoint bundle
        $\mathrm{Ad}(\Bundle)=\Bundle\times_{\mathrm{Ad}}\mathfrak{g}$,
        on whose sections the curvature and covariant derivatives act.
\end{itemize}

Thus the classical geometric data of Section~3.1 can be “read” as a
family of $\mathfrak{g}$-valued tensors satisfying the Lie algebra
identities and the Bianchi identity. The GaoZheng G-Algebra (\GAlgebra)
will generalise this picture by allowing controlled violations of
Jacobi and by encoding higher homotopy data.



\subsubsection*{(b) Generalized noncommutative Lie algebras (GNLA): schematic axioms}

In the PFB-GNLA setting, the algebraic object of interest is no longer
just $\mathfrak{g}$, but a space of \emph{operators} built from
connection, curvature, and their covariant derivatives. The Jacobi
identity does not necessarily hold strictly; instead, its failure is
measured and controlled.



\begin{definition}[Jacobiator]
Let $(A,[\cdot,\cdot])$ be a non-associative, not necessarily
  antisymmetric algebra over $\Bbbk$. The \emph{Jacobiator} is the
  trilinear map
  \[
     \Jac_A:
       A^{\otimes 3}\longrightarrow A
  \]
  defined by
  \[
     \Jac_A(x,y,z)
     :=
     [x,[y,z]] + [y,[z,x]] + [z,[x,y]],
     \qquad x,y,z\in A.
  \]
  When $\Jac_A\equiv 0$ and $[\cdot,\cdot]$ is antisymmetric, $A$
  is a Lie algebra.
\end{definition}



\begin{definition}[Generalized noncommutative Lie algebra (GNLA, schematic)]
A \emph{generalized noncommutative Lie algebra} (GNLA) is a tuple
  \[
     \mathsf{GNLA}
     =
     (A,[\cdot,\cdot],\Jac_A;\mathsf{Str}),
  \]
  where
  \begin{itemize}
    \item $A$ is a $\Bbbk$-vector space (possibly $\mathbb{Z}$-graded);
    \item $[\cdot,\cdot]:A\otimes A\to A$ is a bilinear bracket, not
          necessarily antisymmetric;
    \item $\Jac_A:A^{\otimes 3}\to A$ is the Jacobiator associated
          with $[\cdot,\cdot]$;
    \item $\mathsf{Str}$ is a collection of structural data (grading,
          filtration, distinguished subspaces, pairing with geometric
          objects, etc.),
  \end{itemize}
  subject to the following schematic conditions:
  \begin{enumerate}[label=\roman*)]
    \item (\emph{Controlled noncommutativity.})
          The antisymmetrisation
          \[
             [x,y]_- := \tfrac12\bigl([x,y]-[y,x]\bigr)
          \]
          defines a Lie bracket on a distinguished subspace
          $A_{\mathrm{skew}}\subseteq A$; symmetric parts and higher
          corrections are organised in $\mathsf{Str}$.
    \item (\emph{Controlled Jacobi failure.})
          The Jacobiator $\Jac_A$ takes values in a controlled
          subspace $J\subseteq A$ (e.g.\ a higher-degree or curvature
          sector) and satisfies identities of the form
          \[
             D\Jac_A + \text{(higher Jacobiators)} = 0,
          \]
          where $D$ is a suitable derivation (often a covariant
          differential).
    \item (\emph{Compatibility with geometry or law-space.})
          There exists a functorial assignment
          \[
             \Phi_{\mathrm{alg}}:
               \mathsf{GeomObj}
               \longrightarrow
               \mathsf{GNLAObj},
          \]
          placing $A$ as an operator-level shadow of geometric or
          law-level data (connections, curvatures, flows in \GZLS,
          etc.).
  \end{enumerate}
\end{definition}

The strict PFB-GNLA version will select GNLA’s where the Jacobiator
either vanishes or can be absorbed into a controlled sector compatible
with classical Lie theory. The homotopy and variable functional–operator
versions will keep $\Jac_A$ explicit and promote it to a higher bracket
in an $L_\infty$-structure.



\subsubsection*{(c) GNLA arising from principal bundles (schematic picture)}

For the PFB-GNLA associated with a principal bundle
$(\Bundle,\Base,\Group)$ and a connection \Conn, a typical choice of
$A$ looks like the following:

\[
   A
   \;\sim\;
   \Gamma\bigl(\mathrm{Ad}(\Bundle)\bigr)
   \;\oplus\;
   \Gamma\bigl(\mathrm{End}(E)\bigr)
   \;\oplus\;
   \text{(higher-form / path-space operators)},
\]
where $E$ is an associated bundle and “higher operators” capture
covariant derivatives, curvature insertions, and path-ordered
exponentials.



\begin{definition}[Geometric GNLA (schematic)]
A \emph{geometric GNLA} associated with a principal bundle
  $(\Bundle,\Base,\Group)$ and connection \Conn\ is a GNLA
  $(A,[\cdot,\cdot],\Jac_A;\mathsf{Str})$ together with:
  \begin{itemize}
    \item an embedding
          \[
             \iota:\mathfrak{g}
             \hookrightarrow A,
          \]
          sending Lie algebra elements to operators (for example,
          infinitesimal gauge transformations or curvature generators);
    \item a map
          \[
             \Theta:
               (\Conn,\Curv,
                \nabla,\text{holonomy data})
               \longrightarrow
               \mathsf{Str},
          \]
          which injects connection/curvature information into the
          structural data of $A$;
    \item a compatibility relation
          \[
             \Jac_A
             \sim
             \text{(curvature / higher-homotopy terms derived from }
             \Curv\text{ and }\Theta).
          \]
  \end{itemize}
\end{definition}

Informally, the Jacobiator is no longer forced to vanish; instead, it
records higher-order commutator anomalies and curvature effects. These
will later be organised, in homotopy PFB-GNLA, as components of an
$L_\infty$-structure, with $\Jac_A$ becoming the image of a ternary
bracket.



\subsubsection*{(d) Homotopy Lie structures: the $L_\infty$ framework}

To treat non-vanishing Jacobiators in a principled way, it is natural
to move from strict Lie algebras to \emph{homotopy Lie algebras}
($L_\infty$-algebras). We recall the schematic definition.

\begin{definition}[$L_\infty$-algebra (schematic)]
An \emph{$L_\infty$-algebra} over $\Bbbk$ is a $\mathbb{Z}$-graded
  vector space
  \[
     V = \bigoplus_{n\in\mathbb{Z}} V^n
  \]
  together with a family of multilinear maps
  \[
     \ell_k:
       V^{\otimes k}\longrightarrow V,
     \qquad k\ge1,
  \]
  of degree $2-k$, such that:
  \begin{enumerate}[label=\roman*)]
    \item each $\ell_k$ is graded antisymmetric;
    \item the $\ell_k$ satisfy a system of higher Jacobi identities,
          which for each $n\ge1$ impose
          \[
            \sum_{i+j=n+1}
            \sum_{\sigma\in S(i,n-i)}
            \chi(\sigma)\,
            (-1)^{i(j-1)}
            \ell_j\bigl(
               \ell_i(x_{\sigma(1)},\dots,x_{\sigma(i)}),
               x_{\sigma(i+1)},\dots,x_{\sigma(n)}
            \bigr)
            = 0,
          \]
          for all homogeneous $x_1,\dots,x_n\in V$, where
          $S(i,n-i)$ is the set of $(i,n-i)$-shuffles and $\chi(\sigma)$
          is the Koszul sign.
  \end{enumerate}
\end{definition}

We will only need low-degree cases explicitly:

\begin{itemize}
  \item $\ell_1:V^n\to V^{n+1}$ is a degree-$+1$ differential with
        $\ell_1^2=0$;
  \item $\ell_2:V\otimes V\to V$ is a degree-$0$ graded bracket; the
        $n=2$ identity says that $\ell_1$ is a graded derivation of
        $\ell_2$;
  \item $\ell_3:V^{\otimes3}\to V$ is a degree $-1$ ternary bracket;
        the $n=3$ identity expresses that the failure of $\ell_2$ to
        satisfy the graded Jacobi identity is homotopic to zero via
        $\ell_3$:
        \[
           \Jac_{\ell_2}(x,y,z)
           =
           \ell_1\bigl(\ell_3(x,y,z)\bigr)
           + \text{(terms involving }\ell_3\text{ and }\ell_2).
        \]
\end{itemize}

Thus in an $L_\infty$-algebra, the Jacobiator of the binary bracket
need not vanish, but is \emph{controlled} by higher operations. This is
exactly the situation we wish to encode for PFB-GNLA: the Jacobiator
built from connection- and curvature-derived operators is not zero,
but is constrained by differential identities and higher homotopies.



\subsubsection*{(e) Two-term and curvature-twisted homotopy Lie structures}

For the purposes of PFB-GNLA, we are particularly interested in
\emph{two-term} $L_\infty$-algebras, where only two degrees are
non-zero, and in curvature-twisted versions where a closed 3-form
(e.g.\ $H$ or a law-curvature component) twists the homotopy relations.



\begin{definition}[Two-term $L_\infty$-algebra (schematic)]
A \emph{two-term $L_\infty$-algebra} is an $L_\infty$-algebra
  concentrated in degrees $0$ and $-1$,
  \[
     V = V^{-1}\oplus V^0,
  \]
  with $\ell_k=0$ for $k\ge4$. The non-trivial structure maps are:
  \begin{itemize}
    \item $\ell_1:V^{-1}\to V^0$ (a differential);
    \item $\ell_2:V^i\otimes V^j\to V^{i+j}$ (a degree-$0$ bracket);
    \item $\ell_3:V^0\otimes V^0\otimes V^0\to V^{-1}$ (a degree $-1$
          ternary operation).
  \end{itemize}
  The higher Jacobi identities reduce to a finite set of equations
  linking $\ell_1,\ell_2,\ell_3$.
\end{definition}

In a curvature-twisted situation, one often has a closed 3-form
$H\in\Omega^3(\Base)$ or a law-curvature component in \GZLCurv\ that
enters the $\ell_3$-operation. Schematically, one can think of
\[
   \ell_3(x,y,z)
   \sim
   \int_{\Delta^3} H\wedge \Phi(x,y,z)
\]
for suitable $\Phi$; the precise formula depends on the geometric
context. The resulting structure is sometimes called an $H$-twisted
$L_\infty$-algebra, and provides the algebraic prototype for the
“$H$-twisted two-term $L_\infty$ structures and GZ-Harmony” developed
later in the homotopy version of \GAlgebra.



\subsubsection*{(f) Position of GNLA and homotopy Lie structures in \GZLS\ and \GZGMS}

We now summarise how generalized noncommutative Lie algebras and
homotopy Lie structures sit inside the GaoZheng Law-Space \GZLS\ and
the GaoZheng Generalized Mathematical Structure \GZGMS.



\begin{itemize}
  \item \emph{Law-objects.}
        Each GNLA $(A,[\cdot,\cdot],\Jac_A;\mathsf{Str})$ is regarded
        as a law-object $L\in\GZLS$, whose internal operations and
        Jacobiator are part of the law-level data. Strict Lie algebras
        appear as the special case $\Jac_A=0$.
  \item \emph{Homotopy upgrades.}
        When $A$ carries an $L_\infty$-structure
        $(\ell_1,\ell_2,\ell_3,\dots)$ extending its GNLA bracket,
        the corresponding law-object is enriched to a homotopy
        law-object $L^{\mathrm{homo}}\in\GZLS$, with higher brackets
        encoded by law-curvature components and CH-layering data.
  \item \emph{Embedding into \GZGMS.}
        Inside the GaoZheng Generalized Mathematical Structure
        \[
           \GZGMS =
           \bigl(\mathsf{Obj},\mathsf{Law},\mathsf{Flow},\mathsf{Layer},
                 \GZTLC;\GZLCurv\bigr),
        \]
        the GNLA or $L_\infty$-algebra lives in the $\mathsf{Law}$
        component, while its geometric realisation via PFB data lives
        in $\mathsf{Obj}$. The triple-layer connection \GZTLC\ couples
        them: horizontal flows in $\mathsf{Obj}$ induce homotopy
        flows in $\mathsf{Law}$ and vice versa.
  \item \emph{Gap statements and certificates.}
        Deviations of GNLA structures from strict Lie behaviour
        (non-zero Jacobiator, non-trivial higher brackets) are captured
        by law-curvature \GZLCurv\ and by certificate-type invariants
        such as Jacobiator-gap certificates. The associated Gap
        statements then constrain which geometric PFB configurations
        can realise which law-level GNLA or $L_\infty$ structures.
\end{itemize}

In this way, classical Lie algebras, generalized noncommutative Lie
algebras, and homotopy Lie structures form a hierarchy of algebraic
models for the same underlying idea: that connection, curvature,
higher homotopy, and law-level obstructions can be encoded as operator
algebras with controlled Jacobi failure. The remaining chapters of
Part~II will instantiate this hierarchy concretely in the four
integrated PFB-GNLA constructions, each of which realises a different
balance between strictness, homotopy, and operator-path dynamics.



\section{From Object-Level Geometry to Law-Level Structures: Positioning \GAlgebra in \GFramework}


With the classical principal-bundle geometry and the generalized
(noncommutative, homotopy) Lie structures in place, we now explain how
these data are assembled into law-objects in the GaoZheng Law-Space
\GZLS, and how the GaoZheng G-Algebra (\GAlgebra, PFB-GNLA) is
positioned inside the GaoZheng G-Framework (\GFramework) as the
canonical law-level shadow of object-level geometry.



\subsubsection*{(a) From principal bundles to law-objects in \GZLS}

We first make precise the idea that a principal bundle with connection
and curvature determines a law-object in \GZLS.

Let $\mathsf{GeomObj}$ denote a category whose objects are tuples
\[
   (\Bundle,\Base,\Group;\Conn,\Curv;\text{aux}),
\]
where $(\Bundle,\Base,\Group)$ is a principal bundle, $\Conn$ a
connection, $\Curv$ its curvature, and “aux” encodes auxiliary data
such as associated bundles, covariant derivatives, and holonomy
representations. Morphisms are bundle maps that preserve these
structures up to gauge equivalence.

\begin{definition}[Geometry-to-law functor \texorpdfstring{$\Phi_{\mathrm{geom}}$}{Phi\_geom}]
A \emph{geometry-to-law functor} is a functor
  \[
     \Phi_{\mathrm{geom}}:
       \mathsf{GeomObj}\longrightarrow\GZLS,
  \]
  which assigns to each object
  $(\Bundle,\Base,\Group;\Conn,\Curv;\text{aux})$ a law-object
  \[
     L_{\mathrm{geom}}
     =
     \Phi_{\mathrm{geom}}
       (\Bundle,\Base,\Group;\Conn,\Curv;\text{aux})
     \in \mathfrak{L}\subset\GZLS,
  \]
  and to each morphism a law-transform, such that:
  \begin{enumerate}[label=\roman*)]
    \item curvature, holonomy, and covariant derivatives are encoded as
          composition and metric data in $L_{\mathrm{geom}}$, i.e.\ as
          part of $(J,\mathrm{Loc},\Sigma)$ for \GZLS;
    \item gauge-equivalent geometric data are mapped to isomorphic
          law-objects in \GZLS.
  \end{enumerate}
\end{definition}

Concretely, $L_{\mathrm{geom}}$ may be thought of as a package
containing:

\begin{itemize}
  \item the adjoint bundle $\mathrm{Ad}(\Bundle)$ and its sections;
  \item the connection $\Conn$ and curvature $\Curv$, expressed in a
        gauge-invariant way (e.g.\ via parallel transport and holonomy);
  \item induced covariant derivatives on associated bundles and their
        spectrum (eigenvalues, resolvent bounds, etc.), which feed into
        the metric structure $J$;
  \item locality and threshold data $(\mathrm{Loc},\Sigma)$ specifying
        which deformations are admissible and when regime changes occur.
\end{itemize}



\subsubsection*{(b) From law-objects to operator algebras: \GAlgebra as GNLA}

The next step is to pass from the law-object $L_{\mathrm{geom}}$ to an
operator algebra that will be recognised as an instance of the GaoZheng
G-Algebra (\GAlgebra, PFB-GNLA).

\begin{definition}[Operator layer associated with $L_{\mathrm{geom}}$ (schematic)]
To a law-object $L_{\mathrm{geom}}\in\GZLS$ one associates an operator
  space
  \[
     A(L_{\mathrm{geom}})
     \;\subseteq\;
     \mathrm{End}\bigl(\Gamma(E_1)\oplus\dots\oplus\Gamma(E_r)\bigr)
     \;\oplus\;
     \text{(higher-form / path-space operators)},
  \]
  where the $E_i$ are associated bundles determined by
  $(\Bundle,\Base,\Group;\Conn,\Curv;\text{aux})$, and where the
  operators in $A(L_{\mathrm{geom}})$ are built from covariant
  derivatives, curvature insertions, and holonomy-derived maps along
  admissible paths.
\end{definition}

The space $A(L_{\mathrm{geom}})$ is naturally equipped with a
(noncommutative) bracket $[\cdot,\cdot]$ given by graded commutators,
together with a Jacobiator $\Jac_{A}$ that measures the failure of the
Jacobi identity in the presence of curvature and higher homotopies.

\begin{definition}[GaoZheng G-Algebra (\GAlgebra) at the strict level (schematic)]
The \\ \emph{(strict) GaoZheng G-Algebra} associated with
  $(\Bundle,\Base,\Group;\Conn,\Curv)$ is a geometric GNLA
  \[
     \GAlgebra_{\mathrm{strict}}
     =
     \bigl(
       A_{\mathrm{strict}},
       [\cdot,\cdot]_{\mathrm{strict}},
       \Jac_{\mathrm{strict}};
       \mathsf{Str}_{\mathrm{strict}}
     \bigr),
  \]
  where:
  \begin{itemize}
    \item $A_{\mathrm{strict}}$ is a distinguished subspace of
          $A(L_{\mathrm{geom}})$ containing infinitesimal gauge
          operators, curvature operators, and a controlled class of
          covariant-derivative–type operators;
    \item $[\cdot,\cdot]_{\mathrm{strict}}$ is the graded commutator
          restricted to $A_{\mathrm{strict}}$;
    \item $\Jac_{\mathrm{strict}}$ is the corresponding Jacobiator;
    \item $\mathsf{Str}_{\mathrm{strict}}$ records gradings,
          filtrations, and bounds imposed by $(J,\mathrm{Loc},\Sigma)$.
  \end{itemize}
  The strict regime is characterised by the fact that
  $\Jac_{\mathrm{strict}}$ vanishes or takes values in a curvature
  sector that is either trivialised or made central at the law-level
  (depending on the version).
\end{definition}

Stated more informally, \GAlgebra\ at the strict level is the operator
algebra that encodes “what the connection and curvature can do” to
associated fields, under the constraint that Jacobi-type relations
follow the classical pattern up to controlled curvature effects. This
is the algebraic avatar of the classical geometry in the GaoZheng
G-Framework.



\subsubsection*{(c) Positioning \GAlgebra\ as a law-object in \GZLS}

The preceding constructions can be summarised as a two-step process:

\[
   (\Bundle,\Base,\Group;\Conn,\Curv;\text{aux})
   \xrightarrow{\ \Phi_{\mathrm{geom}}\ }
   L_{\mathrm{geom}}\in\GZLS
   \xrightarrow{\ \mathsf{Op}\ }
   \GAlgebra_{\mathrm{strict}},
\]
where $\mathsf{Op}$ is the (schematic) functor that sends a law-object
to its associated operator GNLA. We now regard
$\GAlgebra_{\mathrm{strict}}$ itself as a law-object.

\begin{definition}[G-Algebra as a law-object]
A (strict) GaoZheng G-Algebra $\GAlgebra_{\mathrm{strict}}$ is said
  to be \emph{realised in \GZLS} if there exists a law-object
  $L_{\GAlgebra}\in\GZLS$ and an isomorphism
  \[
     \Psi:
       \GAlgebra_{\mathrm{strict}}
       \xrightarrow{\ \cong\ }
       L_{\GAlgebra}
  \]
  in the category underlying \GZLS, such that:
  \begin{enumerate}[label=\roman*)]
    \item the composition
          $(\Bundle,\Base,\Group;\Conn,\Curv;\text{aux})
           \mapsto L_{\GAlgebra}$
          agrees with
          $(\Bundle,\Base,\Group;\Conn,\Curv;\text{aux})
           \mapsto L_{\mathrm{geom}}$
          up to a law-transform in \GZLS;
    \item the GNLA structure of \GAlgebra\ (bracket, Jacobiator,
          structural data) is encoded in the law-space composition and
          metric data $(J,\mathrm{Loc},\Sigma)$ of $L_{\GAlgebra}$.
  \end{enumerate}
\end{definition}

In practice, one identifies $L_{\GAlgebra}$ with a region of law-space
in which the admissible law-objects are precisely those that can be
realised by strict PFB-GNLA constructions based on some principal
bundle and connection. This makes \GAlgebra\ both an operator algebra
and a law-object: it sits simultaneously in the algebraic and law-level
layers of the framework.



\subsubsection*{(d) Embedding \GAlgebra\ into the GaoZheng G-Framework \GFramework}

We now recall that the GaoZheng G-Framework is the meta-paradigm that
organises object-level categories, law-space, flows, and layering data,
as summarised schematically by
\[
   \GZGMS =
   \bigl(\mathsf{Obj},\mathsf{Law},\mathsf{Flow},\mathsf{Layer},
         \GZTLC;\GZLCurv\bigr),
\]
with \GFramework\ specifying the “design envelope” of admissible choices.
Within this context, the positioning of \GAlgebra\ is as follows.

\begin{itemize}
  \item \emph{In $\mathsf{Obj}$.}
        The principal bundle $(\Bundle,\Base,\Group)$, connection
        \Conn, curvature \Curv, and associated bundles belong to the
        object-level category $\mathsf{Obj}$.
  \item \emph{In $\mathsf{Law}$.}
        The law-objects $L_{\mathrm{geom}}$ and $L_{\GAlgebra}$ live in
        the law-space component $\mathsf{Law}=\GZLS$, equipped with
        law-transforms \GZLT, law-connection \GZLOC, and law-curvature
        \GZLCurv.
  \item \emph{In $\mathsf{Flow}$ and $\GZTLC$.}
        Object-level deformations of $(\Bundle,\Conn)$, parameter
        changes, and law-level flows are coupled by the triple-layer
        connection \GZTLC, as described in Section~2. In particular,
        horizontal flows in $\mathsf{Obj}$ correspond to horizontally
        admissible flows in $\mathsf{Law}$ that preserve the relation
        “geometry $\leftrightarrow$ \GAlgebra”.
  \item \emph{In $\mathsf{Layer}$.}
        CH-layering and semantic metrics distinguish regimes in which
        \GAlgebra\ may be treated as a strict Lie-type object (e.g.\ at
        extensional layers $l_2$) from regimes in which its homotopy
        refinements must be taken into account (layers $l_{\ge3}$).
\end{itemize}

Thus, inside \GFramework, the GaoZheng G-Algebra is the law-level
representation of how a given class of principal-bundle geometries
behaves under deformations and flows, as constrained by \GZTLC\ and
\GZLCurv.



\subsubsection*{(e) Interface with PL-PI layers and PFB-GNLA variants}

From the perspective of the PL-PI Meta-Mathematical Theory, the
GaoZheng G-Algebra plays a bridging role between the pan-logic (PL)
and pan-iterative (PI) layers:

\begin{itemize}
  \item On the \emph{PL side}, commutators in \GAlgebra\ encode local
        logical update rules: how inference operations, constraint
        propagations, and gauge transformations combine and interact,
        subject to Jacobi-type constraints or their homotopy relaxations.
  \item On the \emph{PI side}, operators in \GAlgebra\ act as generators
        of iterative schemes (e.g.\ evolution operators, update
        kernels, law-space learning steps), whose composition and
        curvature are governed by GNLA/$L_\infty$ identities.
\end{itemize}

In different PFB-GNLA variants, this interface takes different
concrete forms:

\begin{itemize}
  \item \textbf{Strict version.}
        \GAlgebra\ is realised as a GNLA with (effectively) vanishing
        Jacobiator; the PL and PI layers see it as a classical
        Lie-type algebra of generators and constraints.
  \item \textbf{Homotopy version.}
        \GAlgebra\ is upgraded to an $L_\infty$-structure; the pan-logic
        layer incorporates higher-sequent rules, and the pan-iterative
        layer sees homotopy-corrected evolution operators; law-curvature
        \GZLCurv\ and Jacobiator-gap certificates capture the deviation
        from strictness.
  \item \textbf{Variable functional–operator version.}
        \GAlgebra\ becomes an algebra of operators on path and
        functional spaces, directly tied to triple-layer flows in
        \GZTLC; law-space learning and operator-path dynamics are
        encoded algebraically.
  \item \textbf{Meta-mathematical original version.}
        The direction is reversed: PL-PI provenance and law-level flows
        generate both \GAlgebra\ and the underlying geometric category,
        showing that “geometry” can be recovered from law-level GNLA
        data.
\end{itemize}



\subsubsection*{(f) Summary: \GAlgebra as the law-level image of PFB geometry}

To summarise, the positioning of the GaoZheng G-Algebra in the
GaoZheng G-Framework can be encapsulated as:

\begin{itemize}
  \item Start from object-level PFB geometry:
        $(\Bundle,\Base,\Group;\Conn,\Curv;\text{aux})\in\mathsf{Obj}$.
  \item Use $\Phi_{\mathrm{geom}}$ to map it into law-space:
        $L_{\mathrm{geom}}\in\GZLS$.
  \item Construct the operator GNLA $A(L_{\mathrm{geom}})$ and
        specialise to $\GAlgebra_{\mathrm{strict}}$ (or its homotopy /
        operator upgrades) as the canonical operator layer.
  \item Realise \GAlgebra\ itself as a law-object
        $L_{\GAlgebra}\in\GZLS$, compatible with law-transforms and
        law-curvature \GZLCurv.
  \item Use the triple-layer connection \GZTLC\ to couple object-level
        deformations and law-level flows in a way that preserves the
        identification “geometry $\leftrightarrow$ \GAlgebra”, and
        quantify departures from strictness via Gap statements and
        certificate-type invariants.
\end{itemize}

In this sense, \GAlgebra\ is the law-level image of principal-bundle
geometry under the PL-PI meta-paradigm: it is where curvature,
holonomy, and higher homotopies become algebraic generators and
relations in a unified law-space, ready to be deployed in the strict,
homotopy, and variable functional–operator PFB-GNLA constructions that
occupy the remainder of Part~II.



\chapter[Strict Version of the GaoZheng G-Algebra]{Strict Version of the GaoZheng G-Algebra:\\
         PFB-GNLA with Parameterized Connections}


\section{Parameterized Connections and Families of Operators}


In the strict version of the GaoZheng G-Algebra (\GAlgebra), the basic
geometric input is not a single connection on a fixed principal bundle,
but a \emph{smooth family} of connections, organised over a parameter
manifold. This section formalises such parameterized connections and the
associated families of covariant and curvature-type operators, which
will later be assembled into the strict PFB-GNLA operator algebra.



\subsubsection*{(a) Parameter manifold and connection family}

Let $(\Bundle,\Base,\Group)$ be a fixed principal bundle as in
Section~3.1. We now introduce a smooth parameter manifold
$\Lambda_{\mathrm{par}}$ and use it to index a family of connections.

\begin{definition}[Parameter manifold and parameterized connections]
A \emph{parameter manifold} for the strict PFB-GNLA construction is a
  smooth manifold $\Lambda_{\mathrm{par}}$ together with a smooth map
  \[
     \Lambda_{\mathrm{par}}
     \longrightarrow
     \mathsf{Conn}(\Bundle),
     \qquad
     \lambda\longmapsto\Conn_\lambda,
  \]
  where each $\Conn_\lambda\in\Omega^1(\Bundle;\mathfrak{g})$ is a
  connection 1-form on $(\Bundle,\Base,\Group)$. The corresponding
  curvature is denoted
  \[
     \Curv_\lambda
     =
     \mathrm{d}\Conn_\lambda
     + \tfrac12[\Conn_\lambda,\Conn_\lambda]
     \in\Omega^2(\Bundle;\mathfrak{g}).
  \]
\end{definition}

In local coordinates $(\lambda^a)$ on $\Lambda_{\mathrm{par}}$, smooth
dependence means that in any local trivialisation one may write
\[
   \Conn_\lambda
   =
   \Conn_0
   + \sum_a \lambda^a\,K_a
   + \text{(higher-order terms in }\lambda),
\]
with $K_a\in\Omega^1(\Bundle;\mathfrak{g})$ smooth, and similarly for
$\Curv_\lambda$. The $K_a$ encode the infinitesimal response of the
connection to changes in the parameters.



\subsubsection*{(b) Product bundle and total parameterized connection}

It is convenient to package the family $\{\Conn_\lambda\}$ into a single
connection on a product bundle over $\Base\times\Lambda_{\mathrm{par}}$.

\begin{definition}[Product principal bundle and total connection]
Consider the product manifold
  \[
     \widehat{\Base}
     :=
     \Base\times\Lambda_{\mathrm{par}},
  \]
  and the product principal bundle
  \[
     \widehat{\Bundle}
     :=
     \Bundle\times\Lambda_{\mathrm{par}}
     \xrightarrow{\ \widehat{\pi}\ }
     \widehat{\Base},
     \qquad
     \widehat{\pi}(u,\lambda)
     =
     (\pi(u),\lambda),
  \]
  with structure group still $\Group$ acting on the first factor. A
  \emph{total parameterized connection} is a $\mathfrak{g}$-valued
  1-form
  \[
     \widehat{\Conn}
     \in
     \Omega^1(\widehat{\Bundle};\mathfrak{g})
  \]
  which satisfies the usual connection axioms and restricts on each
  slice $\lambda=\lambda_0$ to the pull-back of $\Conn_{\lambda_0}$:
  \[
     \widehat{\Conn}\big|_{\Bundle\times\{\lambda_0\}}
     =
     \mathrm{pr}_1^*(\Conn_{\lambda_0}).
  \]
\end{definition}

Locally one can decompose $\widehat{\Conn}$ into its “geometric” and
“parameter” components. Writing $T_{(u,\lambda)}\widehat{\Bundle}
\cong T_u\Bundle\oplus T_\lambda\Lambda_{\mathrm{par}}$, we may set
\[
   \widehat{\Conn}(X,0) = \Conn_\lambda(X),\qquad
   \widehat{\Conn}(0,v) = B_{\mathrm{par}}(v),
\]
for some $\mathfrak{g}$-valued 1-form
$B_{\mathrm{par}}\in\Omega^1(\Lambda_{\mathrm{par}};\mathfrak{g})$
(the “parameter connection component”). In coordinates, this can be
written schematically as
\[
   \widehat{\Conn}
   =
   \Conn_\lambda
   + B_{\mathrm{par}},
\]
where the first term acts along $\Base$ and the second along
$\Lambda_{\mathrm{par}}$.

The curvature of $\widehat{\Conn}$,
\[
   \widehat{\Curv}
   :=
   \mathrm{d}\widehat{\Conn}
   + \tfrac12[\widehat{\Conn},\widehat{\Conn}]
   \in\Omega^2(\widehat{\Bundle};\mathfrak{g}),
\]
decomposes accordingly into:

\begin{itemize}
  \item the \emph{geometric curvature} component
        $\Curv_\lambda$ along $\Base$;
  \item the \emph{parameter curvature} component
        $\Curv_{\mathrm{par}}$ along $\Lambda_{\mathrm{par}}$;
  \item \emph{mixed curvature} components encoding how variations in
        $\lambda$ affect curvature along $\Base$.
\end{itemize}



\subsubsection*{(c) Families of covariant derivatives and curvature operators}

Let $\rho:\Group\to\mathrm{GL}(V)$ be a finite-dimensional
representation and $E:=\Bundle\times_\rho V$ the associated vector
bundle over $\Base$. For each $\lambda\in\Lambda_{\mathrm{par}}$, the
connection $\Conn_\lambda$ induces a covariant derivative
\[
   \nabla^{(\lambda)}:
      \Gamma(E)\longrightarrow
      \Gamma(T^*\Base\otimes E),
\]
and a curvature operator
\[
   R^{(\lambda)}_\nabla
   =
   (\nabla^{(\lambda)})^2
   :
   \Gamma(E)
   \longrightarrow
   \Gamma(\Lambda^2T^*\Base\otimes E).
\]

\begin{definition}[Spatial covariant-derivative family]
The \emph{spatial covariant-derivative family} associated with the
  parameterized connection is the collection
  \[
     \mathcal{D}_{\Base}
     :=
     \bigl\{
       \nabla^{(\lambda)}_X
       \,\big|\,
       X\in\Gamma(T\Base),\;
       \lambda\in\Lambda_{\mathrm{par}}
     \bigr\},
  \]
  where $\nabla^{(\lambda)}_X$ denotes covariant differentiation along
  $X$ with respect to $\Conn_\lambda$.
\end{definition}

Using the total connection $\widehat{\Conn}$ on
$\widehat{\Bundle}=\Bundle\times\Lambda_{\mathrm{par}}$, we similarly
obtain a covariant derivative
\[
   \widehat{\nabla}:
     \Gamma(\widehat{E})
     \longrightarrow
     \Gamma(T^*\widehat{\Base}\otimes\widehat{E}),
\]
where $\widehat{E}:=E\times\Lambda_{\mathrm{par}}$ and
$\widehat{\Base}=\Base\times\Lambda_{\mathrm{par}}$. Decomposing
$T\widehat{\Base}\cong T\Base\oplus T\Lambda_{\mathrm{par}}$ yields
\[
   \widehat{\nabla}
   =
   \nabla^{(\lambda)}_{\Base}
   \;\oplus\;
   \nabla^{(\lambda)}_{\Lambda},
\]
where:

\begin{itemize}
  \item $\nabla^{(\lambda)}_{\Base}$ acts along $\Base$ and coincides
        with $\nabla^{(\lambda)}$;
  \item $\nabla^{(\lambda)}_{\Lambda}$ acts along $\Lambda_{\mathrm{par}}$
        and captures how sections of $E$ change with $\lambda$ under
        the parameter connection component $B_{\mathrm{par}}$.
\end{itemize}

\begin{definition}[Parameter-direction covariant-derivative family]
The \\ \emph{parameter-direction covariant-derivative family} is the collection
  \[
     \mathcal{D}_{\Lambda}
     :=
     \bigl\{
       \nabla^{(\lambda)}_{v}
       \,\big|\,
       v\in\Gamma(T\Lambda_{\mathrm{par}}),\;
       \lambda\in\Lambda_{\mathrm{par}}
     \bigr\},
  \]
  where $\nabla^{(\lambda)}_{v}$ is induced by the $\Lambda$-component
  of $\widehat{\nabla}$.
\end{definition}

We also record the family of curvature operators
\[
   \mathcal{R}
   :=
   \bigl\{
     R^{(\lambda)}_\nabla
     \,\big|\,
     \lambda\in\Lambda_{\mathrm{par}}
   \bigr\},
\]
together with mixed curvature components derived from $\widehat{\Curv}$.



\subsubsection*{(d) The strict PFB-GNLA operator space: basic construction}

The strict PFB-GNLA operator layer associated with
$(\Bundle,\Base,\Group)$ and the parameter manifold $\Lambda_{\mathrm{par}}$
is built from the families $\mathcal{D}_{\Base}$, $\mathcal{D}_{\Lambda}$
and $\mathcal{R}$.

\begin{definition}[Strict operator space $A_{\mathrm{strict}}$ (schematic)]
Fix $(\Bundle,\Base,\Group)$, a parameter manifold
  $\Lambda_{\mathrm{par}}$, and a parameterized connection
  $\{\Conn_\lambda\}_{\lambda\in\Lambda_{\mathrm{par}}}$ with total
  connection $\widehat{\Conn}$. The \emph{strict operator space}
  associated with this data is the linear space
  \[
     A_{\mathrm{strict}}
     :=
     \mathrm{span}\Bigl(
       \mathcal{D}_{\Base}
       \;\cup\;
       \mathcal{D}_{\Lambda}
       \;\cup\;
       \mathcal{R}
     \Bigr)
     \;\subseteq\;
     \mathrm{End}\bigl(\Gamma(E)\times\Lambda_{\mathrm{par}}\bigr),
  \]
  consisting of finite linear combinations and (when needed) suitable
  closures under composition of:
  \begin{itemize}
    \item spatial covariant derivatives $\nabla^{(\lambda)}_X$;
    \item parameter-direction derivatives $\nabla^{(\lambda)}_{v}$;
    \item curvature operators $R^{(\lambda)}_\nabla$ and mixed
          curvature components.
  \end{itemize}
\end{definition}

We equip $A_{\mathrm{strict}}$ with the graded commutator bracket
\[
   [T_1,T_2]_{\mathrm{strict}}
   :=
   T_1\circ T_2 - T_2\circ T_1,
\]
and, when appropriate, with a grading induced by the degree (number of
derivatives, curvature insertions, etc.). This bracket is noncommutative
in general, but inherits strong constraints from the underlying
differential-geometric identities:

\begin{itemize}
  \item commutators of covariant derivatives reproduce curvature
        operators (e.g.\ $[\nabla^{(\lambda)}_X,
        \nabla^{(\lambda)}_Y] = R^{(\lambda)}_\nabla(X,Y)$ on $E$);
  \item Bianchi identities impose algebraic relations among
        commutators of curvature operators and higher compositions;
  \item parameter-direction derivatives and mixed curvatures record how
        these relations change with $\lambda$.
\end{itemize}

These identities are precisely the source of the “strictness” in the
strict PFB-GNLA regime: they ensure that the Jacobiator, when computed
on the relevant subspace and after appropriate identification of
curvature sectors, behaves in a controlled way (to be detailed in
Section~4.2).



\subsubsection*{(e) Relation to the triple-layer connection \texorpdfstring{\GZTLC}{GZ-TLC}}

Although in the strict version we only use the geometric and parameter
layers explicitly, it is instructive to relate the above construction
to the triple-layer connection \GZTLC\ introduced in Part~I.

Recall that \GZTLC\ is a principal connection on a triple-layer bundle
with base
\[
   X_{\mathrm{tot}}
   =
   \Base \times \Lambda_{\mathrm{par}} \times \mathfrak{L}_0,
\]
whose law-component is the pull-back of the law-connection \GZLOC\ on
\GZLS. In the strict PFB-GNLA regime, we work in a region of law-space
where \GZLOC\ and its curvature \GZLCurv\ are trivial or sufficiently
controlled so that law-level effects can be absorbed into the geometric
and parameter layers.

From this point of view:

\begin{itemize}
  \item the product bundle $\widehat{\Bundle}=\Bundle\times\Lambda_{\mathrm{par}}$
        with its total connection $\widehat{\Conn}$ can be seen as the
        restriction of the triple-layer bundle to a fixed law-state
        (or a trivial law-family);
  \item the operator space $A_{\mathrm{strict}}$ captures the part of
        the triple-layer curvature $\mathcal{F}_{\GZTLC}$ that is
        visible at the geometric and parameter layers, while the law
        layer is kept flat or decoupled;
  \item the strict GaoZheng G-Algebra $\GAlgebra_{\mathrm{strict}}$
        will be built from $A_{\mathrm{strict}}$ as a GNLA in which the
        Jacobiator is either forced to vanish or confined to a trivial
        law-sector, reflecting the absence of non-trivial law-curvature.
\end{itemize}



\subsubsection*{(f) Summary and outlook}

To summarise, the data of a parameter manifold $\Lambda_{\mathrm{par}}$
and a smooth family of connections $\{\Conn_\lambda\}$ on a fixed
principal bundle $(\Bundle,\Base,\Group)$ give rise to:

\begin{itemize}
  \item a total parameterized connection $\widehat{\Conn}$ on the
        product bundle $\widehat{\Bundle}=\Bundle\times\Lambda_{\mathrm{par}}$;
  \item spatial and parameter-direction covariant-derivative families
        $\mathcal{D}_{\Base}$ and $\mathcal{D}_{\Lambda}$, together
        with a curvature family $\mathcal{R}$ and mixed curvature
        components;
  \item a strict operator space $A_{\mathrm{strict}}$ generated by these
        operators, equipped with the graded commutator bracket.
\end{itemize}

In the next section, these ingredients will be assembled into the
strict PFB-GNLA structure data: we will specify the precise GNLA
constraints on $(A_{\mathrm{strict}},[\cdot,\cdot]_{\mathrm{strict}},
\Jac_{\mathrm{strict}})$, explain how the Jacobi identity is enforced
(or its failure controlled) in the strict regime, and position the
resulting GaoZheng G-Algebra $\GAlgebra_{\mathrm{strict}}$ as the
baseline chart of the four integrated PFB-GNLA constructions.



\section{Structural Data of the Strict PFB-GNLA (\GAlgebra)}


The purpose of this section is to assemble the geometric and operator
data of Section~4.1 into a single, well-structured object: the \emph{strict
principal-bundle-based generalized noncommutative Lie algebra} (strict
PFB-GNLA), denoted by \GAlgebra\ in the strict regime. We make explicit
which pieces of data are retained, how they are organised, and which
algebraic constraints implement the “strict” character of this version.



\subsubsection*{(a) Geometric–parameter block}

We keep fixed a principal bundle $(\Bundle,\Base,\Group)$ and a
parameter manifold $\Lambda_{\mathrm{par}}$, together with a smooth
family of connections as in Section~4.1.

\begin{definition}[Geometric–parameter block]
A \emph{geometric–parameter block} for the strict PFB-GNLA consists
  of:
  \[
     \mathsf{GeomPar}
     =
     \bigl(
       \Bundle,\Base,\Group;
       \Lambda_{\mathrm{par}};
       \{\Conn_\lambda\}_{\lambda\in\Lambda_{\mathrm{par}}},
       \widehat{\Conn},\widehat{\Curv}
     \bigr),
  \]
  where:
  \begin{itemize}
    \item $(\Bundle,\Base,\Group)$ is a principal bundle with
          connection family $\{\Conn_\lambda\}$;
    \item $\Lambda_{\mathrm{par}}$ is the parameter manifold;
    \item $\widehat{\Conn}\in\Omega^1(\widehat{\Bundle};\mathfrak{g})$
          is a total parameterized connection on
          $\widehat{\Bundle}=\Bundle\times\Lambda_{\mathrm{par}}$;
    \item $\widehat{\Curv}$ is its curvature, decomposing into
          geometric, parameter, and mixed components.
  \end{itemize}
\end{definition}

Associated to this block are families of covariant derivatives and curvature operators \\ 
$(\mathcal{D}_{\Base},\mathcal{D}_{\Lambda},\mathcal{R})$,
as defined in Section~4.1; these will generate the strict operator
layer.



\subsubsection*{(b) Strict operator block and GNLA structure}

We now package the operator families into a GNLA-type algebraic
structure. The starting point is the strict operator space
$A_{\mathrm{strict}}$ introduced in Definition~4.1.

\begin{definition}[Strict operator block]
Given a geometric–parameter block $\mathsf{GeomPar}$ and a choice of
  associated bundles $E_1,\dots,E_r$, the \emph{strict operator block}
  is the tuple
  \[
     \mathsf{Op}_{\mathrm{strict}}
     =
     \bigl(
       A_{\mathrm{strict}},
       [\cdot,\cdot]_{\mathrm{strict}},
       \Jac_{\mathrm{strict}};
       \mathsf{Gr},\mathsf{Filt}
     \bigr),
  \]
  where:
  \begin{itemize}
    \item $A_{\mathrm{strict}}$ is the operator space generated by
          $\mathcal{D}_{\Base}\cup\mathcal{D}_{\Lambda}\cup\mathcal{R}$
          acting on $\Gamma(E_1)\oplus\dots\oplus\Gamma(E_r)$;
    \item $[\cdot,\cdot]_{\mathrm{strict}}$ is the graded commutator;
    \item $\Jac_{\mathrm{strict}}$ is the Jacobiator of
          $[\cdot,\cdot]_{\mathrm{strict}}$;
    \item $\mathsf{Gr}$ is a grading (e.g.\ by “differential order +
          curvature degree”);
    \item $\mathsf{Filt}$ is a filtration compatible with
          $[\cdot,\cdot]_{\mathrm{strict}}$.
  \end{itemize}
\end{definition}

We impose the following \emph{strict GNLA constraints} on
$\mathsf{Op}_{\mathrm{strict}}$.

\begin{definition}[Strict GNLA constraints (schematic)]
The strict operator block is said to satisfy the \emph{strict GNLA
  constraints} if:
  \begin{enumerate}[label=(S\arabic*)]
    \item \emph{Filtration compatibility.}
          The bracket respects the filtration:
          \[
             [\mathsf{Filt}^pA_{\mathrm{strict}},
              \mathsf{Filt}^qA_{\mathrm{strict}}]_{\mathrm{strict}}
             \subseteq
             \mathsf{Filt}^{p+q}A_{\mathrm{strict}}.
          \]
          In particular, commutators of operators of bounded differential
          order remain of controlled order.
    \item \emph{Curvature reproduction.}
          For first-order spatial covariant derivatives,
          \[
             [\nabla^{(\lambda)}_X,\nabla^{(\lambda)}_Y]_{\mathrm{strict}}
             =
             R^{(\lambda)}_\nabla(X,Y)
             + \text{(lower-order terms)},
          \]
          and analogous identities hold involving
          $\mathcal{D}_{\Lambda}$ and mixed curvature, expressing the usual
          geometric commutation relations.
    \item \emph{Controlled Jacobiator.}
          The Jacobiator $\Jac_{\mathrm{strict}}$ takes values in a
          distinguished “curvature sector”
          $J_{\mathrm{curv}}\subseteq A_{\mathrm{strict}}$, and there
          exists a projection
          $q:A_{\mathrm{strict}}\to A_{\mathrm{strict}}/J_{\mathrm{curv}}$
          such that the induced bracket on
          $A_{\mathrm{strict}}/J_{\mathrm{curv}}$ satisfies the Jacobi
          identity strictly.
    \item \emph{Centralisation in the strict regime.}
          In the strict PFB-GNLA regime, the curvature sector
          $J_{\mathrm{curv}}$ is either:
          \begin{itemize}
            \item trivial (identically zero), or
            \item central up to a prescribed order, i.e.\ for all
                  $T\in A_{\mathrm{strict}}$ and all
                  $C\in J_{\mathrm{curv}}$,
                  \(
                     [T,C]_{\mathrm{strict}}
                     \in \mathsf{Filt}^{>k}A_{\mathrm{strict}}
                  \)
                  for some fixed $k$, making $J_{\mathrm{curv}}$
                  negligible at the law-level scale considered.
          \end{itemize}
  \end{enumerate}
\end{definition}

Under these conditions, $(A_{\mathrm{strict}},[\cdot,\cdot]_{\mathrm{strict}},
\Jac_{\mathrm{strict}};\mathsf{Gr},\mathsf{Filt})$ is a GNLA whose
“non-Lie” behaviour is fully quarantined into a controlled curvature
sector, and whose quotient by that sector is a genuine Lie algebra of
operators.



\subsubsection*{(c) Law-level block and embedding into \GZLS}

The strict PFB-GNLA must be positioned as a law-object in the GaoZheng
Law-Space \GZLS. We therefore add a law-level block that records how
the operator GNLA is seen in law-space.

\begin{definition}[Strict law-level block]
A \emph{strict law-level block} associated with \\
  $(\mathsf{GeomPar},\mathsf{Op}_{\mathrm{strict}})$ is a law-object
  \[
     L_{\mathrm{strict}}
     =
     \bigl(
       \mathfrak{L}_{\mathrm{strict}};
       J_{\mathrm{strict}},
       \mathrm{Loc}_{\mathrm{strict}},
       \Sigma_{\mathrm{strict}}
     \bigr)
     \in \GZLS,
  \]
  together with:
  \begin{itemize}
    \item an injection
          \(
             \iota_{\mathrm{law}}:
             A_{\mathrm{strict}}\hookrightarrow
             \mathfrak{L}_{\mathrm{strict}}
          \),
          embedding operators as law-objects;
    \item metric and localisation data
          $(J_{\mathrm{strict}},\mathrm{Loc}_{\mathrm{strict}})$, and
          threshold set $\Sigma_{\mathrm{strict}}$, compatible with
          the grading and filtration on $A_{\mathrm{strict}}$;
    \item a law-connection
          \(
             {\GZLOC}_{\mathrm{strict}}
          \)
          on $\mathfrak{L}_{\mathrm{strict}}$ with curvature
          \(
             {\GZLCurv}_{\mathrm{strict}},
          \)
          such that ${\GZLCurv}_{\mathrm{strict}}$ is either flat or
          confined to a controlled sector corresponding to
          $J_{\mathrm{curv}}$.
  \end{itemize}
\end{definition}

In the strict regime, we are in a region of \GZLS\ where law-curvature
can be neglected or treated as central, so that the primary law-level
structure is the Lie-type algebra obtained by projecting
$A_{\mathrm{strict}}$ modulo $J_{\mathrm{curv}}$. This matches the
intuition that strict PFB-GNLA describes “near-flat” or “controlled
curvature” law-regimes, serving as a baseline chart in law-space.



\subsubsection*{(d) Strict PFB-GNLA datum: integrated definition}

We can now collect the geometric–parameter block, strict operator block,
and strict law-level block into a single integrated datum.

\begin{definition}[Strict PFB-GNLA datum]
A \emph{strict PFB-GNLA datum} (or strict GaoZheng G-Algebra datum)
  is a tuple
  \[
    \mathsf{PFB\text- GNLA}_{\mathrm{strict}}
    =
    \bigl(
      \mathsf{GeomPar},
      \mathsf{Op}_{\mathrm{strict}},
      L_{\mathrm{strict}};
      \Phi_{\mathrm{geom}},\mathcal{C}_{\mathrm{strict}}
    \bigr),
  \]
  where:
  \begin{itemize}
    \item $\mathsf{GeomPar}$ is a geometric–parameter block;
    \item $\mathsf{Op}_{\mathrm{strict}}$ is a strict operator block
          satisfying (S1)–(S4);
    \item $L_{\mathrm{strict}}\in\GZLS$ is a strict law-level block as
          above;
    \item $\Phi_{\mathrm{geom}}$ is the geometry-to-law functor of
          Section~3.3, restricted to the strict regime;
    \item $\mathcal{C}_{\mathrm{strict}}$ is a strict law–geometry
          coupling (as in Section~2.2) satisfying:
          \[
             \Phi_{\mathrm{geom}}
               (\Bundle,\Base,\Group;\Conn_\lambda,\Curv_\lambda;\text{aux})
             \simeq
             L_{\mathrm{strict}}
             \quad\text{and}\quad
             \iota_{\mathrm{law}}(A_{\mathrm{strict}})
             \subseteq
             \mathfrak{L}_{\mathrm{strict}}.
          \]
  \end{itemize}
\end{definition}

Informally, a strict PFB-GNLA datum says:

\begin{quote}
  “We have a principal bundle with a parameterized family of
   connections, we form the operator algebra of covariant derivatives
   and curvature operators, we embed this algebra as a law-object in
   law-space, and the coupling between geometry and law is such that
   all non-Lie behaviour is confined to a controlled curvature sector
   that is negligible or central in the strict regime.”
\end{quote}



\subsubsection*{(e) Strictness as a baseline for homotopy and operator-path upgrades}

Finally, we clarify the role of “strictness” in relation to the later
PFB-GNLA variants.

\begin{remark}[Strictness as a baseline regime]
The strict PFB-GNLA regime is characterised by:
  \begin{itemize}
    \item law-curvature \GZLCurv\ being flat or confined to a central,
          controlled sector;
    \item Jacobi failure being either absent or quarantined in
          $J_{\mathrm{curv}}$, which is negligible at the law-level
          scale considered;
    \item operator dynamics being describable entirely within
          $A_{\mathrm{strict}}$ and its Lie quotient \\
          $A_{\mathrm{strict}}/J_{\mathrm{curv}}$.
  \end{itemize}
  In later chapters:
  \begin{itemize}
    \item the \emph{homotopy version} will relax these constraints by
          promoting $\Jac_{\mathrm{strict}}$ to a genuine homotopy
          operation $\ell_3$ in an $L_\infty$-structure, thus turning
          curvature-induced Jacobi failure into controlled homotopies;
    \item the \emph{variable functional–operator version} will enlarge
          $A_{\mathrm{strict}}$ to operator-path and functional spaces,
          with \GZTLC\ acting directly on the operator layer;
    \item the \emph{meta-mathematical original version} will reconstruct
          $\mathsf{GeomPar}$ from law-level and homotopy data, reversing
          the direction of the strict construction.
  \end{itemize}
\end{remark}

From the vantage point of the GaoZheng G-Framework, the strict
PFB-GNLA data defined in this section provide the “coordinate chart”
in which principal-bundle geometry is encoded as a near-Lie operator
algebra in law-space. The next section will make explicit how this
strict chart is matched with the GZ Law-Space \GZLS\ and the triple-layer
connection \GZTLC, preparing the ground for the integrated four-way
PFB-GNLA comparison in Part~III.



\section[Matching the Strict Version with \GZLStxt\ and \GZTLCtxt]{Matching the Strict Version with \GZLS\ (\GZLStxt) and \GZTLC\ (\GZTLCtxt)}


The strict PFB-GNLA datum of Section~4.2 consists of a geometric–parameter
block, a strict operator block, and a strict law-level block. In this
section we make precise how this datum is embedded into the GaoZheng
Law-Space \GZLS\ and how it is carried by a triple-layer connection
\GZTLC\ in the sense of Part~II, thereby identifying the strict PFB-GNLA
as a distinguished “flat or near-flat slice” of the full
GaoZheng G-Framework.



\subsubsection*{(a) Strict slice of the GaoZheng Law-Space \texorpdfstring{\GZLS}{GZ-LS}}

Recall that the GaoZheng Law-Space is schematically written as
\[
   \GZLS
   =
   (\mathfrak{L};J,\mathrm{Loc},\Sigma),
\]
where $\mathfrak{L}$ is the semantic universe of law-objects, $J$ is a
family of metrics or energy-like quantities, $\mathrm{Loc}$ encodes
localisation, and $\Sigma$ records thresholds for regime changes. The
strict law-level block $L_{\mathrm{strict}}$ of Section~4.2 picks out a
distinguished region of \GZLS.

\begin{definition}[Strict GZ-LS slice]
A \emph{strict GZ-LS slice} is a pair
  \[
     {\GZLS}_{\mathrm{strict}}
     :=
     \bigl(
       \mathfrak{L}_{\mathrm{strict}};
       J_{\mathrm{strict}},
       \mathrm{Loc}_{\mathrm{strict}},
       \Sigma_{\mathrm{strict}}
     \bigr),
  \]
  together with an embedding
  \[
     \iota_{\mathrm{LS}}:
        \mathfrak{L}_{\mathrm{strict}}\hookrightarrow\mathfrak{L},
  \]
  such that:
  \begin{enumerate}[label=\roman*)]
    \item $(J_{\mathrm{strict}},\mathrm{Loc}_{\mathrm{strict}},
           \Sigma_{\mathrm{strict}})$ are restrictions of
          $(J,\mathrm{Loc},\Sigma)$ to $\mathfrak{L}_{\mathrm{strict}}$;
    \item the law-connection ${\GZLOC}_{\mathrm{strict}}$ on
          $\mathfrak{L}_{\mathrm{strict}}$ is the restriction of the
          global law-connection \GZLOC, and its law-curvature
          ${\GZLCurv}_{\mathrm{strict}}$ takes values in a controlled
          sector corresponding to $J_{\mathrm{curv}}$ in
          $A_{\mathrm{strict}}$ (trivial or central in the strict regime);
    \item the strict law-level block $L_{\mathrm{strict}}$ of
          Section~4.2 is a law-object in $\mathfrak{L}_{\mathrm{strict}}$.
  \end{enumerate}
\end{definition}

In other words, ${\GZLS}_{\mathrm{strict}}$ is the region of law-space
where law-curvature is flat or effectively central at the resolution
of interest, and where the law-objects coming from strict PFB-GNLA
data live. The embedding $\iota_{\mathrm{LS}}$ situates this slice
inside the full \GZLS.



\subsubsection*{(b) Extending the strict datum to a triple-layer setting}

The strict PFB-GNLA datum
\[
   \mathsf{PFB\text- GNLA}_{\mathrm{strict}}
   =
   \bigl(
     \mathsf{GeomPar},
     \mathsf{Op}_{\mathrm{strict}},
     L_{\mathrm{strict}};
     \Phi_{\mathrm{geom}},\mathcal{C}_{\mathrm{strict}}
   \bigr)
\]
contains a geometric–parameter block and a strict law-level block.
To match it with the triple-layer connection \GZTLC\ of Part~II, we
need to place it into a triple-layer principal bundle setting.

Recall that in Section~2.1 the total base of a triple-layer system was
taken as
\[
   X_{\mathrm{tot}}
   =
   \Base\times\Lambda_{\mathrm{par}}\times\mathfrak{L}_0,
\]
for a suitable local slice $\mathfrak{L}_0\subset\mathfrak{L}$ of
law-space, and that $\GZTLC$ is a principal connection on a
$\widehat{G}$-bundle over $X_{\mathrm{tot}}$ whose law-component is the
pull-back of \GZLOC.

In the strict regime we choose $\mathfrak{L}_0$ so that
$\mathfrak{L}_{\mathrm{strict}}\subseteq\mathfrak{L}_0$ and restrict
attention to the strict slice:

\[
   X_{\mathrm{strict}}
   :=
   \Base\times\Lambda_{\mathrm{par}}\times
   \mathfrak{L}_{\mathrm{strict}}.
\]

\begin{definition}[Strict triple-layer bundle and connection]
A \emph{strict triple-layer bundle} associated with a strict
  PFB-GNLA datum is a principal $\widehat{G}$-bundle
  \[
     \widehat{\pi}_{\mathrm{strict}}:
       \widehat{\Bundle}_{\mathrm{strict}}
       \longrightarrow
       X_{\mathrm{strict}},
  \]
  together with a triple-layer connection
  \[
     \GZTLC_{\mathrm{strict}}
     \in
     \Omega^1(\widehat{\Bundle}_{\mathrm{strict}};\widehat{\mathfrak{g}}),
  \]
  such that:
  \begin{enumerate}[label=\roman*)]
    \item its restriction to the geometric–parameter directions
          $(\Base,\Lambda_{\mathrm{par}})$ coincides with the total
          parameterized connection $\widehat{\Conn}$ from Section~4.1;
    \item its law-component is the pull-back of
          ${\GZLOC}_{\mathrm{strict}}$ along the projection to
          $\mathfrak{L}_{\mathrm{strict}}$;
    \item its curvature satisfies
          \[
             {\GZLCurv}_{\mathrm{strict}}
             =
             \pi_{\mathrm{law}}
               \bigl(\mathcal{F}_{\GZTLC_{\mathrm{strict}}}\bigr)
             \in
             J_{\mathrm{curv}}^{\mathrm{law}},
          \]
          where $J_{\mathrm{curv}}^{\mathrm{law}}$ is the law-level
          image of the curvature sector $J_{\mathrm{curv}}\subseteq
          A_{\mathrm{strict}}$, and is trivial or central at the scale
          considered.
  \end{enumerate}
\end{definition}

Thus $\GZTLC_{\mathrm{strict}}$ is a genuine triple-layer connection
whose geometric and parameter components reproduce the strict PFB-GNLA
geometry, and whose law-component is the restriction of the global
law-connection to the strict slice. The “strictness” assumption ensures
that the law-curvature produced by $\GZTLC_{\mathrm{strict}}$ is
harmless from the viewpoint of the strict GNLA structure.



\subsubsection*{(c) Matching diagrams: geometry, operators, law-space, and \GZTLC}

We can now summarise the matching between the strict PFB-GNLA datum,
the strict slice of \GZLS, and the strict triple-layer connection via a
commutative diagrammatic picture.

On the \emph{geometry–operator–law} side, we have:
\[
   (\Bundle,\Base,\Group;\{\Conn_\lambda\})
   \xrightarrow{\ \Phi_{\mathrm{geom}}\ }
   L_{\mathrm{strict}}
   \xrightarrow{\ \iota_{\mathrm{law}}\ }
   {\GZLS}_{\mathrm{strict}}\subset\GZLS,
\]
and
\[
   L_{\mathrm{strict}}
   \xrightarrow{\ \mathsf{Op}\ }
   \GAlgebra_{\mathrm{strict}}
   =
   (A_{\mathrm{strict}},[\cdot,\cdot]_{\mathrm{strict}},
    \Jac_{\mathrm{strict}};\mathsf{Gr},\mathsf{Filt}),
\]
with $A_{\mathrm{strict}}$ embedded in
$\mathfrak{L}_{\mathrm{strict}}$ via $\iota_{\mathrm{law}}$.

On the \emph{triple-layer} side, we have:
\[
   (\Bundle,\Base,\Group;\Lambda_{\mathrm{par}};\mathfrak{L}_{\mathrm{strict}})
   \longrightarrow
   \widehat{\Bundle}_{\mathrm{strict}}
   \xrightarrow{\ \GZTLC_{\mathrm{strict}}\ }
   \text{horizontal flows in }X_{\mathrm{strict}},
\]
and the following compatibility:

\begin{proposition}[Matching of horizontal flows (schematic)]
Let $\gamma(t)=(x(t),\lambda(t),L(t))$ be a smooth curve in
  $X_{\mathrm{strict}}$ whose law-component is horizontal for
  ${\GZLOC}_{\mathrm{strict}}$ and whose geometric–parameter component
  is horizontal for $\widehat{\Conn}$. Then:
  \begin{enumerate}[label=\roman*)]
    \item the $\GZTLC_{\mathrm{strict}}$-horizontal lift
          $\widehat{\gamma}(t)$ induces a family of connections
          $\Conn_{\lambda(t)}$ and associated operators in
          $A_{\mathrm{strict}}$ such that
          \[
             \Phi_{\mathrm{geom}}
               (\Bundle,\Base,\Group;\Conn_{\lambda(t)})
             \simeq
             L(t)
          \]
          as law-objects in ${\GZLS}_{\mathrm{strict}}$;
    \item conversely, any strict PFB-GNLA deformation of
          $(\Bundle,\Conn_\lambda)$ along $\Lambda_{\mathrm{par}}$
          that preserves the strict GNLA constraints is realised as the
          projection of a $\GZTLC_{\mathrm{strict}}$-horizontal flow
          in $X_{\mathrm{strict}}$.
  \end{enumerate}
  A complete proof is given in Technical Appendix~I
  (Section~\ref{app:I-P4}).
\end{proposition}

\begin{proof}[Proof sketch]
Condition (i) follows from the definition of a GZ-LS–compatible
  triple-layer connection (Section~2.2) specialised to the strict
  slice, together with the fact that $L_{\mathrm{strict}}$ and
  $\GAlgebra_{\mathrm{strict}}$ are linked by
  $\Phi_{\mathrm{geom}}$ and $\mathsf{Op}$. Condition (ii) uses the
  existence of strict law–geometry couplings $\mathcal{C}_{\mathrm{strict}}$
  and the fact that strict GNLA constraints ensure that any admissible
  deformation can be lifted to a triple-layer horizontal flow with
  law-curvature confined to $J_{\mathrm{curv}}^{\mathrm{law}}$. Details
  are standard in the triple-layer framework and omitted.
\end{proof}



\subsubsection*{(d) Strict PFB-GNLA as a chart of \GZGMS}

From the vantage point of the GaoZheng Generalized Mathematical
Structure
\[
   \GZGMS =
   \bigl(
     \mathsf{Obj},\mathsf{Law},\mathsf{Flow},\mathsf{Layer},
     \GZTLC;\GZLCurv
   \bigr),
\]
the strict PFB-GNLA datum and its matching with
${\GZLS}_{\mathrm{strict}}$ and $\GZTLC_{\mathrm{strict}}$ can be
rephrased as follows.

\begin{itemize}
  \item In $\mathsf{Obj}$, we restrict to those object-level geometries
        that are principal bundles with parameterized connections
        satisfying the strict GNLA constraints.
  \item In $\mathsf{Law}$, we restrict to the strict slice
        ${\GZLS}_{\mathrm{strict}}$ and to law-objects realised by strict
        PFB-GNLA data.
  \item In $\mathsf{Flow}$, we restrict to flows that are horizontal
        for $\GZTLC_{\mathrm{strict}}$ and remain inside
        $X_{\mathrm{strict}}$; law-curvature along these flows is flat
        or central in the sense of $J_{\mathrm{curv}}^{\mathrm{law}}$.
  \item In $\mathsf{Layer}$, we work primarily at extensional or
        low-homotopy layers (schematically $l_2$), where the strict Lie
        quotient of $A_{\mathrm{strict}}$ suffices to describe all
        relevant law-level invariants.
\end{itemize}

In this sense, the strict PFB-GNLA together with
$({\GZLS}_{\mathrm{strict}},\GZTLC_{\mathrm{strict}})$ defines a
\emph{chart} of \GZGMS in which principal-bundle geometry appears as a
Lie-type law-space structure with controlled curvature, and where the
homotopy and operator-path refinements are not yet activated.



\subsubsection*{(e) Interface to non-strict regimes and future upgrades}

Matching the strict version with \GZLS\ and \GZTLC\ also clarifies how
non-strict regimes arise: they correspond precisely to moving out of
the strict slice ${\GZLS}_{\mathrm{strict}}$ and to turning on
non-trivial law-curvature components.

\begin{remark}[Leaving the strict slice]
When one perturbs $\GZTLC_{\mathrm{strict}}$ so that the law-curvature
  $\GZLCurv$ develops components outside $J_{\mathrm{curv}}^{\mathrm{law}}$,
  or when one allows deformations that violate the strict GNLA
  constraints (S1)–(S4), the law-objects $L(t)$ move out of
  ${\GZLS}_{\mathrm{strict}}$ into more general regions of \GZLS. In
  algebraic terms, the Jacobiator of the operator algebra ceases to be
  confined to a negligible or central sector; homotopy corrections
  (higher brackets) must be introduced, and the appropriate framework
  becomes that of homotopy PFB-GNLA and variable functional–operator
  PFB-GNLA, treated in subsequent chapters.
\end{remark}

From the perspective of Gap statements and certificate-type invariants
(Section~2.3), the strict slice ${\GZLS}_{\mathrm{strict}}$ is
characterised by vanishing or minimal values of certain certificates,
such as curvature-based invariants and Jacobiator-gap certificates.
Matching the strict version with \GZLS\ and \GZTLC\ makes this precise:
strict PFB-GNLA corresponds to those triple-layer systems for which
\[
   \mathsf{Cert}_{\mathrm{curv}}(S_{\mathrm{strict}}) \approx 0,
   \qquad
   \mathsf{JacGap}(S_{\mathrm{strict}}) = 0,
\]
up to the controlled curvature sector. Moving away from the strict
slice is then equivalent to turning on non-zero certificate components,
which the homotopy and operator-path versions are designed to handle.



\medskip

In summary, matching the strict PFB-GNLA with the GaoZheng Law-Space
\GZLS\ and the triple-layer connection \GZTLC\ identifies the strict
version as:

\begin{itemize}
  \item a \emph{law-space chart} ${\GZLS}_{\mathrm{strict}}$ where law
        curvature is flat or central and GNLA reduces to a near-Lie
        structure; and
  \item a \emph{triple-layer chart} $(X_{\mathrm{strict}},
        \widehat{\Bundle}_{\mathrm{strict}},\GZTLC_{\mathrm{strict}})$
        where horizontal flows implement all admissible strict
        PFB-GNLA deformations.
\end{itemize}

This chart serves as the baseline against which the homotopy and
variable functional–operator PFB-GNLA constructions in the next
chapters will be compared: they will be obtained by relaxing the
strictness conditions and by allowing non-trivial law-curvature and
homotopy data to propagate along \GZTLC, thereby enriching the law-level
image of principal-bundle geometry in the GaoZheng G-Framework.



\chapter[Homotopy Version of the GaoZheng G-Algebra]{Homotopy Version of the GaoZheng G-Algebra:\\
         From the Jacobiator to \texorpdfstring{$H$}{H}-Twisted Structures}


\section{The Jacobiator and Higher Homotopies}


The strict PFB-GNLA version organises curvature and parameter
dependence into an operator GNLA whose Jacobiator is confined to a
controlled curvature sector and becomes negligible or central at the
law-level scale of interest. The homotopy version abandons this
“negligible” assumption: the Jacobiator is promoted to a primary piece
of data, and higher homotopies are introduced so that failures of
strict Jacobi identities are absorbed into an $L_\infty$-type
structure. This section formalises this transition, preparing for the
$H$-twisted two-term structures and GZ-Harmony discussed in
Section~5.2.



\subsubsection*{(a) Jacobiator in the strict operator algebra}

We recall the strict operator space
$A_{\mathrm{strict}}$ and its bracket
$[\cdot,\cdot]_{\mathrm{strict}}$ from Section~4.2. For any operators
$T_1,T_2,T_3\in A_{\mathrm{strict}}$, the Jacobiator is
\[
  \Jac_{\mathrm{strict}}(T_1,T_2,T_3)
  :=
  [T_1,[T_2,T_3]_{\mathrm{strict}}]_{\mathrm{strict}}
 +[T_2,[T_3,T_1]_{\mathrm{strict}}]_{\mathrm{strict}}
 +[T_3,[T_1,T_2]_{\mathrm{strict}}]_{\mathrm{strict}}.
\]

By construction, $\Jac_{\mathrm{strict}}$ takes values in a curvature
sector $J_{\mathrm{curv}}\subseteq A_{\mathrm{strict}}$; the strict GNLA
constraints ensure that either $J_{\mathrm{curv}}=0$ or that it is
central/negligible up to a prescribed order. Concretely, as soon as
curvature is non-zero and parameter dependence is present, the
Jacobiator is typically non-vanishing on operators built from
covariant derivatives and curvature insertions.

A schematic example is given by commutators of covariant derivatives with parameter \\ dependence:
\[
  \bigl[\nabla^{(\lambda)}_X,\,
        [\nabla^{(\lambda)}_Y,\nabla^{(\lambda)}_Z]\bigr]
 +\text{cyclic permutations in }(X,Y,Z),
\]
which, once resolved by standard differential-geometric identities,
contains \\ derivative-of-curvature terms, mixed curvature components, and
higher-order contributions in $\widehat{\Curv}$. These contributions
constitute a typical source of Jacobiator terms that cannot be ignored
once higher-order effects become relevant (e.g.\ in strongly coupled or
multi-scale regimes).



\subsubsection*{(b) Jacobiator as a shadow of law-curvature}

From the law-level viewpoint, the Jacobiator in
$A_{\mathrm{strict}}$ is not arbitrary: it is tightly related to the
law-curvature \GZLCurv\ of the law-connection \GZLOC\ on \GZLS. In the
triple-layer picture, the curvature of \GZTLC\ decomposes into
geometric, parameter, and law components; projecting to the law-layer
one obtains
\[
   \GZLCurv
   =
   \pi_{\mathrm{law}}\bigl(\mathcal{F}_{\GZTLC}\bigr).
\]

In the homotopy regime, we regard $\Jac_{\mathrm{strict}}$ as the
operator-level shadow of certain components of \GZLCurv: schematically,
there exists an assignment
\[
   \mathcal{K}:
     \GZLCurv
     \longrightarrow
     \mathrm{Hom}\bigl(
       A_{\mathrm{strict}}^{\otimes3},A_{\mathrm{strict}}
     \bigr)
\]
such that
\[
   \Jac_{\mathrm{strict}}
   \approx
   \mathcal{K}(\GZLCurv),
\]
up to reparametrisation and lower-order terms. Intuitively, when the
law-connection \GZLOC\ has non-trivial curvature along certain
directions in law-space, transporting operator bases around small
triangles in the triple-layer base produces a “holonomy defect” that
shows up as a non-zero Jacobiator for the induced operator algebra.

This observation motivates the homotopy upgrade: instead of forcing
$\Jac_{\mathrm{strict}}$ to vanish (or live in a negligible sector),
one includes it as part of a higher operation (a ternary bracket
$\ell_3$), and tracks its relation to \GZLCurv\ via homotopy identities.



\subsubsection*{(c) Promoting the Jacobiator to a homotopy operation}

Let $A_{\mathrm{homo}}$ denote an enlargement of $A_{\mathrm{strict}}$
suitable for hosting an $L_\infty$-type structure (for example by
adding auxiliary degrees or by splitting curvature sectors into lower
and higher grades). We set up a graded vector space
\[
   V
   =
   V^{-1}\oplus V^0
   \;\oplus\;\text{(optional higher degrees)},
\]
with the following identifications in the simplest two-term setting:

\begin{itemize}
  \item $V^0$ is (a completion of) $A_{\mathrm{strict}}$, the space of
        degree-zero operators (covariant derivatives, curvature
        insertions, etc.);
  \item $V^{-1}$ is a “defect space” that receives Jacobiator terms
        and other controlled failures (e.g.\ curvature differences
        between law- and geometry-induced structures).
\end{itemize}

We then introduce the following structure maps:
\begin{itemize}
  \item $\ell_1:V^{-1}\to V^0$ encodes the way defects can be “pushed
        back” into the operator layer via differentiation or
        cancellation;
  \item $\ell_2:V^i\otimes V^j\to V^{i+j}$ extends the strict
        commutator $[\cdot,\cdot]_{\mathrm{strict}}$ (possibly with
        curvature corrections) to a graded bracket;
  \item $\ell_3:V^0\otimes V^0\otimes V^0\to V^{-1}$ is defined so
        that
        \[
           \ell_3(T_1,T_2,T_3)
           \approx
           \Jac_{\mathrm{strict}}(T_1,T_2,T_3)
        \]
        up to terms in the image of $\ell_1$ and law-curvature
        projections.
\end{itemize}

\begin{definition}[Homotopy realisation of the Jacobiator (schematic)]
A \emph{homotopy realisation} of the strict Jacobiator is a choice
  of $(V,\ell_1,\ell_2,\ell_3)$ as above such that:
  \begin{enumerate}[label=\roman*)]
    \item $\ell_2$ restricts to $[\cdot,\cdot]_{\mathrm{strict}}$ on
          $A_{\mathrm{strict}}\subseteq V^0$;
    \item the graded Jacobi identity for $\ell_2$ is satisfied up to
          a boundary
          \[
             \Jac_{\ell_2}(T_1,T_2,T_3)
             =
             \ell_1\bigl(\ell_3(T_1,T_2,T_3)\bigr),
          \]
          for all $T_i\in V^0$, possibly plus higher corrections if
          degrees $>0$ are present;
    \item $\ell_3$ is related to the law-curvature via a
          curvature–homotopy map
          \[
             \ell_3
             \sim
             \mathcal{K}(\GZLCurv)
             \mod \mathrm{Im}(\ell_1),
          \]
          as in the previous subsection.
  \end{enumerate}
\end{definition}

The condition $\Jac_{\ell_2}=\ell_1\circ\ell_3$ says that the failure
of $\ell_2$ to satisfy the Jacobi identity is “homotopically
trivial”: it is the boundary of $\ell_3$. Operationally, this provides
a consistent way to keep track of Jacobiator effects without breaking
the higher coherence of the algebraic structure.



\subsubsection*{(d) Higher homotopies and iterated curvature effects}

While a two-term structure $(V^{-1},V^0)$ and a single $\ell_3$ already
captures the basic Jacobiator, more refined regimes may require higher
homotopies $\ell_k$ for $k\ge4$. Schematically, these encode:
\begin{itemize}
  \item higher-order curvature effects (e.g.\ curvature of law-curvature
        itself along law-space loops);
  \item interactions between different sectors of the operator algebra
        (e.g.\ mixing between geometric, parameter, and law-induced
        operators);
  \item corrections arising from CH-layering transitions, where the
        relevant metric shifts from extensional to semantic/homotopy
        based.
\end{itemize}

The general $L_\infty$ identities then impose a tower of coherence
conditions of the form
\[
   \sum_{i+j=n+1}
   \sum_{\sigma\in S(i,n-i)}
   \chi(\sigma)(-1)^{i(j-1)}
   \ell_j\bigl(\ell_i(\cdots),\cdots\bigr)
   = 0,
   \qquad n\ge1,
\]
which ensure that all these curvature and mixing effects fit together
consistently. In the homotopy PFB-GNLA version, we will typically keep
only a small finite number of non-trivial $\ell_k$ (often up to
$\ell_3$ or $\ell_4$), reflecting the fact that only a bounded set of
higher interactions are relevant for the law-level phenomena of
interest.



\subsubsection*{(e) Jacobiator-gap and homotopy completion (preview)}

The homotopy interpretation of the Jacobiator leads naturally to the
Jacobiator-gap certificates and the GZ-OHU theorem sketched in Part~III.
At a high level:

\begin{itemize}
  \item the magnitude and structure of $\Jac_{\mathrm{strict}}$ (or
        equivalently of $\ell_3$) determine a \emph{Jacobiator-gap}
        vector $\mathsf{JacGap}(L)$ for a given law-level operator
        system $L$;
  \item if $\mathsf{JacGap}(L)=0$, the homotopy data can be “collapsed”
        to yield a strictly Jacobi-satisfying structure without changing
        low-order observables;
  \item if $\mathsf{JacGap}(L)\neq0$, any attempt to deform $L$ into a
        strict regime must cross a gap-region in law-space where
        curvature and homotopy invariants remain bounded away from zero.
\end{itemize}

The homotopy PFB-GNLA construction thus serves a dual purpose:

\begin{enumerate}[label=\alph*)]
  \item it provides a \emph{positive} framework in which non-trivial
        Jacobiators and higher law-curvature effects are consistently
        encoded as homotopies $(\ell_3,\ell_4,\dots)$, rather than
        treated as “defects”;
  \item it provides, via $\mathsf{JacGap}$ and related certificates,
        a \emph{quantitative} framework for stating Gap theorems about
        which law-level evolutions can or cannot be completed to strict
        PFB-GNLA regimes.
\end{enumerate}



\subsubsection*{(f) Summary: from strict Jacobi to homotopy-controlled Jacobi}

To summarise, the role of the Jacobiator and higher homotopies in the
homotopy version of the GaoZheng G-Algebra is as follows:

\begin{itemize}
  \item The Jacobiator $\Jac_{\mathrm{strict}}$ of the strict operator
        algebra is recognised as an operator-level imprint of
        law-curvature \GZLCurv\ in \GZLS, arising from triple-layer
        holonomy effects.
  \item Instead of forcing $\Jac_{\mathrm{strict}}$ to vanish (or
        hiding it in a negligible sector), we enlarge the operator
        space to $V$ and promote $\Jac_{\mathrm{strict}}$ to a ternary
        homotopy operation $\ell_3$, accompanied by a differential
        $\ell_1$ and a graded bracket $\ell_2$.
  \item The higher homotopies $(\ell_k)_{k\ge3}$ organise iterated
        curvature effects and multi-sector couplings into an
        $L_\infty$-type structure, ensuring global coherence of the
        homotopy PFB-GNLA.
  \item Jacobiator-gap certificates and related invariants measure
        how far a homotopy PFB-GNLA lies from the strict regime and
        constrain possible deformations between homotopy and strict
        law-level structures.
\end{itemize}

The next section will specialise this general homotopy picture to the
case of $H$-twisted two-term $L_\infty$ structures and introduce the
notion of GZ-Harmony: a compatibility condition between curvature-type
3-forms (such as $H$ and law-curvature components) and the homotopy
brackets $(\ell_1,\ell_2,\ell_3)$ that will play a central role in the
homotopy version of the GaoZheng G-Algebra.



\section[H-Twisted Two-Term L-infinity Structures and GZ-Harmony]{\texorpdfstring{$H$}{H}-Twisted Two-Term $L_\infty$ Structures and GZ-Harmony}


The homotopy interpretation of the Jacobiator in Section~5.1 motivates
a concrete class of $L_\infty$-structures in which the ternary bracket
$\ell_3$ is controlled by a closed 3-form $H$ and by law-curvature
components in \GZLCurv. These are the $H$-twisted two-term
$L_\infty$-structures. In the context of the GaoZheng G-Framework, such
structures become meaningful only when they are compatible with the
triple-layer connection \GZTLC\ and with the semantic layering encoded
in \GZGMS. The name \emph{GZ-Harmony} will be used for this compatibility
condition.



\subsubsection*{(a) Geometric $H$-data and its law-level shadow}

We begin by specifying the geometric $H$-data that will twist the
two-term $L_\infty$-structure. Let $(\Bundle,\Base,\Group)$ be a
principal bundle with parameterized connection as in Part~IV, and let
$H$ denote a closed 3-form on $\Base$ (or more generally a closed
3-form on a suitable configuration space built from $\Base$ and
$\Lambda_{\mathrm{par}}$):

\[
   H \in \Omega^3(\Base),
   \qquad
   \mathrm{d}H = 0.
\]

In many concrete situations, $H$ may originate from:

\begin{itemize}
  \item a Chern–Simons-type 3-form constructed from $\Conn_\lambda$ and
        $\Curv_\lambda$;
  \item a law-curvature 3-form extracted from \GZLCurv\ via the
        curvature–homotopy map $\mathcal{K}$ of Section~5.1;
  \item a semantic or CH-layering 3-form that measures transitions
        between layers $l_2$ and $l_{\ge3}$.
\end{itemize}

At the law-level, the same data is registered as a component
\[
   H_{\mathrm{law}}
   \in
   \Omega^3(\mathfrak{L}_{\mathrm{strict}})
\]
(or on a suitable law-bundle) which enters the law-curvature
\GZLCurv\ and thus feeds directly into the Jacobiator and higher
homotopies in the operator algebra.



\subsubsection*{(b) $H$-twisted two-term $L_\infty$-structures}

We now specialise the general homotopy picture of Section~5.1 to
two-term $L_\infty$-structures twisted by $H$.

\begin{definition}[$H$-twisted two-term $L_\infty$-structure (schematic)]
An \emph{$H$-twisted two-term $L_\infty$-structure} consists of:
  \begin{itemize}
    \item a $\mathbb{Z}$-graded vector space
          \(
             V = V^{-1}\oplus V^0
          \),
          with all other degrees trivial;
    \item a degree $+1$ differential
          \(
             \ell_1:V^{-1}\to V^0
          \);
    \item a degree $0$ graded antisymmetric bracket
          \(
             \ell_2:V^i\otimes V^j\to V^{i+j}
          \)
          extending a given commutator on $V^0$;
    \item a degree $-1$ trilinear map
          \(
             \ell_3^H:V^0\otimes V^0\otimes V^0\to V^{-1},
          \)
          called the \emph{$H$-twisted Jacobiator},
  \end{itemize}
  such that the $L_\infty$ identities reduce to:
  \begin{enumerate}[label=\roman*)]
    \item $\ell_1^2=0$;
    \item $\ell_1$ is a derivation of $\ell_2$;
    \item the Jacobi identity for $\ell_2$ is twisted by $\ell_3^H$:
          for all $x,y,z\in V^0$,
          \[
             \Jac_{\ell_2}(x,y,z)
             =
             \ell_1\bigl(\ell_3^H(x,y,z)\bigr);
          \]
    \item $\ell_3^H$ satisfies a coherence condition which depends on
          $H$, typically of the schematic form
          \[
             \mathrm{d}H = 0
             \quad\Longrightarrow\quad
             \text{higher } n=4 \text{ identity for }
             \ell_2,\ell_3^H \text{ holds.}
          \]
  \end{enumerate}
\end{definition}

The last condition encodes that the twist by $H$ does not break the
$L_\infty$ coherence: the closedness of $H$ is mirrored by a higher
homotopy identity involving $\ell_3^H$ and $\ell_2$. Concrete versions
of this phenomenon appear in many known examples (string Lie 2-algebras,
Courant algebroids, $H$-twisted Poisson structures), and the
PFB-GNLA homotopy version can be seen as a law-level generalisation of
this pattern.

In our setting, $V^0$ will be an operator layer (an enlargement of
$A_{\mathrm{strict}}$), while $V^{-1}$ records controlled defects
(Jacobiator and curvature differences) that are “resolved” by
$\ell_1$.



\subsubsection*{(c) Coupling $H$ with law-curvature and Jacobiator}

The $H$-dependence of $\ell_3^H$ is not arbitrary: it must be
compatible with the law-curvature \GZLCurv\ and with the Jacobiator of
the underlying operator algebra. Schematically, we require a map
\[
   \mathcal{K}_H:
     \bigl(H,\GZLCurv\bigr)
     \longrightarrow
     \mathrm{Hom}\bigl(
       V^{0\otimes3},V^{-1}
     \bigr)
\]
such that:

\begin{enumerate}[label=(C\arabic*)]
  \item \emph{Jacobiator realisation.}
        For $x,y,z\in V^0$ corresponding to operators in
        $A_{\mathrm{strict}}$, we have
        \[
           \ell_3^H(x,y,z)
           \approx
           \mathcal{K}_H\bigl(H,\GZLCurv\bigr)(x,y,z),
        \]
        and
        \[
           \Jac_{\ell_2}(x,y,z)
           =
           \ell_1\bigl(\ell_3^H(x,y,z)\bigr),
        \]
        recovering the Jacobiator as a boundary term.
  \item \emph{Curvature compatibility.}
        Under triple-layer parallel transport by \GZTLC, the
        combination
        \[
           (H,\GZLCurv,\ell_3^H)
        \]
        transforms covariantly; in particular, if $H$ and \GZLCurv\ are
        transported horizontally along a closed loop in
        $X_{\mathrm{tot}}$, the change in $\ell_3^H$ is measured by
        higher homotopies $\ell_k$ for $k\ge4$, consistent with the
        $L_\infty$ identities.
\end{enumerate}

Condition (C1) ties $\ell_3^H$ to both the geometric 3-form $H$ and the
law-curvature \GZLCurv; (C2) ensures that this coupling behaves well
under the triple-layer geometry, so that homotopy structures are stable
under admissible flows.



\subsubsection*{(d) GZ-Harmony: compatibility of $H$, \GZTLC, and the $L_\infty$-structure}

We are now ready to formulate the notion of \emph{GZ-Harmony}. This is a
compatibility condition between:

\begin{itemize}
  \item the geometric $H$-data and its law-level shadow;
  \item the triple-layer connection \GZTLC\ (and its law-curvature
        \GZLCurv);
  \item the $H$-twisted two-term $L_\infty$-structure
        $(V^{-1}\oplus V^0,\ell_1,\ell_2,\ell_3^H)$.
\end{itemize}

\begin{definition}[GZ-Harmony (schematic)]
An $H$-twisted two-term $L_\infty$-structure
  $(V^{-1}\oplus V^0,\ell_1,\ell_2,\ell_3^H)$ on an operator layer
  associated with a PFB-GNLA is said to be in \emph{GZ-Harmony} with
  $(H,\GZTLC)$ if:
  \begin{enumerate}[label=\roman*)]
    \item (\emph{Curvature–homotopy matching})
          There exists a curvature–homotopy map
          \[
             \mathcal{K}_H:
               (H,\GZLCurv)
               \longrightarrow
               \ell_3^H
          \]
          such that conditions (C1)–(C2) above hold, and the closedness
          $\mathrm{d}H=0$ implies the $n=4$ $L_\infty$ identity for
          $(\ell_2,\ell_3^H)$.
    \item (\emph{Layer compatibility})
          The CH-layering data $\mathsf{Layer}$ in \GZGMS\ is such that:
          \begin{itemize}
            \item at extensional layers $l_2$, the contribution of
                  $\ell_3^H$ is invisible in the effective law-metrics
                  $J$ (i.e.\ it contributes only to higher-order
                  corrections);
            \item at higher layers $l_{\ge3}$, the metrics include
                  contributions from $\ell_3^H$ and from $H$, so that
                  semantic/homotopy metrics are sensitive to the
                  $H$-twist.
          \end{itemize}
    \item (\emph{Triple-layer horizontality})
          For any $\GZTLC$-horizontal flow in $X_{\mathrm{tot}}$, the
          evolution of $(H,\GZLCurv,\ell_3^H)$ is governed by the
          $L_\infty$ identities: infinitesimal changes in $H$ along
          law-level directions correspond to higher homotopies
          involving $\ell_3^H$ and possibly $\ell_4$, but do not
          generate inconsistencies in the underlying operator algebra.
  \end{enumerate}
\end{definition}

In short, GZ-Harmony states that “what $H$ and \GZLCurv\ do
geometrically” is exactly matched by “what $\ell_3^H$ and the higher
homotopies do algebraically”, and that this match respects both the
triple-layer geometry and the CH-layering semantics.



\subsubsection*{(e) Role of GZ-Harmony in the homotopy PFB-GNLA}

Within the homotopy version of the GaoZheng G-Algebra, the existence
of GZ-Harmony has several consequences.

\begin{itemize}
  \item \emph{Well-posed homotopy upgrade.}
        Starting from a strict PFB-GNLA datum, turning on a non-trivial
        3-form $H$ and law-curvature \GZLCurv\ that satisfy
        GZ-Harmony ensures that the resulting operator algebra admits
        an $H$-twisted two-term $L_\infty$-upgrade; Jacobiator effects
        are consistently captured by $\ell_3^H$.
  \item \emph{Controlled deviation from strictness.}
        The size and structure of $\ell_3^H$ (or equivalently of
        $\mathsf{JacGap}$) measure the deviation from the strict
        PFB-GNLA regime; GZ-Harmony guarantees that this deviation is
        geometrically meaningful (controlled by $H$ and \GZLCurv) and
        algebraically coherent.
  \item \emph{Compatibility with law-space flows.}
        Because GZ-Harmony is formulated in terms of \GZTLC\ and
        \GZLCurv, the homotopy PFB-GNLA is stable under law-space flows
        and CH-layer transitions: $H$-twisted homotopy structures are
        transported consistently along admissible curves in \GZLS.
\end{itemize}

From the viewpoint of the overall monograph, GZ-Harmony is the key
bridge between:

\begin{enumerate}[label=\alph*)]
  \item the “geometric side”: principal bundles, connections, 3-forms
        $H$, and triple-layer curvature; and
  \item the “homotopy–algebraic side”: $H$-twisted two-term
        $L_\infty$-structures on operator layers, with Jacobiator and
        higher homotopies as primary data.
\end{enumerate}



\subsubsection*{(f) Summary and preparation for the homotopy version of \texorpdfstring{\GAlgebra}{G-Algebra}}

To summarise:

\begin{itemize}
  \item $H$-twisted two-term $L_\infty$-structures provide a concrete
        homotopy framework in which the Jacobiator is encoded by a
        ternary operation $\ell_3^H$ controlled by a closed 3-form $H$
        and by law-curvature components \GZLCurv.
  \item GZ-Harmony is the compatibility condition ensuring that the
        triple $(H,\GZTLC,\ell_\bullet)$ is coherent across geometry,
        law-space, and homotopy algebra; it enforces curvature–homotopy
        matching, CH-layer compatibility, and triple-layer horizontality.
  \item Under GZ-Harmony, the homotopy version of the GaoZheng
        G-Algebra can be constructed as an $H$-twisted two-term
        $L_\infty$-algebra on an operator layer associated with
        PFB geometry, with Jacobiator-gap certificates and Gap
        statements becoming natural invariants of the law-level
        structure.
\end{itemize}

The next section will build on this foundation to present the homotopy
version of \GAlgebra\ as an integrated PFB-GNLA construction: starting
from strict data, turning on $H$ and law-curvature in a GZ-Harmonised
way, and deriving a homotopy GaoZheng G-Algebra that sits inside the
GaoZheng G-Framework as the “Jacobiator-aware” counterpart of the
strict version.



\section[Homotopy Version of the G-Algebra: Integrated Construction]{Homotopy Version of \GAlgebra as an Integrated Construction}


The homotopy version of the GaoZheng G-Algebra upgrades the strict
PFB-GNLA by turning non-trivial Jacobiators and law-curvature effects
into primary homotopy data. Concretely, this means passing from a
near-Lie operator algebra $(A_{\mathrm{strict}},[\cdot,\cdot]_{\mathrm{strict}})$
to an $H$-twisted two-term $L_\infty$-structure
$(V^{-1}\oplus V^0,\ell_1,\ell_2,\ell_3^H)$ that is in GZ-Harmony with
the triple-layer connection \GZTLC\ and the law-curvature \GZLCurv. In
this section we assemble these elements into an integrated homotopy
PFB-GNLA datum and position it inside the GaoZheng G-Framework.



\subsubsection*{(a) Homotopy geometric–parameter–$H$ block}

We start by extending the geometric–parameter block of Section~4.2 to
include $H$-data relevant for the homotopy twist.

\begin{definition}[Homotopy geometric–parameter–$H$ block]
A \emph{homotopy geometric–parameter–$H$ block} for the homotopy
  PFB-GNLA consists of
  \[
     \mathsf{GeomPar}_H
     =
     \bigl(
       \Bundle,\Base,\Group;
       \Lambda_{\mathrm{par}};
       \{\Conn_\lambda\}_{\lambda\in\Lambda_{\mathrm{par}}},
       \widehat{\Conn},\widehat{\Curv};
       H
     \bigr),
  \]
  where:
  \begin{itemize}
    \item $(\Bundle,\Base,\Group)$ and
          $\Lambda_{\mathrm{par}}$, $\{\Conn_\lambda\}$,
          $\widehat{\Conn}$, $\widehat{\Curv}$ are as in
          $\mathsf{GeomPar}$ (Section~4.2);
    \item $H\in\Omega^3(\Base)$ (or in a suitable extension such as
          $\Omega^3(\Base\times\Lambda_{\mathrm{par}})$) is a closed
          3-form,
          \[
             \mathrm{d}H = 0,
          \]
          whose law-level shadow contributes to the law-curvature
          \GZLCurv\ and to the homotopy bracket $\ell_3^H$.
  \end{itemize}
\end{definition}

In practice, $H$ will often be constructed from the connection and
curvature (e.g.\ Chern–Simons-type forms, secondary characteristic
classes, or law-space analogues derived from \GZLCurv). For the
homotopy PFB-GNLA, however, we only require that $H$ be closed and
GZ-Harmonised with \GZTLC\ and law-curvature in the sense of
Section~5.2.



\subsubsection*{(b) Homotopy operator block: $H$-twisted two-term structure}

We now build a homotopy operator block that extends the strict operator
block $\mathsf{Op}_{\mathrm{strict}}$ into an $H$-twisted two-term
$L_\infty$-structure.

\begin{definition}[Homotopy operator block]
A \emph{homotopy operator block} associated with \\
  $\mathsf{GeomPar}_H$ is a tuple
  \[
     \mathsf{Op}_{\mathrm{homo}}
     =
     \bigl(
       V^{-1}\oplus V^0,
       \ell_1,\ell_2,\ell_3^H;
       \mathsf{Gr}_{\mathrm{homo}},\mathsf{Filt}_{\mathrm{homo}}
     \bigr),
  \]
  such that:
  \begin{enumerate}[label=\roman*)]
    \item (\emph{Embedding of the strict operator space})
          There is an injective map
          \[
             \iota_{\mathrm{strict}}:
               A_{\mathrm{strict}}
               \hookrightarrow
               V^0,
          \]
          identifying the strict operator space with a degree-zero
          subspace of $V^0$.
    \item (\emph{$H$-twisted two-term $L_\infty$ structure})
          $(V^{-1}\oplus V^0,\ell_1,\ell_2,\ell_3^H)$ is an
          $H$-twisted two-term $L_\infty$-structure in the sense of
          Section~5.2, with:
          \[
             \ell_2\big|_{A_{\mathrm{strict}}}
             =
             [\cdot,\cdot]_{\mathrm{strict}},
          \]
          up to curvature-controlled corrections compatible with
          $\mathsf{Gr}_{\mathrm{homo}}$ and $\mathsf{Filt}_{\mathrm{homo}}$.
    \item (\emph{Jacobiator realisation})
          For $T_1,T_2,T_3\in A_{\mathrm{strict}}\subseteq V^0$,
          the Jacobiator of $\ell_2$ satisfies
          \[
             \Jac_{\ell_2}(T_1,T_2,T_3)
             =
             \ell_1\bigl(
               \ell_3^H(T_1,T_2,T_3)
             \bigr),
          \]
          and $\ell_3^H(T_1,T_2,T_3)$ agrees, up to $\mathrm{Im}(\ell_1)$,
          with the Jacobiator of the strict commutator and the
          curvature–homotopy map $\mathcal{K}_H(H,\GZLCurv)$.
    \item (\emph{Homotopy filtration compatibility})
          The grading and filtration
          $(\mathsf{Gr}_{\mathrm{homo}},\mathsf{Filt}_{\mathrm{homo}})$
          refine $(\mathsf{Gr},\mathsf{Filt})$ and are compatible with
          the $L_\infty$ operations: for instance,
          \[
             \ell_k\bigl(
               \mathsf{Filt}^pV^{\bullet},\dots,\mathsf{Filt}^qV^{\bullet}
             \bigr)
             \subseteq
             \mathsf{Filt}^{p+\dots+q}V^{\bullet}.
          \]
  \end{enumerate}
\end{definition}

Intuitively, $V^0$ plays the role of an “extended operator layer”
containing the strict operators and their law-driven deformations,
while $V^{-1}$ collects controlled defect modes (Jacobiator components,
curvature mismatches, etc.) that are resolved homotopically via
$\ell_1$ and $\ell_3^H$.



\subsubsection*{(c) Homotopy law-level block and homotopy GZ-LS slice}

On the law-level, the homotopy version lives in a region of law-space
where law-curvature and CH-layering make the $H$-twist and higher
homotopies visible. We denote this region by ${\GZLS}_{\mathrm{homo}}$.

\begin{definition}[Homotopy GZ-LS slice and law-level block]
A \emph{homotopy GZ-LS slice} is a region
  \[
     {\GZLS}_{\mathrm{homo}}
     =
     \bigl(
       \mathfrak{L}_{\mathrm{homo}};
       J_{\mathrm{homo}},
       \mathrm{Loc}_{\mathrm{homo}},
       \Sigma_{\mathrm{homo}}
     \bigr),
  \]
  together with an embedding
  \(
     \iota_{\mathrm{LS}}^{\mathrm{homo}}:
     \mathfrak{L}_{\mathrm{homo}}\hookrightarrow\mathfrak{L}
  \),
  such that:
  \begin{enumerate}[label=\roman*)]
    \item $\mathfrak{L}_{\mathrm{strict}}\subseteq\mathfrak{L}_{\mathrm{homo}}$,
          and $(J_{\mathrm{homo}},\mathrm{Loc}_{\mathrm{homo}},
          \Sigma_{\mathrm{homo}})$ extend the strict metrics and
          localisation data to include homotopy contributions;
    \item the law-connection \GZLOC\ restricted to
          $\mathfrak{L}_{\mathrm{homo}}$ has law-curvature
          ${\GZLCurv}_{\mathrm{homo}}$ whose components feed into the
          $H$-twisted homotopy structure via $\mathcal{K}_H$;
    \item there is a law-object
          \[
             L_{\mathrm{homo}}
             =
             \bigl(
               \mathfrak{L}_{\mathrm{homo}}^{(0)},
               \text{homotopy data}
             \bigr)
             \in{\GZLS}_{\mathrm{homo}},
          \]
          equipped with an embedding
          \[
             \iota_{\mathrm{law}}^{\mathrm{homo}}:
               V^0\hookrightarrow\mathfrak{L}_{\mathrm{homo}}^{(0)},
          \]
          such that $(V^{-1}\oplus V^0,\ell_1,\ell_2,\ell_3^H)$ is
          realised as homotopy law-data for $L_{\mathrm{homo}}$.
  \end{enumerate}
\end{definition}

In effect, $L_{\mathrm{homo}}$ is the homotopy law-object whose
internal algebraic structure is given by the $H$-twisted two-term
$L_\infty$-algebra $\mathsf{Op}_{\mathrm{homo}}$, and
${\GZLS}_{\mathrm{homo}}$ is the region of law-space where such homotopy
objects naturally live.



\subsubsection*{(d) Homotopy PFB-GNLA datum: integrated definition}

We can now define the homotopy version of the PFB-GNLA as an integrated
datum, paralleling the strict definition of Section~4.2.

\begin{definition}[Homotopy PFB-GNLA datum]
A \emph{homotopy PFB-GNLA datum} (homotopy GaoZheng G-Algebra datum)
  is a tuple
  \[
    \mathsf{PFB\text- GNLA}_{\mathrm{homo}}
    =
    \bigl(
      \mathsf{GeomPar}_H,
      \mathsf{Op}_{\mathrm{homo}},
      L_{\mathrm{homo}};
      \Phi_{\mathrm{geom}}^{\mathrm{homo}},
      \mathcal{C}_{\mathrm{homo}}
    \bigr),
  \]
  where:
  \begin{itemize}
    \item $\mathsf{GeomPar}_H$ is a homotopy geometric–parameter–$H$
          block;
    \item $\mathsf{Op}_{\mathrm{homo}}$ is a homotopy operator block
          as above, with $(V^{-1}\oplus V^0,\ell_1,\ell_2,\ell_3^H)$
          in GZ-Harmony with $(H,\GZTLC)$;
    \item $L_{\mathrm{homo}}\in{\GZLS}_{\mathrm{homo}}$ is a homotopy
          law-level block with internal $L_\infty$-structure
          identified with $\mathsf{Op}_{\mathrm{homo}}$;
    \item $\Phi_{\mathrm{geom}}^{\mathrm{homo}}$ is a homotopy
          geometry-to-law functor extending $\Phi_{\mathrm{geom}}$
          such that
          \[
            \Phi_{\mathrm{geom}}^{\mathrm{homo}}
              (\Bundle,\Base,\Group;
               \{\Conn_\lambda\},H;\text{aux})
            \simeq
            L_{\mathrm{homo}}
          \]
          as homotopy law-objects;
    \item $\mathcal{C}_{\mathrm{homo}}$ is a homotopy law–geometry
          coupling that inverts $\Phi_{\mathrm{geom}}^{\mathrm{homo}}$
          up to homotopy and realises the $H$-twist in \GZTLC.
  \end{itemize}
\end{definition}

Informally, a homotopy PFB-GNLA datum packages:

\begin{itemize}
  \item a PFB geometry with parameter dependence and a closed 3-form $H$;
  \item an $H$-twisted two-term $L_\infty$ operator layer on the
        associated bundles;
  \item a homotopy law-object in \GZLS capturing this structure;
  \item a bidirectional (up to homotopy) correspondence between
        geometry and law, consistent with \GZTLC.
\end{itemize}



\subsubsection*{(e) Relation to the strict version and Jacobiator-gap}

The homotopy datum is designed to extend the strict PFB-GNLA datum in a
controlled way. This relationship can be formalised via “upgrade” and
“strictification” maps.

\begin{proposition}[Homotopy upgrade and strictification (schematic)]
Under appropriate regularity and boundedness assumptions, there exist
  functorial constructions:
  \[
     \mathsf{Upgrade}:
       \mathsf{PFB\text- GNLA}_{\mathrm{strict}}
       \longrightarrow
       \mathsf{PFB\text- GNLA}_{\mathrm{homo}},
  \]
  and
  \[
     \mathsf{Strictify}:
       \mathsf{PFB\text- GNLA}_{\mathrm{homo}}
       \dashrightarrow
       \mathsf{PFB\text- GNLA}_{\mathrm{strict}},
  \]
  satisfying:
  \begin{enumerate}[label=\roman*)]
    \item $\mathsf{Upgrade}$ embeds strict data as homotopy data with
          trivial or controlled $H$ and $\ell_3^H$ (e.g.\ $H=0$,
          $\ell_3^H=0$), recovering the strict PFB-GNLA as a special
          case of the homotopy version.
    \item $\mathsf{Strictify}$ is defined when the Jacobiator-gap
          certificate $\mathsf{JacGap}(L_{\mathrm{homo}})$ vanishes:
          in that case, there exists a homotopy equivalence between
          $L_{\mathrm{homo}}$ and a strict law-object
          $L_{\mathrm{strict}}$, and the corresponding
          $\mathsf{PFB\text- GNLA}_{\mathrm{homo}}$ can be collapsed to
          a strict PFB-GNLA datum.
    \item In general, if $\mathsf{JacGap}(L_{\mathrm{homo}})\neq0$,
          any strictification homotopy must cross regions of law-space
          where curvature and homotopy invariants remain bounded away
          from zero, as encoded by Gap statements.
  \end{enumerate}
  Further details and a complete proof are given in Technical Appendix~II
  (Section~\ref{app:II-P5}).
\end{proposition}

The details of these constructions, and the precise statement of the
GZ-OHU theorem, will be developed in Part~III. For the present section,
the key point is that the homotopy PFB-GNLA is not an arbitrary
relaxation: it sits in a controlled envelope around the strict regime,
with Jacobiator-gap certificates measuring how far one has moved away
from strictness.



\subsubsection*{(f) Embedding into \texorpdfstring{\GZLS}{GZ-LS}, \texorpdfstring{\GZTLC}{GZ-TLC}, and \texorpdfstring{\GZGMS}{GZ-GMS}}

Finally, we summarise how the homotopy PFB-GNLA datum integrates into
the GaoZheng G-Framework.

\begin{itemize}
  \item In \textbf{$\mathsf{Obj}$}, homotopy PFB-GNLA restricts to
        principal bundles with parameterised connections and an
        attached closed 3-form $H$, satisfying geometric conditions
        required by GZ-Harmony.
  \item In \textbf{$\mathsf{Law}$}, it occupies the slice
        ${\GZLS}_{\mathrm{homo}}$ of law-space, where law-objects are
        endowed with internal $H$-twisted $L_\infty$ structures and
        where law-curvature \GZLCurv\ has non-trivial components
        relevant for homotopy.
  \item In \textbf{$\mathsf{Flow}$}, it uses the triple-layer
        connection \GZTLC\ to transport $(H,\GZLCurv,\ell_\bullet)$
        along admissible flows, ensuring that homotopy structures are
        preserved up to higher homotopies along the flow.
  \item In \textbf{$\mathsf{Layer}$}, it lives primarily at higher
        CH-layers $l_{\ge3}$, where semantic/homotopy metrics are
        sensitive to $\ell_3^H$ and to the $H$-twist; the strict regime
        corresponds to a lower-layer projection where $\ell_3^H$ is
        effectively invisible.
\end{itemize}

Within \GZGMS, the homotopy PFB-GNLA thus provides a refined chart
above the strict chart of Section~4.3: the same geometric substrate is
now viewed through a homotopy-sensitive law-space lens, where
Jacobiator and higher curvature effects are modelled as $L_\infty$
operations subject to GZ-Harmony.



\medskip

In conclusion, the homotopy version of the GaoZheng G-Algebra is an
integrated construction that:

\begin{itemize}
  \item starts from PFB geometry with parameter dependence and a closed
        3-form $H$;
  \item builds an $H$-twisted two-term $L_\infty$ operator algebra in
        GZ-Harmony with \GZTLC\ and \GZLCurv;
  \item realises this algebra as a homotopy law-object in a dedicated
        slice of \GZLS;
  \item connects continuously to the strict PFB-GNLA regime via
        Jacobiator-gap–controlled upgrades and strictifications.
\end{itemize}

The subsequent chapter will extend this pattern to the variable
functional–operator homotopy version, where path spaces and operator
flows become explicit objects, and the triple-layer connection \GZTLC\
acts directly on law–operator dynamics in a way compatible with the
homotopy structures introduced here.



\chapter[Variable Functional–Operator Homotopy Version]{Variable Functional–Operator Homotopy Version:\\
         From Path Spaces to Law–Operator Dynamics}


\section{Operator Paths and Connections on Variable Functionals}


The homotopy version of \GAlgebra\ treats operators as static elements
of a graded space $(V^{-1}\oplus V^0,\ell_1,\ell_2,\ell_3^H)$. In the
variable functional–operator homotopy version, operators become
\emph{paths} and even \emph{histories}: the basic objects are
time-parameterised (or scale-parameterised) operator curves and
functionals on these curves. The role of \GZTLC\ is then to provide a
connection not only on the triple-layer base
$(\Base,\Lambda_{\mathrm{par}},\GZLS)$, but also on the induced bundles
over operator path spaces. This section sets up the path-space and
variable functional geometry that will be used in Sections~6.2 and~6.3.



\subsubsection*{(a) Operator paths as basic dynamic objects}

Let $V^0$ be the degree-zero operator layer of the homotopy PFB-GNLA,
containing (completions of) strict operators and their law-driven
deformations. For concreteness one may keep in mind
\[
   V^0
   \subseteq
   \mathrm{End}\bigl(\Gamma(E_1)\oplus\cdots\oplus\Gamma(E_r)\bigr)
   \;\oplus\; \text{(higher operator sectors)},
\]
but the discussion will remain abstract.

\begin{definition}[Operator path]
Fix a compact parameter interval $I:=[0,1]$. An \emph{operator path}
  is a map
  \[
     \Gamma:I\longrightarrow V^0,
  \]
  such that:
  \begin{enumerate}[label=\roman*)]
    \item $\Gamma$ is continuous in an appropriate operator topology
          (e.g.\ norm, strong, or weak topology, depending on the
          context);
    \item on each finite-dimensional test subspace of the underlying
          state space, the map $s\mapsto\Gamma(s)$ is $C^1$.
  \end{enumerate}
  The space of such paths is denoted
  \[
     \mathsf{Path}(V^0)
     :=
     C^0(I;V^0)
     \quad\text{(with the chosen regularity)}.
  \]
\end{definition}

In many applications, $\Gamma(s)$ is itself an \emph{evolution
generator} or \emph{propagator} indexed by a physical or algorithmic
time $s$ (e.g.\ an RG scale, training step, or control horizon). The
path space $\mathsf{Path}(V^0)$ is thus the natural “dynamic stage” for
law–operator dynamics.



\begin{definition}[Concatenation and reversal]
For two operator paths $\Gamma_1,\Gamma_2\in\mathsf{Path}(V^0)$ with
  matching endpoints $\Gamma_1(1)=\Gamma_2(0)$, their concatenation
  $\Gamma_2\ast\Gamma_1$ is defined by:
  \[
    (\Gamma_2\ast\Gamma_1)(s)
    :=
    \begin{cases}
      \Gamma_1(2s),
        & 0\le s\le \tfrac12,\\[0.2em]
      \Gamma_2(2s-1),
        & \tfrac12\le s\le 1.
    \end{cases}
  \]
  The reversal of $\Gamma$ is the path
  \[
     \Gamma^\vee(s) := \Gamma(1-s).
  \]
\end{definition}

These operations will be compatible with the law-level composition and
with the triple-layer parallel transport implemented by \GZTLC.

\subsubsection*{(b) Variable functionals on operator paths}

We now formalise the notion of \emph{variable functionals} on operator
paths. These are the “observables of observables”: maps that take an
entire operator trajectory as input and output numerical data, law-data,
or new operators. They are the natural carriers of cost, action,
certificate, and law-learning objectives.

\begin{definition}[Variable functional on operator paths]
A \emph{variable functional} on operator paths is a map
  \[
     \mathcal{F}:\mathsf{Path}(V^0)\longrightarrow W,
  \]
  where $W$ is a target space (typically $\mathbb{R}$, $\mathbb{C}$,
  a finite-dimensional vector space, or another operator/law space),
  such that $\mathcal{F}$ satisfies suitable regularity conditions
  (e.g.\ continuity, Fréchet differentiability) with respect to the
  chosen topology on $\mathsf{Path}(V^0)$.
\end{definition}

In practice we will often restrict to \emph{cylindrical} or
\emph{locally supported} functionals.

\begin{definition}[Cylindrical and local functionals]
A variable functional $\mathcal{F}$ is called:
  \begin{itemize}
    \item \emph{cylindrical} if there exist evaluation points
          $0\le s_1<\dots<s_k\le1$ and a smooth map
          $f:V^0\times\cdots\times V^0\to W$ such that
          \[
             \mathcal{F}(\Gamma)
             =
             f\bigl(\Gamma(s_1),\dots,\Gamma(s_k)\bigr).
          \]
    \item \emph{local} if, informally, its value on $\Gamma$ depends
          only on the values of $\Gamma$ and a finite number of its
          derivatives in a neighbourhood of each $s$, and can be
          represented as an integral of a local density:
          \[
             \mathcal{F}(\Gamma)
             =
             \int_0^1 \mathcal{L}\bigl(s,\Gamma(s),
                                      \dot{\Gamma}(s),\dots\bigr)
                     \,\mathrm{d}s,
          \]
          for some Lagrangian-type map $\mathcal{L}$.
  \end{itemize}
\end{definition}

Cylindrical functionals are convenient for algebraic manipulations and
for defining path-integral–type objects; local functionals are natural
for action principles and cost functionals in law-space learning.



\subsubsection*{(c) Bundles over path space and connections on variable functionals}

To discuss \emph{connections} on variable functionals, we introduce
bundles over the path space $\mathsf{Path}(V^0)$ whose fibers carry the
values of functionals and their law-level structure.

Let $W$ be a fixed target space as above. Consider the trivial bundle
\[
   \mathcal{E}_0
   :=
   \mathsf{Path}(V^0)\times W
   \xrightarrow{\ \mathrm{pr}_1\ }
   \mathsf{Path}(V^0).
\]
A variable functional $\mathcal{F}:\mathsf{Path}(V^0)\to W$ can be
regarded as a section of $\mathcal{E}_0$. However, in the PFB-GNLA
context, $W$ will often carry extra structure (e.g.\ law-level GNLA or
$L_\infty$ structure), and the bundle will be non-trivial due to
law-space and triple-layer effects.

\begin{definition}[Law-structured path bundle (schematic)]
A \emph{law-structured path bundle} is a fiber bundle
  \[
     \mathcal{E}
     \xrightarrow{\ \pi_{\mathcal{E}}\ }
     \mathsf{Path}(V^0),
  \]
  whose fiber over $\Gamma$ is a law-object
  $L_\Gamma\in\GZLS$ (for example, an incarnation of
  $L_{\mathrm{strict}}$ or $L_{\mathrm{homo}}$ specialised to the
  dynamic regime corresponding to $\Gamma$), together with:
  \begin{itemize}
    \item a family of law-connections on the fibers compatible with
          \GZLOC;
    \item a projection from each $L_\Gamma$ to a “numerical” part
          $W_\Gamma$ in which values of functionals live:
          \[
             \pi_W:L_\Gamma\to W_\Gamma.
          \]
  \end{itemize}
\end{definition}

Sections of $\mathcal{E}$, i.e.\ maps $\Gamma\mapsto\sigma(\Gamma)\in
L_\Gamma$, represent \emph{law-valued functionals} on operator paths.
Composing with $\pi_W$ yields “ordinary” functionals with values in
$W_\Gamma$.

To differentiate such sections along variations of $\Gamma$, we need a
notion of connection on $\mathcal{E}$.

\begin{definition}[Connection on a law-structured path bundle]
A \emph{connection} on $\mathcal{E}$ is a rule that assigns to any
  smooth curve
  \[
     t\mapsto\Gamma_t
     \in\mathsf{Path}(V^0),
  \]
  and any section $\sigma$ of $\mathcal{E}$, a covariant derivative
  along the curve,
  \[
     \nabla^{\mathrm{path}}_{\dot{\Gamma}_t}\sigma
     \in\Gamma(\mathcal{E}),
  \]
  satisfying:
  \begin{enumerate}[label=\roman*)]
    \item linearity in $\dot{\Gamma}_t$ and derivation properties in
          $\sigma$;
    \item compatibility with the law-connections on fibers $L_\Gamma$
          induced by \GZLOC;
    \item reduction to the trivial derivative on $\mathcal{E}_0$ when
          the law-structure is trivial and the bundle is a product.
  \end{enumerate}
\end{definition}

Operationally, if $\sigma(\Gamma_t)$ is parallel-transported along a
path of operator paths $\Gamma_t$ in the triple-layer sense (using
\GZTLC on the underlying base plus law-space), then
$\nabla^{\mathrm{path}}_{\dot{\Gamma}_t}\sigma=0$ encodes the condition
that “the law attached to $\Gamma_t$ is moved without introducing
spurious distortions”.



\subsubsection*{(d) Infinitesimal variations of operator paths}

To express $\nabla^{\mathrm{path}}$ more concretely, we introduce
infinitesimal variations of operator paths. A tangent vector at
$\Gamma\in\mathsf{Path}(V^0)$ can be represented by a variation
$\delta\Gamma\in T_\Gamma\mathsf{Path}(V^0)$, i.e.\ a map
\[
   \delta\Gamma:I\to V^0
\]
of the same regularity class as $\Gamma$, thought of as the derivative
$\delta\Gamma(s) = \left.\frac{\mathrm{d}}{\mathrm{d}t}
\right|_{t=0}\Gamma_t(s)$ for a curve $t\mapsto\Gamma_t$.

Given a functional $\mathcal{F}:\mathsf{Path}(V^0)\to W$, its
\emph{Gateaux derivative} at $\Gamma$ in the direction $\delta\Gamma$
is
\[
   (\mathrm{D}\mathcal{F})_\Gamma[\delta\Gamma]
   :=
   \left.
     \frac{\mathrm{d}}{\mathrm{d}t}
   \right|_{t=0}
   \mathcal{F}(\Gamma_t),
\]
where $\Gamma_t$ is any curve with $\Gamma_0=\Gamma$ and
$\left.\frac{\mathrm{d}}{\mathrm{d}t}\right|_{t=0}\Gamma_t
=\delta\Gamma$.

If $\mathcal{F}$ is local and admits a density $\mathcal{L}$, this
derivative often decomposes as
\[
   (\mathrm{D}\mathcal{F})_\Gamma[\delta\Gamma]
   =
   \int_0^1
     \Bigl\langle
       \frac{\delta\mathcal{L}}{\delta\Gamma}(s),
       \delta\Gamma(s)
     \Bigr\rangle
   \,\mathrm{d}s,
\]
where $\frac{\delta\mathcal{L}}{\delta\Gamma}$ is a functional
derivative and $\langle\cdot,\cdot\rangle$ is a pairing between
$V^0$ and an appropriate dual space.

A connection $\nabla^{\mathrm{path}}$ on $\mathcal{E}$ refines this
raw derivative by “subtracting” the contribution coming from the
parallel transport of law-data along $\Gamma_t$ according to \GZTLC.
In particular, if $\sigma$ is a law-valued functional, then
\[
   \nabla^{\mathrm{path}}_{\delta\Gamma}\sigma(\Gamma)
   =
   (\mathrm{D}\sigma)_\Gamma[\delta\Gamma]
   - \text{(law-connection term induced by \GZLOC/\GZTLC)},
\]
in exact analogy with how covariant derivatives differ from ordinary
derivatives in bundle geometry.



\subsubsection*{(e) Coupling with the triple-layer connection \texorpdfstring{\GZTLC}{GZ-TLC}}

The variable functional–operator homotopy version does not live on
operator path space in isolation: it is coupled to the triple-layer
base via \GZTLC. To make this explicit, consider a family of operator
paths
\[
   (t,s)\longmapsto \Gamma_{t}(s)
\]
parametrised by $t$ (external evolution, e.g.\ law-learning time) and
$s\in I$ (internal path parameter), together with a corresponding
family of triple-layer base points
\[
   t\longmapsto
   (x(t),\lambda(t),L(t))
   \in \Base\times\Lambda_{\mathrm{par}}\times\GZLS,
\]
such that the pair
\[
   t\longmapsto
   \bigl(
     x(t),\lambda(t),L(t);\,
     \Gamma_t(\cdot)
   \bigr)
\]
is an admissible trajectory in the combined base
\[
   X_{\mathrm{vf-op}}
   :=
   \bigl(\Base\times\Lambda_{\mathrm{par}}\times\GZLS\bigr)
   \times \mathsf{Path}(V^0).
\]

\begin{definition}[Triple-layer–path coupling (schematic)]
A \emph{triple-layer–path coupling} is a connection
  \[
     \widetilde{\GZTLC}
  \]
  on a suitable bundle over $X_{\mathrm{vf-op}}$ such that:
  \begin{enumerate}[label=\roman*)]
    \item its restriction to $\Base\times\Lambda_{\mathrm{par}}\times
          \GZLS$ coincides with the original \GZTLC;
    \item its restriction to $\mathsf{Path}(V^0)$ induces the
          path-wise connection $\nabla^{\mathrm{path}}$ on
          law-structured path bundles as above;
    \item the induced curvature $\widetilde{\mathcal{F}}$ decomposes
          into:
          \[
             \widetilde{\mathcal{F}}
             =
             \mathcal{F}_{\GZTLC}
             + \mathcal{F}_{\mathrm{path}}
             + \text{(mixed terms)},
          \]
          where $\mathcal{F}_{\mathrm{path}}$ measures the failure of
          path-wise parallel transport to commute with triple-layer
          parallel transport, and the mixed terms encode law–operator
          dynamical curvature.
  \end{enumerate}
\end{definition}

In Section~6.2, we will show how to construct $\widetilde{\GZTLC}$ in a
way compatible with the homotopy operator structure of \GAlgebra, and
how variable functionals become the natural observables for this
extended connection.



\subsubsection*{(f) Summary: from static operators to path-level connection geometry}

To summarise, the passage from the homotopy version to the variable
functional–operator homotopy version of \GAlgebra\ involves the
following conceptual upgrades:

\begin{itemize}
  \item static degree-zero operators $V^0$ are replaced (or enriched)
        by operator paths $\mathsf{Path}(V^0)$;
  \item scalar- or law-valued observables are replaced by variable
        functionals on operator paths, realised as sections of
        law-structured path bundles $\mathcal{E}\to\mathsf{Path}(V^0)$;
  \item the triple-layer connection \GZTLC\ is extended to a
        path-sensitive connection $\widetilde{\GZTLC}$ whose path-wise
        component defines a covariant derivative
        $\nabla^{\mathrm{path}}$ on variable functionals;
  \item infinitesimal variations of operator paths are treated via
        functional derivatives, corrected by law-connection terms to
        yield path-space covariant derivatives.
\end{itemize}

This sets the stage for Section~6.2, where \GZTLC\ will be implemented
explicitly on operator path spaces, and for Section~6.3, where the
variable functional–operator homotopy version of \GAlgebra\ will be
constructed as a full PFB-GNLA version living on
$(\Base,\Lambda_{\mathrm{par}},\GZLS,\mathsf{Path}(V^0))$ with
operator-path–sensitive law-connections and homotopies.



\section[Implementing \GZTLCtxt\ on Operator Path Spaces]{Implementing \GZTLC\ (\GZTLCtxt) on Operator Path Spaces}


In the variable functional–operator homotopy version, the triple-layer
connection \GZTLC\ must act not only on the base
$\Base\times\Lambda_{\mathrm{par}}\times\GZLS$, but also on operator
paths $\Gamma\in\mathsf{Path}(V^0)$ and on law-structured bundles over
this path space. The goal of this section is to describe how \GZTLC\ is
\emph{lifted} to operator path spaces and how its curvature acquires
operator-path components. This will provide the geometric backbone for
the variable functional–operator PFB-GNLA in Section~6.3.



\subsubsection*{(a) Evaluation maps and pull-back of \GZTLC}

Recall that the triple-layer base is
\[
   X_{\mathrm{tot}}
   :=
   \Base\times\Lambda_{\mathrm{par}}\times\GZLS,
\]
with a principal $\widehat{G}$-bundle
\[
   \widehat{\pi}:\widehat{\Bundle}\longrightarrow X_{\mathrm{tot}}
\]
carrying the triple-layer connection
\(
   \GZTLC\in\Omega^1(\widehat{\Bundle};\widehat{\mathfrak{g}})
\).
Let $V^0$ be the degree-zero operator layer (Section~6.1) and
\[
   \mathsf{Path}(V^0) = C^0(I;V^0)
\]
its path space.

For each $s\in I$, we have an evaluation map
\[
   \mathrm{ev}_s:
     \mathsf{Path}(V^0)
     \longrightarrow
     V^0,\qquad
   \Gamma\longmapsto \Gamma(s).
\]
Similarly, at the triple-layer base level we consider a family of base
paths
\[
   \Xi:I\longrightarrow X_{\mathrm{tot}},
   \qquad
   s\longmapsto \Xi(s)
   :=\bigl(x(s),\lambda(s),L(s)\bigr),
\]
and we write
\[
   \mathrm{Ev}_s:X_{\mathrm{tot}}^I\longrightarrow X_{\mathrm{tot}},
   \qquad
   \Xi\longmapsto \Xi(s)
\]
for the evaluation on base paths.

We use these evaluation maps to pull back the triple-layer bundle and
connection to various function spaces.

\begin{definition}[Path-lifted triple-layer bundle (schematic)]
The \emph{path-lifted triple-layer bundle} is the pull-back bundle
  \[
     \widehat{\Bundle}^{I}
     :=
     \bigl\{
       \widehat{\Xi}:I\to\widehat{\Bundle}
       \,\big|\,
       \widehat{\pi}\circ\widehat{\Xi}=\Xi\in X_{\mathrm{tot}}^I
     \bigr\},
  \]
  equipped with the evaluation maps
  \[
     \widehat{\mathrm{Ev}}_s:
       \widehat{\Bundle}^{I}\longrightarrow \widehat{\Bundle},
       \qquad
       \widehat{\Xi}\longmapsto\widehat{\Xi}(s).
  \]
  The triple-layer connection \GZTLC\ on $\widehat{\Bundle}$ can be
  pulled back along $\widehat{\mathrm{Ev}}_s$ to yield a family of
  connections on $\widehat{\Bundle}^I$.
\end{definition}

Heuristically, $\widehat{\Bundle}^{I}$ collects pointwise lifts
$\widehat{\Xi}(s)$ of base paths $\Xi(s)$, and the pull-backs
$\widehat{\mathrm{Ev}}_s^*(\GZTLC)$ describe how to transport triple-layer
data at each “internal time” $s$ along the path direction.



\subsubsection*{(b) Path-averaged and pointwise lifted connections}

There are two natural ways to extend \GZTLC\ to operator path spaces:
\emph{pointwise lifting} and \emph{path averaging}.

\paragraph{Pointwise lifting.}

For each fixed $s\in I$, the pull-back
\[
   \GZTLC^{(s)}
   :=
   \widehat{\mathrm{Ev}}_s^*(\GZTLC)
\]
is a connection on the pull-back bundle
$\widehat{\Bundle}^{(s)}:=\widehat{\mathrm{Ev}}_s^*(\widehat{\Bundle})$
over the function space $X_{\mathrm{tot}}^I$. Restricting to those base
paths that correspond to admissible law–operator dynamics, we obtain a
connection that acts on infinitesimal variations “at time $s$”.

\paragraph{Path-averaged lifting.}

To encode cumulative effects along the path, we define a
\emph{path-averaged} connection by integrating these pointwise
connections along $I$. Schematicallly, for a tangent vector
$\delta\widehat{\Xi}\in T_{\widehat{\Xi}}\widehat{\Bundle}^I$, we set
\[
   \GZTLC^{\mathrm{path}}_{\widehat{\Xi}}
     (\delta\widehat{\Xi})
   :=
   \int_0^1
     \alpha(s)\cdot
     \GZTLC_{\widehat{\Xi}(s)}
       \bigl(\delta\widehat{\Xi}(s)\bigr)
   \,\mathrm{d}s,
\]
where $\alpha:I\to\mathbb{R}$ is a weighting function (for example
$\alpha\equiv1$ or a smooth bump function emphasising a portion of the
path). The resulting 1-form $\GZTLC^{\mathrm{path}}$ is a connection on
an appropriate associated bundle over path space.

\begin{remark}[Choice of weighting]
The choice of $\alpha$ reflects how strongly we want to weight
  different segments of the path in the effective connection. In many
  cases one takes $\alpha\equiv 1$, leading to a uniform averaging of
  triple-layer effects along the entire path. Other choices can be used
  to emphasise early-time or late-time behaviour (e.g.\ in law-space
  learning or control problems).
\end{remark}



\subsubsection*{(c) Horizontal lifts of operator paths under \texorpdfstring{\GZTLC}{GZ-TLC}}

We next specify what it means for an operator path to be \emph{horizontal}
with respect to the lifted connection. Let
\[
   t\longmapsto\Gamma_t\in\mathsf{Path}(V^0)
\]
be a one-parameter family of operator paths, and
\[
   t\longmapsto
   \bigl(x(t),\lambda(t),L(t)\bigr)\in X_{\mathrm{tot}}
\]
the corresponding triple-layer base trajectory.

\begin{definition}[Horizontal operator-path evolution]
A family $t\mapsto(\Gamma_t,x(t),\lambda(t),L(t))$ is said to be
  \emph{horizontal for \GZTLC on operator path spaces} if:
  \begin{enumerate}[label=\roman*)]
    \item for each fixed $s\in I$, the triple-layer base path
          $t\mapsto(x(t,s),\lambda(t,s),L(t,s))$ associated with
          $\Gamma_t(s)$ is horizontal for \GZTLC, i.e.\ its lift to
          $\widehat{\Bundle}$ has tangent vector in the horizontal
          subbundle;
    \item the variation of $\Gamma_t$ with respect to $t$ is compatible
          with the law-level parallel transport dictated by \GZTLC,
          in the sense that for any law-structured path bundle
          $\mathcal{E}\to\mathsf{Path}(V^0)$ and any section
          $\sigma$, one has
          \[
             \nabla^{\mathrm{path}}_{\partial_t\Gamma_t}
               \sigma(\Gamma_t)
             = 0
          \]
          whenever $\sigma$ is defined to be globally parallel with
          respect to \GZTLC along the underlying triple-layer base
          trajectory.
  \end{enumerate}
\end{definition}

Intuitively, a horizontal evolution of operator paths is one in which
all changes in $\Gamma_t$ can be “explained” by triple-layer geometry
without introducing extra twisting in the law-layer: the associated
law-objects $L_\Gamma$ are moved by pure parallel transport.



\subsubsection*{(d) Curvature on operator path spaces and law–operator dynamics}

The curvature of the lifted connection $\widetilde{\GZTLC}$ on
operator path spaces decomposes into:

\[
   \widetilde{\mathcal{F}}
   =
   \mathcal{F}_{\GZTLC}
   + \mathcal{F}_{\mathrm{path}}
   + \mathcal{F}_{\mathrm{mix}},
\]
where:

\begin{itemize}
  \item $\mathcal{F}_{\GZTLC}$ is the original triple-layer curvature, 
        encoding geometric, parameter, \\ and law-curvature;
  \item $\mathcal{F}_{\mathrm{path}}$ is the curvature purely along
        path directions in $\mathsf{Path}(V^0)$; it reflects
        path-dependent law–operator effects that arise even when the
        base point $(x,\lambda,L)$ is held fixed;
  \item $\mathcal{F}_{\mathrm{mix}}$ consists of mixed components
        coupling triple-layer directions with path-space directions;
        it measures how base evolution and operator-path evolution
        fail to commute.
\end{itemize}

\begin{definition}[Law–operator dynamical curvature (schematic)]
The \emph{law–operator dynamical curvature} is the collection of
  mixed curvature components
  \[
     \mathcal{F}_{\mathrm{dyn}}
     :=
     \mathcal{F}_{\mathrm{mix}}
     \;\subseteq\;
     \widetilde{\mathcal{F}},
  \]
  valued in the Lie algebra $\widehat{\mathfrak{g}}$ of the lifted
  structure group. It quantifies the obstruction to viewing operator
  dynamics as mere triple-layer parallel transport; non-zero
  $\mathcal{F}_{\mathrm{dyn}}$ signals genuinely new law–operator
  interaction effects.
\end{definition}

In particular, the homotopy data $(\ell_1,\ell_2,\ell_3^H)$ of the
operator layer can be read as algebraic manifestations of
$\mathcal{F}_{\mathrm{dyn}}$ and of law-curvature components seen in
the path-space setting. For example, higher-order commutator anomalies
in time-ordered exponentials of $\Gamma_t(s)$ correspond to non-trivial
holonomy of $\widetilde{\GZTLC}$ along loops in the combined space of
$(x,\lambda,L,\Gamma)$.



\subsubsection*{(e) Compatibility with homotopy \texorpdfstring{\GAlgebra}{G-Algebra}}

The lifted connection must respect the homotopy structure of the
operator layer. Concretely, if $(V^{-1}\oplus V^0,\ell_1,\ell_2,\ell_3^H)$
is in GZ-Harmony with $(H,\GZTLC)$ (Section~5.2), then:

\begin{itemize}
  \item $\nabla^{\mathrm{path}}$ acts on sections valued in $V^0$ in a
        way that is compatible with $\ell_1$ and $\ell_2$; in
        particular, the covariant derivative of the bracket is equal
        to the bracket of covariant derivatives, up to controlled
        homotopy terms governed by $\ell_3^H$ and higher $\ell_k$;
  \item the operator-path curvature $\mathcal{F}_{\mathrm{path}}$
        contributes to $\ell_3^H$ and to higher homotopies $\ell_k$
        in a manner consistent with the $L_\infty$ identities;
  \item the law–operator dynamical curvature
        $\mathcal{F}_{\mathrm{dyn}}$ is mapped to Jacobiator-gap and
        related certificate-type invariants, measuring the deviation
        of operator-path dynamics from strict law-level parallel
        transport.
\end{itemize}

Formally, one expects identities of the form
\[
   \nabla^{\mathrm{path}}_{\delta\Gamma}
     \ell_2(T_1,T_2)
   -
   \ell_2\bigl(
     \nabla^{\mathrm{path}}_{\delta\Gamma}T_1,\,
     T_2
   \bigr)
   -
   \ell_2\bigl(
     T_1,\,
     \nabla^{\mathrm{path}}_{\delta\Gamma}T_2
   \bigr)
   \sim
   \ell_1\bigl(
     \ell_3^H(\delta\Gamma,T_1,T_2)
   \bigr),
\]
for appropriate interpretations of $\delta\Gamma$ as an element of
$V^0$ along the path. This expresses that covariant differentiation
along operator-path directions fails to be a strict derivation of
$\ell_2$ by a homotopy term controlled by $\ell_3^H$.



\subsubsection*{(f) Summary: \GZTLC\ as a path-level engine for law–operator dynamics}

We can now summarise the implementation of \GZTLC\ on operator path
spaces:

\begin{itemize}
  \item The triple-layer connection \GZTLC\ on
        $\widehat{\Bundle}\to X_{\mathrm{tot}}$ is pulled back along
        evaluation maps to path-lifted bundles and then combined, via
        pointwise lifting and path-averaging, into a connection
        $\widetilde{\GZTLC}$ on a bundle over
        $X_{\mathrm{tot}}\times\mathsf{Path}(V^0)$.
  \item Horizontal operator-path evolutions are those for which
        operator dynamics and triple-layer base evolution are jointly
        governed by $\widetilde{\GZTLC}$, with no additional twisting
        in the law-layer beyond what is dictated by the connection.
  \item The curvature of $\widetilde{\GZTLC}$ splits into triple-layer,
        pure path, and mixed (law–operator dynamical) components; the
        latter encode genuinely new law–operator interaction effects.
  \item Compatibility with the homotopy structure of \GAlgebra\ is
        enforced by GZ-Harmony and by the requirement that
        $\nabla^{\mathrm{path}}$ respect the $L_\infty$ operations
        $(\ell_1,\ell_2,\ell_3^H,\dots)$ up to controlled homotopies.
\end{itemize}

This completes the geometric preparation needed to formulate the
variable functional–operator homotopy version of the GaoZheng
G-Algebra. In Section~6.3, these ideas will be assembled into a full
PFB-GNLA construction in which operator paths, variable functionals,
and the lifted connection $\widetilde{\GZTLC}$ jointly define a law–
operator dynamics that lives natively on
$\Base\times\Lambda_{\mathrm{par}}\times\GZLS\times\mathsf{Path}(V^0)$.



\section[Variable Functional--Operator Homotopy Version of the G-Algebra]{Variable Functional--Operator Homotopy Version of \\ \GAlgebra}


The variable functional--operator homotopy version of the GaoZheng
G-Algebra extends the homotopy PFB-GNLA by making operator paths and
functionals on these paths into first-class citizens. Instead of
working with static operators in $V^0$, we work on the path space
$\mathsf{Path}(V^0)$ and on law-structured bundles over it, equipped
with the path-lifted triple-layer connection $\widetilde{\GZTLC}$. The
goal of this section is to assemble these ingredients into an
integrated PFB-GNLA datum: the \emph{variable functional--operator
homotopy version} of \GAlgebra.



\subsubsection*{(a) Variable functional--operator geometric block}

We begin by extending the homotopy geometric--parameter--$H$ block
(Section~5.3) with explicit operator path and functional data.

\begin{definition}[Variable functional--operator geometric block]
A \\ \emph{variable functional--operator geometric block} is a tuple
  \[
     \mathsf{GeomPar}_{\mathrm{vf\text- op}}
     =
     \bigl(
       \Bundle,\Base,\Group;
       \Lambda_{\mathrm{par}};
       \{\Conn_\lambda\}_{\lambda\in\Lambda_{\mathrm{par}}},
       \widehat{\Conn},\widehat{\Curv};
       H;
       V^0,\mathsf{Path}(V^0)
     \bigr),
  \]
  where:
  \begin{itemize}
    \item $(\Bundle,\Base,\Group)$,
          $\Lambda_{\mathrm{par}}$,
          $\{\Conn_\lambda\}$,
          $\widehat{\Conn}$,
          $\widehat{\Curv}$,
          $H$ are as in the homotopy block $\mathsf{GeomPar}_H$;
    \item $V^0$ is the degree-zero operator layer of the homotopy
          PFB-GNLA (containing operators built from covariant
          derivatives, curvature, etc.);
    \item $\mathsf{Path}(V^0)$ is the operator path space
          $\Gamma:I\to V^0$ introduced in Section~6.1.
  \end{itemize}
\end{definition}

In addition, we assume that the triple-layer connection \GZTLC\ on
$X_{\mathrm{tot}}=\Base\times\Lambda_{\mathrm{par}}\times\GZLS$ has
been lifted to a path-level connection
$\widetilde{\GZTLC}$ on a suitable bundle over
\[
   X_{\mathrm{tot}}\times\mathsf{Path}(V^0),
\]
as described in Section~6.2. This lifted connection will control the
parallel transport of law-objects attached to operator paths.



\subsubsection*{(b) Variable functional--operator block}

We now introduce the algebraic layer built on operator paths and
variable functionals. This is where homotopy operators and path-level
connections meet.

\begin{definition}[Variable functional--operator block]
A \emph{variable functional--operator block} associated with
  $\mathsf{GeomPar}_{\mathrm{vf\text- op}}$ is a tuple
  \[
     \mathsf{VF\text- Op}
     =
     \bigl(
       \mathsf{Path}(V^0),
       \{\mathcal{F}_\alpha\}_{\alpha\in A};
       \mathcal{E},
       \nabla^{\mathrm{path}};
       \ell_\bullet^{\mathrm{vf}}
     \bigr),
  \]
  with:
  \begin{itemize}
    \item $\mathsf{Path}(V^0)$ the operator path space;
    \item $\{\mathcal{F}_\alpha\}_{\alpha\in A}$ a distinguished
          family of variable functionals
          $\mathcal{F}_\alpha:\mathsf{Path}(V^0)\to W_\alpha$,
          where each $W_\alpha$ is a law- or number-valued target
          space, typically representing cost, action, certificate, or
          law-learning objectives;
    \item $\mathcal{E}\to\mathsf{Path}(V^0)$ a law-structured path
          bundle whose fiber over $\Gamma$ is a law-object
          $L_\Gamma\in\GZLS$ (Section~6.1);
    \item $\nabla^{\mathrm{path}}$ a path-wise connection on
          $\mathcal{E}$, induced from the lifted triple-layer
          connection $\widetilde{\GZTLC}$;
    \item $\ell_\bullet^{\mathrm{vf}}$ a collection of multilinear
          operations on sections of $\mathcal{E}$ (or on an associated
          graded space) extending the homotopy operations
          $(\ell_1,\ell_2,\ell_3^H,\dots)$ to the variable
          functional--operator setting.
  \end{itemize}
\end{definition}

Here $\ell_\bullet^{\mathrm{vf}}$ plays the role of an
$L_\infty$-type structure not just on static operators but on
\emph{operator-path–dependent} objects. For instance, one may think of
\[
  \ell_2^{\mathrm{vf}}(\sigma_1,\sigma_2)(\Gamma)
  \sim
  \text{(graded commutator of the law-valued functionals
  $\sigma_1(\Gamma),\sigma_2(\Gamma)$)},
\]
and similarly for higher brackets, with $\nabla^{\mathrm{path}}$
ensuring that these operations transform covariantly under path-wise
parallel transport.



\subsubsection*{(c) Variable functional--operator law-level block}

On the law-level, we need a region of \GZLS where law-objects are
equipped not only with internal homotopy structure, but also with
path-level data and functionals.

\begin{definition}[Variable functional--operator GZ-LS slice]
A \emph{variable functional--operator GZ-LS slice} is a region
  \[
     {\GZLS}_{\mathrm{vf\text- op}}
     =
     \bigl(
       \mathfrak{L}_{\mathrm{vf\text- op}};
       J_{\mathrm{vf\text- op}},
       \mathrm{Loc}_{\mathrm{vf\text- op}},
       \Sigma_{\mathrm{vf\text- op}}
     \bigr),
  \]
  with:
  \begin{enumerate}[label=\roman*)]
    \item $\mathfrak{L}_{\mathrm{homo}}
           \subseteq \mathfrak{L}_{\mathrm{vf\text- op}}
           \subseteq \mathfrak{L}$, i.e.\ it contains the homotopy
          slice as a subregion;
    \item metrics and localisation data
          $(J_{\mathrm{vf\text- op}},\mathrm{Loc}_{\mathrm{vf\text- op}},
            \Sigma_{\mathrm{vf\text- op}})$ that extend the homotopy
          metrics to include contributions from operator-path data and
          variable functionals (e.g.\ action integrals, cost
          functionals, certificate trajectories);
    \item for each operator path $\Gamma$, a law-object $L_\Gamma$
          in $\mathfrak{L}_{\mathrm{vf\text- op}}$ whose internal
          homotopy structure is given by $\ell_\bullet^{\mathrm{vf}}$,
          and whose metric/threshold data reflect the values of
          selected functionals $\{\mathcal{F}_\alpha(\Gamma)\}$.
  \end{enumerate}
\end{definition}

We then record the law-level block:

\begin{definition}[Variable functional--operator law-level block]
A \emph{variable functional--operator law-level block} is a law-object
  \[
     L_{\mathrm{vf\text- op}}
     \in {\GZLS}_{\mathrm{vf\text- op}},
  \]
  together with an embedding
  \[
     \iota_{\mathrm{law}}^{\mathrm{vf\text- op}}:
       \mathsf{Path}(V^0)
       \hookrightarrow
       \mathfrak{L}_{\mathrm{vf\text- op}},
  \]
  such that:
  \begin{enumerate}[label=\roman*)]
    \item the image of $\iota_{\mathrm{law}}^{\mathrm{vf\text- op}}$
          consists of law-objects of the form $L_\Gamma$;
    \item the homotopy and path-level structure on $L_\Gamma$ is
          encoded by the variable functional--operator block
          $\mathsf{VF\text- Op}$ via $\ell_\bullet^{\mathrm{vf}}$ and
          $\nabla^{\mathrm{path}}$;
    \item the law-connection on
          $\mathfrak{L}_{\mathrm{vf\text- op}}$ induced from \GZLOC\
          coincides with the law-structured part of the lifted
          connection $\widetilde{\GZTLC}$ acting on $\mathcal{E}$.
  \end{enumerate}
\end{definition}



\subsubsection*{(d) Variable functional--operator PFB-GNLA datum}

We now assemble the geometric, operator-path, functional, and law-level
blocks into a single integrated PFB-GNLA datum for the variable
functional--operator homotopy version.

\begin{definition}[Variable functional--operator PFB-GNLA datum]
A \emph{variable functional--operator PFB-GNLA datum} is a tuple
  \[
    \mathsf{PFB\text- GNLA}_{\mathrm{vf\text- op}}
    =
    \bigl(
      \mathsf{GeomPar}_{\mathrm{vf\text- op}},
      \mathsf{VF\text- Op},
      L_{\mathrm{vf\text- op}};
      \Phi_{\mathrm{geom}}^{\mathrm{vf\text- op}},
      \mathcal{C}_{\mathrm{vf\text- op}}
    \bigr),
  \]
  where:
  \begin{itemize}
    \item $\mathsf{GeomPar}_{\mathrm{vf\text- op}}$ is the variable
          functional--operator geometric block;
    \item $\mathsf{VF\text- Op}$ is the variable functional--operator
          block;
    \item $L_{\mathrm{vf\text- op}}\in{\GZLS}_{\mathrm{vf\text- op}}$ is
          the corresponding law-level block;
    \item $\Phi_{\mathrm{geom}}^{\mathrm{vf\text- op}}$ is a
          geometry-to-law functor at the path level, assigning to
          each choice of PFB geometry, parameter path, $H$, and
          operator path family an element of
          $\mathfrak{L}_{\mathrm{vf\text- op}}$ that matches
          $L_{\mathrm{vf\text- op}}$;
    \item $\mathcal{C}_{\mathrm{vf\text- op}}$ is a law–geometry–path
          coupling that reconstructs operator paths and variable
          functionals (up to homotopy) from law-level data and the
          lifted connection $\widetilde{\GZTLC}$.
  \end{itemize}
\end{definition}

Informally, a variable functional--operator PFB-GNLA datum says:

\begin{quote}
  “We have a PFB geometry with parameters and a closed 3-form $H$; we
   consider operator paths on the associated bundles, equip them with
   variable functionals and a path-wise connection induced by
   $\widetilde{\GZTLC}$; we encode the resulting law–operator dynamics
   as homotopy law-objects in \GZLS; and we maintain a bi-directional
   (up to homotopy) correspondence between geometry, operator paths,
   functionals, and law-space.”
\end{quote}



\subsubsection*{(e) Relation to strict and homotopy versions}

The variable functional--operator version encompasses the strict and
homotopy versions as special cases, corresponding to progressively
forgetting path-level structure.

\begin{proposition}[Projection to homotopy and strict regimes (schematic)]
There exist natural projection functors
  \[
     \Pi_{\mathrm{homo}}:
       \mathsf{PFB\text- GNLA}_{\mathrm{vf\text- op}}
       \longrightarrow
       \mathsf{PFB\text- GNLA}_{\mathrm{homo}},
  \]
  \[
     \Pi_{\mathrm{strict}}:
       \mathsf{PFB\text- GNLA}_{\mathrm{homo}}
       \longrightarrow
       \mathsf{PFB\text- GNLA}_{\mathrm{strict}},
  \]
  such that:
  \begin{enumerate}[label=\roman*)]
    \item $\Pi_{\mathrm{homo}}$ forgets operator-path and functional
          dependence (collapsing $\mathsf{Path}(V^0)$ to $V^0$ and
          retaining only the pointwise homotopy structure
          $(\ell_1,\ell_2,\ell_3^H,\dots)$ and law-curvature
          \GZLCurv;
    \item $\Pi_{\mathrm{strict}}$ further forgets homotopy data,
          replacing the $L_\infty$ structure by its strict GNLA
          quotient (when the Jacobiator-gap certificate permits),
          thus recovering the strict PFB-GNLA.
  \end{enumerate}
  Further details and a complete proof are given in Technical Appendix~II
  (Section~\ref{app:II-P6}).
\end{proposition}

Conversely, the variable functional--operator version can be viewed as
a controlled extension of the homotopy version:

\begin{itemize}
  \item It adds explicit path-time or scale-time structure to the
        operator layer.
  \item It attaches variable functionals that encode objectives,
        actions, and certificates along operator paths.
  \item It extends \GZTLC\ to $\widetilde{\GZTLC}$ on path space,
        making law–operator dynamics a geometric object in its own
        right.
\end{itemize}



\subsubsection*{(f) Position inside \texorpdfstring{\GZLS}{GZ-LS}, \texorpdfstring{\GZTLC}{GZ-TLC}, and \texorpdfstring{\GZGMS}{GZ-GMS}}

Within the GaoZheng G-Framework, the variable functional--operator
PFB-GNLA occupies the following position:

\begin{itemize}
  \item In $\mathsf{Obj}$, it considers principal bundles with
        parameterised connections, closed 3-forms $H$, and operator
        path families, together with their induced law–operator
        dynamics.
  \item In $\mathsf{Law}$, it lives in the extended slice
        ${\GZLS}_{\mathrm{vf\text- op}}$, where law-objects carry
        internal homotopy structures and path-level functional data.
  \item In $\mathsf{Flow}$, it uses the lifted connection
        $\widetilde{\GZTLC}$ as the engine governing the joint
        evolution of geometry, parameters, laws, operator paths, and
        variable functionals.
  \item In $\mathsf{Layer}$, it primarily belongs to higher CH-layers
        where semantic, homotopy, and dynamical metrics (including
        those derived from variable functionals and certificates) are
        fully visible.
\end{itemize}

From this vantage point, the four versions of PFB-GNLA can be viewed
as nested charts of \GZGMS:

\[
   \mathsf{PFB\text- GNLA}_{\mathrm{strict}}
   \;\subset\;
   \mathsf{PFB\text- GNLA}_{\mathrm{homo}}
   \;\subset\;
   \mathsf{PFB\text- GNLA}_{\mathrm{vf\text- op}}
   \;\subset\;
   \mathsf{PFB\text- GNLA}_{\mathrm{orig}},
\]
where the meta-mathematical original version (Part~III) reconstructs
these structures from PL-PI provenance and law-level flows.



\medskip

In conclusion, the variable functional--operator homotopy version of
the GaoZheng \\
G-Algebra provides the most dynamic and expressive
realisation of PFB-GNLA within the present monograph. It is the level
at which:

\begin{itemize}
  \item law-level homotopy structures,
  \item operator-path dynamics,
  \item variable functionals (actions, costs, certificates),
  \item and triple-layer geometry
\end{itemize}

are unified into a single, path-sensitive law–operator geometry, ready
to be deployed in applications to physics, AI, and life/complex systems in Part~V
and to be compared with the meta-mathematical original version in
Part~III.





\part[Original G-Algebra and Four-Way Comparison]{Meta-Mathematical Original GaoZheng G-Algebra and Four-Way Comparison}


\chapter[Original G-Algebra under PL-PI MMT]{Original GaoZheng G-Algebra under the PL-PI MMT Provenance}


\section{Law-Level Framework from Pan-Logic and \ Pan-Iterative Analysis}


In this chapter we reverse the direction of the constructions in
Parts~II and~III: instead of starting from principal bundles and
connections and then encoding them as law-objects, we start from the
PL-PI Meta-Mathematical Theory (\PLPIMMT) and reconstruct, at the
law-level, the framework in which the original GaoZheng G-Algebra
(\GAlgebra) lives. The present section isolates the pure law-level
framework induced by Pan-Logic Analysis (PL) and Pan-Iterative Analysis
(PI); Section~7.2 will then realise the original PFB-GNLA as a
law-level \GAlgebra\ built inside this PL-PI law-space.



\subsubsection*{(a) Pan-Logic Analysis: law-objects for inference schemes}

Pan-Logic Analysis views each concrete logical or inferential scheme as
a law-object. We make this precise by working with a category of
\emph{logic profiles}.

\begin{definition}[Logic profile and logic law-object]
A \emph{logic profile} is a tuple
  \[
    \mathsf{LP}
    =
    \bigl(
      \mathsf{Form},\ \vdash,\ 
      \mathsf{Mod},\ \mathrm{Sem}
    \bigr),
  \]
  where:
  \begin{itemize}
    \item $\mathsf{Form}$ is a space (or category) of formulas,
          judgments, or types;
    \item $\vdash\ \subseteq \mathsf{Form}^{k+1}$ is a derivability
          relation (possibly multi-ary) encoding inference rules;
    \item $\mathsf{Mod}$ is a class of models / valuations;
    \item $\mathrm{Sem}$ is a satisfaction or semantics map
          $\mathrm{Sem}:\mathsf{Mod}\times\mathsf{Form}\to\{\text{truth
          values / degrees}\}$ compatible with $\vdash$.
  \end{itemize}
  A \emph{logic law-object} is a law-object in \GZLS\ of the form
  \[
    L_{\mathrm{PL}}
    =
    \Phi_{\mathrm{PL}}(\mathsf{LP})
    \in\GZLS,
  \]
  obtained by a functor
  \(
    \Phi_{\mathrm{PL}}:\mathsf{LogicProf}\to\GZLS
  \)
  that encodes $\mathsf{LP}$ into law-space data
  $(\mathfrak{L};J,\mathrm{Loc},\Sigma)$.
\end{definition}

The category $\mathsf{LogicProf}$ has objects $\mathsf{LP}$ and
morphisms given by structure-preserving translations between logics
(embeddings, conservative extensions, modality additions, etc.). Pan-
Logic Analysis is then the study of the induced law-space region
\[
   \mathfrak{L}_{\mathrm{PL}}
   :=
   \Phi_{\mathrm{PL}}(\mathsf{LogicProf})
   \subseteq\mathfrak{L},
\]
equipped with law-transforms $\GZLT^{\alpha\beta}$ corresponding to
logic translations and with \\ a law-connection \GZLOC\ that organises
continuous deformations between logics.



\subsubsection*{(b) Pan-Iterative Analysis: law-objects for procedures}

Pan-Iterative Analysis plays the analogous role for iterative
procedures: optimisation algorithms, learning rules, control loops,
renormalisation flows, and so forth.

\begin{definition}[Iterative profile and iterative law-object]
An \emph{iterative profile} is a tuple
  \[
    \mathsf{IP}
    =
    \bigl(
      \mathsf{State},\ \mathsf{Update},\ \mathsf{Obs},\ \mathsf{Eval}
    \bigr),
  \]
  where:
  \begin{itemize}
    \item $\mathsf{State}$ is a state space (possibly a manifold or a
          space of measures / distributions);
    \item $\mathsf{Update}$ is a family of maps
          $U_\theta:\mathsf{State}\to\mathsf{State}$ indexed by
          design parameters $\theta$ (step size, learning rate, etc.);
    \item $\mathsf{Obs}$ is a space of observables on
          $\mathsf{State}$;
    \item $\mathsf{Eval}$ is an evaluation functional assigning
          convergence, stability, cost, or performance scores to
          trajectories generated by $\mathsf{Update}$.
  \end{itemize}
  An \emph{iterative law-object} is a law-object in \GZLS\ of the form
  \[
    L_{\mathrm{PI}}
    =
    \Phi_{\mathrm{PI}}(\mathsf{IP})
    \in\GZLS,
  \]
  where
  \(
    \Phi_{\mathrm{PI}}:\mathsf{IterProf}\to\GZLS
  \)
  encodes $\mathsf{IP}$ as law-space data with composition and metric
  structure.
\end{definition}

Pan-Iterative Analysis is then the study of the region
\[
   \mathfrak{L}_{\mathrm{PI}}
   :=
   \Phi_{\mathrm{PI}}(\mathsf{IterProf})
   \subseteq\mathfrak{L},
\]
together with law-transforms corresponding to allowed modifications of
iterative schemes \\ (changing loss functions, adding noise, changing
schedules, etc.) and a law-connection describing smooth deformations
between such schemes.



\subsubsection*{(c) PL-PI law-systems and their coupling}

The original PFB-GNLA does not arise from PL or PI alone, but from
their \emph{coupling}: a law-level system that specifies both which
logic governs inference and which iterative schemes implement
procedural evolution, together with constraints linking the two.

\begin{definition}[PL-PI law-system (schematic)]
A \emph{PL-PI law-system} is a tuple
  \[
    \mathsf{LS}_{\mathrm{PL\text- PI}}
    =
    \bigl(
      L_{\mathrm{PL}},\ L_{\mathrm{PI}};\ \Theta_{\mathrm{PL\text- PI}}
    \bigr),
  \]
  where:
  \begin{itemize}
    \item $L_{\mathrm{PL}}\in\mathfrak{L}_{\mathrm{PL}}$ is a logic
          law-object;
    \item $L_{\mathrm{PI}}\in\mathfrak{L}_{\mathrm{PI}}$ is an
          iterative law-object;
    \item $\Theta_{\mathrm{PL\text- PI}}$ is a coupling datum
          encoding:
          \begin{enumerate}[label=\alph*)]
            \item how inference rules in $L_{\mathrm{PL}}$ constrain
                  admissible updates in $L_{\mathrm{PI}}$;
            \item how evaluation functionals in $L_{\mathrm{PI}}$
                  determine or update logical commitments in
                  $L_{\mathrm{PL}}$;
            \item how both are embedded into the global law-space
                  \GZLS\ with a common law-connection \GZLOC.
          \end{enumerate}
  \end{itemize}
\end{definition}

A simple example is a scheme where $L_{\mathrm{PL}}$ specifies a
probabilistic logic with certain conditional independences, and
$L_{\mathrm{PI}}$ specifies a stochastic gradient procedure; the
coupling $\Theta_{\mathrm{PL\text- PI}}$ enforces that updates respect
the conditional independence structure and that inference rules are
adapted to the noise and cost landscape seen by the iterative
procedure.

From the \PLPIMMT\ perspective, PL-PI law-systems are the \emph{basic
law-level objects} from which both object-level geometries (principal
bundles, connections) and operator-level algebras (\GAlgebra) will be
reconstructed.



\subsubsection*{(d) Law-space of PL-PI systems and flows}

We gather all PL-PI law-systems into a law-space category and introduce
flows between them.

\begin{definition}[PL-PI law-space and flows]
The \emph{PL-PI law-space} is the category
  \[
    \mathsf{Law}_{\mathrm{PL\text- PI}}
    :=
    \mathrm{Obj}(\mathsf{Law}_{\mathrm{PL\text- PI}}),
    \quad
    \mathrm{Mor}(\mathsf{Law}_{\mathrm{PL\text- PI}}),
  \]
  whose objects are PL-PI law-systems
  $\mathsf{LS}_{\mathrm{PL\text- PI}}$ and whose morphisms are
  triples
  \[
    \Phi
    =
    \bigl(
      \Phi_{\mathrm{PL}},\Phi_{\mathrm{PI}};
      \Phi_\Theta
    \bigr),
  \]
  where $\Phi_{\mathrm{PL}}$ and $\Phi_{\mathrm{PI}}$ are law-transforms
  in $\mathfrak{L}_{\mathrm{PL}}$ and $\mathfrak{L}_{\mathrm{PI}}$
  respectively, and $\Phi_\Theta$ is a consistency map between
  couplings, compatible with \GZLOC.
  
  A \emph{PL-PI law-flow} is a smooth path
  \[
    \Gamma_{\mathrm{law}}:I\to\mathsf{Law}_{\mathrm{PL\text- PI}},
    \qquad
    t\longmapsto
    \mathsf{LS}_{\mathrm{PL\text- PI}}(t)
  \]
  such that the associated curve in $\GZLS$ is horizontal for the
  law-connection \GZLOC.
\end{definition}

The law-connection \GZLOC\ thus induces a notion of parallel transport
on $\mathsf{Law}_{\mathrm{PL\text- PI}}$, with a law-curvature
component \GZLCurv\ that measures obstructions to path-independence of
PL-PI deformations.



\subsubsection*{(e) Law-level configuration space for the original \texorpdfstring{\GAlgebra}{G-Algebra}}

From the viewpoint of the original GaoZheng G-Algebra, only some
PL-PI law-systems are admissible: those whose structure is rich enough
to support a PFB-GNLA reconstruction. We denote this admissible
subspace by
\[
   \mathsf{Law}_{\mathrm{PL\text- PI}}^{\mathrm{GZ}}
   \subseteq
   \mathsf{Law}_{\mathrm{PL\text- PI}}.
\]

\begin{definition}[GZ-admissible PL-PI systems (schematic)]
A PL-PI law-system
  $\mathsf{LS}_{\mathrm{PL\text- PI}}=
   (L_{\mathrm{PL}},L_{\mathrm{PI}};\Theta_{\mathrm{PL\text- PI}})$
  is \emph{GZ-admissible} if:
  \begin{enumerate}[label=\roman*)]
    \item (Principal bundle realisability)
          There exists a category $\mathsf{GeomObj}^{\mathrm{PFB}}$
          of principal bundles with connection and curvature, and a
          reconstruction functor
          \[
            \mathcal{R}_{\mathrm{PFB}}:
              \mathsf{Law}_{\mathrm{PL\text- PI}}^{\mathrm{GZ}}
              \longrightarrow
              \mathsf{GeomObj}^{\mathrm{PFB}},
          \]
          such that logical structure in $L_{\mathrm{PL}}$ and
          iterative structure in $L_{\mathrm{PI}}$ determine (up to
          gauge) the data of a PFB geometry.
    \item (PFB-GNLA realisability)
          The reconstructed PFB geometry admits a strict/homotopy/
          variable functional–operator PFB-GNLA structure whose
          operator algebra \GAlgebra\ is recovered as a law-object in
          \GZLS, with internal structure consistent with
          $\Theta_{\mathrm{PL\text- PI}}$.
    \item (CH-layer compatibility)
          The PL-PI metrics and stratifications induced by
          $L_{\mathrm{PL}}$ and $L_{\mathrm{PI}}$ align with the CH-
          layering encoded in $\mathsf{Layer}$, so that law-level
          metrics and invariants match those used in the PFB-GNLA
          charts.
  \end{enumerate}
\end{definition}

The law-level configuration space for the original \GAlgebra\ is then
the full subcategory
\[
   \mathsf{Law}_{\mathrm{PL\text- PI}}^{\mathrm{GZ}}
   \subset
   \mathsf{Law}_{\mathrm{PL\text- PI}},
\]
equipped with the restricted law-connection and law-curvature. It is
this space in which the original PFB-GNLA will be constructed purely
from PL-PI data in Section~7.2.



\subsubsection*{(f) Summary: PL-PI MMT as the provenance of the law-level stage}

At this point, the role of the PL-PI Meta-Mathematical Theory in the
original version of \GAlgebra\ can be summarised as follows:

\begin{itemize}
  \item Pan-Logic Analysis (PL) supplies a law-space region
        $\mathfrak{L}_{\mathrm{PL}}$ of logic law-objects, with
        law-transforms and law-connections encoding admissible changes
        of inference scheme.
  \item Pan-Iterative Analysis (PI) supplies a law-space region
        $\mathfrak{L}_{\mathrm{PI}}$ of iterative law-objects, with
        law-transforms and law-connections encoding admissible changes
        of procedure.
  \item PL-PI law-systems $(L_{\mathrm{PL}},L_{\mathrm{PI}};
        \Theta_{\mathrm{PL\text- PI}})$ form a law-space
        $\mathsf{Law}_{\mathrm{PL\text- PI}}$ whose flows and
        curvature are governed by \GZLOC\ and \GZLCurv.
  \item The GZ-admissible subspace
        $\mathsf{Law}_{\mathrm{PL\text- PI}}^{\mathrm{GZ}}$ consists of
        those PL-PI systems that can be reconstructed as PFB geometries
        and PFB-GNLA structures; this is the law-level arena in which
        the original GaoZheng G-Algebra is defined.
\end{itemize}

From the viewpoint of the monograph, this section identifies the
law-level “stage” on which the original version of \GAlgebra\ will be
enacted: a PL-PI law-space endowed with connections and curvatures,
stratified by CH-layering and filtered by GZ-admissibility. The next
section will construct the original PFB-GNLA directly on this stage,
making explicit how object-level geometry and operator algebras arise
as emergent structures from PL-PI law-flows.



\section{Original PFB-GNLA as Law-Level \GAlgebra}


We now put the PL-PI law-level framework of Section~7.1 to use and
describe the \emph{original} version of the GaoZheng G-Algebra
(\GAlgebra, PFB-GNLA) as a \emph{purely law-level} construction. In
contrast to the strict, homotopy, and variable functional--operator
versions (which start from principal bundles and then climb up to
law-space), the original version starts from GZ-admissible PL-PI
law-systems in $\mathsf{Law}_{\mathrm{PL\text- PI}}^{\mathrm{GZ}}$ and
reconstructs PFB-GNLA structures as emergent objects inside \GZLS.



\subsubsection*{(a) From GZ-admissible PL-PI law-systems to law-geometry}

Let
\[
   \mathsf{LS}_{\mathrm{PL\text- PI}}
   =
   (L_{\mathrm{PL}},L_{\mathrm{PI}};\Theta_{\mathrm{PL\text- PI}})
   \in
   \mathsf{Law}_{\mathrm{PL\text- PI}}^{\mathrm{GZ}}
\]
be a GZ-admissible PL-PI law-system (Definition~7.1). Intuitively:

\begin{itemize}
  \item $L_{\mathrm{PL}}$ specifies which inferential laws are allowed;
  \item $L_{\mathrm{PI}}$ specifies which iterative/procedural laws are
        allowed;
  \item $\Theta_{\mathrm{PL\text- PI}}$ specifies how they cohere as a
        single law-system.
\end{itemize}

The first step in the original PFB-GNLA construction is to extract from
$\mathsf{LS}_{\mathrm{PL\text- PI}}$ a \emph{law-geometry} object, i.e.\
a law-level analogue of a principal bundle with connection.

\begin{definition}[Law-geometry associated with a PL-PI law-system]
A \emph{law-geometry} associated with a GZ-admissible PL-PI
  law-system $\mathsf{LS}_{\mathrm{PL\text- PI}}$ is a tuple
  \[
    \mathsf{LG}
    =
    \bigl(
      \mathfrak{L}_{\mathrm{base}},
      \mathfrak{L}_{\mathrm{fib}},
      \GZLOC^{\mathrm{tot}},
      \GZLCurv^{\mathrm{tot}}
    \bigr),
  \]
  where:
  \begin{itemize}
    \item $\mathfrak{L}_{\mathrm{base}}$ is a “base” region of law-space
          describing admissible PL-PI configurations (e.g.\ the image
          of $\mathsf{Law}_{\mathrm{PL\text- PI}}^{\mathrm{GZ}}$ under
          a forgetful functor into \GZLS);
    \item $\mathfrak{L}_{\mathrm{fib}}$ is a “fiber” region describing
          internal law-degrees of freedom (e.g.\ choices of local
          logics, local procedures, or internal law parameters);
    \item $\GZLOC^{\mathrm{tot}}$ is a law-connection on a suitable
          law-bundle
          \[
            \pi_{\mathrm{law}}:
              \mathcal{P}_{\mathrm{law}}
              \longrightarrow
              \mathfrak{L}_{\mathrm{base}},
          \]
          with fiber $\mathfrak{L}_{\mathrm{fib}}$;
    \item $\GZLCurv^{\mathrm{tot}}$ is the associated law-curvature,
          recording law-level analogues of “field strength” and
          “holonomy defects” in PL-PI space.
  \end{itemize}
\end{definition}

This $\mathsf{LG}$ should be thought of as the \emph{law-level
principal-bundle surrogate}: its base encodes admissible PL-PI
configurations, its fiber encodes internal law-degrees of freedom, and
its connection encodes what it means to “move law-consistently” between
different PL-PI configurations.



\subsubsection*{(b) Law-level operator layer and original \GAlgebra}

Given a law-geometry $\mathsf{LG}$, we introduce a purely law-level
operator layer. Instead of starting from differential operators on
sections of associated bundles, we start from operators generated by
law-transforms and law-curvature data, acting on law-objects in
$\mathcal{P}_{\mathrm{law}}$.

\begin{definition}[Law-level operator layer]
Let $\mathsf{LG}$ be as above. The \emph{law-level operator layer}
  associated with $\mathsf{LG}$ is a vector space
  \[
     A_{\mathrm{law}}
     \subseteq
     \mathrm{End}\bigl(
       \Gamma(\mathcal{P}_{\mathrm{law}})
     \bigr),
  \]
  generated (schematically) by:
  \begin{itemize}
    \item law-parallel-transport operators along curves in
          $\mathfrak{L}_{\mathrm{base}}$, derived from
          $\GZLOC^{\mathrm{tot}}$;
    \item law-curvature insertions derived from
          $\GZLCurv^{\mathrm{tot}}$;
    \item PL-PI induced transformations corresponding to changes in
          logic profile or iterative profile within
          $\mathsf{LS}_{\mathrm{PL\text- PI}}$.
  \end{itemize}
\end{definition}

Equipped with the commutator bracket $[\cdot,\cdot]$ (and its homotopy
upgrades when needed), $A_{\mathrm{law}}$ carries a natural GNLA or
$L_\infty$ structure. The \emph{original law-level \GAlgebra} is then
defined by isolating the appropriate GNLA data inside $A_{\mathrm{law}}$.

\begin{definition}[Original law-level \GAlgebra\ (schematic)]
The \emph{original law-level GaoZheng G-Algebra} associated with a
  GZ-admissible PL-PI law-system is a GNLA (or $L_\infty$-algebra)
  \[
     \GAlgebra_{\mathrm{orig}}
     =
     \bigl(
       A_{\mathrm{law}},
       [\cdot,\cdot]_{\mathrm{law}},
       \Jac_{\mathrm{law}};\mathsf{Str}_{\mathrm{law}}
     \bigr),
  \]
  where:
  \begin{itemize}
    \item $A_{\mathrm{law}}$ is a distinguished subspace of the
          law-operator layer constructed above;
    \item $[\cdot,\cdot]_{\mathrm{law}}$ is the commutator (or graded
          commutator) restricted to $A_{\mathrm{law}}$;
    \item $\Jac_{\mathrm{law}}$ is the corresponding Jacobiator;
    \item $\mathsf{Str}_{\mathrm{law}}$ collects gradings, filtrations,
          and law-metric data induced from $J,\mathrm{Loc},\Sigma$ and
          from the PL-PI evaluation functionals.
  \end{itemize}
\end{definition}

By construction, $\GAlgebra_{\mathrm{orig}}$ is \emph{defined entirely
inside \GZLS}: it is a law-operator algebra built from PL-PI law space
and law-geometry, prior to any explicit reference to object-level
bundles. The principal-bundle interpretations of Parts~II and~III are
then recovered as \emph{representations} of
$\GAlgebra_{\mathrm{orig}}$.



\subsubsection*{(c) Representation theorem: from law-level \GAlgebra\ to PFB-GNLA}

The key statement connecting the original law-level \GAlgebra\ to the
PFB-GNLA charts of Parts~II and~III is a representation theorem: any
GZ-admissible PL-PI law-system gives rise to a law-level
$\GAlgebra_{\mathrm{orig}}$, and conversely any such law-level
\GAlgebra\ admits a family of PFB-GNLA representations.

\begin{proposition}[Law-level \GAlgebra\ and PFB-GNLA representations (schematic)]
Let
  $\mathsf{LS}_{\mathrm{PL\text- PI}}
   \in\mathsf{Law}_{\mathrm{PL\text- PI}}^{\mathrm{GZ}}$
  be GZ-admissible, with associated law-geometry $\mathsf{LG}$ and
  original law-level \GAlgebra\ $\GAlgebra_{\mathrm{orig}}$.
  Then:
  \begin{enumerate}[label=\roman*)]
    \item (Existence of PFB-GNLA representations)
          There exists a non-empty class of PFB-GNLA data
          (strict/homotopy/variable functional--operator) such that
          each datum defines a representation
          \[
             \rho:
               \GAlgebra_{\mathrm{orig}}
               \longrightarrow
               \mathrm{End}(\mathcal{H}),
          \]
          where $\mathcal{H}$ denotes an appropriate space of
          sections, functionals, or operator paths, and such that the
          representation preserves the GNLA/$L_\infty$ structure up to
          controlled homotopies.
    \item (Equivalence up to law-level gauge)
          Two PFB-GNLA representations \\ $\rho_1,\rho_2$ of
          $\GAlgebra_{\mathrm{orig}}$ that arise from the same PL-PI
          law-system are equivalent up to law-level gauge
          transformations: there exists an invertible law-operator
          $U$ such that
          \[
             \rho_2(T)
             =
             U\circ\rho_1(T)\circ U^{-1},
             \qquad
             \forall T\in \GAlgebra_{\mathrm{orig}},
          \]
          and $U$ respects PL-PI coupling and CH-layering.
  \end{enumerate}
  A complete proof is given in Technical Appendix~I
  (Section~\ref{app:I-P7}).
\end{proposition}

We do not prove this proposition in detail here; it is a schematic
summary of the reconstruction results scattered across the earlier
parts of the monograph. The main point is that the original PFB-GNLA
is fundamentally a law-level GNLA; PFB geometry enters as a convenient
“strict chart” on which to realise its representations.



\subsubsection*{(d) Law-level curvature, Jacobiator, and CH-layering in the original version}

Within the original law-level \GAlgebra, curvature and Jacobiator are
both encoded directly in $\mathsf{LG}$ and $A_{\mathrm{law}}$.

\begin{itemize}
  \item The law-curvature $\GZLCurv^{\mathrm{tot}}$ determines
        components of $\Jac_{\mathrm{law}}$ via the same curvature–
        homotopy mechanism used in the homotopy version, but now at
        the level of PL-PI law-space.
  \item CH-layering (Section~1.4 and Section~4.1 of the CH-layering
        part) stratifies the law-level stage so that:
        \begin{itemize}
          \item at extensional layers $l_2$, $\GAlgebra_{\mathrm{orig}}$
                looks approximately like a strict Lie algebra;
          \item at higher layers $l_{\ge3}$, the semantic/homotopy
                contributions of $\Jac_{\mathrm{law}}$ and higher
                brackets become visible, and law-space metrics depend
                on homotopy invariants (e.g.\ Jacobiator-gap).
        \end{itemize}
\end{itemize}

Thus, in the original version, the role of law-curvature and CH-layering
is not derivative of PFB geometry; instead, \emph{PFB geometry is
interpreted as a law-level reparametrisation} in which these law-level
objects acquire a more familiar geometric form.



\subsubsection*{(e) Original \GAlgebra\ as the common source of the four PFB-GNLA versions}

We can now describe the relationship between the original law-level
\GAlgebra\ and the four PFB-GNLA versions developed earlier.

\begin{proposition}[Four PFB-GNLA versions as law-level charts (schematic)]
Let \\ $\GAlgebra_{\mathrm{orig}}$ be the original law-level \GAlgebra\
  associated with a GZ-admissible PL-PI law-system. Then there exist
  embeddings
  \[
    \iota_{\mathrm{strict}},
    \iota_{\mathrm{homo}},
    \iota_{\mathrm{vf\text- op}}
  \]
  such that:
  \begin{enumerate}[label=\roman*)]
    \item (\emph{Strict chart})
          $\iota_{\mathrm{strict}}$ realises a sub-structure of
          $\GAlgebra_{\mathrm{orig}}$ as the strict PFB-GNLA
          $\GAlgebra_{\mathrm{strict}}$:
          \[
            \iota_{\mathrm{strict}}:
               \GAlgebra_{\mathrm{strict}}
               \hookrightarrow
               \GAlgebra_{\mathrm{orig}},
          \]
          corresponding to a region of PL-PI law-space where
          law-curvature and Jacobiator are negligible/central.
    \item (\emph{Homotopy chart})
          $\iota_{\mathrm{homo}}$ realises a homotopy upgrade of
          $\GAlgebra_{\mathrm{strict}}$ as a sub-structure of
          $\GAlgebra_{\mathrm{orig}}$ with non-trivial homotopy
          brackets $(\ell_1,\ell_2,\ell_3^H,\dots)$, corresponding to
          a PL-PI region with controlled but non-negligible
          law-curvature.
    \item (\emph{Variable functional--operator chart})
          $\iota_{\mathrm{vf\text- op}}$ realises the variable
          functional--operator PFB-GNLA
          $\GAlgebra_{\mathrm{vf\text- op}}$ as a path-level chart of
          $\GAlgebra_{\mathrm{orig}}$, where operator-path dynamics
          and variable functionals encode law-flows in PL-PI space.
  \end{enumerate}
  Further details and a complete proof are given in Technical Appendix~II
  (Section~\ref{app:II-P8}).
\end{proposition}

Conceptually, the original \GAlgebra\ is the “global law-level object”;
the other three versions are different coordinate systems and
truncations adapted to:

\begin{itemize}
  \item near-flat / near-Lie behaviour (strict version),
  \item homotopy-rich behaviour (homotopy version),
  \item dynamical/path-dependent behaviour (variable functional--operator
        version).
\end{itemize}



\subsubsection*{(f) Summary: Original PFB-GNLA as a law-level \GAlgebra\ derived from PL-PI MMT}

Summarising the construction:

\begin{itemize}
  \item Start from a GZ-admissible PL-PI law-system
        $\mathsf{LS}_{\mathrm{PL\text- PI}}\in
         \mathsf{Law}_{\mathrm{PL\text- PI}}^{\mathrm{GZ}}$;
  \item Build a law-geometry $\mathsf{LG}$ inside \GZLS\ from PL-PI
        data and the global law-connection \GZLOC;
  \item Construct a law-level operator layer $A_{\mathrm{law}}$ from
        law-parallel transport, law-curvature, and PL-PI transformations;
  \item Equip $A_{\mathrm{law}}$ with a GNLA/$L_\infty$ structure to 
        obtain the original law-level \\ \GAlgebra\
        $\GAlgebra_{\mathrm{orig}}$;
  \item Represent $\GAlgebra_{\mathrm{orig}}$ via PFB-GNLA structures
        (strict, homotopy, variable functional--operator), yielding the
        four versions of PFB-GNLA studied in earlier parts.
\end{itemize}

In this sense, the original PFB-GNLA is not “another” version on top of
the previous three; it is the \emph{source} of which the others are
geometric and operator-level projections. The PL-PI Meta-Mathematical
Theory provides the provenance of this law-level \GAlgebra; PFB
geometry and operator paths provide its concrete realisations in
physics, AI, and life/complex systems, which will be the focus of the application
chapters in Part~V.



\chapter[Four Versions of the G-Algebra: Comparison]{Four Versions of the GaoZheng G-Algebra: Comparison and Alignment}


\section{Overview of Structural Data and a Quick Map}


This chapter is devoted to comparing the four integrated PFB-GNLA
constructions developed in Parts~II and~III. Before entering into
detailed comparison statements and Gap-type theorems, it is useful to
collect, in one place, the structural data carried by each version of
the GaoZheng G-Algebra (\GAlgebra), and to provide a “quick map” that
shows how these versions embed into one another and into the law-space
\GZLS\ and the GaoZheng Generalized Mathematical Structure \GZGMS. This
section provides such an overview.



\subsubsection*{(a) Structural data for the four PFB-GNLA versions}

For reference, we summarise here the core structural data of the four
versions of PFB-GNLA, using a common schematic pattern.

\paragraph{(1) Strict version.}

The strict version is built from a parameterised principal-bundle
geometry and yields a near-Lie operator GNLA:

\begin{itemize}
  \item \emph{Geometric--parameter block}
    \[
      \mathsf{GeomPar}
      =
      \bigl(
        \Bundle,\Base,\Group;
        \Lambda_{\mathrm{par}};
        \{\Conn_\lambda\},
        \widehat{\Conn},\widehat{\Curv}
      \bigr).
    \]
  \item \emph{Strict operator block}
    \[
      \mathsf{Op}_{\mathrm{strict}}
      =
      \bigl(
        A_{\mathrm{strict}},
        [\cdot,\cdot]_{\mathrm{strict}},
        \Jac_{\mathrm{strict}};
        \mathsf{Gr},\mathsf{Filt}
      \bigr),
    \]
    with GNLA constraints (S1)--(S4), Jacobiator confined to a
    curvature sector $J_{\mathrm{curv}}$.
  \item \emph{Strict law-level block}
    \[
      L_{\mathrm{strict}}
      \in
      {\GZLS}_{\mathrm{strict}}
      \subset\GZLS,
    \]
    with law-connection ${\GZLOC}_{\mathrm{strict}}$ and law-curvature
    \GZLCurv\ either trivial or central in the strict regime.
  \item \emph{Strict PFB-GNLA datum}
    \[
      \mathsf{PFB\text- GNLA}_{\mathrm{strict}}
      =
      \bigl(
        \mathsf{GeomPar},
        \mathsf{Op}_{\mathrm{strict}},
        L_{\mathrm{strict}};
        \Phi_{\mathrm{geom}},
        \mathcal{C}_{\mathrm{strict}}
      \bigr).
    \]
\end{itemize}

\paragraph{(2) Homotopy version.}

The homotopy version adds a closed 3-form $H$ and promotes the strict
GNLA to an $H$-twisted two-term $L_\infty$-structure in GZ-Harmony:

\begin{itemize}
  \item \emph{Homotopy geometric--parameter--$H$ block}
    \[
      \mathsf{GeomPar}_H
      =
      \bigl(
        \Bundle,\Base,\Group;
        \Lambda_{\mathrm{par}};
        \{\Conn_\lambda\},
        \widehat{\Conn},\widehat{\Curv};
        H
      \bigr),
      \quad \mathrm{d}H=0.
    \]
  \item \emph{Homotopy operator block}
    \[
      \mathsf{Op}_{\mathrm{homo}}
      =
      \bigl(
        V^{-1}\oplus V^0,
        \ell_1,\ell_2,\ell_3^H;
        \mathsf{Gr}_{\mathrm{homo}},
        \mathsf{Filt}_{\mathrm{homo}}
      \bigr),
    \]
    where $(V^{-1}\oplus V^0,\ell_1,\ell_2,\ell_3^H)$ is an
    $H$-twisted two-term $L_\infty$-structure in GZ-Harmony with
    $(H,\GZTLC)$.
  \item \emph{Homotopy law-level block}
    \[
      L_{\mathrm{homo}}
      \in
      {\GZLS}_{\mathrm{homo}}
      \subset\GZLS,
    \]
    with internal homotopy data induced from
    $(\ell_1,\ell_2,\ell_3^H)$.
  \item \emph{Homotopy PFB-GNLA datum}
    \[
      \mathsf{PFB\text- GNLA}_{\mathrm{homo}}
      =
      \bigl(
        \mathsf{GeomPar}_H,
        \mathsf{Op}_{\mathrm{homo}},
        L_{\mathrm{homo}};
        \Phi_{\mathrm{geom}}^{\mathrm{homo}},
        \mathcal{C}_{\mathrm{homo}}
      \bigr).
    \]
\end{itemize}

\paragraph{(3) Variable functional--operator homotopy version.}

The variable functional--operator version works on operator paths and
their functionals, with a lifted connection $\widetilde{\GZTLC}$:

\begin{itemize}
  \item \emph{Variable functional--operator geometric block}
    \[
      \mathsf{GeomPar}_{\mathrm{vf\text- op}}
      =
      \bigl(
        \Bundle,\Base,\Group;
        \Lambda_{\mathrm{par}};
        \{\Conn_\lambda\},
        \widehat{\Conn},\widehat{\Curv};
        H;
        V^0,\mathsf{Path}(V^0)
      \bigr),
    \]
    with \GZTLC\ lifted to $\widetilde{\GZTLC}$ on path space.
  \item \emph{Variable functional--operator block}
    \[
      \mathsf{VF\text- Op}
      =
      \bigl(
        \mathsf{Path}(V^0),
        \{\mathcal{F}_\alpha\};
        \mathcal{E},
        \nabla^{\mathrm{path}};
        \ell_\bullet^{\mathrm{vf}}
      \bigr),
    \]
    capturing operator paths, variable functionals, path-wise
    connection, and path-level homotopy operations.
  \item \emph{Variable functional--operator law-level block}
    \[
      L_{\mathrm{vf\text- op}}
      \in
      {\GZLS}_{\mathrm{vf\text- op}}
      \subset\GZLS.
    \]
  \item \emph{Variable functional--operator PFB-GNLA datum}
    \[
      \mathsf{PFB\text- GNLA}_{\mathrm{vf\text- op}}
      =
      \bigl(
        \mathsf{GeomPar}_{\mathrm{vf\text- op}},
        \mathsf{VF\text- Op},
        L_{\mathrm{vf\text- op}};
        \Phi_{\mathrm{geom}}^{\mathrm{vf\text- op}},
        \mathcal{C}_{\mathrm{vf\text- op}}
      \bigr).
    \]
\end{itemize}

\paragraph{(4) Original law-level version.}

The original version is defined directly in the PL-PI law-space and
\GZLS, with PFB-GNLA structures recovered as representations:

\begin{itemize}
  \item \emph{GZ-admissible PL-PI law-system}
    \[
      \mathsf{LS}_{\mathrm{PL\text- PI}}
      =
      (L_{\mathrm{PL}},L_{\mathrm{PI}};
       \Theta_{\mathrm{PL\text- PI}})
      \in
      \mathsf{Law}_{\mathrm{PL\text- PI}}^{\mathrm{GZ}}.
    \]
  \item \emph{Law-geometry}
    \[
      \mathsf{LG}
      =
      \bigl(
        \mathfrak{L}_{\mathrm{base}},
        \mathfrak{L}_{\mathrm{fib}},
        \GZLOC^{\mathrm{tot}},
        \GZLCurv^{\mathrm{tot}}
      \bigr),
    \]
    fully inside \GZLS.
  \item \emph{Law-level operator layer and original \GAlgebra}
    \[
      \GAlgebra_{\mathrm{orig}}
      =
      \bigl(
        A_{\mathrm{law}},
        [\cdot,\cdot]_{\mathrm{law}},
        \Jac_{\mathrm{law}};
        \mathsf{Str}_{\mathrm{law}}
      \bigr),
    \]
    defined purely at the law-level.
\end{itemize}

\subsubsection*{(b) A quick map between the four versions}

At a schematic level, the relationship between the four versions can be
summarised by the following “tower” of structures:
\[
   \GAlgebra_{\mathrm{strict}}
   \;\rightsquigarrow\;
   \GAlgebra_{\mathrm{homo}}
   \;\rightsquigarrow\;
   \GAlgebra_{\mathrm{vf\text- op}}
   \;\subset\;
   \GAlgebra_{\mathrm{orig}},
\]
where:

\begin{itemize}
  \item $\GAlgebra_{\mathrm{strict}}$ is a near-Lie GNLA realised on
        PFB geometry with controlled curvature, corresponding to a
        near-flat region of \GZLS;
  \item $\GAlgebra_{\mathrm{homo}}$ adds $H$-twisted homotopy data
        $(\ell_1,\ell_2,\ell_3^H,\dots)$, corresponding to PL-PI law-
        systems with controlled but non-negligible law-curvature;
  \item $\GAlgebra_{\mathrm{vf\text- op}}$ further promotes operators
        to operator paths and equips them with variable functionals and
        a path-lifted connection $\widetilde{\GZTLC}$;
  \item $\GAlgebra_{\mathrm{orig}}$ is the global law-level object in
        PL-PI law-space, from which the other three versions are
        recovered as coordinate charts and truncations.
\end{itemize}

Each arrow is to be understood as:

\begin{itemize}
  \item an \emph{embedding} when read from left to right (strict
        structures extend to homotopy and path-level structures);
  \item a \emph{projection} or \emph{strictification} (when permitted
        by Jacobiator-gap and curvature constraints) when read from
        right to left.
\end{itemize}



\subsubsection*{(c) Quick map by questions: which version for what?}

For later reference, it is convenient to index the four versions by the
type of questions they are best adapted to.

\begin{itemize}
  \item \textbf{Strict version}:
    \begin{itemize}
      \item Questions about “pure PFB geometry + GNLA” where curvature
            is small or can be centralised;
      \item alignment with classical differential geometry and gauge
            theory language;
      \item baseline for Gap statements: “what happens in the ideal,
            near-flat regime?”
    \end{itemize}
  \item \textbf{Homotopy version}:
    \begin{itemize}
      \item Questions where Jacobi failure and higher law-curvature
            effects matter but can still be encoded in a controlled
            $L_\infty$ structure;
      \item design and interpretation of Jacobiator-gap certificates
            and GZ-OHU-type theorems;
      \item law-level counterparts of “topologically non-trivial”
            situations.
    \end{itemize}
      \item \textbf{Variable functional--operator version}:
    \begin{itemize}
      \item Questions about law–operator dynamics along paths (renormalisation
            flows, learning trajectories, control loops);
      \item use of variable functionals as actions, costs, or
            certificates along operator paths;
      \item explicit implementation of $\widetilde{\GZTLC}$ in dynamic
            settings (physics, AI, life/complex systems).
    \end{itemize}
  \item \textbf{Original law-level version}:
    \begin{itemize}
      \item Questions where PL-PI provenance is primary: “given a
            meta-level specification of inference and procedure, what
            PFB-GNLA structures can emerge?”;
      \item design of law-level models before choosing any particular
            geometric realisation;
      \item global comparison between different PFB-GNLA charts and
            their law-level equivalence classes.
    \end{itemize}
\end{itemize}



\subsubsection*{(d) Placement inside \texorpdfstring{\GZLS}{GZ-LS} and \texorpdfstring{\GZGMS}{GZ-GMS}}

Finally, we record the placement of the four versions inside the law-
space and the generalised structure:

\begin{itemize}
  \item In \GZLS, we have nested slices
    \[
      {\GZLS}_{\mathrm{strict}}
      \subset
      {\GZLS}_{\mathrm{homo}}
      \subset
      {\GZLS}_{\mathrm{vf\text- op}}
      \subset
      \mathsf{Law}_{\mathrm{PL\text- PI}}^{\mathrm{GZ}}
      \subset
      \GZLS,
    \]
    with law-curvature and homotopy visibility increasing along the
    chain.
  \item In \GZGMS, the four versions correspond to nested choices in \\ 
        $(\mathsf{Obj},\mathsf{Law},\mathsf{Flow},\mathsf{Layer},
          \GZTLC;\GZLCurv)$:
    \begin{itemize}
      \item \emph{Strict}: object-level PFB geometries with controlled
            curvature, low CH-layers ($l_2$);
      \item \emph{Homotopy}: same geometries plus $H$-twisted homotopy
            law-objects, higher CH-layers;
      \item \emph{Variable functional--operator}: previous data plus
            operator paths and variable functionals, with lifted
            connection $\widetilde{\GZTLC}$;
      \item \emph{Original}: PL-PI law-space as primary, with all other
            levels reconstructed as emergent representations.
    \end{itemize}
\end{itemize}

This quick structural overview will be used repeatedly in the remainder
of Part~III as a reference frame: comparison theorems, Gap statements,
and certificate-based results will be phrased in terms of how they
behave across these four nested versions and the corresponding slices of
\GZLS\ and charts of \GZGMS.



\section[From Strict Geometry to Law--Operator Homotopy]{From Strict Geometry to Law--Operator Homotopy \\
         and Meta-Mathematical Original}


The purpose of this section is to make explicit how the four PFB-GNLA
constructions are connected \emph{as a process} rather than as four
static objects. Starting from strict principal-bundle geometry and its
near-Lie operator algebra, we trace the controlled deformations that
lead first to law--operator homotopy (homotopy and variable
functional--operator versions), and then to the meta-mathematical
original version where all PFB and operator data are reinterpreted as
emergent from PL-PI law-systems.



\subsubsection*{(a) From strict PFB geometry to strict law--operator structures}

We begin with the strict PFB-GNLA datum
\[
  \mathsf{PFB\text- GNLA}_{\mathrm{strict}}
  =
  \bigl(
    \mathsf{GeomPar},
    \mathsf{Op}_{\mathrm{strict}},
    L_{\mathrm{strict}};
    \Phi_{\mathrm{geom}},
    \mathcal{C}_{\mathrm{strict}}
  \bigr),
\]
as summarised in Section~8.1. Conceptually, this datum performs three
translations:

\begin{enumerate}[label=\roman*)]
  \item \emph{From geometry to operators.}
        The geometric--parameter block \\
        $\mathsf{GeomPar}=(\Bundle,\Base,\Group;
           \Lambda_{\mathrm{par}};\{\Conn_\lambda\},
           \widehat{\Conn},\widehat{\Curv})$
        is converted to an operator layer
        $A_{\mathrm{strict}}$ via covariant derivatives, curvature
        operators, and their commutators.
  \item \emph{From operators to law-space.}
        The operator block \\
        $\mathsf{Op}_{\mathrm{strict}}=
          (A_{\mathrm{strict}},
           [\cdot,\cdot]_{\mathrm{strict}},
           \Jac_{\mathrm{strict}};\mathsf{Gr},\mathsf{Filt})$
        is realised as a law-object \\
        $L_{\mathrm{strict}}\in{\GZLS}_{\mathrm{strict}}\subset\GZLS$
        via the embedding $\iota_{\mathrm{law}}$ and the functor
        $\Phi_{\mathrm{geom}}$.
  \item \emph{From law-space back to geometry (consistency).}
        The strict coupling $\mathcal{C}_{\mathrm{strict}}$ ensures
        that parallel transport and curvature in law-space are
        compatible with the original PFB geometry; in particular,
        strict GNLA constraints (S1)--(S4) guarantee that Jacobi
        failure is either absent or confined to a negligible curvature
        sector $J_{\mathrm{curv}}$.
\end{enumerate}

The upshot is a \emph{near-Lie} law--operator structure:
\[
   \GAlgebra_{\mathrm{strict}}
   =
   \bigl(
     A_{\mathrm{strict}},
     [\cdot,\cdot]_{\mathrm{strict}},
     \Jac_{\mathrm{strict}};\mathsf{Str}_{\mathrm{strict}}
   \bigr)
   \;\in\;
   {\GZLS}_{\mathrm{strict}},
\]
realised concretely on PFB geometry in a region of law-space where
law-curvature is flat or central and CH-layering is effectively at
$l_2$.



\subsubsection*{(b) Turning on controlled homotopy: from strict to law--operator homotopy}

The first deformation away from strictness is to \emph{stop} treating
the Jacobiator and law-curvature as negligible. Instead, we promote
them to first-class data and encode them in an $H$-twisted two-term
$L_\infty$-structure.

Formally, we move from $\mathsf{GeomPar}$ to
$\mathsf{GeomPar}_H$ by introducing a closed 3-form
\[
  H\in\Omega^3(\Base)
  \quad\text{with}\quad
  \mathrm{d}H=0,
\]
whose law-level shadow enters \GZLCurv\ and generates a homotopy
Jacobiator $\ell_3^H$ in GZ-Harmony with \GZTLC. On the operator side,
$A_{\mathrm{strict}}$ is embedded into $V^0$, and a defect space
$V^{-1}$ is introduced to host Jacobiator and curvature mismatch
modes.

The deformation
\[
   \GAlgebra_{\mathrm{strict}}
   \leadsto
   \GAlgebra_{\mathrm{homo}}
\]
is then implemented by
\[
  (A_{\mathrm{strict}},[\cdot,\cdot]_{\mathrm{strict}},
   \Jac_{\mathrm{strict}})
  \;\rightsquigarrow\;
  (V^{-1}\oplus V^0,\ell_1,\ell_2,\ell_3^H),
\]
subject to:

\begin{itemize}
  \item $\ell_2$ extends the strict commutator on $A_{\mathrm{strict}}$;
  \item $\Jac_{\ell_2}=\ell_1\circ\ell_3^H$, so Jacobi failure is a
        homotopy boundary;
  \item $\ell_3^H$ is tied to $(H,\GZLCurv)$ via the curvature–homotopy
        map $\mathcal{K}_H$ and GZ-Harmony.
\end{itemize}

Geometrically, this corresponds to opening up law-space from the
strict slice ${\GZLS}_{\mathrm{strict}}$ into the homotopy slice
${\GZLS}_{\mathrm{homo}}$, where law-curvature is non-trivial but
controlled, and CH-layering has moved from $l_2$ to $l_{\ge3}$.



\subsubsection*{(c) Adding dynamics: from homotopy \texorpdfstring{\GAlgebra}{G-Algebra} to law--operator dynamics}

The second deformation concerns \emph{dynamics}: instead of static
operators in $V^0$, we consider operator paths $\Gamma\in\mathsf{Path}(V^0)$
and variable functionals on these paths, governed by a lifted
connection $\widetilde{\GZTLC}$.

In law--operator terms, the transition
\[
   \GAlgebra_{\mathrm{homo}}
   \leadsto
   \GAlgebra_{\mathrm{vf\text- op}}
\]
has three main components:

\begin{enumerate}[label=\roman*)]
  \item \emph{Path upgrade.}
        Replace $V^0$ by $\mathsf{Path}(V^0)$ and enrich the law-space
        slice ${\GZLS}_{\mathrm{homo}}$ to
        ${\GZLS}_{\mathrm{vf\text- op}}$, where law-objects carry path-
        dependent structure $L_\Gamma$.
  \item \emph{Functional layer.}
        Introduce variable functionals
        $\mathcal{F}_\alpha:\mathsf{Path}(V^0)\to W_\alpha$ and
        law-structured path bundles
        $\mathcal{E}\to\mathsf{Path}(V^0)$ with connection
        $\nabla^{\mathrm{path}}$ induced from $\widetilde{\GZTLC}$.
  \item \emph{Dynamic homotopy.}
        Extend the homotopy operations to
        $\ell_\bullet^{\mathrm{vf}}$ acting on sections of
        $\mathcal{E}$, so that covariant differentiation along path
        directions respects the $L_\infty$ structure up to controlled
        homotopies:
        \[
          \nabla^{\mathrm{path}}(\ell_2)
          \sim
          \ell_1\circ\ell_3^H,
          \quad\text{etc.}
        \]
\end{enumerate}

The curvature of $\widetilde{\GZTLC}$ splits into triple-layer,
path, and mixed components
$\mathcal{F}_{\GZTLC}+\mathcal{F}_{\mathrm{path}}
 +\mathcal{F}_{\mathrm{dyn}}$.
The mixed part $\mathcal{F}_{\mathrm{dyn}}$ is precisely what we
interpret as \emph{law--operator dynamical curvature}: the obstruction
to seeing operator dynamics as mere triple-layer parallel transport.



\subsubsection*{(d) Inverting the perspective: from law--operator dynamics to meta-mathematical original}

The third step is conceptually the most radical: instead of regarding
PL-PI law-systems as background data for PFB geometries and operator
dynamics, we regard \emph{PFB geometries and law--operator structures
as representations of a law-level \GAlgebra\ that is defined directly
in the PL-PI law-space}.

We summarised this in Section~7.2 by introducing the original
\GAlgebra:
\[
   \GAlgebra_{\mathrm{orig}}
   =
   (A_{\mathrm{law}},
    [\cdot,\cdot]_{\mathrm{law}},
    \Jac_{\mathrm{law}};\mathsf{Str}_{\mathrm{law}}),
\]
built from:

\begin{itemize}
  \item a GZ-admissible PL-PI law-system
        $\mathsf{LS}_{\mathrm{PL\text- PI}}\in
         \mathsf{Law}_{\mathrm{PL\text- PI}}^{\mathrm{GZ}}$,
  \item an associated law-geometry
        $\mathsf{LG}=(\mathfrak{L}_{\mathrm{base}},
                      \mathfrak{L}_{\mathrm{fib}},
                      \GZLOC^{\mathrm{tot}},
                      \GZLCurv^{\mathrm{tot}})$,
  \item a law-operator layer $A_{\mathrm{law}}$ generated by law-
        parallel transport, law-curvature, and PL-PI transforms.
\end{itemize}

The key conceptual inversion is:

\begin{quote}
  In the strict / homotopy / path-level versions, PL-PI structure is
  “forgotten” or only indirectly visible; PFB geometry and operators
  are primary. In the meta-mathematical original version, PL-PI
  structure is primary; PFB geometry and operators are \emph{secondary}
  realisations or coordinate expressions of $\GAlgebra_{\mathrm{orig}}$.
\end{quote}

Concretely, this inversion is captured by reconstruction functors of
the form
\[
   \mathcal{R}_{\mathrm{PFB}}:
     \mathsf{Law}_{\mathrm{PL\text- PI}}^{\mathrm{GZ}}
     \longrightarrow
     \mathsf{GeomObj}^{\mathrm{PFB}},
\]
and by representation maps
\[
   \rho_{\mathrm{strict}},
   \rho_{\mathrm{homo}},
   \rho_{\mathrm{vf\text- op}}:
     \GAlgebra_{\mathrm{orig}}
     \longrightarrow
     \mathrm{End}(\mathcal{H}_{\bullet}),
\]
where each $\mathcal{H}_\bullet$ is a space of sections, operators, or
operator paths appropriate to the strict, homotopy, or path-level
version.



\subsubsection*{(e) A commuting picture: geometry, homotopy, dynamics, provenance}

We can now summarise the situation by a schematic commuting picture.
Let us denote, at a very high level:

\begin{tikzcd}[column sep=3.8em,row sep=2.8em]
  & \GAlgebra_{\mathrm{homo}}
    \arrow[dr,hook,"{\text{path upgrade}}"] & \\
  \GAlgebra_{\mathrm{strict}}
    \arrow[ur,hook,"{\text{turn on }H,\ell_3^H}"]
    \arrow[dr,hook',"{\text{embed}}"'] & &
  \GAlgebra_{\mathrm{vf\text{-}op}}
    \arrow[dl,hook,"{\text{representation}}"'] \\
  & \GAlgebra_{\mathrm{orig}} &
\end{tikzcd}

Here:

\begin{itemize}
  \item the upper left arrow corresponds to switching on controlled
        law-curvature and homotopy brackets (homotopy upgrade);
  \item the upper right arrow corresponds to promoting static operators
        to operator paths with variable functionals (dynamic upgrade);
  \item the lower arrows correspond to embedding these PFB-GNLA
        versions as representations or charts of the original law-level
        \GAlgebra.
\end{itemize}

From the perspective of the overall monograph, this diagram encodes two
orthogonal directions:

\begin{itemize}
  \item a \emph{vertical direction} (bottom to top) capturing the
        passage from PL-PI provenance and law-level structure to
        concrete PFB geometries and operators;
  \item a \emph{horizontal direction} (left to right) capturing the
        passage from static, near-Lie structures to homotopy-rich,
        path-dependent law--operator dynamics.
\end{itemize}



\subsubsection*{(f) Summary: from strict geometry to original meta-mathematical law}

To conclude, the phrase “From Strict Geometry to Law–Operator Homotopy and \\ Meta-Mathematical Original” can be unpacked as the following
layered narrative:

\begin{enumerate}[label=\arabic*)]
  \item \textbf{Strict geometry.}
        Start with principal bundles, connections, and curvature; build
        a near-Lie operator algebra $A_{\mathrm{strict}}$ and embed it
        into a near-flat law-space slice ${\GZLS}_{\mathrm{strict}}$.
  \item \textbf{Law–operator homotopy.}
        Turn on controlled law-curvature and a closed 3-form $H$;
        promote $A_{\mathrm{strict}}$ to an $H$-twisted two-term
        $L_\infty$-algebra $(V^{-1}\oplus V^0,\ell_1,\ell_2,\ell_3^H)$ \\
        in GZ-Harmony; extend to operator-path dynamics and variable
        functionals via $\widetilde{\GZTLC}$.
  \item \textbf{Meta-mathematical original.}
        Recognise that all these geometric and operator constructions
        are representations of an original law-level \GAlgebra\
        $\GAlgebra_{\mathrm{orig}}$ living in the PL-PI law-space,
        whose curvature, homotopy, and CH-layering structure are
        primary; PFB geometries and law--operator dynamics are its
        various charts and slices.
\end{enumerate}

Seen in this way, the four integrated PFB-GNLA constructions are not
four unrelated theories but four \emph{windows} onto a single meta-
mathematical object: the law-level GaoZheng G-Algebra emergent from the
PL-PI Meta-Mathematical Theory. The remaining sections of Part~III will
exploit this unified picture to formulate precise comparison theorems,
Gap statements, and certificate-based results that apply simultaneously
across all four versions.



\section{GZ-OHU: \ Operator--Homotopy Universality and Failure-Free Construction}


The GZ-OHU (GaoZheng Operator--Homotopy Universality) framework is the
point where the four PFB-GNLA versions meet. At an intuitive level, it
makes two claims:

\begin{itemize}
  \item \emph{Universality:} once the Jacobiator and higher homotopy
        data of a law-level operator system are brought under control,
        there is a \emph{universal} homotopy completion
        $L^{\mathrm{OHU}}$ that realises all operator--homotopy types
        compatible with those constraints.
  \item \emph{Failure-free construction:} there exist constructive
        schemes that build $L^{\mathrm{OHU}}$ without ever getting
        stuck in an un-certified failure mode: any intermediate
        breakdown of strict Jacobi behaviour is either explicitly
        detected by a finite set of invariants or repaired by a finite
        sequence of homotopy moves.
\end{itemize}

This section formalises these ideas in terms of Jacobiator-gap
certificates and states the schematic GZ-OHU theorem, in a way that
integrates and extends the earlier informal uses of these notions.

\subsubsection*{(a) Law-level operator systems and homotopy data}

We work at the level of law-space. A \emph{law-level operator system}
$L$ should be understood as any of the following, depending on the
context:

\begin{itemize}
  \item a homotopy version of \GAlgebra\ living in
        ${\GZLS}_{\mathrm{homo}}$, with underlying
        $H$-twisted two-term $L_\infty$-structure
        $(V^{-1}\oplus V^0,\ell_1,\ell_2,\ell_3^H)$;
  \item a variable functional--operator version living in
        ${\GZLS}_{\mathrm{vf\text- op}}$, where the same homotopy
        structure is enriched by operator paths and variable
        functionals, and transported by $\widetilde{\GZTLC}$;
  \item the law-level original version $\GAlgebra_{\mathrm{orig}}$
        restricted to an appropriate homotopy sector.
\end{itemize}

In all cases, $L$ carries at least:

\begin{itemize}
  \item a binary operation (graded commutator)
        $\ell_2:[\cdot,\cdot]$;
  \item a higher homotopy operation $\ell_3$ that encodes the
        Jacobiator via relations of the form
        $\Jac_{\ell_2} = \ell_1\circ\ell_3$;
  \item a law-connection (either \GZLOC\ or a lifted version
        $\widetilde{\GZTLC}$) that governs how these algebraic
        structures are transported along law-space or operator-path
        flows.
\end{itemize}

The GZ-OHU framework asks: \emph{how far} is $L$ from being strict, and
\emph{how much homotopy completion} is needed so that all admissible
operator--homotopy types compatible with its “would-be strict” part are
present? The answer is measured by Jacobiator-gap certificates.

\subsubsection*{(b) Jacobiator-gap certificates}

A Jacobiator-gap certificate is a finite family of law-level invariants
that quantifies the controlled failure of strict Jacobi behaviour and
tracks how this failure behaves under law-connections and homotopies.

\begin{definition}[Jacobiator-gap certificate (schematic)]
  Let $L$ be a law-level operator system (for example, a homotopy
  version of \GAlgebra). A \emph{Jacobiator-gap certificate} for $L$
  is a finite collection of invariants
  \[
    \mathsf{JacGap}(L) = \bigl(c_1(L),\dots,c_k(L)\bigr),
  \]
  constructed from the higher brackets and homotopy data of $L$
  (typically from $\ell_2,\ell_3$ and their interaction with
  law-curvature), such that:
  \begin{enumerate}[label=\roman*)]
    \item \emph{Quantitative Jacobi failure.}
          Each $c_i(L)$ quantitatively measures a controlled failure
          of the strict Jacobi identity; for instance, $c_i(L)$ may be
          defined as a norm or spectral radius of certain components
          of the Jacobiator projected to distinguished defect sectors.
    \item \emph{Vanishing characterisation.}
          $\mathsf{JacGap}(L)=0$ if and only if $L$ can be promoted to
          a strictly Jacobi-satisfying law-object
          $L_{\mathrm{strict}}$ (possibly after enlarging the space and
          adding auxiliary variables) \emph{without} changing its
          observable low-order behaviour (e.g.\ low-degree brackets,
          low-order correlators, or a prescribed set of evaluation
          functionals).
    \item \emph{Stability under law-connections.}
          $\mathsf{JacGap}(L)$ is stable under the deformations
          induced by the triple-layer connection \GZTLC\ and its
          path-lifted version $\widetilde{\GZTLC}$, in the sense that
          along any horizontal law-flow
          $t\mapsto L(t)$ one has controlled variation:
          \[
            \frac{\mathrm{d}}{\mathrm{d}t}\mathsf{JacGap}(L(t))
            \ \text{is bounded by curvature and homotopy terms of }
            \GZTLC/\widetilde{\GZTLC}.
          \]
  \end{enumerate}
\end{definition}

The precise choice of invariants $c_i$ depends on the level (strict,
homotopy, variable functional--operator, original) and on which
components of the Jacobiator are declared “observable” or “protected”.
The important point is that \emph{finitely many} numbers (or tensors)
encode all the obstructions that matter for passing to a strict regime
without altering the visible behaviour of $L$.

\subsubsection*{(c) GZ-OHU extension: operator--homotopy universality}

Given a law-level operator system $L$ with a Jacobiator-gap certificate,
the GZ-OHU framework asserts the existence of a homotopy-complete
extension $L^{\mathrm{OHU}}$ that is universal for operator-homotopy
types compatible with $\mathsf{JacGap}(L)=0$.

\begin{theorem}[GZ-OHU theorem (schematic)]
  Under suitable finiteness, regularity, and \\ boundedness assumptions on
  a law-level operator system $L$ that admits a Jacobiator-gap
  certificate $\mathsf{JacGap}(L)$, there exists a homotopy-complete
  extension $L^{\mathrm{OHU}}$ and a morphism
  \[
    \iota : L \longrightarrow L^{\mathrm{OHU}},
  \]
  satisfying:
  \begin{enumerate}[label=\roman*)]
    \item \emph{Operator--homotopy universality.}
          $L^{\mathrm{OHU}}$ realises all operator--homotopy types that
          are compatible with vanishing Jacobiator-gap. More precisely,
          if $L'$ is any law-level operator system such that
          $\mathsf{JacGap}(L')=0$ and whose low-order observables agree
          with those of $L$, then there exists a homotopy equivalence
          between $L^{\mathrm{OHU}}$ and $L'$ respecting the law-
          connection and CH-layering structure.
    \item \emph{Failure-free constructive schemes.}
          There exist constructive schemes (e.g.\ iterative extensions
          in the sense of $L_\infty$-envelopes, or law-space learning
          procedures guided by \GZTLC) for building $L^{\mathrm{OHU}}$
          from $L$ such that any intermediate violation of the strict
          Jacobi identity is either:
          \begin{itemize}
            \item detected by a non-zero component of
                  $\mathsf{JacGap}$, or
            \item resolved by a finite sequence of homotopy moves
                  (adjustments of $\ell_3$ and higher brackets) that
                  strictly decrease an appropriate “defect potential”
                  derived from $\mathsf{JacGap}$.
          \end{itemize}
  \end{enumerate}
  A complete statement and proof are given in Technical Appendix~I
  (Section~\ref{app:I-T1}).
\end{theorem}

Informally, $L^{\mathrm{OHU}}$ is the “maximally completed” homotopy
version of $L$ that still respects the constraint
$\mathsf{JacGap}(L)=0$. Any other such completion factors through it
up to homotopy. This is the operator--homotopy universality property.

\subsubsection*{(d) Failure-free construction across the four PFB-GNLA versions}

The “failure-free” aspect of GZ-OHU is best understood in the context
of the four PFB-GNLA versions:

\begin{itemize}
  \item In the \textbf{strict} version, one attempts to maintain
        strict Jacobi behaviour while exploring parameter and geometric
        deformations. GZ-OHU says that when strictness is about to
        break, this is signalled by $\mathsf{JacGap}$ becoming
        non-zero, and that a controlled homotopy upgrade is available.
  \item In the \textbf{homotopy} version, one works directly with
        $(V^{-1}\oplus V^0,\ell_1,\ell_2,\ell_3^H)$; GZ-OHU ensures
        that such homotopy data can be extended step by step to a
        homotopy-complete structure $L^{\mathrm{OHU}}$ without
        encountering unsignalled inconsistencies: any obstruction is
        either encoded in $\mathsf{JacGap}$ or removable via finitely
        many homotopy moves.
  \item In the \textbf{variable functional--operator} version,
        operator-path dynamics and variable functionals are added; the
        path-lifted connection $\widetilde{\GZTLC}$ may generate
        additional dynamical Jacobi defects. GZ-OHU asserts that
        these too can be classified by a finite-dimensional
        $\mathsf{JacGap}$ and compensated by path-level homotopies,
        yielding a failure-free upgrade of law--operator dynamics.
  \item In the \textbf{original law-level} version, the entire
        picture is seen directly in PL-PI law-space: GZ-OHU becomes a
        statement about the existence of universal law-level completions
        of GZ-admissible PL-PI systems and of failure-free law-space
        deformation procedures.
\end{itemize}

In all cases, the guiding principle is the same: \emph{do not attempt
to enforce strict identities blindly}. Instead, track deviations via
$\mathsf{JacGap}$, interpret them as signals of missing homotopy data,
and enlarge the law-level operator system along the directions indicated
by GZ-OHU until a homotopy-complete regime is reached.

\subsubsection*{(e) Relation to Gap theorems and CH-layering}

Jacobiator-gap certificates also underpin Gap-type statements and CH-
layering behaviour:

\begin{itemize}
  \item \emph{Gap theorems.} Roughly speaking, if
        $\mathsf{JacGap}(L)$ is bounded away from zero along any
        path connecting $L$ to a strict regime, then there is a
        “Jacobiator gap”: no continuous deformation within the
        admissible class can make $L$ strictly Jacobi without crossing
        a region of law-space where certain curvature/homotopy
        invariants are non-negligible. GZ-OHU identifies the
        universal homotopy-complete object $L^{\mathrm{OHU}}$ that
        sits just at the boundary of this gap.
  \item \emph{CH-layering.} At lower CH-layers (e.g.\ $l_2$), the
        contribution of $\mathsf{JacGap}$ may be invisible, and
        $L^{\mathrm{OHU}}$ and $L_{\mathrm{strict}}$ may have identical
        observables. At higher layers ($l_{\ge3}$), semantic and
        homotopy metrics become sensitive to $\mathsf{JacGap}$ and to
        the full homotopy structure of $L^{\mathrm{OHU}}$, revealing
        the “hidden” operator-homotopy universality behind apparently
        strict behaviours.
\end{itemize}

Thus, GZ-OHU sits at the intersection of three themes of this part of
the monograph: \emph{operator homotopy}, \emph{Gap phenomena}, and
\emph{CH-layering of law-space}.

\subsubsection*{(f) Summary}

The GZ-OHU framework can be summarised in three key points:

\begin{enumerate}[label=\arabic*)]
  \item \textbf{Measurement:} Jacobiator-gap certificates
        $\mathsf{JacGap}(L)$ provide a finite measurement of how far a
        law-level operator system is from being strictly Jacobi, in a
        way that is stable under law-connections and compatible with
        the homotopy structure of \GAlgebra.
  \item \textbf{Completion:} For any such system, there exists a
        homotopy-complete extension $L^{\mathrm{OHU}}$ that is
        universal for operator-homotopy types compatible with
        $\mathsf{JacGap}(L)=0$; all other completions factor through
        it up to homotopy.
  \item \textbf{Failure-free construction:} There are constructive
        schemes, aligned with \GZTLC\ and its path-level lift, that
        build $L^{\mathrm{OHU}}$ from $L$ without unsignalled
        breakdowns of Jacobi behaviour: every obstruction is either
        detected or repaired in finitely many steps.
\end{enumerate}

In this sense, GZ-OHU provides the operator-homotopy backbone of the
four-way PFB-GNLA comparison: it explains how strict, homotopy,
variable functional--operator, and original law-level versions of
\GAlgebra\ can be navigated without conceptual dead-ends, and how
their interrelation is governed by a single family of law-level
invariants, $\mathsf{JacGap}$.



\part{CH-Layering and the Breakdown/Reconstruction of Continuum Semantics}


\chapter[Overview of CH-Layering]{Overview of CH-Layering:\\
         A Structural Analogy with the Fifth Postulate Revolution}


\section{Classical CH and Independence Results: A Brief Review}


The aim of this section is to recall, in a compact but conceptually
coherent way, the classical story surrounding the Continuum Hypothesis
(CH) and its independence from ZFC. This will serve as the baseline
against which the CH-layering framework of the later sections is to be
positioned: what exactly is being “layered” when we speak of CH, and
where do the usual independence results leave genuine room for a
law-level reinterpretation.



\subsubsection*{(a) The classical formulation of CH and related hypotheses}

We start with the standard set-theoretic background.

\begin{definition}[Cardinalities, continuum, and initial alephs]
  Let $\mathbb{N}=\{0,1,2,\dots\}$, $\mathbb{R}$ the real line, and
  $|\cdot|$ denote cardinality. The \emph{cardinality of the continuum}
  is
  \[
    \mathfrak{c} := |\mathbb{R}|.
  \]
  The infinite initial ordinals are denoted
  \[
    \aleph_0,\ \aleph_1,\ \aleph_2,\ \dots
  \]
  where $\aleph_0=|\mathbb{N}|$ and each $\aleph_\alpha$ is the least
  cardinal strictly greater than all $\aleph_\beta$ with $\beta<\alpha$.
\end{definition}

The Continuum Hypothesis can then be stated in several equivalent ways
in the presence of ZFC.

\begin{definition}[Continuum Hypothesis (CH)]
  The \emph{Continuum Hypothesis} is the assertion that there is no
  cardinal strictly between $\aleph_0$ and $\mathfrak{c}$, i.e.
  \[
    \text{CH:}\qquad
    \mathfrak{c} = \aleph_1.
  \]
  Equivalently: every infinite subset of $\mathbb{R}$ has cardinality
  either $\aleph_0$ or $\mathfrak{c}$.
\end{definition}

More generally, the \emph{Generalized Continuum Hypothesis} (GCH)
asserts that the power-set of each infinite cardinal is the next
cardinal:

\begin{definition}[Generalized Continuum Hypothesis (GCH)]
  The \emph{Generalized Continuum Hypothesis} states that for every
  infinite cardinal $\kappa$,
  \[
    2^\kappa = \kappa^+,
  \]
  where $\kappa^+$ denotes the immediate successor cardinal of $\kappa$.
  In particular, at $\kappa=\aleph_0$ one recovers CH as
  $2^{\aleph_0}=\aleph_1$.
\end{definition}

At the purely combinatorial level, CH is a statement about the
cardinal arithmetic of the continuum; its subtlety arises from the fact
that ZFC does not settle this arithmetic. The modern view of this
subtlety is phrased in terms of \emph{models} of ZFC.

\begin{definition}[Models of ZFC and truth of CH]
  A (set-theoretic) \emph{model of ZFC} is a structure
  \[
     M = (|M|,\in^M)
  \]
  satisfying the axioms of ZFC, where $\in^M$ is a binary relation on
  the domain $|M|$ interpreting membership. One writes
  \[
    M \models \text{ZFC}
  \]
  to indicate that ZFC holds in $M$, and
  \[
    M \models \text{CH}
    \quad\text{or}\quad
    M \models \neg\text{CH}
  \]
  to record the truth value of CH \emph{inside} $M$.
\end{definition}



\subsubsection*{(b) Gödel's constructible universe: consistency of CH with ZFC}

The first major milestone is due to Gödel (1940): CH is \emph{not}
refutable from ZFC, assuming ZFC itself is consistent. The proof goes
via the constructible universe $L$.

\begin{definition}[Constructible universe $L$ (very briefly)]
  The \emph{constructible universe} $L$ is the class obtained by
  transfinite recursion
  \[
    L_0 := \emptyset,\qquad
    L_{\alpha+1} := \mathrm{Def}(L_\alpha),
  \]
  where $\mathrm{Def}(X)$ denotes the collection of subsets of $X$
  definable over $(X,\in)$ with parameters from $X$, and for limit
  ordinals
  \[
    L_\lambda := \bigcup_{\alpha<\lambda} L_\alpha.
  \]
  The class $L := \bigcup_{\alpha} L_\alpha$ is transitive and can be
  shown to satisfy ZFC.
\end{definition}

Gödel's analysis established:

\begin{theorem}[Gödel, 1940: ZFC + CH is consistent if ZFC is]
  If ZFC is consistent, then so is ZFC + CH. More precisely (see, for example, \cite{Godel1940,Jech2003}):
  \begin{enumerate}[label=\roman*)]
    \item $L$ is a model of ZFC: $L\models\text{ZFC}$;
    \item In $L$, the Generalized Continuum Hypothesis holds:
          $L\models\text{GCH}$, hence in particular $L\models\text{CH}$.
  \end{enumerate}
  A clean restatement and further bibliographical pointers are provided
  in Technical Appendix~III (Section~\ref{app:III-T2}); for full proofs,
  see, for example, \cite{Godel1940,Jech2003}.
\end{theorem}

The upshot is that ZFC cannot prove $\neg\text{CH}$, unless ZFC itself
is inconsistent. This is usually expressed as:

\[
  \text{Con(ZFC)} \ \Longrightarrow\ \text{Con(ZFC + CH)}.
\]



\subsubsection*{(c) Cohen forcing: consistency of $\neg\text{CH}$ with ZFC}

The second milestone, due to Cohen (1963), is that CH is also \emph{not
provable} from ZFC: there are models of ZFC where CH fails.

Cohen introduced the method of \emph{forcing}, which can be thought of
as a way of adjoining new sets (e.g.\ new reals) to a model of ZFC in a
controlled way.

\begin{definition}[Very schematic forcing extension]
  Let $M\models\text{ZFC}$ be a countable transitive model. A
  \emph{forcing notion} (or poset) $\mathbb{P}\in M$ is used to define
  $M$-generic filters $G\subseteq\mathbb{P}$; the \emph{forcing
  extension} $M[G]$ is a new model of ZFC, built from $M$ and the
  generic filter $G$, in which additional sets (often, new subsets of
  $\omega$) exist.
\end{definition}

Cohen showed that by forcing with a suitable poset that “adds new
reals”, one can make the continuum strictly larger than $\aleph_1$:

\begin{theorem}[Cohen, 1963: ZFC + $\neg\text{CH}$ is consistent if ZFC is]
  If ZFC is consistent, then so is ZFC + $\neg\text{CH}$. More precisely (see, for example, \cite{Cohen1963,Jech2003}):
  \begin{enumerate}[label=\roman*)]
    \item Starting from a ground model $M\models\text{ZFC}$ with CH
          (such as $L$), there exists a forcing notion $\mathbb{P}$ and
          an $M$-generic filter $G\subseteq\mathbb{P}$ such that the
          extension $M[G]$ satisfies:
          \[
             M[G]\models\text{ZFC} + \neg\text{CH}.
          \]
    \item In particular, one can arrange
          $M[G]\models 2^{\aleph_0} = \aleph_2$
          (and many other patterns of continuum cardinality).
  \end{enumerate}
  A clean restatement and further bibliographical pointers are provided
  in Technical Appendix~III (Section~\ref{app:III-T3}); for full proofs,
  see, for example, \cite{Cohen1963,Jech2003}.
\end{theorem}

Combining Gödel and Cohen, one obtains:

\[
  \text{Con(ZFC)}
  \ \Longrightarrow\
  \text{Con(ZFC + CH)}\ \wedge\ \text{Con(ZFC + $\neg$CH)}.
\]

Thus ZFC alone cannot decide CH; CH is \emph{independent} of ZFC.



\subsubsection*{(d) Beyond CH: continuum patterns, forcing axioms, and large cardinals}

Subsequent work has greatly refined our understanding of what the
continuum can look like in models of ZFC, and which additional
principles tend to push the continuum in one direction or another.

\paragraph{Continuum patterns and Easton’s theorem.}

Easton (1967) showed (see \cite{Easton1970,Jech2003}) that, subject to certain monotonicity and
coherence constraints, the function
\[
   \kappa \mapsto 2^\kappa
\]
at regular cardinals can be made to follow a wide variety of
prescriptions via forcing. In particular, even if one settles CH at
$\aleph_0$, the behaviour of $2^{\aleph_\alpha}$ for higher $\alpha$
can be arbitrarily wild, subject to these constraints. This illustrates
that CH is just the first instance of a much broader “continuum
pattern” problem.

\paragraph{Forcing axioms.}

Strong forcing axioms such as Martin’s Maximum (MM) or the Proper
Forcing Axiom (PFA), consistent relative to large cardinals, often
entail $\neg\text{CH}$ and impose strong regularity and structural
properties on subsets of $\mathbb{R}$ (e.g.\ projective determinacy in
certain fragments, strong reflection principles, etc.). In this sense,
for many set theorists, “natural” strengthenings of ZFC coming from
forcing axioms tend to favour $\neg\text{CH}$ (see, e.g., \cite{Jech2003,Kanamori2003}).

\paragraph{Large cardinals and inner models.}

On the other side, large cardinal axioms (measurable, supercompact,
Woodin cardinals, etc.) and inner model theory provide deep structural
constraints on the universe of sets (see, for example, \cite{Kanamori2003}). Some research programs (e.g.\ the
“Ultimate-$L$” program) aim to find canonical inner models that might
once again favour CH; others (e.g.\ Woodin’s work on $\Omega$-logic and
determinacy) suggest that strong regularity properties of reals push
towards $\neg\text{CH}$.

The resulting picture is highly non-trivial: \emph{different} “natural”
axiomatic extensions of ZFC point in different directions regarding CH.



\subsubsection*{(e) Conceptual lessons: why a law-level reinterpretation is meaningful}

For the purposes of this monograph, there are three conceptual lessons
from the classical CH story:

\begin{enumerate}[label=\arabic*)]
  \item \textbf{Extensional metrics are not enough.}
        Classical CH is formulated entirely in terms of cardinalities
        and membership-based structure. All models that differ only in
        “how many” reals there are, but keep the same extensional
        language of sets, look equally legitimate from the standpoint
        of ZFC.
  \item \textbf{Model pluralism is forced by ZFC.}
        Once one accepts ZFC as a baseline, Gödel and Cohen show that
        there will inevitably be multiple non-isomorphic models of ZFC
        disagreeing on CH. From a purely extensional viewpoint, there
        is no canonical way within ZFC to declare one continuum
        structure “the” true one.
  \item \textbf{Additional principles are meta-level and law-like.}
        Forcing axioms, large cardinal axioms, determinacy principles,
        and canonical inner model hypotheses all live at a “law-level”
        above ZFC: they prescribe global patterns of behaviour for
        sets and reals. Their influence on CH suggests that what is at
        stake is \emph{not} merely the size of $\mathbb{R}$, but the
        law-level regime governing the continuum (which regularities,
        which reflection principles, which determinacy patterns are
        deemed “lawful”).
\end{enumerate}

From the vantage point of the GaoZheng G-Framework, this motivates the
CH-layering programme:

\begin{itemize}
  \item treat classical CH as a \emph{layer-$l_2$} statement about
        extensional cardinalities;
  \item treat forcing axioms, large cardinal axioms, and related
        principles as part of a higher law-space structure governing
        the continuum, visible only in layers $l_{\ge3}$ where
        homotopy, curvature, and semantic metrics play a role;
  \item ask how the “breakdown” and “reconstruction” of continuum
        semantics can be modelled as a change of layer in a unified
        law-space, rather than as an absolute yes/no verdict on CH
        inside a single extensional background.
\end{itemize}



The next section will develop this idea explicitly, by introducing
CH-layering as a stratification of law-space: CH and its variants are
re-positioned as statements whose meaning depends on which layer
(extensional vs.\ semantic/homotopy) one is working in, and the
classical independence results become manifestations of the fact that
extensional layers alone cannot fully determine the law-level regime of
the continuum.



\section{Layering CH: From Cardinal Metrics to Semantic Metrics}


The classical treatment of the Continuum Hypothesis (CH) is entirely
\emph{extensional}: it asks about the existence of intermediate
cardinalities between $|\mathbb{N}|$ and $|\mathbb{R}|$. From the point
of view of the GaoZheng G-Framework and CH-layering, this corresponds
to working exclusively at a \emph{binary-law, extensional layer}
(denoted schematically by $l_2$), where “how many elements” is the
only meaningful metric. At higher layers $l_{\ge3}$, where law-space
curvature, homotopy, and generative semantics become visible, the
meaning of “size of the continuum” is no longer adequately captured by
cardinals alone.

This section formalises this idea by introducing a layered view of CH:
CH as an $l_2$-statement about cardinal metrics, contrasted with
CH-like statements at $l_{\ge3}$ where semantic/generative metrics
take over. The goal is not to “refute” classical independence results,
but to reposition them as statements about \emph{one particular layer}
of a richer law-level structure.



\subsubsection*{(a) The extensional layer $l_2$: cardinal metrics as the only visible structure}

We begin by isolating the layer in which classical CH naturally lives.
In the CH-layering language, this is the \emph{binary-law extensional
layer} $l_2$, characterised by:

\begin{itemize}
  \item objects are sets (or Polish spaces like $\mathbb{R}$) up to
        bijection;
  \item morphisms are (measurable/Borel) injections, surjections, and
        bijections, but only insofar as they witness cardinal
        inequalities;
  \item the primary metric is cardinality, or measures built “on top
        of” cardinality (counting measure, Lebesgue measure, etc.),
        which depend only on how many elements sit in a region, not on
        how they are generated or coupled.
\end{itemize}

Formally, one can think of a forgetful functor
\[
  U_{l_2} : \mathsf{Law}_{\mathrm{Cont}}
            \longrightarrow
            \mathsf{Set}_{\mathrm{card}},
\]
from some law-category of continuum-like objects to a category whose
only invariants are cardinalities and basic measure-theoretic
statistics.

\begin{definition}[Extensional CH at layer $l_2$]
  The \emph{layer-$l_2$ version} of the Continuum Hypothesis,
  denoted $\mathrm{CH}^{(l_2)}$, is simply the classical CH:
  \[
    \mathrm{CH}^{(l_2)} :\quad
    \mathfrak{c} = 2^{\aleph_0} = \aleph_1,
  \]
  or equivalently: there is no set $X$ such that
  \[
    |\mathbb{N}| < |X| < |\mathbb{R}|.
  \]
  All statements are interpreted in a model of ZFC, and all invariants
  are computed after applying the forgetful functor $U_{l_2}$.
\end{definition}

At this layer, two continuum-like systems $X$ and $Y$ that differ only
by “how they are generated” but agree in cardinality are indistinguish\-able:
$U_{l_2}(X)\cong U_{l_2}(Y)$.

From the CH-layering perspective, this is the layer where Gödel and
Cohen’s independence results live: they show that in the realm where
only $|\cdot|$ is visible, ZFC cannot determine whether
$\mathfrak{c}=\aleph_1$ or not.



\subsubsection*{(b) Higher layers $l_{\ge3}$: semantic and generative metrics}

At higher layers, one enriches the description of continuum-like
objects by tracking \emph{how} they are generated, how they are coupled
across scales, and what homotopy and law-curvature data they carry.

We denote by $l_{\ge3}$ a collection of layers in which at least some
of the following structures are visible:

\begin{itemize}
  \item \emph{generative complexity}: a measure
        $\sigma_{\ge3}(X)$ of how complex the process is that generates
        the continuum-like object $X$ (e.g.\ minimal description length
        of laws, homotopy depth of generating systems, law-curvature
        spectrum);
  \item \emph{semantic continuity}: a measure $\vartheta(X)$ of how
        coherent the semantic or structural relations are across the
        continuum (e.g.\ regularity of definable sets, determinacy
        patterns, reflection properties);
  \item \emph{coupling metrics}: measures of how strong the couplings
        are between different “faces” of the continuum (topology,
        measure, category, descriptive-set-theoretic hierarchy, etc.).
\end{itemize}

Schematically, we may package this as a “semantic metric vector”
\[
  \mathbf{S}(X)
  :=
  \bigl(
    \sigma_{\mathrm{gen}}(X),\,
    \vartheta_{\mathrm{sem}}(X),\,
    \kappa_{\mathrm{coup}}(X),\dots
  \bigr),
\]
which is invisible at $l_2$ (only cardinal $|X|$ is seen) but becomes
visible at $l_{\ge3}$.

\begin{definition}[n-th law layer and n-th order law-bundle space (schematic)]
  Let $\mathsf{Layer}$ denote the semantic layering data in \GZGMS.
  For each integer $n\ge2$, an \emph{$n$-th law layer} $l_n$ is a
  stratum in $\mathsf{Layer}$ which fixes:
  \begin{itemize}
    \item the class of admissible law-objects and flows that are
          visible at that level; and
    \item the family of law-space metrics and invariants that may be
          used to distinguish them.
  \end{itemize}
  The second law layer $l_2$ is the purely extensional,
  cardinality-based layer, while the higher layers $l_n$ for $n\ge3$
  are semantic/homotopy layers that incorporate curvature, homotopy,
  and generative data.

  Given a base region $B\subset\mathsf{Obj}$ and a choice of law
  layer $l_n$, an \emph{$n$-th order law-bundle space} is a
  principal-bundle--type object
  \[
    \pi_n : \mathcal{E}_n \to B
  \]
  equipped with a law-space fibre over each point of $B$ and a
  restriction of the law-connection and CH-layering data to layer
  $\le n$, so that horizontal lifts in $\mathcal{E}_n$ respect the
  $l_n$-visible metrics and invariants. Intuitively, $\mathcal{E}_2$
  records only the extensional law-data, whereas $\mathcal{E}_n$ with
  $n\ge3$ also carries the higher homotopy and curvature information.
\end{definition}


\begin{definition}[Semantic CH data at layer $l_{\ge3}$ (schematic)]
  For a continuum-like object $X$ (e.g.\ a model of reals in a given
  law-space), its \emph{semantic CH data} at layer $l_{\ge3}$ is the
  pair
  \[
    \bigl(|X|,\mathbf{S}(X)\bigr),
  \]
  where $|X|$ is the extensional cardinality and $\mathbf{S}(X)$
  collects homotopy/semantic metrics. Two objects $X,Y$ that are
  indistinguish\-able at $l_2$ (same cardinal) may be sharply
  distinguishable at $l_{\ge3}$ via $\mathbf{S}(X)\neq\mathbf{S}(Y)$.
\end{definition}

At these higher layers, the question “what is the continuum?” is not
merely “how many reals are there” but also “how are the reals
organised by law-level structure?” Classical CH answers only the first
half.



\subsubsection*{(c) Law-space layering: a structured view via \texorpdfstring{\GZLS}{GZ-LS}}

Within the GaoZheng Law-Space $\GZLS=(\mathfrak{L};J,\mathrm{Loc},\Sigma)$
and the GaoZheng Generalized Mathematical Structure
\[
  \GZGMS =
  \bigl(
    \mathsf{Obj},\mathsf{Law},\mathsf{Flow},\mathsf{Layer},
    \GZTLC;\GZLCurv
  \bigr),
\]
CH-layering is implemented by the \emph{layer component}
$\mathsf{Layer}$.

At a schematic level, one can think of $\mathsf{Layer}$ as assigning to
each law-object $L\in\mathfrak{L}$ a finite or countable sequence of
“visible aspects”
\[
  \mathsf{Layer}(L)
  =
  \bigl(
    \text{extensional data at }l_2,\,
    \text{homotopy/semantic data at }l_3,\,
    \dots
  \bigr).
\]

\begin{definition}[Layer projection functors]
  For each layer index $\Layer$ (e.g.\ $\Layer=l_2,l_3,l_4,\dots$),
  there is a projection
  \[
    \Pi_\Layer :
      \mathsf{Law} \longrightarrow \mathsf{Law}_\Layer,
  \]
  where $\mathsf{Law}_\Layer$ is the category of law-objects as seen
  at layer $\Layer$ (i.e.\ after discarding all structure invisible at
  that layer). In particular:
  \begin{itemize}
    \item $\Pi_{l_2}$ forgets all semantic/homotopy data and retains
          only cardinal and measure-theoretic structure;
    \item $\Pi_{l_{\ge3}}$ retains additional homotopy and semantic
          invariants as encoded by $J,\mathrm{Loc},\Sigma$ and
          \GZLCurv.
  \end{itemize}
\end{definition}

The classical CH story is entirely about $\Pi_{l_2}(L)$ for law-objects
$L$ representing continuum-like structures. CH-layering asks: what
happens to CH-like statements when we replace $\Pi_{l_2}$ by
$\Pi_{l_{\ge3}}$?



\subsubsection*{(d) Rewriting CH at different layers}

With these projections in hand, we can define layer-dependent forms of
CH. Let $L_{\mathrm{cont}}\in\mathsf{Law}$ be a law-object representing
a continuum structure (e.g.\ the reals in a given model of ZFC plus
additional law-level data).

\begin{definition}[Layered CH statements]
  For a continuum law-object $L_{\mathrm{cont}}$, we define:
  \begin{itemize}
    \item \emph{Extensional CH} (layer-$l_2$):
      \[
        \mathrm{CH}^{(l_2)}(L_{\mathrm{cont}})
        \quad\text{means}\quad
        \text{CH holds in }\Pi_{l_2}(L_{\mathrm{cont}}),
      \]
      i.e.\ the underlying model of sets has $2^{\aleph_0}=\aleph_1$.
    \item \emph{Semantic CH} at layer $\Layer\ge3$ (schematic):
      \[
        \mathrm{CH}^{(\Layer)}(L_{\mathrm{cont}})
        \quad\text{is a constraint of the form}\quad
        \mathbf{S}_\Layer(L_{\mathrm{cont}}) \in \mathcal{R}_\Layer,
      \]
      where $\mathbf{S}_\Layer$ is the semantic metric vector visible
      at layer $\Layer$ and $\mathcal{R}_\Layer$ is a region in
      parameter space encoding law-level regularities one desires
      (e.g.\ certain determinacy patterns or reflection principles).
  \end{itemize}
\end{definition}

In particular, one may have situations where:

\begin{itemize}
  \item $\mathrm{CH}^{(l_2)}(L_{\mathrm{cont}})$ holds (cardinal CH
        is true in the underlying set-theoretic model),
  \item but $\mathrm{CH}^{(l_3)}(L_{\mathrm{cont}})$ fails, because the
        semantic metrics at $l_3$ violate certain coherence constraints
        (e.g.\ they reveal irregularities in definable sets or law-
        curvature patterns incompatible with a “minimal continuum”),
  \item or conversely, $\mathrm{CH}^{(l_2)}$ fails while
        $\mathrm{CH}^{(l_3)}$ holds in a semantic sense (e.g.\ certain
        determinacy and regularity principles hold that treat the
        continuum as if it were “minimal” from a law-level perspective,
        despite cardinal excess).
\end{itemize}

The precise form of $\mathrm{CH}^{(\Layer)}$ at higher layers will be
developed in later subsections (e.g.\ via GZ-CBT: continuum breakdown
and law-level rewriting). For this section, the key point is that CH is
no longer a single, absolute yes/no statement, but a \emph{family} of
layer-dependent constraints.



\subsubsection*{(e) Independence as a projection phenomenon}

Gödel--Cohen independence results can now be re-read as the statement
that \emph{the projection $\Pi_{l_2}$ carries too little structure to
fix the law-level continuum regime}. More concretely:

\begin{itemize}
  \item For any law-object $L_{\mathrm{cont}}$ whose extensional layer
        $\Pi_{l_2}(L_{\mathrm{cont}})$ is a model of ZFC, there exist
        two law-objects
        $L^{\mathrm{CH}}_{\mathrm{cont}},
         L^{\neg\mathrm{CH}}_{\mathrm{cont}}$
        in the same PL-PI law-space fibre such that:
        \[
          \Pi_{l_2}\bigl(L^{\mathrm{CH}}_{\mathrm{cont}}\bigr)
          \models \mathrm{CH},
          \qquad
          \Pi_{l_2}\bigl(L^{\neg\mathrm{CH}}_{\mathrm{cont}}\bigr)
          \models \neg\mathrm{CH},
        \]
        while their higher-layer data
        $\Pi_{l_{\ge3}}(\cdot)$ may be drastically different.
  \item Thus, \emph{extensional} CH is undecidable in ZFC because
        $l_2$ forgets too much of the law-level structure; higher
        layers may still distinguish law-regimes that look identical
        from the $l_2$ viewpoint.
\end{itemize}

In this light, the classical independence of CH is not an obstacle but
a \emph{hint}: CH is not the appropriate type of invariant to capture
the true law-level semantics of the continuum. A law-level theory must
work with semantic metrics that survive passage to higher layers, not
with cardinals alone.



\subsubsection*{(f) Summary and outlook}

We can now summarise the CH-layering perspective as follows:

\begin{enumerate}[label=\arabic*)]
  \item At the extensional layer $l_2$, CH is a statement about the
        cardinality of the continuum, and Gödel--Cohen show that this
        cardinal statement is independent of ZFC.
  \item At higher layers $l_{\ge3}$, continuum-like objects are
        endowed with semantic and generative metrics
        $\mathbf{S}(X)$ capturing law-curvature, homotopy, and
        regularity patterns; CH-like statements become constraints on
        these metrics, not just on cardinalities.
  \item Law-space layering via $\mathsf{Layer}$ and projections
        $\Pi_\Layer$ makes precise how the same law-object
        $L_{\mathrm{cont}}$ can have different “faces” at $l_2$ and
        $l_{\ge3}$; independence at $l_2$ does not preclude strong
        structural statements at $l_{\ge3}$.
\end{enumerate}

In this sense, CH-layering is not a solution to CH \emph{within} ZFC,
but a reframing: it asserts that the correct “location” for continuum
questions is not the extensional layer alone, but the full law-space,
with CH as the $l_2$ shadow and more refined semantic CH constraints
living at $l_{\ge3}$. The next section (GZ-CBT: Continuum Breakdown and
the Rewriting of Law-Level Structures) will make this reframing
concrete by analysing how extensional CH can “break down” when lifted
to higher layers, and how law-level dynamics can “reconstruct”
continuum semantics in a controlled, certificate-based way.



\section{Coupling with the GaoZheng G-Framework and \GZGMS}


CH-layering, as introduced in Sections~9.1--9.2, is not an isolated
construction: it is designed to live \emph{inside} the GaoZheng
G-Framework (\GFramework) and the GaoZheng Generalized Mathematical
Structure (\GZGMS). In this section we make this coupling explicit. The
central point is that CH-layering is a specific instance of the general
“layer” component $\mathsf{Layer}$ in \GZGMS, and that the breakdown
and reconstruction of continuum semantics correspond to controlled
changes of layer along triple-layer flows governed by \GZTLC.



\subsubsection*{(a) CH-layering as a component of \texorpdfstring{\GZLS}{GZ-LS} and \texorpdfstring{\GZGMS}{GZ-GMS}}

Recall that the law-space and generalized structure were schematically
given by
\[
  \GZLS = (\mathfrak{L};J,\mathrm{Loc},\Sigma),
  \qquad
  \GZGMS =
  \bigl(
    \mathsf{Obj},\mathsf{Law},\mathsf{Flow},\mathsf{Layer},
    \GZTLC;\GZLCurv
  \bigr),
\]
where $\mathsf{Law}$ is a suitable category based on $\GZLS$, and
$\mathsf{Layer}$ encodes which invariants are visible at which “semantic
levels” (extensional vs.\ homotopy/semantic).

From this viewpoint, CH-layering is the following specialisation:

\begin{itemize}
  \item we select a subcategory $\mathsf{Law}_{\mathrm{Cont}}
        \subseteq\mathsf{Law}$ of \emph{continuum law-objects}
        (law-level incarnations of $\mathbb{R}$ and its variants in
        different models of ZFC and PL-PI extensions);
  \item on this subcategory, we restrict the general $\mathsf{Layer}$
        mechanism to two distinguished strata:
        \[
           l_2 \ \text{(extensional/cardinal layer)},\qquad
           l_{\ge3} \ \text{(semantic/homotopy layers)}.
        \]
  \item the resulting pair
        \[
          \mathsf{Layer}_{\mathrm{CH}}
          : \mathsf{Law}_{\mathrm{Cont}}
          \longrightarrow
          \mathsf{Law}_{l_2}\times\mathsf{Law}_{l_{\ge3}}
        \]
        is the CH-layering structure: it associates to each
        continuum law-object its $l_2$ projection
        $\Pi_{l_2}(L_{\mathrm{cont}})$ (classical cardinal face) and its
        $l_{\ge3}$ projection $\Pi_{l_{\ge3}}(L_{\mathrm{cont}})$
        (semantic/homotopy face).
\end{itemize}

In other words, CH-layering is simply the restriction of the general
layering mechanism of \GZGMS to a very specific family of law-objects
and two specific layers ($l_2$ and $l_{\ge3}$), together with the
bridge between them.



\subsubsection*{(b) Triple-layer connection \texorpdfstring{\GZTLC}{GZ-TLC} and layer transitions}

Triple-layer connections were introduced as
\[
  \GZTLC \in \Omega^1(\widehat{\Bundle};\widehat{\mathfrak{g}}),
\]
on a principal bundle over a base of the form
\[
  X_{\mathrm{tot}} =
  \Base\times\Lambda_{\mathrm{par}}\times\GZLS,
\]
with curvature
\[
  \mathcal{F}_{\GZTLC}
  =
  \mathrm{d}\GZTLC
  + \tfrac{1}{2}[\GZTLC,\GZTLC].
\]

CH-layering couples to \GZTLC in two ways:

\begin{enumerate}[label=\roman*)]
  \item \emph{Horizontal layer transitions.}
        Along a \GZTLC-horizontal flow
        \[
          t\longmapsto
          \bigl(x(t),\lambda(t),L_{\mathrm{cont}}(t)\bigr)
          \in X_{\mathrm{tot}},
        \]
        the projections
        \[
          \Pi_{l_2}\bigl(L_{\mathrm{cont}}(t)\bigr),
          \quad
          \Pi_{l_{\ge3}}\bigl(L_{\mathrm{cont}}(t)\bigr)
        \]
        evolve in a coupled way. Law-curvature components in
        $\GZLCurv$ decide when a purely extensional deformation in
        $l_2$ (changing $2^{\aleph_0}$, for instance by forcing) must
        be accompanied by changes in the semantic/homotopy invariants
        at $l_{\ge3}$ (e.g.\ changes in determinacy patterns,
        reflection principles, or PL-PI regimes).
  \item \emph{Layer-change operators.}
        Certain components of \GZTLC\ (or of its curvature) can be
        interpreted as \emph{layer-change operators}: they measure the
        failure of a law-flow to stay within a fixed layer. For
        continuum law-objects, such components act as generators of
        transitions between $\mathrm{CH}^{(l_2)}$ and
        $\mathrm{CH}^{(l_{\ge3})}$ regimes, in the sense that when
        these components are non-zero, the extensional and semantic
        CH faces of $L_{\mathrm{cont}}$ cannot be kept in fixed
        relation along the flow.
\end{enumerate}

This provides the geometric backbone for CH-breakdown and
reconstruction: flows that keep certain layer-change components of
\GZLCurv\ small preserve a “near-consistency” between CH at $l_2$ and
semantic CH at $l_{\ge3}$; large curvature in these components signals
a breakdown.



\subsubsection*{(c) Coupling CH-layering with PL-PI law-systems}

In Part~III we introduced PL-PI law-systems
\[
  \mathsf{LS}_{\mathrm{PL\text- PI}}
  =
  (L_{\mathrm{PL}},L_{\mathrm{PI}};\Theta_{\mathrm{PL\text- PI}})
  \in\mathsf{Law}_{\mathrm{PL\text- PI}},
\]
and the GZ-admissible subspace
\(
  \mathsf{Law}_{\mathrm{PL\text- PI}}^{\mathrm{GZ}}
\subseteq \mathsf{Law}_{\mathrm{PL\text- PI}}.
\)

For continuum questions, PL-PI law-systems determine:

\begin{itemize}
  \item \emph{logical} regimes affecting definability and structural
        properties of subsets of the continuum (via $L_{\mathrm{PL}}$);
  \item \emph{procedural} regimes affecting how continuum objects are
        generated or approximated (via $L_{\mathrm{PI}}$).
\end{itemize}

CH-layering couples to PL-PI in the following way:

\begin{itemize}
  \item The extensional face $\Pi_{l_2}(L_{\mathrm{cont}})$ is sensitive
        to the underlying ZFC + additional set-theoretic axioms
        encoded by $L_{\mathrm{PL}}$ (e.g.\ forcing axioms, large
        cardinals).
  \item The semantic face $\Pi_{l_{\ge3}}(L_{\mathrm{cont}})$ is also
        sensitive to the \emph{iterative} structure in $L_{\mathrm{PI}}$:
        for example, how measure, category, or game-theoretic
        procedures are allowed to act on reals (determinacy, regularity
        of projective sets, etc.).
  \item The coupling $\Theta_{\mathrm{PL\text- PI}}$ ensures that
        certain PL and PI choices move together. In CH-layering terms,
        this means that certain regions in PL-PI law-space correspond
        to “semantic CH-like” regimes (high regularity, near-minimal
        continuum structure) even when extensional CH fails, and
        conversely.
\end{itemize}

Thus, CH-layering is tightly interwoven with PL-PI law-systems: the
location of a continuum law-object in PL-PI law-space largely
determines its higher-layer semantic CH characteristics, whereas its
lower-layer extensional CH status is controlled by the underlying
ZFC-based face.



\subsubsection*{(d) Interaction with PFB-GNLA charts: strict, homotopy, and path levels}

The four PFB-GNLA versions provide different “charts” on \GZGMS; CH-
layering appears differently in each:

\begin{itemize}
  \item \textbf{Strict version.}
        In the strict PFB-GNLA chart
        $\mathsf{PFB\text- GNLA}_{\mathrm{strict}}$, one works
        essentially at $l_2$: continuum-related quantities enter only
        via cardinalities and classical measures; law-curvature is
        controlled so that $L_{\mathrm{cont}}$ looks “almost flat”
        from the CH perspective.
  \item \textbf{Homotopy version.}
        In the homotopy chart
        $\mathsf{PFB\text- GNLA}_{\mathrm{homo}}$, $H$-twisted
        homotopy data and non-trivial \GZLCurv\ are visible; they
        influence semantic CH metrics at $l_{\ge3}$ (e.g.\ through
        Jacobiator-gap behaviour in continuum-related operator
        systems).
  \item \textbf{Variable functional--operator version.}
        In the variable functional--operator chart
        $\mathsf{PFB\text- GNLA}_{\mathrm{vf\text- op}}$, operator
        paths trace law-space deformations of continuum structures;
        variable functionals can be chosen to measure, along these
        paths, how semantic CH metrics evolve. Law--operator dynamical
        curvature $\mathcal{F}_{\mathrm{dyn}}$ then encodes CH-breakdown
        and reconstruction along these flows.
  \item \textbf{Original law-level version.}
        In the original \GAlgebra\ chart, CH-layering is seen directly
        in PL-PI law-space: $\mathrm{CH}^{(l_2)}$ and
        $\mathrm{CH}^{(l_{\ge3})}$ become properties of
        $\GAlgebra_{\mathrm{orig}}$’s law-geometry, without any need
        to pass through PFB representations.
\end{itemize}

This shows that CH-layering is not a mere “add-on” but an integral part
of how the four PFB-GNLA versions carve up \GZGMS.



\subsubsection*{(e) CH breakdown and reconstruction as law-level dynamics}

Within this coupled picture, CH breakdown and reconstruction acquire a
precise law-level meaning:

\begin{itemize}
  \item \emph{Breakdown.}
        A CH breakdown event is a region in $X_{\mathrm{tot}}$ where
        layer-change components of \GZLCurv\ and law--operator
        dynamical curvature $\mathcal{F}_{\mathrm{dyn}}$ become large,
        causing $\mathrm{CH}^{(l_2)}$ and semantic CH data at
        $l_{\ge3}$ to drift apart. From the PL-PI point of view, this
        corresponds to moving between PL-PI regimes (e.g.\ switching
        from a CH-favouring to a $\neg\mathrm{CH}$-favouring extension)
        without preserving higher-layer regularities.
  \item \emph{Reconstruction.}
        A CH reconstruction path is a \GZTLC-horizontal (or
        approximately horizontal) flow which returns the continuum
        law-object to a regime where extensional and semantic CH
        constraints are “re-aligned” within acceptable tolerances
        (e.g.\ Jacobiator-gap-based bounds on semantic deviations). In
        path-level charts, such reconstructions can be described by
        law-space learning procedures that drive variable functionals
        towards minima constrained by CH-layer metrics.
\end{itemize}

These notions will be made more concrete in the GZ-CBT (Continuum
Breakdown and Rewriting) framework, but the structural coupling with
\GFramework and \GZGMS is already clear: CH breakdown/reconstruction is
just one particular family of law-level phenomena among many that
\GZTLC\ and \GZLCurv\ can organise.



\subsubsection*{(f) Summary}

Coupling CH-layering with the GaoZheng G-Framework and \GZGMS yields
the following picture:

\begin{enumerate}[label=\arabic*)]
  \item CH-layering is implemented by the general $\mathsf{Layer}$
        component of \GZGMS, restricted to continuum law-objects and
        to the pair of layers $(l_2,l_{\ge3})$.
  \item The triple-layer connection \GZTLC\ and its curvature provide
        the geometric mechanism for layer transitions: they control
        how extensional and semantic CH data co-evolve along law
        flows.
  \item PL-PI law-systems determine where a continuum law-object sits
        in PL-PI law-space; this position governs both its extensional
        CH status and its higher-layer semantic CH profile.
  \item The four PFB-GNLA versions realise CH-layering in different
        charts: strict (near-$l_2$), homotopy (visible $l_{\ge3}$),
        variable functional--operator (dynamic CH-layer flows), and
        original (direct PL-PI law-level).
\end{enumerate}

In this sense, CH-layering is a test case for the broader philosophy of
the GaoZheng G-Framework: extensional independence phenomena (like CH in
ZFC) are reinterpreted as projections of a richer law-level geometry,
in which connections, curvature, homotopy, and semantic metrics jointly
govern the behaviour of continuum and other complex structures.



\chapter[GZ-CBT: Continuum Breakdown]{GZ-CBT: Continuum Breakdown and the Rewriting of Law-Level Structures}


The CH-layering framework of Chapter~9 reinterprets the Continuum
Hypothesis as a layer-dependent family of constraints on continuum
law-objects inside the GaoZheng G-Framework: extensional CH at layer
$l_2$, semantic CH at layers $l_{\ge3}$. The aim of this chapter is to
move from \emph{static} CH-layering to \emph{dynamic} CH-layer
transitions, viewed as law-space flows governed by the triple-layer
connection \GZTLC. The acronym \emph{GZ-CBT} (GaoZheng Continuum
Breakdown and Rewriting) refers to this dynamical picture:

\begin{itemize}
  \item \emph{Continuum breakdown} (CBT) = loss of coherence between
        extensional and semantic CH data along law-level flows, driven
        by specific components of law-curvature \GZLCurv\ and
        law--operator dynamical curvature.
  \item \emph{Rewriting of law-level structures} = systematic
        reconstruction of continuum semantics by adjusting law-space
        connections, homotopy data, and PL-PI configurations so that
        CH-layer constraints are restored (or stabilised in a new
        regime), in a certificate-based, failure-free manner.
\end{itemize}

In other words, CH-layering answers “what CH looks like at different
layers”; GZ-CBT answers “how CH-like semantics can fail and be
rewritten along law-level evolutions”.



\section{Conceptual Architecture of GZ-CBT}


We begin with a high-level description of the objects and notions that
enter the GZ-CBT framework. The technical constructions (precise
definitions of breakdown events, rewriting schemes, and their relation
to certificates) will be developed in the subsequent sections of this
chapter.

\subsubsection*{(a) Continuum law-objects and CH-layer data}

Fix a continuum law-object
\[
  L_{\mathrm{cont}} \in \mathsf{Law}_{\mathrm{Cont}}\subseteq\mathsf{Law},
\]
representing, at the law-level, a real-line--like structure in a given
PL-PI regime. CH-layering associates to $L_{\mathrm{cont}}$ two faces:

\begin{itemize}
  \item the \emph{extensional face} at layer $l_2$,
        \[
           \Pi_{l_2}(L_{\mathrm{cont}}),
        \]
        whose invariants are cardinalities and traditional measure-
        theoretic data (e.g.\ whether
        $2^{\aleph_0}=\aleph_1,\aleph_2,\dots$);
  \item the \emph{semantic/homotopy face} at layers $l_{\ge3}$,
        \[
           \Pi_{l_{\ge3}}(L_{\mathrm{cont}}),
        \]
        whose invariants include semantic metrics
        $\mathbf{S}(L_{\mathrm{cont}})$ (regularity of definable sets,
        determinacy patterns, law-curvature spectral data, etc.).
\end{itemize}

For each such law-object we can thus speak of a \emph{CH-profile}
\[
   \mathsf{CHProfile}(L_{\mathrm{cont}})
   :=
   \bigl(
     \mathrm{CH}^{(l_2)}(L_{\mathrm{cont}}),\,
     \mathrm{CH}^{(l_{\ge3})}(L_{\mathrm{cont}})
   \bigr),
\]
where $\mathrm{CH}^{(l_2)}$ is the classical CH truth value (or a more
refined specification of $2^{\aleph_0}$), and
$\mathrm{CH}^{(l_{\ge3})}$ is a semantic CH constraint (or family of
constraints) on $\mathbf{S}(L_{\mathrm{cont}})$.



\subsubsection*{(b) Law-level flows and continuum trajectories}

Continuum law-objects evolve along law-space flows governed by
\GZTLC. Concretely, consider a law-space trajectory
\[
   t \longmapsto L_{\mathrm{cont}}(t)
   \in \mathsf{Law}_{\mathrm{Cont}},
   \qquad t\in I=[0,1],
\]
lifted to the triple-layer bundle via \GZTLC. At each time $t$, CH-
layering yields
\[
  \mathsf{CHProfile}\bigl(L_{\mathrm{cont}}(t)\bigr)
  =
  \bigl(
    \mathrm{CH}^{(l_2)}(t),\,
    \mathrm{CH}^{(l_{\ge3})}(t)
  \bigr),
\]
and we may track these as functions of $t$.

In dynamic settings (e.g.\ law-space learning, RG flows, evolving PL-PI
regimes), it is natural to consider:

\begin{itemize}
  \item \emph{continuum trajectories}:
        time-indexed or scale-indexed families of continuum
        law-objects $L_{\mathrm{cont}}(t)$, often coupled to PFB-GNLA
        data (strict/homotopy/path-level \GAlgebra);
  \item \emph{CH-layer observables}:
        functionals $\mathcal{O}_{\mathrm{CH}}$ that assign to each
        trajectory a numerical or law-valued summary of its CH-profile
        evolution, for example:
        \[
           \mathcal{O}_{\mathrm{CH}}\bigl(L_{\mathrm{cont}}(\cdot)\bigr)
           :=
           \int_0^1
             F\Bigl(
               \mathrm{CH}^{(l_2)}(t),\,
               \mathrm{CH}^{(l_{\ge3})}(t)
             \Bigr)\,\mathrm{d}t,
        \]
        where $F$ penalises certain misalignments between layers.
\end{itemize}

GZ-CBT will use such trajectories and observables to formalise what it
means for continuum semantics to “break down” and to be “rewritten”
along law-level flows.



\subsubsection*{(c) Continuum breakdown: definition and intuition}

Informally, a \emph{continuum breakdown} is a region along a trajectory
where the extensional and semantic CH faces of $L_{\mathrm{cont}}(t)$
are no longer mutually compatible in a prescribed sense, and this
incompatibility is driven by specific law-curvature and homotopy
effects.

At a schematic level, we fix a tolerance region
\[
   \mathcal{R}_{\mathrm{align}}
   \subset
   \{\bigl(\mathrm{CH}^{(l_2)},\mathrm{CH}^{(l_{\ge3})}\bigr)\},
\]
encoding what it means for the two CH faces to be “sufficiently
aligned” (e.g.\ extensional CH holds and semantic CH metrics lie within
a prescribed range, or extensional $\neg\mathrm{CH}$ is compensated by
semantic regularities of a certain form).

\begin{definition}[Continuum breakdown event (schematic)]
  Let $t\mapsto L_{\mathrm{cont}}(t)$ be a continuum law-space
  trajectory. A \emph{continuum breakdown event} on an interval
  $[t_1,t_2]\subseteq I$ is a maximal connected region in which
  \begin{enumerate}[label=\roman*)]
    \item the CH-profile leaves the alignment region:
          \[
             \bigl(
               \mathrm{CH}^{(l_2)}(t),\,
               \mathrm{CH}^{(l_{\ge3})}(t)
             \bigr)
             \notin
             \mathcal{R}_{\mathrm{align}},
          \]
    \item and certain layer-change curvature components of \GZLCurv\
          and law--operator dynamical curvature $\mathcal{F}_{\mathrm{dyn}}$
          exceed prescribed thresholds:
          \[
             \|\mathcal{F}_{\mathrm{layer}}(t)\|
             > \varepsilon_{\mathrm{layer}},
             \qquad
             \|\mathcal{F}_{\mathrm{dyn}}(t)\|
             > \varepsilon_{\mathrm{dyn}}.
          \]
  \end{enumerate}
\end{definition}

Condition (i) is purely semantic (“CH faces no longer match within the
desired regime”), while (ii) ties this semantic failure to concrete
geometric data in \GZTLC\ and \GZLCurv. The thresholds
$\varepsilon_{\mathrm{layer}}$, $\varepsilon_{\mathrm{dyn}}$ are
part of the design of GZ-CBT; they control how “sensitive” one wants
to be to deviations.



\subsubsection*{(d) Rewriting: law-level reconstruction of continuum semantics}

Once a breakdown has been detected, GZ-CBT prescribes a \emph{rewriting}
procedure: a controlled deformation of law-level structures that brings
the CH-profile back into an acceptable region, or stabilises it in a
new, explicitly characterised regime. This rewriting can occur at
different levels:

\begin{itemize}
  \item by modifying PL-PI law-systems (changing $L_{\mathrm{PL}}$,
        $L_{\mathrm{PI}}$, or their coupling $\Theta_{\mathrm{PL\text- PI}}$);
  \item by adjusting law-connection \GZLOC\ and associated curvature
        in targeted directions;
  \item by enriching homotopy data (e.g.\ passing from a strict to a
        homotopy or path-level PFB-GNLA chart) so that previously
        invisible structure becomes explicit and manageable;
  \item by redefining variable functionals that measure continuum
        behaviour along operator paths, thereby changing the “cost
        landscape” that drives law-space learning or RG flows.
\end{itemize}

Formally, GZ-CBT looks for a \emph{rewriting path}
\[
  \gamma_{\mathrm{rew}} :
     [t_2,t_3]\longrightarrow\mathsf{Law}_{\mathrm{Cont}},
\]
with $\gamma_{\mathrm{rew}}(t_2)=L_{\mathrm{cont}}(t_2)$, such that:

\begin{enumerate}[label=\roman*)]
  \item $\gamma_{\mathrm{rew}}$ is (approximately) horizontal for
        \GZTLC\ outside a controlled “surgery region” where explicit
        law-modifications occur;
  \item the CH-profile returns to alignment:
        \[
          \bigl(
            \mathrm{CH}^{(l_2)}\bigl(\gamma_{\mathrm{rew}}(t)\bigr),\,
            \mathrm{CH}^{(l_{\ge3})}\bigl(\gamma_{\mathrm{rew}}(t)\bigr)
          \bigr)
          \in
          \mathcal{R}_{\mathrm{align}},
          \quad
          t\approx t_3;
        \]
  \item the rewriting respects certificate-based constraints (e.g.\
        Jacobiator-gap bounds, curvature-energy budgets) so that no
        uncontrolled singularities or hidden inconsistencies are
        introduced.
\end{enumerate}

The last condition indicates the tight coupling between GZ-CBT and the
GZ-OHU framework: just as GZ-OHU ensures failure-free homotopy
completion for operator systems, GZ-CBT seeks failure-free law-level
rewriting of continuum semantics.



\subsubsection*{(e) GZ-CBT as a template for law-level crisis management}

Beyond CH itself, the conceptual architecture of GZ-CBT provides a
template for dealing with \emph{any} law-level semantic crisis:

\begin{itemize}
  \item identify a layered profile for the structure of interest
        (extensional vs.\ semantic/homotopy layers);
  \item detect breakdown regions where layer projections drift apart
        beyond acceptable tolerances;
  \item measure the dynamical curvature driving this drift via
        \GZTLC, \GZLCurv, and law--operator dynamical curvature;
  \item design rewriting paths that restore or relocate the system in
        a layer-aligned regime, constrained by certificates (such as
        JacGap) that guarantee failure-free evolution.
\end{itemize}

In this chapter, the continuum serves as a particularly vivid and
historically charged example: classical CH independence provides the
“stress test” for any attempt to give law-level meaning to continuum
semantics. GZ-CBT aims to show that, within the GaoZheng G-Framework,
such meaning can be assigned in a way that is compatible with
classical independence but goes beyond it, by leveraging the full
\GZLS/\GZGMS machinery.



\section{CH-Breakdown Indicators and Curvature-Driven Diagnostics}


The conceptual architecture of GZ-CBT treats continuum breakdown as a
misalignment between the extensional CH face and the semantic/homotopy
CH face of a continuum law-object $L_{\mathrm{cont}}(t)$ along a
law-space trajectory, driven by specific components of the triple-layer
curvature. In this section we introduce quantitative indicators for
this misalignment and relate them to curvature data of \GZTLC\ and its
path-level lift. These indicators will play, for CH, the role that
Jacobiator-gap certificates play for operator homotopy.



\subsubsection*{(a) CH-misalignment indicators along law-space trajectories}

Let
\[
   t \longmapsto L_{\mathrm{cont}}(t)
   \in \mathsf{Law}_{\mathrm{Cont}},
   \qquad t\in I=[0,1],
\]
be a continuum law-space trajectory as in Section~10.1. At each time
$t$ we have a CH-profile
\[
   \mathsf{CHProfile}\bigl(L_{\mathrm{cont}}(t)\bigr)
   =
   \Bigl(
     \mathrm{CH}^{(l_2)}(t),\,
     \mathrm{CH}^{(l_{\ge3})}(t)
   \Bigr),
\]
where $\mathrm{CH}^{(l_2)}(t)$ is an extensional CH datum (e.g.\ the
identity of $2^{\aleph_0}$) and $\mathrm{CH}^{(l_{\ge3})}(t)$ is a
semantic CH datum (e.g.\ a vector of semantic metrics).

To measure misalignment between these faces we introduce a
\emph{CH-misalignment indicator}.

\begin{definition}[CH-misalignment indicator]
Fix a tolerance region
  \[
     \mathcal{R}_{\mathrm{align}}
     \subset
     \mathcal{X}_{l_2}\times\mathcal{X}_{l_{\ge3}},
  \]
  where $\mathcal{X}_{l_2}$ and $\mathcal{X}_{l_{\ge3}}$ are the
  spaces of possible extensional and semantic CH data. A
  \emph{CH-misalignment indicator} is a function
  \[
    \Delta_{\mathrm{CH}}:
      \mathcal{X}_{l_2}\times\mathcal{X}_{l_{\ge3}}
      \longrightarrow [0,+\infty),
  \]
  such that:
  \begin{enumerate}[label=\roman*)]
    \item $\Delta_{\mathrm{CH}}(\xi_2,\xi_{\ge3}) = 0$ if and only if
          $(\xi_2,\xi_{\ge3})\in\mathcal{R}_{\mathrm{align}}$;
    \item $\Delta_{\mathrm{CH}}$ is continuous (or piecewise smooth)
          with respect to a chosen topology on
          $\mathcal{X}_{l_2}\times\mathcal{X}_{l_{\ge3}}$;
    \item $\Delta_{\mathrm{CH}}$ is invariant under law-level gauge
          transformations that preserve both CH faces (e.g.\ internal
          reshufflings in $L_{\mathrm{cont}}$ that do not change
          $\mathrm{CH}^{(l_2)}$ or $\mathrm{CH}^{(l_{\ge3})}$).
  \end{enumerate}
\end{definition}

For a continuum trajectory, we define the time-dependent misalignment
indicator
\[
   \Delta_{\mathrm{CH}}(t)
   :=
   \Delta_{\mathrm{CH}}\Bigl(
     \mathrm{CH}^{(l_2)}(t),\,
     \mathrm{CH}^{(l_{\ge3})}(t)
   \Bigr).
\]

A typical example is obtained by embedding both faces into a common
metric space and taking a distance to $\mathcal{R}_{\mathrm{align}}$,
for instance:
\[
   \Delta_{\mathrm{CH}}(t)
   =
   \mathrm{dist}\Bigl(
     \bigl(\xi_2(t),\xi_{\ge3}(t)\bigr),\,
     \mathcal{R}_{\mathrm{align}}
   \Bigr),
\]
where $\xi_2(t)$ encodes $2^{\aleph_0}(t)$ and $\xi_{\ge3}(t)$ encodes
semantic metrics, and $\mathrm{dist}$ is a suitable metric.



\subsubsection*{(b) Curvature decompositions and layer-change drivers}

We now identify the curvature components of the triple-layer connection that drive \\ CH-misalignment. Recall that the triple-layer connection
\GZTLC\ on
\[
  X_{\mathrm{tot}} = \Base\times\Lambda_{\mathrm{par}}\times\GZLS
\]
has curvature
\[
  \mathcal{F}_{\GZTLC}
  =
  \mathrm{d}\GZTLC + \tfrac12[\GZTLC,\GZTLC],
\]
and, when lifted to operator path spaces, the extended connection
$\widetilde{\GZTLC}$ has curvature
\[
  \widetilde{\mathcal{F}}
  =
  \mathcal{F}_{\GZTLC}
  +\mathcal{F}_{\mathrm{path}}
  +\mathcal{F}_{\mathrm{dyn}}.
\]

In the CH-layering context we isolate two types of curvature drivers:

\begin{itemize}
  \item \emph{Layer-change curvature} $\mathcal{F}_{\mathrm{layer}}$:
        components of $\mathcal{F}_{\GZTLC}$ that mix directions in
        which $l_2$ and $l_{\ge3}$ projections differ, i.e.\ the
        obstruction to $\Pi_{l_2}$ and $\Pi_{l_{\ge3}}$ commuting with
        parallel transport. Schematically,
        \[
          \mathcal{F}_{\mathrm{layer}}
          =
          \Pi_{l_2\leftrightarrow l_{\ge3}}
          \bigl(\mathcal{F}_{\GZTLC}\bigr).
        \]
  \item \emph{Law--operator dynamical curvature} $\mathcal{F}_{\mathrm{dyn}}$:
        mixed components in $\widetilde{\mathcal{F}}$ that couple
        law-space directions to operator-path directions, encoding how
        dynamic law–operator effects distort CH-layer data along
        trajectories.
\end{itemize}

Both types contribute to the time derivative of $\Delta_{\mathrm{CH}}$.

\begin{definition}[Curvature-based CH diagnostics]
Along a continuum trajectory $t\mapsto L_{\mathrm{cont}}(t)$, we
  define the curvature-driven CH diagnostic quantities:
  \[
    B_{\mathrm{layer}}(t)
    :=
    \bigl\|\mathcal{F}_{\mathrm{layer}}(t)\bigr\|,
    \qquad
    B_{\mathrm{dyn}}(t)
    :=
    \bigl\|\mathcal{F}_{\mathrm{dyn}}(t)\bigr\|,
  \]
  where $\|\cdot\|$ denotes a norm or seminorm on curvature valued in
  $\widehat{\mathfrak{g}}$ (or an associated invariant such as a
  spectral radius or trace norm). A combined curvature diagnostic is
  then
  \[
    B_{\mathrm{CH}}(t)
    :=
    \sqrt{
      B_{\mathrm{layer}}(t)^2
      +
      B_{\mathrm{dyn}}(t)^2
    }.
  \]
\end{definition}

Intuitively, $B_{\mathrm{layer}}(t)$ measures how strongly a small
loop in law-space moves $L_{\mathrm{cont}}(t)$ between different
layers; $B_{\mathrm{dyn}}(t)$ measures how strongly operator-path
dynamics pushes law-objects out of layer-aligned regimes.



\subsubsection*{(c) Local breakdown index and threshold criteria}

We now combine misalignment and curvature diagnostics into a single
\emph{local breakdown index}.

\begin{definition}[Local CH breakdown index]
The \emph{local CH breakdown index} at time $t$ is defined by
  \[
    \mathcal{B}_{\mathrm{CH}}(t)
    :=
    w_{\Delta}\,\Delta_{\mathrm{CH}}(t)
    +
    w_{B}\,B_{\mathrm{CH}}(t),
  \]
  where $w_{\Delta},w_{B}>0$ are fixed weights determining the relative
  importance of misalignment versus curvature. A time $t$ is said to
  be \emph{locally CH-regular} if
  \[
     \mathcal{B}_{\mathrm{CH}}(t) \le \theta_{\mathrm{CH}},
  \]
  and \emph{locally CH-singular} if
  \[
     \mathcal{B}_{\mathrm{CH}}(t) > \theta_{\mathrm{CH}},
  \]
  for a chosen threshold $\theta_{\mathrm{CH}}>0$.
\end{definition}

A \emph{CH breakdown event} on an interval $[t_1,t_2]$ may then be
rephrased as the maximal interval on which
$\mathcal{B}_{\mathrm{CH}}(t)>\theta_{\mathrm{CH}}$. This blends the
two conditions of Definition~10.1 (semantic misalignment and curvature
excitation) into a single index.

Note that $\mathcal{B}_{\mathrm{CH}}$ is designed to be:

\begin{itemize}
  \item \emph{law-gauge invariant}, in the sense that internal
        symmetries of \GZLS\ that preserve $\mathsf{CHProfile}$ and
        the norms of relevant curvature components do not change
        $\mathcal{B}_{\mathrm{CH}}$;
  \item \emph{reparametrisation invariant}, in the sense that if
        $t\mapsto L_{\mathrm{cont}}(\tau(t))$ is a reparametrisation
        with monotone $\tau$, then
        $\mathcal{B}_{\mathrm{CH}}(t)$ changes only by a controlled
        scaling if at all (depending on how the weights are normalised).
\end{itemize}

These invariance properties ensure that GZ-CBT diagnostics are
well-defined at the level of law-space trajectories, not tied to an
arbitrary choice of coordinates.



\subsubsection*{(d) Coupling with Jacobiator-gap and operator-side certificates}

Continuum breakdown is not only a law-space phenomenon; it is also
reflected on the operator side, in the behaviour of PFB-GNLA versions
along the trajectory. In particular, the Jacobiator-gap certificate
$\mathsf{JacGap}$ introduced in Section~8.3 provides a complementary
set of invariants for operator homotopy.

For continuum-related operator systems $L^{\mathrm{op}}_{\mathrm{cont}}(t)$
(e.g.\ GZ-GMS operator layers acting on continuum observables), we can
attach a Jacobiator-gap vector
\[
   \mathsf{JacGap}\bigl(L^{\mathrm{op}}_{\mathrm{cont}}(t)\bigr)
   =
   \bigl(
     c_1(t),\dots,c_k(t)
   \bigr),
\]
and define an operator-side Jacobiator index
\[
   J_{\mathrm{op}}(t)
   :=
   \Bigl\|
     \mathsf{JacGap}\bigl(L^{\mathrm{op}}_{\mathrm{cont}}(t)\bigr)
   \Bigr\|.
\]

We then define a \emph{combined continuum--operator breakdown index}
\[
   \mathcal{B}_{\mathrm{CH+op}}(t)
   :=
   \mathcal{B}_{\mathrm{CH}}(t)
   +
   w_{\mathrm{op}}\,J_{\mathrm{op}}(t),
\]
for some weight $w_{\mathrm{op}}>0$.

The rationale is:

\begin{itemize}
  \item large $\Delta_{\mathrm{CH}}(t)$ and $B_{\mathrm{CH}}(t)$ signal
        law-level CH misalignment and curvature excitation;
  \item large $J_{\mathrm{op}}(t)$ signals that the corresponding
        operator system is far from a strict regime and may require
        homotopy completion via GZ-OHU;
  \item their combined index
        $\mathcal{B}_{\mathrm{CH+op}}(t)$ gives an integrated picture
        of how severe the continuum crisis is at both law and operator
        levels.
\end{itemize}



\subsubsection*{(e) Breakdown profiles and energy-like functionals}

For a given trajectory, one may want to summarise breakdown behaviour
over an interval, not just at individual times. For this we introduce
\emph{breakdown profiles} and associated “energy-like” functionals.

\begin{definition}[CH breakdown profile and energy]
Given a trajectory $t\mapsto L_{\mathrm{cont}}(t)$, its \emph{CH
  breakdown profile} on $I$ is the function
  \[
     t \longmapsto \mathcal{B}_{\mathrm{CH}}(t).
  \]
  The associated \emph{CH breakdown energy} on $I$ is defined by
  \[
     \mathcal{E}_{\mathrm{CH}}[L_{\mathrm{cont}}(\cdot)]
     :=
     \int_0^1
       \Phi\bigl(\mathcal{B}_{\mathrm{CH}}(t)\bigr)\,\mathrm{d}t,
  \]
  where $\Phi:[0,+\infty)\to[0,+\infty)$ is a convex, increasing
  function (e.g.\ $\Phi(x)=x^2$). Similarly, the combined
  continuum--operator breakdown energy is
  \[
     \mathcal{E}_{\mathrm{CH+op}}[L_{\mathrm{cont}}(\cdot)]
     :=
     \int_0^1
       \Phi\bigl(\mathcal{B}_{\mathrm{CH+op}}(t)\bigr)\,\mathrm{d}t.
  \]
\end{definition}

These energies are not physical in the usual sense, but they play an
analogous role in law-level optimisation or learning problems: GZ-CBT
reconstruction procedures will aim to find law-space trajectories (or
deformations of given trajectories) that minimise
$\mathcal{E}_{\mathrm{CH}}$ or $\mathcal{E}_{\mathrm{CH+op}}$ subject
to constraints (e.g.\ preserving certain PL-PI features, staying within
a given PFB-GNLA chart, etc.).



\subsubsection*{(f) Summary: diagnostic layer of GZ-CBT}

To summarise, Section~10.2 has introduced a diagnostic layer for GZ-CBT
consisting of:

\begin{itemize}
  \item CH-misalignment indicators $\Delta_{\mathrm{CH}}(t)$ capturing
        the semantic gap between \\ $\mathrm{CH}^{(l_2)}$ and $\mathrm{CH}^{(l_{\ge3})}$;
  \item curvature-driven diagnostics
        $B_{\mathrm{layer}}(t)$, $B_{\mathrm{dyn}}(t)$ and their
        combination $B_{\mathrm{CH}}(t)$ capturing the geometric forces
        that drive layer transitions;
  \item a local CH breakdown index
        $\mathcal{B}_{\mathrm{CH}}(t)$ and its continuum--operator
        extension $\mathcal{B}_{\mathrm{CH+op}}(t)$, blending semantic
        misalignment, curvature excitation, and Jacobiator-gap
        behaviour;
  \item energy-like functionals
        $\mathcal{E}_{\mathrm{CH}}$ and $\mathcal{E}_{\mathrm{CH+op}}$
        summarising breakdown severity along trajectories.
\end{itemize}

In the next section, these diagnostics will be used to define and
analyse explicit GZ-CBT rewriting schemes: law-level procedures that
modify PL-PI law-systems, law-connections, and homotopy data so as to
steer continuum law-objects away from breakdown regions and towards
layer-aligned CH regimes, in a manner compatible with the GZ-OHU
failure-free philosophy.



\section{Law-Level Rewriting Schemes for Continuum Semantics}


The diagnostic layer of GZ-CBT (Section~10.2) provides quantitative
indicators of continuum breakdown: CH-misalignment
$\Delta_{\mathrm{CH}}(t)$, curvature diagnostics $B_{\mathrm{layer}}(t)$
and $B_{\mathrm{dyn}}(t)$, and the combined breakdown indices
$\mathcal{B}_{\mathrm{CH}}(t)$ and $\mathcal{B}_{\mathrm{CH+op}}(t)$.
In this section we introduce \emph{rewriting schemes}: law-level
procedures that deform continuum law-objects and their associated
operator systems so as to reduce these indices, ideally driving the
system back into a CH-aligned regime. The guiding philosophy is
failure-free construction in the sense of GZ-OHU: every step is
either certifiably safe or accompanied by an explicit certificate
indicating what homotopy completion or law-level adjustment is needed.



\subsubsection*{(a) The law-level rewriting problem as constrained optimisation}

We formalise the rewriting problem as a constrained optimisation problem
in law-space. Let $\mathcal{C}_{\mathrm{adm}}$ be the class of
GZ-admissible continuum trajectories
\[
  \gamma(\cdot): I=[0,1]\to\mathsf{Law}_{\mathrm{Cont}},
  \qquad
  t\mapsto L_{\mathrm{cont}}(t),
\]
satisfying structural constraints (e.g.\ PL-PI admissibility, staying
within a given PFB-GNLA chart, bounds on curvature, etc.).

Given diagnostic functionals
$\mathcal{E}_{\mathrm{CH}}[\gamma]$ and
$\mathcal{E}_{\mathrm{CH+op}}[\gamma]$ as in Section~10.2, we define:

\begin{definition}[GZ-CBT rewriting problem (schematic)]
  Fix an initial trajectory
  $\gamma_0\in\mathcal{C}_{\mathrm{adm}}$ and admissible constraint set
  $\mathcal{K}$ (encoding invariants we wish to preserve: PL-PI features,
  law-level symmetries, bounds on Jacobiator-gap, etc.). The
  \emph{GZ-CBT rewriting problem} is to find a new trajectory
  $\gamma^\star\in\mathcal{C}_{\mathrm{adm}}$ and a homotopy
  \[
    H:[0,1]\times I\to\mathsf{Law}_{\mathrm{Cont}},
    \quad
    (s,t)\mapsto H(s,t),
  \]
  such that:
  \begin{enumerate}[label=\roman*)]
    \item $H(0,\cdot)=\gamma_0(\cdot)$,\quad $H(1,\cdot)=\gamma^\star(\cdot)$;
    \item $H(s,\cdot)\in\mathcal{C}_{\mathrm{adm}}$ for all $s$;
    \item the constraints $\mathcal{K}$ are satisfied along $H$;
    \item the breakdown energy is reduced:
          \[
            \mathcal{E}_{\mathrm{CH+op}}[\gamma^\star]
            \le
            \mathcal{E}_{\mathrm{CH+op}}[\gamma_0],
          \]
          with strict inequality unless $\gamma_0$ is already
          CH-aligned in a suitable sense.
  \end{enumerate}
\end{definition}

Thus, rewriting is a \emph{law-space homotopy} of trajectories that
respects admissibility constraints and decreases the breakdown energy.



\subsubsection*{(b) Gradient-like flows in law-space}

A natural class of rewriting schemes is obtained by defining a
gradient-like flow of the breakdown energy on law-space, using the
metric information encoded in $J$ (the law-metric component of \GZLS)
and the connection \GZLOC.

Assume that $\mathsf{Law}_{\mathrm{Cont}}$ carries a Riemannian (or
pseudo-Riemannian) structure induced from $J$, so that for smooth
functionals $\mathcal{F}$ one can define a gradient
$\nabla\mathcal{F}(L_{\mathrm{cont}})$ in the tangent space at
$L_{\mathrm{cont}}$.

\begin{definition}[CH-gradient flow (schematic)]
  A \emph{CH-gradient flow} is a law-space evolution
  \[
    \frac{\mathrm{d}}{\mathrm{d}s}L_{\mathrm{cont}}(s)
    =
    -\,\mathrm{Proj}_{\mathcal{C}_{\mathrm{adm}}}
      \Bigl(
        \nabla\mathcal{E}_{\mathrm{CH+op}}
          \bigl[L_{\mathrm{cont}}(\cdot)\bigr](s)
      \Bigr),
  \]
  where:
  \begin{itemize}
    \item $s$ is an “algorithmic time” parameter for rewriting;
    \item the gradient is computed with respect to the law-metric;
    \item $\mathrm{Proj}_{\mathcal{C}_{\mathrm{adm}}}$ is a projection
          onto admissible directions (enforcing constraints $\mathcal{K}$);
    \item parallel transport along $s$ is taken with respect to
          \GZLOC\ (or \GZTLC\ when geometry and parameters are included).
  \end{itemize}
\end{definition}

In discrete time, this gives a scheme
\[
  L_{\mathrm{cont}}^{(n+1)}
  =
  \mathrm{Exp}_{L_{\mathrm{cont}}^{(n)}}
   \Bigl(
     -\eta_n\,
      \mathrm{Proj}_{\mathcal{C}_{\mathrm{adm}}}
        \bigl(\nabla\mathcal{E}_{\mathrm{CH+op}}\bigr)
   \Bigr),
\]
where $\mathrm{Exp}$ is the exponential map induced by the connection
and $\eta_n$ is a step size.

Under suitable conditions (smoothness, convexity in directions allowed
by $\mathcal{C}_{\mathrm{adm}}$), one expects monotone decrease of
$\mathcal{E}_{\mathrm{CH+op}}$ along such flows, at least locally.



\subsubsection*{(c) Discrete CBT-OHU schemes: interleaving law rewriting and operator homotopy}

Gradient flows provide a continuous picture, but in practice one often
works with \emph{discrete} schemes that interleave:

\begin{itemize}
  \item law-level rewrites of PL-PI data and law-connections
        (GZ-CBT steps);
  \item operator-level homotopy completions via GZ-OHU (GZ-OHU steps).
\end{itemize}

We call such schemes \emph{CBT-OHU schemes}.

\begin{definition}[CBT-OHU rewriting scheme (schematic)]
  A \emph{CBT-OHU rewriting scheme} for an initial pair
  $(L_{\mathrm{cont}}^{(0)},L_{\mathrm{op}}^{(0)})$ consists of
  sequences
  \[
    \bigl(L_{\mathrm{cont}}^{(n)},L_{\mathrm{op}}^{(n)}\bigr)_{n\ge0}
  \]
  and operations of two types:

  \begin{itemize}
    \item \emph{CBT step (law-level)}:
      \[
        \bigl(L_{\mathrm{cont}}^{(n)},L_{\mathrm{op}}^{(n)}\bigr)
        \mapsto
        \bigl(L_{\mathrm{cont}}^{(n+1)},L_{\mathrm{op}}^{(n)}\bigr),
      \]
      where $L_{\mathrm{cont}}^{(n+1)}$ is obtained from
      $L_{\mathrm{cont}}^{(n)}$ via a small law-space move guided by
      $\nabla\mathcal{E}_{\mathrm{CH}}$ (or its projection), and
      $L_{\mathrm{op}}^{(n)}$ is pulled back/forward along \GZLOC/\GZTLC\;
      \item \emph{OHU step (operator-level)}:
      \[
        \bigl(L_{\mathrm{cont}}^{(n+1)},L_{\mathrm{op}}^{(n)}\bigr)
        \mapsto
        \bigl(L_{\mathrm{cont}}^{(n+1)},L_{\mathrm{op}}^{(n+1)}\bigr),
      \]
      where $L_{\mathrm{op}}^{(n+1)}$ is obtained from
      $L_{\mathrm{op}}^{(n)}$ via a GZ-OHU homotopy completion that
      reduces the operator-side Jacobiator-gap
      $J_{\mathrm{op}}^{(n)}$.
  \end{itemize}

  At each combined step, the total breakdown index
  $\mathcal{B}_{\mathrm{CH+op}}^{(n)}$ is required to be non-increasing,
  with strict decrease whenever a breakdown threshold is exceeded.
\end{definition}

A CBT step targets law-level CH misalignment and curvature; an OHU
step targets \\ operator-level Jacobi defects. By construction, such
schemes inherit the failure-free guarantees of GZ-OHU on the operator
side, and—under suitable design of the law-level moves—avoid
uncontrolled law-space singularities.



\subsubsection*{(d) Schematic convergence and failure-free properties}

We record a schematic statement capturing the intended behaviour of
CBT-OHU rewriting schemes.

\begin{proposition}[Monotonicity and failure-free behaviour (schematic)]
  Let $\bigl(L_{\mathrm{cont}}^{(n)},L_{\mathrm{op}}^{(n)}\bigr)_{n\ge0}$
  be produced by a CBT-OHU scheme satisfying:
  \begin{enumerate}[label=\roman*)]
    \item each CBT step is chosen so that
          \[
            \mathcal{E}_{\mathrm{CH}}
              \bigl[L_{\mathrm{cont}}^{(n+1)}(\cdot)\bigr]
            \le
            \mathcal{E}_{\mathrm{CH}}
              \bigl[L_{\mathrm{cont}}^{(n)}(\cdot)\bigr],
          \]
          with strict inequality whenever
          $\mathcal{B}_{\mathrm{CH}}^{(n)}>\theta_{\mathrm{CH}}$;
    \item each OHU step applies a GZ-OHU completion that reduces or
          preserves $J_{\mathrm{op}}^{(n)}$ and is failure-free in the
          sense of the GZ-OHU theorem (every intermediate Jacobi
          violation is certified or resolved);
    \item step sizes and curvature are controlled so that the sequence
          remains in a compact region of law-space where the diagnostics
          are continuous.
  \end{enumerate}
  Then:
  \begin{enumerate}[label=\alph*)]
    \item the combined breakdown energy
          $\mathcal{E}_{\mathrm{CH+op}}^{(n)}$ is monotonically
          non-increasing in $n$;
    \item any accumulation point of the sequence is CH-regular in the
          sense that $\mathcal{B}_{\mathrm{CH+op}}\le\theta_{\mathrm{CH}}$
          and $\mathsf{JacGap}$ is minimal within the admissible class;
    \item along the scheme there are no unsignalled breakdowns:
          whenever a Jacobi or CH-related invariant crosses a critical
          threshold, this is reflected in
          $\mathcal{B}_{\mathrm{CH+op}}^{(n)}$ and triggers a remedial
          step prescribed by the scheme.
  \end{enumerate}
  A complete proof is given in Technical Appendix~I
  (Section~\ref{app:I-P9}).
\end{proposition}

This proposition is schematic; a fully rigorous statement would specify
precise regularity assumptions, compactness properties, and the form of
$\Phi$ in the energies. However, it captures the conceptual content of
“failure-free CH rewriting”: the process is guided and monitored by
diagnostics that prevent hidden divergences.



\subsubsection*{(e) Compatibility with the four PFB-GNLA charts}

Each of the four PFB-GNLA versions offers a different “implementation
layer” for CBT-OHU schemes:

\begin{itemize}
  \item \textbf{Strict chart:}
        rewriting amounts to adjusting PFB connections and parameters
        while keeping Jacobi strict; CBT steps dominate, OHU steps are
        rarely invoked (JacGap $\approx0$ by design).
  \item \textbf{Homotopy chart:}
        rewriting includes explicit modification of homotopy brackets
        ($\ell_3^H$ and possibly higher $\ell_k$) in tandem with
        law-space moves; OHU steps are central.
  \item \textbf{Variable functional--operator chart:}
        rewriting operates on operator paths and variable functionals
        with $\widetilde{\GZTLC}$ controlling both law and operator
        dynamics; CH breakdown and reconstruction can be described in
        terms of path-space flows minimising
        $\mathcal{E}_{\mathrm{CH+op}}$.
  \item \textbf{Original law-level chart:}
        rewriting is phrased entirely in PL-PI law-space, with PFB and
        operator manifestations recovered as representations; CBT-OHU
        becomes a statement about deforming PL-PI systems themselves.
\end{itemize}

In all cases, the same diagnostic quantities
$\Delta_{\mathrm{CH}},B_{\mathrm{CH}},\mathsf{JacGap}$ and the same
structural constraints (PL-PI admissibility, CH-layering, CH-profile
alignment) inform the design of rewriting steps.



\subsubsection*{(f) Summary}

GZ-CBT rewriting schemes provide the constructive half of the
continuum breakdown story:

\begin{itemize}
  \item Diagnostics from Section~10.2 quantify where and how continuum
        semantics misalign between layers and how strongly curvature
        drives this misalignment.
  \item Gradient-like flows in law-space and discrete CBT-OHU schemes
        give concrete procedures for deforming continuum law-objects
        and their associated operator systems to reduce CH breakdown
        energy.
  \item The GZ-OHU framework guarantees that operator-side homotopy
        completions can be performed failure-free, while \GZTLC\ and
        CH-layering constrain law-level moves to avoid uncontrolled
        singularities.
\end{itemize}

In the remaining sections of this chapter, these rewriting schemes will
be illustrated in specific scenarios (e.g.\ transitions between CH and
$\neg\mathrm{CH}$ regimes under different PL-PI extensions), and their
interaction with certificate-based AI and law-space learning will be
sketched, demonstrating how continuum semantics can be treated as a
dynamic, law-level object within the GaoZheng G-Framework.



\section[Scenario Sketches: CH → ¬CH, Forcing Axioms, and Large Cardinals]{Scenario Sketches: CH\texorpdfstring{$\to\neg\mathrm{CH}$}{ → ¬CH}, Forcing Axioms, and Large Cardinals}


The previous sections developed GZ-CBT as a general diagnostic and
rewriting framework for continuum semantics in law-space, coupled to
the PFB-GNLA charts and the PL-PI law-level picture. In this section we
outline several canonical scenarios in which continuum breakdown and
rewriting are expected to occur, without attempting a fully detailed
model in each case. The intent is to show how classical set-theoretic
phenomena---CH$\to\neg\mathrm{CH}$ transitions, forcing axioms,
large cardinal assumptions---fit into the GZ-CBT template.



\subsubsection*{(a) Scenario I: Forcing from CH to $\neg\mathrm{CH}$}

Consider a ground model $M\models\text{ZFC+CH}$ and a forcing
extension $M[G]\models\text{ZFC}+\neg\mathrm{CH}$ obtained by adding
Cohen reals or other suitable generic subsets of $\omega$. At the
extensional layer $l_2$, this is a classical CH$\to\neg\mathrm{CH}$
transition. We now recast this as a GZ-CBT scenario.

\paragraph{Law-level description.}

Let $L_{\mathrm{cont}}^{\mathrm{ground}}$ and
$L_{\mathrm{cont}}^{\mathrm{ext}}$ be continuum law-objects
corresponding to the reals in $M$ and $M[G]$, respectively. There is a
law-space path
\[
  \gamma_{\mathrm{forc}}(t),
  \qquad
  t\in[0,1],
\]
interpolating between these law-objects in $\mathsf{Law}_{\mathrm{Cont}}$,
with the property that
\[
  \gamma_{\mathrm{forc}}(0)
    = L_{\mathrm{cont}}^{\mathrm{ground}},
  \qquad
  \gamma_{\mathrm{forc}}(1)
    = L_{\mathrm{cont}}^{\mathrm{ext}}.
\]

At layer $l_2$ we have
\[
  \mathrm{CH}^{(l_2)}\bigl(\gamma_{\mathrm{forc}}(0)\bigr)=\text{true},
  \qquad
  \mathrm{CH}^{(l_2)}\bigl(\gamma_{\mathrm{forc}}(1)\bigr)=\text{false}.
\]
The evolution of $\mathrm{CH}^{(l_{\ge3})}$ along the path depends on
how forcing interacts with semantic metrics (regularity of definable
sets, determinacy, etc.).

\paragraph{CBT diagnostics.}

Assume that the forcing has been chosen so that semantic regularities
are significantly altered (e.g.\ adding many Cohen reals tends to
destroy some regularity of the Baire/measurable hierarchy). Then along
$\gamma_{\mathrm{forc}}$:

\begin{itemize}
  \item $\Delta_{\mathrm{CH}}(t)$ will typically increase around the
        “forcing region”, as extensional CH flips while semantic
        metrics respond in a less controlled way;
  \item $B_{\mathrm{layer}}(t)$ captures the fact that small loops in
        law-space near the forcing step do not preserve the relative
        alignment of $l_2$ and $l_{\ge3}$ faces;
  \item if continuum-related operator systems (e.g.\ those encoding
        descriptive-set-theoretic hierarchies) are included, their
        Jacobiator-gap $J_{\mathrm{op}}(t)$ may also spike as the
        law-regime is abruptly changed.
\end{itemize}

This is a prototypical CH-breakdown event: a discrete change in
extensional CH status induces a cascade of semantic disturbances, and
the GZ-CBT indices register a significant excursion in
$\mathcal{B}_{\mathrm{CH+op}}(t)$.

\paragraph{Rewriting sketch.}

A CBT-OHU rewriting scheme could attempt to “heal” part of this
breakdown by:

\begin{itemize}
  \item adjusting PL-PI law-systems in the extension $M[G]$ (e.g.\
        adding forcing axioms that restore strong regularity properties
        for definable sets);
  \item homotopy-completing operator systems to incorporate new
        symmetry or determinacy patterns;
  \item designing variable functionals whose minima correspond to
        law-regimes where semantic CH metrics regain a near-minimal
        character, even though $\mathrm{CH}^{(l_2)}$ now fails.
\end{itemize}

The result is not a return to CH at $l_2$, but a transition to a new
semantic CH regime at $l_{\ge3}$ in which the continuum behaves “as
regularly as possible” given the extensional change.



\subsubsection*{(b) Scenario II: Forcing axioms and $\neg\mathrm{CH}$-favouring regimes}

Strong forcing axioms such as PFA or MM typically imply $\neg\mathrm{CH}$
and enforce robust regularity properties on $\mathbb{R}$ (and on
definable sets more generally). From the CH-layering perspective, these
axioms define a PL-PI subregion in which:

\begin{itemize}
  \item $\mathrm{CH}^{(l_2)}$ is false (the continuum is “too large” in
        cardinal terms);
  \item $\mathrm{CH}^{(l_{\ge3})}$ may be regarded as \emph{true} in a
        semantic sense: the continuum is highly regular, many
        definability hierarchies collapse, and game-theoretic
        determinacy holds in wide fragments.
\end{itemize}

Within GZ-CBT, this corresponds to trajectories that \emph{start} in a
CH-like extensional regime, move into a forcing-axiom regime, and then
stabilise in a law-level region where:

\[
  \Delta_{\mathrm{CH}}(t) \text{ is small for } t\gg0,
\]
despite $\mathrm{CH}^{(l_2)}$ being permanently false.

Rewriting schemes in this scenario are less about “repairing damage”
and more about \emph{navigating} into a desired forcing-axiom regime:

\begin{itemize}
  \item CBT steps drive the PL-PI system towards law-level fixed points
        characterised by \\ PFA/MM-like constraints;
  \item OHU steps ensure that operator systems encoding continuum
        regularity (e.g.\ projective hierarchies) are homotopy-complete
        and failure-free under those axioms;
  \item the end state has high $\mathcal{E}_{\mathrm{CH}}$ reduction
        at $l_{\ge3}$, even though the cardinal continuum is large.
\end{itemize}



\subsubsection*{(c) Scenario III: Large cardinals and inner-model guided regimes}

Large cardinal axioms and inner model theory provide a different way of
shaping continuum behaviour. Depending on the program (e.g.\ inner
model with Woodin cardinals, Ultimate-$L$), one may:

\begin{itemize}
  \item favour CH via canonical inner models resembling $L$ but enriched
        with large cardinals;
  \item or favour $\neg\mathrm{CH}$ via determinacy-inspired regimes
        where projective sets of reals are exceptionally regular.
\end{itemize}

In GZ-CBT terms:

\begin{itemize}
  \item the PL-PI law-space contains regions corresponding to specific
        large-cardinal/inner-model hypotheses;
  \item continuum law-objects $L_{\mathrm{cont}}$ in these regions have
        CH-profiles that may be sharply constrained at $l_{\ge3}$, even
        though $l_2$ may still be free (in the sense that various
        $2^{\aleph_0}$ are compatible).
\end{itemize}

A CH-rewriting path in this scenario might consist of:

\begin{itemize}
  \item moving from a generic forcing-based $\neg\mathrm{CH}$ regime
        with poor semantic regularity to an inner-model regime where
        large cardinals impose strong structure on reals;
  \item adjusting law-connections so that law-curvature aligns with
        inner-model hierarchies (e.g.\ canonical scales, extender
        sequences);
  \item performing OHU completions on operator systems that model
        these hierarchies, thus reducing both
        $\mathcal{B}_{\mathrm{CH}}$ and $J_{\mathrm{op}}$.
\end{itemize}

The resulting law-level picture can be seen as a “semantic
stabilisation” of the continuum guided by large cardinal structure.



\subsubsection*{(d) Scenario IV: Law-space learning and continuum-aware complex-systems models based on total-operator principal bundles}

While the previous scenarios are purely set-theoretic, the same
GZ-CBT machinery can be applied to continuum-like structures in applied
domains, especially where continuum approximations or infinite-horizon
limits appear (e.g.\ in models of total-operator principal bundles for complex-systems modelling
(developed in separate applied volumes for complex-systems law-objects), or in
RL/AI systems that operate on spaces of trajectories).

In such contexts:

\begin{itemize}
  \item $L_{\mathrm{cont}}$ may represent a law-level model of a
        continuum state space used by a total-operator bundle-type model for complex-systems modelling or by an RL agent;
  \item $\mathrm{CH}^{(l_2)}$ is replaced by coarse “size” or capacity
        notions (e.g.\ covering numbers, \\ VC-dimensions, entropy rates);
  \item $\mathrm{CH}^{(l_{\ge3})}$ is replaced by semantic metrics
        describing how well the model captures relevant regularities
        (e.g.\ Lipschitz continuity of value functions, stability
        of optimal policies).
\end{itemize}

Continuum breakdown in this setting corresponds to failure of
approximation regimes (e.g.\ infinite-horizon limits misbehaving,
function classes ceasing to generalise). GZ-CBT rewriting then becomes
a law-space learning process:

\begin{itemize}
  \item diagnostics $\Delta_{\mathrm{CH}}$, $B_{\mathrm{CH}}$ are
        instantiated by application-specific metrics (e.g.\ violation
        of uniform convergence bounds, curvature of value landscapes);
  \item CBT steps adjust the law-level model class, regularisation
        schemes, or structural parameters of total-operator principal bundles used for complex-systems modelling (developed in separate applied volumes);
  \item OHU steps ensure that the induced operator systems (e.g.\ Bellman
        operators, gradient flow generators) remain homotopy-complete
        and free of hidden degeneracies.
\end{itemize}

Although the vocabulary differs, the underlying template is the same:
layered semantics, curvature-driven breakdown, certificate-guided
rewriting.



\subsubsection*{(e) Scenario V: Mixed transitions and multi-stage CBT}

Realistic law-level evolutions may combine several of the above
scenarios:

\begin{itemize}
  \item starting in a CH model with weak semantic regularity;
  \item forcing to $\neg\mathrm{CH}$ to achieve certain combinatorial
        benefits;
  \item adding forcing axioms or large cardinal hypotheses to regain
        regularity at $l_{\ge3}$;
  \item embedding the resulting continuum law-object into PFB-GNLA
        charts used by physical or AI modelling frameworks.
\end{itemize}

GZ-CBT treats such multi-stage transitions as concatenations of
CBT-OHU phases, each with its own diagnostic profile and rewriting
goals. The design principle is to ensure that:

\begin{itemize}
  \item $\mathcal{B}_{\mathrm{CH+op}}$ does not blow up uncontrollably
        in any phase;
  \item each phase ends in a law-region where the next phase can begin
        without violating GZ-admissibility;
  \item accumulated changes in CH semantics are traceable via
        CH-profile histories and certificate logs (e.g.\ time series
        of $\mathsf{JacGap}$ and curvature diagnostics).
\end{itemize}



\subsubsection*{(f) Summary}

The scenarios in this section illustrate how:

\begin{itemize}
  \item classical CH$\leftrightarrow\neg\mathrm{CH}$ transitions,
        forcing axioms, and large cardinal hypotheses can be seen as
        specific regions and flows in PL-PI law-space;
  \item GZ-CBT diagnostics and rewriting schemes organise these flows,
        turning continuum breakdown into a quantifiable, controllable
        law-level phenomenon;
  \item the same template extends to continuum-like structures in
        applied domains (total-operator principal bundles for complex-systems
        modelling, developed in separate applied volumes, and RL/AI), where “CH” becomes a
        metaphor for more general breakdown of infinite-dimensional
        semantics.
\end{itemize}

In this sense, Chapter~10 shows that CH-layering plus GZ-CBT does not
attempt to “solve” CH in the classical sense, but to \emph{manage} the
semantics of continuum-like structures across layers and regimes
within the GaoZheng G-Framework, providing a unified language for
detecting, diagnosing, and rewriting law-level breakdowns of
continuum behaviour.






\part{Outlook: From the GaoZheng G-Framework to Physics}


\chapter[Applications in Physics]{Applications of the GaoZheng G-Algebra and Law-Connections in Physics}


The preceding parts developed the GaoZheng G-Algebra (\GAlgebra) and
the associated law-space structures (\GZLS, \GZTLC, \GZGMS) in a
largely “theory-internal” manner: principal-bundle geometry, GNLA and
$L_\infty$ structures, law-space flows, CH-layering, and PL-PI
provenance. The aim of this chapter is to sketch how these structures
interface with mainstream physics, in a way that preserves the full
law-level perspective rather than merely reproducing known models.

Very roughly, we will see \GAlgebra\ and law-connections appear in
three interconnected roles:

\begin{itemize}
  \item as a \emph{unifying language} for gauge theories and
        gravitation, lifting classical connections and curvatures in
        space-time to law-connections and law-curvatures in \GZLS;
  \item as a \emph{meta-level scaffold} for quantum field theories and
        renormalisation, with law-space flows and operator-path
        structures encoding RG and effective field theory hierarchies;
  \item as a \emph{continuum-law control system} for complex media and
        statistical/condensed-matter models, where CH-layering and
        GZ-CBT–type ideas govern how continuum approximations are
        maintained, broken, and repaired at the law-level.
\end{itemize}

This chapter does \emph{not} attempt to derive any specific physical
theory from \GAlgebra. Rather, it proposes a set of mapping principles:
how to read classical physical data (fields, couplings, symmetries,
effective actions) as projections of law-level objects in \GZLS; how to
lift familiar flows (RG flows, background deformations, phase
transitions) to \GZTLC-horizontal or curvature-driven trajectories;
and how the four PFB-GNLA charts provide different “views” on the same
physical law-space.



\section{From Principal Bundles to Law-Space in Physics}


We begin by recalling how principal bundles already pervade modern
physics and then indicate how the GaoZheng G-Algebra upgrades this
geometry to a law-space framework. This section sets the stage for the
more concrete application sketches in later sections (gauge theories,
gravitation, quantum field theory and RG, complex systems).

\subsubsection*{(a) Principal bundles in gauge theories and gravitation}

In classical and quantum field theory, principal bundles appear in at
least two canonical ways:

\begin{itemize}
  \item \emph{Gauge theories.}
        A principal $\Group$-bundle $(\Bundle,\Base,\Group)$ over
        space-time $\Base$ (or space-time plus extra dimensions) with
        connection $\Conn$ encodes a gauge potential $A$; its curvature
        $\Curv$ encodes field strength $F$. Gauge equivalence classes
        of connections correspond to physical configurations modulo
        gauge redundancy.
  \item \emph{Gravitation.}
        In general relativity and its extensions, one may work with the
        frame bundle or orthonormal frame bundle over $\Base$ (space-
        time manifold), equipped with a Levi-Civita connection or more
        general metric-compatible connections. The curvature is (a
        representation of) the Riemann tensor; geodesic motion and
        Einstein equations can be expressed in this language.
\end{itemize}

In both cases, the principal bundle is an \emph{object-level} device:
it describes how physical fields and particles respond to given
connections and curvatures. What it does not, by itself, encode is the
\emph{space of laws} that govern which connections, couplings, and
fields are admissible, nor how these laws flow under changes of scale,
background, or meta-level constraints.



\subsubsection*{(b) Lifting to law-space: physical connections as slices of \texorpdfstring{\GZLS}{GZ-LS}}

In the GaoZheng G-Framework, a physical principal bundle is embedded
into a larger law-space picture as follows:

\begin{enumerate}[label=\roman*)]
  \item \emph{Object-level geometry.}
        A physical configuration is described by a principal bundle \\ 
        $(\Bundle,\Base,\Group)$ with connection $\Conn$ and curvature
        $\Curv$, together with associated bundles carrying matter
        fields.
  \item \emph{Law-objects in \GZLS.}
        The geometry is mapped, via the functor
        $\Phi_{\mathrm{geom}}$ (Section~3.3), to a law-object
        \[
          L_{\mathrm{phys}}
          =
          \Phi_{\mathrm{geom}}(\Bundle,\Base,\Group;\Conn,\Curv;\text{aux})
          \in\GZLS.
        \]
        This law-object packages not only the raw $(\Conn,\Curv)$ but
        also composition rules, metric data $J$, localisation data
        $\mathrm{Loc}$, and thresholds $\Sigma$ reflecting the
        admissible ranges of couplings, field strengths, and symmetry
        breaking patterns.
  \item \emph{Law-connections and law-curvature.}
        A family of physical configurations (e.g.\ a one-parameter
        family of background fields, an RG trajectory, a family of
        effective theories) is lifted to a law-space trajectory
        $t\mapsto L_{\mathrm{phys}}(t)$, equipped with a law-connection
        \GZLOC\ (or triple-layer connection \GZTLC) and law-curvature
        \GZLCurv. These encode “how laws change” when we vary a
        physical parameter, energy scale, or background.
\end{enumerate}

From the law-space vantage point, each familiar physical configuration
is just a \emph{slice} of a much larger structure: a region in \GZLS
where certain law-curvature and homotopy invariants are small or fall
into a particular pattern (e.g.\ classical, weakly coupled, or
effective-theory regime).



\subsubsection*{(c) The GaoZheng G-Algebra as operator layer for physical laws}

Given a law-object $L_{\mathrm{phys}}\in\GZLS$, the GaoZheng G-Algebra
provides the associated law-level operator layer:

\begin{itemize}
  \item In the \emph{strict} regime, $\GAlgebra_{\mathrm{strict}}$
        encodes operators built from covariant derivatives, curvature
        insertions, and symmetry generators (including gauge
        transformations and diffeomorphisms), with Jacobi behaviour
        approximating that of a Lie algebra.
  \item In the \emph{homotopy} regime, $\GAlgebra_{\mathrm{homo}}$
        encodes higher-bracket structures that arise from anomalies,
        topological sectors, and non-perturbative effects; the ternary
        bracket $\ell_3^H$ captures controlled Jacobi failure linked
        to 3-form data such as Chern–Simons terms or fluxes.
  \item In the \emph{variable functional--operator} regime,
        $\GAlgebra_{\mathrm{vf\text- op}}$ organises operator paths
        (e.g.\ time-evolution generators, RG generators) and variable
        functionals (e.g.\ actions, effective actions, correlation
        functionals) under the path-lifted connection
        $\widetilde{\GZTLC}$.
\end{itemize}

Physically:

\begin{itemize}
  \item commutators in \GAlgebra\ encode not only field-strength
        algebra and symmetry brackets, but also law-level commutation
        relations between “change-of-law” operations (e.g.\ changing
        couplings, adding counterterms, modifying gauge-fixing);
  \item law-curvature in \GZLCurv\ measures obstructions to trivial
        reparametrisations of laws: for example, whether an RG flow can
        be flattened by a change of scheme, or whether a certain
        anomaly can be removed by local counterterms.
\end{itemize}

In this way, \GAlgebra\ is not merely a rephrasing of known Lie
algebras of symmetries; it is an operator-level avatar of the law-space
geometry in which those symmetries live.



\subsubsection*{(d) Triple-layer connection as meta-gauge structure}

The triple-layer connection \GZTLC\ decomposes into:

\begin{itemize}
  \item a geometric component affecting space-time geometry (e.g.\
        connections on $\Bundle$ and its associated bundles);
  \item a parameter component affecting couplings, scales, and other
        physical parameters (e.g.\ energy scale in RG, temperature in
        thermal field theory);
  \item a law component \GZLOC\ affecting law-objects in \GZLS.
\end{itemize}

In physical language, \GZTLC\ acts as a \emph{meta-gauge structure}:

\begin{itemize}
  \item geometric gauge transformations (ordinary gauge symmetry,
        diffeomorphisms) appear as portions of the structure group
        $\widehat{G}$ acting on the geometric component;
  \item parameter reparametrisations (change of renormalisation scheme,
        coupling redefinitions) appear in the parameter component;
  \item law reparametrisations (changes of effective action, choice of
        regulator, choice of “fundamental vs effective” description)
        appear in the law component.
\end{itemize}

Horizontal flows of \GZTLC\ correspond to transformations in which
physical predictions are preserved while geometry, parameters, and laws
are reshuffled in a controlled way. Curvature components measure
obstructions to such reshufflings: anomalies, scheme dependence,
irreducible non-locality, or intrinsic path-dependence of RG flows.



\subsubsection*{(e) Organization of subsequent sections}

The remainder of this chapter will sketch three broad classes of
applications, each seen through the four PFB-GNLA charts and the
underlying law-space:

\begin{itemize}
  \item \emph{Gauge theories and Yang–Mills type models.}
        We will interpret gauge fields and matter couplings as
        principal-bundle data, and show how strict and homotopy
        versions of \GAlgebra\ capture both classical symmetry
        structure and controlled anomalies. Law-connections encode
        flows between different gauge-theory regimes (e.g.\ UV/IR
        phases, Higgs phases) at the law-level.
  \item \emph{Gravitation and space-time geometry.}
        We will map connections on frame bundles to law-objects
        describing gravitational laws, analyse how law-curvature
        reflects different gravitational regimes (GR vs modified
        gravity), and sketch how CH-layering and continuum breakdown
        ideas manifest in the context of continuum space-time models.
  \item \emph{Quantum field theory, renormalisation, and effective
        field theory.}
        We will interpret RG flows, effective actions, and scheme
        dependence as law-space trajectories; variable
        functional--operator structures will be used to encode
        functional integrals and counterterm flows as operator paths
        under $\widetilde{\GZTLC}$, and we will briefly connect this
        to Jacobiator-gap and universality ideas in field theory.
\end{itemize}

Each of these sketches is deliberately high-level: the goal is to
illustrate the mapping from physical objects to law-space objects and
from physical flows to \GZTLC-horizontal or curvature-driven flows,
and to show how the same \GAlgebra/\GZLS/\GZTLC machinery recurs across
apparently very different physical domains.



\section{Gauge Theories and Yang--Mills Type Models}


Gauge theories are the natural first arena for testing the GaoZheng
G-Algebra and \\ law-connection machinery: they are already formulated in
terms of principal bundles, connections, and curvatures, and they
exhibit rich phenomena (symmetry breaking, anomalies, dualities) that
benefit from a law-level reinterpretation. In this section we outline
how Yang--Mills type models fit into the four PFB-GNLA charts and into
the law-space picture.



\subsubsection*{(a) Strict chart: classical Yang--Mills as a PFB-GNLA}

We start from the strict PFB-GNLA perspective, where gauge theories
take their familiar principal-bundle form.

\paragraph{Geometric data.}

Let $\Base$ be a space-time manifold (possibly with Euclidean signature
for path-integral treatments), $\Group$ a compact Lie group (e.g.\
$\mathrm{SU}(N)$), and $\Bundle\to\Base$ a principal $\Group$-bundle
equipped with a connection 1-form
\[
  \Conn \in \Omega^1(\Bundle;\mathfrak{g}),
  \qquad \mathfrak{g}=\mathrm{Lie}(\Group),
\]
with curvature
\[
  \Curv
  = \mathrm{d}\Conn + \tfrac12[\Conn,\Conn]
  \in \Omega^2(\Bundle;\mathfrak{g}).
\]
Matter fields live in associated bundles $E=\Bundle\times_\rho V$ for
representations $\rho:\Group\to\mathrm{GL}(V)$, with covariant
derivative $\nabla$ and curvature $R_\nabla=\rho_*(\Curv)$.

\paragraph{Strict operator layer.}

At the strict level, the GaoZheng G-Algebra associated with this data
is built from the operator space
\[
  A_{\mathrm{strict}}
  =
  \mathrm{span}\Bigl(
    \{\nabla_X\}_{X\in\Gamma(T\Base)},\,
    \{R_\nabla(X,Y)\},\,
    \text{gauge generators},\dots
  \Bigr),
\]
acting on spaces of fields (sections of $E$) and on composite operators.
The commutator
\[
  [\nabla_X,\nabla_Y]
  =
  R_\nabla(X,Y)
\]
and its variants encode the standard gauge-covariant identities; the
strict GNLA constraints ensure that Jacobi failure is confined to a
curvature sector that is either central or negligible in the regime of
interest.

From the law-space point of view, the functor $\Phi_{\mathrm{geom}}$
sends $(\Bundle,\Base,\Group;\Conn,\Curv;\text{matter})$ to a law-object
\[
  L_{\mathrm{YM,strict}}
  \in {\GZLS}_{\mathrm{strict}}\subset\GZLS,
\]
with $\GAlgebra_{\mathrm{strict}}$ as its operator layer. Different
choices of background gauge field $(\Conn,\Curv)$ and of coupling
constants correspond to moving within ${\GZLS}_{\mathrm{strict}}$, with
law-connection \GZLOC\ organising admissible deformations.



\subsubsection*{(b) Homotopy chart: anomalies, topological sectors, and $H$-twists}

Real gauge theories exhibit phenomena that go beyond strict Lie
structure: global anomalies, theta terms, Chern--Simons actions,
Wess--Zumino--Witten terms, etc. These effects naturally live in the
homotopy PFB-GNLA chart.

\paragraph{Three-form data and $H$-twists.}

Many anomaly-related structures are controlled by closed 3-forms $H$
(or their higher-form analogues). Examples include:

\begin{itemize}
  \item Chern--Simons 3-forms
        $\mathrm{CS}(A)\sim \mathrm{Tr}(A\wedge\mathrm{d}A
          +\tfrac23 A\wedge A\wedge A)$ on 3-manifolds;
  \item Wess--Zumino terms in sigma models, defined via integrals of
        3-forms over extension manifolds;
  \item fluxes in higher-form gauge theories and string-inspired
        backgrounds.
\end{itemize}

In the homotopy chart, such 3-forms become the $H$ in an
$H$-twisted two-term $L_\infty$-structure
$(V^{-1}\oplus V^0,\ell_1,\ell_2,\ell_3^H)$ associated with the gauge
theory. The ternary bracket $\ell_3^H$ encodes controlled Jacobi
failure linked to $H$ and to law-curvature components in \GZLCurv.

\paragraph{Anomalies as Jacobiator effects.}

Gauge and gravitational anomalies, at the operator level, often manifest
as failures of conserved currents or of Ward identities. In the
homotopy PFB-GNLA picture:

\begin{itemize}
  \item the strict commutator bracket $\ell_2$ encodes baseline symmetry
        algebra relations;
  \item the Jacobiator $\Jac_{\ell_2}$ captures anomaly-like obstructions
        to making these relations strictly associative/consistent;
  \item the homotopy $\ell_3^H$ represents the “repair data” that makes
        the overall $L_\infty$ structure coherent, even when $\ell_2$
        alone fails;
  \item Jacobiator-gap certificates $\mathsf{JacGap}$ classify which
        anomalies can be absorbed (via homotopy completions or
        counterterms) and which represent genuine obstructions.
\end{itemize}

From this angle, an “anomaly cancellation condition” can be phrased as
$\mathsf{JacGap}=0$ in the relevant homotopy sector, enabling a
strictification or a move to a more benign homotopy regime. Non-zero
$\mathsf{JacGap}$ corresponds to intrinsically anomalous phases.



\subsubsection*{(c) Variable functional--operator chart: actions, path integrals, and RG}

The variable functional--operator chart is the natural home for
action functionals, path integrals, and renormalisation group (RG)
flows.

More concretely, if $S[\Gamma]$ denotes an action functional on an
operator path $\Gamma$, then classical or ``optimal'' paths arise as
stationary points of $S$, i.e.\ as solutions of the variational
equation
\[
  \frac{\delta S[\Gamma]}{\delta \Gamma} = 0.
\]
In the semiclassical regime, the corresponding path integrals are
dominated by contributions from neighbourhoods of such stationary
paths, so that operator paths solving the Euler--Lagrange equations
provide the backbone of the quantum theory.


\paragraph{Operator paths and effective actions.}

In this chart:

\begin{itemize}
  \item operator paths $\Gamma:I\to V^0$ can represent, for example,
        RG trajectories of effective Hamiltonians or Lagrangians:
        $\Gamma(s)$ is the effective operator at scale $s$;
  \item variable functionals
        $\mathcal{F}[\Gamma]$ can be action functionals,
        generating functionals, or free energies evaluated along such
        trajectories;
  \item the lifted connection $\widetilde{\GZTLC}$ encodes how these
        operators and functionals are transported under simultaneous
        changes of space-time geometry, couplings, and law-objects.
\end{itemize}

Law--operator dynamical curvature $\mathcal{F}_{\mathrm{dyn}}$
measures the obstruction to trivialising RG flows by reparametrisation:
if $\mathcal{F}_{\mathrm{dyn}}=0$, an RG trajectory is “pure gauge” in
law-space; non-zero $\mathcal{F}_{\mathrm{dyn}}$ indicates intrinsic,
scheme-independent structure (fixed points, universality classes,
irreducible irrelevant operators, etc.).

\paragraph{GZ-OHU and universality.}

The GZ-OHU framework applies directly to operator families in QFT:

\begin{itemize}
  \item Jacobiator-gap certificates capture controlled failures of
        strict Lie-type relations among composite operators or
        effective generators;
  \item GZ-OHU completions provide homotopy-universal operator
        systems that realise all \\ operator-homotopy types compatible
        with low-energy observables;
  \item fixed points in the RG flow correspond to law-space regions
        where $\mathsf{JacGap}$ and curvature diagnostics stabilise,
        and where different microscopic theories represent the same
        $\GAlgebra_{\mathrm{orig}}$ up to representation.
\end{itemize}

From this point of view, universality in critical phenomena can be
read as a special case of operator--homotopy universality: many
distinct microscopic laws share the same law-level \GAlgebra\ in the
infrared.



\subsubsection*{(d) Original PL-PI law-level chart: gauge theories as PL-PI systems}

Finally, in the meta-mathematical original chart, gauge theories are
seen directly as PL-PI law-systems:

\begin{itemize}
  \item $L_{\mathrm{PL}}$ encodes logical structure: which kinds of
        field configurations, constraints, and gauge equivalences are
        allowed; how propositions about fields are inferred (e.g.\
        locality, causality, analyticity assumptions).
  \item $L_{\mathrm{PI}}$ encodes iterative procedures: RG algorithms,
        variational principles, numerical \\ schemes, Monte Carlo updates,
        etc.
  \item $\Theta_{\mathrm{PL\text- PI}}$ encodes consistency conditions
        (Ward identities, reflection positivity, locality, unitarity)
        linking inference and procedure.
\end{itemize}

The original law-level \GAlgebra\ $\GAlgebra_{\mathrm{orig}}$ then
emerges as a GNLA/$L_\infty$ structure on the space of “change-of-law”
operations: modifying gauge groups, changing representations and
couplings, adding or removing terms in the action, changing
renormalisation prescriptions, etc.

PFB-GNLA charts (strict, homotopy, variable functional--operator) are
various realisations of $\GAlgebra_{\mathrm{orig}}$ on object-level
geometries and operator-path spaces, adapted to specific physical
questions (classical solutions, anomalies, RG flows).



\subsubsection*{(e) Summary: gauge theories as testbeds for law-space physics}

To summarise the role of Yang--Mills type models in the GaoZheng
G-Framework:

\begin{itemize}
  \item They provide a canonical example of PFB geometry, making the
        strict PFB-GNLA chart immediately recognisable to physicists:
        $(\Bundle,\Base,\Group;\Conn,\Curv)$ maps to
        $\GAlgebra_{\mathrm{strict}}$.
  \item They naturally exhibit homotopy phenomena (anomalies, topological
        terms) that demand the homotopy chart
        $\GAlgebra_{\mathrm{homo}}$ and $H$-twisted $L_\infty$
        structures.
  \item Their RG and effective-action structures are cleanly expressed
        in the variable functional--operator chart
        $\GAlgebra_{\mathrm{vf\text- op}}$, where operator paths and
        variable functionals live under $\widetilde{\GZTLC}$.
  \item The entire family of gauge-theory regimes (different groups,
        couplings, anomalies, effective descriptions) can be unified as
        PL-PI law-systems whose original law-level \GAlgebra\
        $\GAlgebra_{\mathrm{orig}}$ governs the space of admissible
        “changes of law”.
\end{itemize}

The next section will apply the same logic to gravitation and
space-time geometry, interpreting frame bundles and metric connections
as slices of law-space and exploring how CH-layering and continuum
breakdown ideas naturally appear when “space-time” itself is treated as
a law-level continuum object within the GaoZheng G-Framework.



\section{Gravitation and Space--Time Geometry}


Gravitation is the second canonical arena where principal bundles and
connections have long been central: general relativity and its
extensions are naturally phrased in terms of frame bundles, affine
connections, curvature tensors, and local Lorentz symmetry. In this
section we outline how these structures sit inside the GaoZheng
G-Framework, and how the four PFB-GNLA charts offer distinct but
compatible views on gravitational laws and space--time geometry.



\subsubsection*{(a) Strict chart: classical GR as a PFB-GNLA}

We begin with the strict PFB-GNLA description of classical general
relativity (GR) and closely related metric theories.

\paragraph{Geometric data: frame bundles and metric connections.}

Let $\Base$ be a 4-dimensional Lorentzian manifold representing
space--time, endowed with a metric $g$ of signature $(-,+,+,+)$. The
orthonormal frame bundle $\Bundle\to\Base$ has structure group
$\Group=\mathrm{SO}(1,3)$ (or its double cover $\mathrm{Spin}(1,3)$).
A space--time connection can be viewed as:

\begin{itemize}
  \item the Levi-Civita connection $\nabla$ on $(\Base,g)$,
        torsion-free and metric-compatible;
  \item or, in gauge-theoretic language, as a connection 1-form
        $\Conn\in\Omega^1(\Bundle;\mathfrak{so}(1,3))$ on the frame
        bundle, with curvature $\Curv$ representing the Riemann tensor
        in suitable bases.
\end{itemize}

Matter fields (e.g.\ spinors, tensors) live in associated bundles
$E=\Bundle\times_\rho V$ for representations $\rho$ of
$\mathrm{SO}(1,3)$ or $\mathrm{Spin}(1,3)$, with covariant derivatives
and curvature acting as usual.

\paragraph{Strict operator layer for gravity.}

From the strict PFB-GNLA perspective, we form an operator space
\[
  A_{\mathrm{strict}}^{\mathrm{grav}}
  =
  \mathrm{span}\Bigl(
    \{\nabla_X\}_{X\in\Gamma(T\Base)},\,
    \{R(X,Y)\},\,
    \text{local frame rotations},\dots
  \Bigr),
\]
acting on sections of $E$ and on composite observables (stress--energy
tensors, curvature scalars, etc.). Commutators reproduce familiar
identities:
\[
  [\nabla_X,\nabla_Y]
  =
  R(X,Y),
\]
and Bianchi identities constrain $R$ and its derivatives.

The strict GNLA constraints (S1)--(S4) from Part~II are in this setting
essentially classical differential-geometric facts: Jacobi failure is
localised in curvature sectors that, for purely geometric operators,
are either centralised or tightly constrained by Bianchi identities.

Via $\Phi_{\mathrm{geom}}$, the data
$(\Bundle,\Base,\Group;\Conn,\Curv;\text{matter})$ define a law-object
\[
  L_{\mathrm{GR,strict}}
  \in {\GZLS}_{\mathrm{strict}}\subset\GZLS,
\]
with $\GAlgebra_{\mathrm{strict}}^{\mathrm{grav}}$ as its operator
layer. Different choices of metric (solutions of Einstein equations,
backgrounds with different curvature, etc.) correspond to different
points or regions in ${\GZLS}_{\mathrm{strict}}$.



\subsubsection*{(b) Homotopy chart: torsion, topological terms, and extended gravity}

Beyond classical GR lie numerous extensions: torsionful connections
(Cartan geometry), topological terms (e.g.\ Gauss--Bonnet, Pontryagin
densities), higher-derivative gravity, and quantum anomalies in
curved backgrounds. These effects are naturally captured in the
homotopy PFB-GNLA chart.

\paragraph{Torsion and non-metricity as homotopy data.}

In a general metric-affine gravity, one allows:

\begin{itemize}
  \item torsion $T\neq0$ (antisymmetric part of connection coefficients);
  \item non-metricity $Q\neq0$ (failure of $\nabla g=0$).
\end{itemize}

At the operator level, this introduces extra commutator and Jacobiator
terms when acting on fields: covariant derivatives no longer satisfy
the same symmetries, and additional defect terms appear in curvature
identities. In the homotopy chart:

\begin{itemize}
  \item the binary bracket $\ell_2$ encodes extended covariant
        derivatives and curvature operators;
  \item the ternary bracket $\ell_3^H$ encodes torsion/non-metricity
        contributions and relations stemming from closed 3-forms $H$
        built from $(T,Q,R)$ and their contractions;
  \item higher brackets $\ell_k$ can encode yet more subtle geometric
        structures (e.g.\ Lovelock terms, higher-curvature invariants).
\end{itemize}

From this angle, torsion and non-metricity are not ad hoc additions,
but controlled components of the $L_\infty$ structure associated with
$L_{\mathrm{GR,homo}}$.

\paragraph{Topological terms and gravitational anomalies.}

Topological invariants in gravity (Euler, Pontryagin, Nieh--Yan
classes) often appear as integrals of differential forms of degree
4 or higher, many of which can be written as $\mathrm{d}$ of 3-forms.
These 3-forms provide natural candidates for the $H$ in $H$-twisted
$L_\infty$ structures.

Similarly, gravitational anomalies (e.g.\ in chiral theories on curved
backgrounds) manifest themselves as failures of certain Ward
identities or index-theoretic balances. In the homotopy chart:

\begin{itemize}
  \item such anomalies show up as non-trivial components of the
        Jacobiator or higher homotopy brackets;
  \item Jacobiator-gap certificates $\mathsf{JacGap}$ measure whether
        these anomalies can be absorbed by extending the operator
        algebra (e.g.\ adding auxiliary fields) or whether they
        represent intrinsic obstructions.
\end{itemize}

Anomaly cancellation conditions (for gravity or mixed gauge--gravity)
can thus be rephrased as conditions for the existence of a strict or
“smaller-gap” homotopy version of $\GAlgebra^{\mathrm{grav}}$.



\subsubsection*{(c) Variable functional--operator chart: actions, path integrals, and effective gravity}

The variable functional--operator chart provides a natural home for:

\begin{itemize}
  \item gravitational action functionals (Einstein--Hilbert, higher-
        curvature corrections, cosmological constant terms);
  \item path integrals over metrics (or over connections and tetrads),
        at least in a formal or effective sense;
  \item effective field theory (EFT) flows of gravitational couplings
        under coarse-graining (RG-like flows in gravity).
\end{itemize}

\paragraph{Operator paths of gravitational couplings.}

In this chart, we may consider:

\begin{itemize}
  \item an operator path $\Gamma(s)$ representing an EFT Hamiltonian or
        Lagrangian at scale $s$, including gravitational terms with
        scale-dependent couplings (Newton’s constant, cosmological
        constant, higher-curvature coefficients);
  \item variable functionals $\mathcal{F}[\Gamma]$ representing, for
        example, effective actions or generating functionals after
        integrating out matter fields up to scale $s$;
  \item the path-lifted connection $\widetilde{\GZTLC}$ encoding how
        these operators and functionals transform when we vary both
        the scale $s$ and the law-level background (e.g.\ choice of
        EFT scheme, choice of background metric family).
\end{itemize}

Non-trivial dynamical curvature $\mathcal{F}_{\mathrm{dyn}}$ signals
that gravitational RG flows have genuine, scheme-independent fixed
points or attractors (as in asymptotic safety scenarios), whereas
vanishing $\mathcal{F}_{\mathrm{dyn}}$ would correspond to flows that
are pure reparametrisations of the same law.

\paragraph{CH-layering and continuum breakdown for space--time.}

Space--time itself is a continuum law-object; CH-layering and GZ-CBT
ideas can be applied here as well:

\begin{itemize}
  \item $\mathrm{CH}^{(l_2)}$ for space--time refers to cardinal
        properties of its continuum (e.g.\ underlying set of events,
        or its embedding into models of ZFC);
  \item $\mathrm{CH}^{(l_{\ge3})}$ refers to semantic metrics on
        space--time: regularity of causal structure, smoothness,
        global hyperbolicity, behaviour of geodesic completeness, etc.
\end{itemize}

Continuum breakdown events (singularities, topology change, failure of
global hyperbolicity) can thus be seen as GZ-CBT events for the
space--time law-object. Rewriting schemes then correspond to law-level
procedures that replace or extend the classical GR description in
regions where CH-like semantics fail (e.g.\ near singularities, at
Planck scales), guided by curvature diagnostics and homotopy
completions.



\subsubsection*{(d) Original PL-PI law-level chart: gravity as a continuum-law system}

From the PL-PI meta-mathematical perspective, gravitational theories
are law-systems specifying:

\begin{itemize}
  \item which space--time structures are admissible (Lorentzian
        manifolds, tetrad fields, connections, causal structures);
  \item which inference schemes are allowed for physical predictions
        (field equations, constraint equations, boundary conditions,
        variational principles);
  \item which iterative schemes generate solutions or effective laws
        (numerical GR algorithms, perturbative expansions, RG flows,
        coarse-graining of degrees of freedom, etc.).
\end{itemize}

In this chart:

\begin{itemize}
  \item $L_{\mathrm{PL}}^{\mathrm{grav}}$ encodes logical rules for
        constructing admissible space--times and interpreting physical
        statements about them;
  \item $L_{\mathrm{PI}}^{\mathrm{grav}}$ encodes iterative procedures
        and algorithms used to solve Einstein equations, perform
        EFT expansions, or implement quantum-gravity-inspired
        discretisations;
  \item $\Theta_{\mathrm{PL\text- PI}}^{\mathrm{grav}}$ enforces
        meta-level coherence (e.g.\ consistency of constraints, Bianchi
        identities, covariance of solution spaces).
\end{itemize}

The original law-level \GAlgebra\ $\GAlgebra_{\mathrm{orig}}^{\mathrm{grav}}$
then governs admissible “changes of gravitational law”: modifying the
action, changing the space of admissible metrics, adding new fields,
imposing different causal or CH-layering constraints, etc.
PFB-GNLA charts provide various representations of this original
\GAlgebra\ on concrete geometric and operator-path spaces.



\subsubsection*{(e) Summary: space--time as a law-level continuum}

To summarise, gravitation and space--time geometry fit into the
GaoZheng G-Framework as follows:

\begin{itemize}
  \item Strict PFB-GNLA reproduces classical GR as a principal-bundle
        geometry with a near-Lie operator algebra of covariant
        derivatives and curvature operators.
  \item Homotopy PFB-GNLA captures torsion, non-metricity, topological
        terms, and gravitational anomalies as controlled homotopy
        corrections in an $H$-twisted $L_\infty$ structure.
  \item Variable functional--operator PFB-GNLA organises gravitational
        actions, path integrals, and EFT/RG flows as operator paths and
        variable functionals under a lifted connection
        $\widetilde{\GZTLC}$.
  \item The original PL-PI law-level chart treats space--time itself as
        a continuum law-object whose semantics can break down and be
        rewritten (e.g.\ near singularities), with CH-layering and
        GZ-CBT providing a quantitative language for such breakdown and
        reconstruction.
\end{itemize}

Thus, in the GaoZheng G-Framework, gravitation is not merely another
field theory on a given background, but a paradigmatic example of how a
continuum law-object (space--time) and its governing laws (Einstein
equations and their extensions) can be unified into a single law-space
geometry, analysed via \GAlgebra, \GZLS, \GZTLC, and the associated
homotopy and CH-layer structures.



\section{Quantum Field Theory, Renormalisation, and Effective Laws}


Quantum Field Theory (QFT) is the natural meeting point of gauge
theories, gravitation, and statistical mechanics (see, for example, \cite{Weinberg1995,Haag1996}). It is also the domain
where renormalisation, effective field theory (EFT), anomalies, and
universality are most prominently displayed. In this section we sketch
how QFT and renormalisation fit into the GaoZheng G-Framework: strict
PFB-GNLA for local field algebras, homotopy PFB-GNLA for anomalies and
BV/BRST structures, variable functional--operator PFB-GNLA for RG flows
and effective actions, and the original PL-PI chart for QFT as a
meta-level law-system of inference and procedures.



\subsubsection*{(a) Strict chart: local QFT as a near-Lie operator algebra}

At the strict level, we view QFT through the lens of local algebras of
fields or observables. Two standard formalisms illustrate this:

\begin{itemize}
  \item \emph{Wightman axioms / Haag--Kastler framework.}
        QFT is specified by a net of local *-algebras
        $\mathcal{A}(\mathcal{O})$ associated with space--time regions
        $\mathcal{O}\subseteq\Base$, with covariance, locality, and
        spectrum conditions.
  \item \emph{Canonical quantisation.}
        Fields $\Phi(x)$ and conjugate momenta $\Pi(x)$ obey equal-time
        commutation relations; smeared operators $\Phi(f)$ generate an
        algebra acting on a Hilbert space.
\end{itemize}

The strict PFB-GNLA chart organises these structures as follows:

\begin{itemize}
  \item geometric data $(\Bundle,\Base,\Group;\Conn)$ describe space--time
        geometry and internal symmetries as in Sections~11.1--11.2;
  \item matter fields and gauge fields define local operator algebras
        (smeared fields, currents, Hamiltonians);
  \item the strict operator layer $A_{\mathrm{strict}}^{\mathrm{QFT}}$
        consists of generators of these algebras and their composites
        (local operators, charges, symmetry generators, Hamiltonian
        densities).
\end{itemize}

The commutator
\[
  [\mathcal{O}_1,\mathcal{O}_2]_{\mathrm{strict}}
\]
approximates canonical or graded commutation relations (microcausality,
current algebras, etc.). Strict GNLA constraints ensure that Jacobi
identities hold in the relevant sector, or that their failures are
centralised (e.g.\ Kac--Moody central extensions) and thus compatible
with a near-Lie description.

Via $\Phi_{\mathrm{geom}}$, a given QFT on a fixed background
(geometry + couplings) determines a law-object
\[
  L_{\mathrm{QFT,strict}}
  \in {\GZLS}_{\mathrm{strict}},
\]
with $\GAlgebra_{\mathrm{strict}}^{\mathrm{QFT}}$ as its operator
layer. Different bare actions and couplings correspond to moving within
this slice.



\subsubsection*{(b) Homotopy chart: BV/BRST structures, anomalies, and OPE}

Strict Lie-type structure is insufficient for a full description of
modern QFT: gauge fixing, ghosts, BRST symmetry, BV formalism, and OPE
cohomology naturally lead to higher-bracket and homotopy structures.

\paragraph{BV/BRST as $L_\infty$ structures.}

The Batalin--Vilkovisky (BV) and BRST formalisms encode gauge symmetry
and its redundancies via:

\begin{itemize}
  \item a graded space of fields, ghosts, antifields, etc.;
  \item a BRST differential $s$ and/or BV operator $\Delta$;
  \item higher-order brackets and homotopies enforcing the master
        equation and its deformations.
\end{itemize}

From the homotopy PFB-GNLA viewpoint:

\begin{itemize}
  \item $V^0$ contains physical and ghost fields (and associated
        operators);
  \item $V^{-1}$ records gauge identities, Noether identities, and
        deformation directions;
  \item $\ell_1$ is the BRST differential;
  \item $\ell_2$ encodes graded commutators;
  \item $\ell_3^H$ and higher $\ell_k$ encode higher-order constraints,
        gauge-for-gauge structure, and anomalies.
\end{itemize}

The $H$ in the $H$-twisted structure can come from 3-form data built from the \\
gauge/gravitational background or from topological terms in
the action (e.g.\ Wess--Zumino terms). Jacobiator-gap certificates
$\mathsf{JacGap}$ measure whether a given BV/BRST complex admits a
strictification (e.g.\ anomaly-free, consistent quantisation) or
requires genuinely higher-form or higher-bracket structure.

\paragraph{OPE and homotopy algebras.}

Operator product expansions (OPE) in QFT generate infinite families of
composite operators with associativity constraints; these constraints
are often naturally expressed in terms of homotopy algebras (e.g.\
$A_\infty$, $L_\infty$). In the homotopy chart:

\begin{itemize}
  \item $\ell_2$ captures leading OPE brackets;
  \item $\ell_3,\ell_4,\dots$ encode homotopy-associativity and higher
        coherence relations between OPE coefficients;
  \item law-curvature \GZLCurv\ records obstructions to simplifying or
        reparametrising OPE structures (e.g.\ non-trivial CFT data).
\end{itemize}

Thus, anomalies, gauge fixing, BV/BRST, and OPE constraints all appear
as manifestations of homotopy PFB-GNLA structure in QFT.



\subsubsection*{(c) Variable functional--operator chart: RG flows and effective field theories}

The variable functional--operator chart is tailored for renormalisation
and effective field theories: actions and couplings evolve under
coarse-graining, and path integrals define scale-dependent functionals.

In particular, when one associates to a family of effective actions
$S_\Lambda$ an appropriate path-integral weight, those RG trajectories
that extremise $S_\Lambda$ (in the sense of $\delta S_\Lambda=0$) play
the role of semiclassical or ``optimal'' flows whose neighbourhoods
give the leading contributions to scale-dependent functionals.


\paragraph{Wilsonian RG as operator paths.}

In Wilsonian RG, one considers a family of effective actions
$S_\Lambda$ or Hamiltonians $H_\Lambda$ as the cutoff scale $\Lambda$
changes. In the variable functional--operator chart:

\begin{itemize}
  \item the parameter $s$ (or $t$) can be taken as RG “time” (logarithm
        of the scale);
  \item operator paths
        $\Gamma(s)\in V^0$ represent effective generators (e.g.\
        Hamiltonians or action functionals) at each scale;
  \item variable functionals $\mathcal{F}[\Gamma]$ represent, for
        instance, generating functionals, free energies, or
        coarse-grained observables.
\end{itemize}

The lifted connection $\widetilde{\GZTLC}$ encodes how these operators
and functionals change when both the scale and law-level background
(e.g.\ choice of RG scheme, field redefinitions) are varied.
Law--operator dynamical curvature $\mathcal{F}_{\mathrm{dyn}}$ measures
the intrinsic, scheme-independent content of the RG flow: fixed points,
critical exponents, universality classes.

\paragraph{Effective field theory as law-space localisation.}

Effective field theory (EFT) can be rephrased as working in a region of
$\GZLS$ where:

\begin{itemize}
  \item law-curvature is small or controlled in directions associated
        with irrelevant operators;
  \item law-connections in relevant/marginal directions (couplings of
        interest) dominate the flow;
  \item higher-homotopy data encode systematic corrections and
        effective interactions (higher-dimensional operators).
\end{itemize}

In GZ-OHU language:

\begin{itemize}
  \item many different microscopic theories represent the same
        homotopy-complete law-object $L^{\mathrm{OHU}}$ in the IR;
  \item Jacobiator-gap certificates classify which deformations are
        “irrelevant” from the law-level viewpoint (flow into the same
        $L^{\mathrm{OHU}}$) and which lead to distinct universality
        classes (different law-level homotopy types with non-trivial
        $\mathsf{JacGap}$).
\end{itemize}

RG equations (beta functions) then appear as coordinate expressions of
\GZTLC-horizontal or curvature-driven flows in $\GZLS$, and EFT
hierarchies as local trivializations of law-curvature in appropriate
directions.



\subsubsection*{(d) Original PL-PI law-level chart: QFT as a law-system of inference and procedures}

From the PL-PI perspective, a QFT is not just a Lagrangian; it is a
paired system of:

\begin{itemize}
  \item \emph{inference laws} $L_{\mathrm{PL}}^{\mathrm{QFT}}$:
        axioms and logical rules governing correlation functions,
        amplitudes, and observables (locality, causality, analyticity,
        crossing symmetry, cluster decomposition, axioms of Wightman or
        Osterwalder--Schrader type);
  \item \emph{iterative laws} $L_{\mathrm{PI}}^{\mathrm{QFT}}$:
        computational/constructive procedures (perturbation theory,
        lattice discretisations, Monte Carlo methods, tensor network
        algorithms, RG schemes) that produce approximations to these
        observables.
\end{itemize}

The coupling $\Theta_{\mathrm{PL\text- PI}}^{\mathrm{QFT}}$ enforces
coherence between these two aspects: it encodes renormalisability,
Ward identities, reflection positivity, and consistency conditions
ensuring that iterative procedures truly approximate the intended
inference laws.

The original law-level \GAlgebra\ $\GAlgebra_{\mathrm{orig}}^{\mathrm{QFT}}$
then becomes:

\begin{itemize}
  \item a GNLA/$L_\infty$ structure acting on “change-of-law”
        operations: modifying Lagrangians, adding new fields,
        changing renormalisation prescriptions, switching between
        different PL-PI regimes;
  \item a meta-level object whose representations on PFB geometries and
        operator-path spaces are precisely the various QFTs and EFTs
        physicists study.
\end{itemize}

Thus, the diversity of QFT models (different local interactions,
dimensions, field content) is organised as a law-space family of
representations of a smaller number of law-level operator structures
and PL-PI patterns.



\subsubsection*{(e) Summary: renormalisation and universality as law-space phenomena}

Across the four PFB-GNLA charts, QFT, renormalisation, and effective
laws take on a unified law-space interpretation:

\begin{itemize}
  \item Strict PFB-GNLA: local field algebras and bare couplings form a
        near-Lie operator algebra on a fixed background geometry;
  \item Homotopy PFB-GNLA: BV/BRST, anomalies, and OPE homotopies are
        encoded as $H$-twisted $L_\infty$ structures on operator
        spaces, with Jacobiator-gap certificates classifying anomaly
        and coherence patterns;
  \item Variable functional--operator PFB-GNLA: RG flows and EFT
        hierarchies are operator paths in $\GZLS$, acted upon by
        $\widetilde{\GZTLC}$; universality classes correspond to
        homotopy-complete law-objects $L^{\mathrm{OHU}}$ in the IR;
  \item Original PL-PI law-level chart: QFTs are PL-PI systems of
      inference and procedures; \(\GAlgebra_{\mathrm{orig}}^{\mathrm{QFT}}\)
      encapsulates the meta-level algebra of law changes.
\end{itemize}

In this sense, renormalisation and universality in QFT are not just
properties of particular Lagrangians, but manifestations of the
geometry and homotopy of law-space encoded by \GZLS, \GZTLC, and
\GAlgebra. The GaoZheng G-Framework thus offers a structural bridge
between abstract PFB-GNLA theory and the concrete, multi-scale
behaviour of quantum fields.



\section{Law-Space Dynamics: Differential and\texorpdfstring{\\}{ }Path-Integral Quantum Mechanisms}


The preceding sections have illustrated how gauge theories, gravitation,
and quantum field theory can be read inside the GaoZheng G-Framework:
principal bundles with connection and curvature live in the
object-level component $\mathsf{Obj}$; law-space \GZLS, law-connection
\GZLOC, and triple-layer geometry \GZTLC\ organise admissible flows and
homotopies; and the four integrated PFB-GNLA charts provide different
but compatible views of the same law-level landscape.

From the dynamical point of view, the core mechanisms of the
GaoZheng G-Framework can be summarised as follows. On the classical
side, law-level dynamics are represented by trajectories
\[
  \Gamma:[0,T]\to\GZLS,
\]
together with horizontal lifts to $\mathsf{Obj}$ with respect to the
triple-layer connection \GZTLC. The evolution of such trajectories is
governed schematically by law-level equations of motion
\[
  \frac{d\Gamma}{dt} = X_{\mathrm{law}}(\Gamma(t)),
\]
where $X_{\mathrm{law}}$ is a law-vector field determined by
\GZLOC\ and \GZLCurv, encoding how changes in space--time geometry,
fields, and couplings must remain compatible with the underlying
law-space constraints.

On the quantum side, the same framework admits a generalized Feynman
path integral on law-space (the GaoZheng law-space path integral,
GZ-LSPI). A law-action functional $S_{\mathrm{law}}[\Gamma]$ is
assigned to each admissible law-level history $\Gamma$, and one
considers the formal law-space partition function
\[
  \mathcal{Z}_{\mathrm{law}}
  = \int_{\mathrm{Hist}(\GZLS)}
    \exp\!\Bigl(\frac{i}{\hbar}S_{\mathrm{law}}[\Gamma]\Bigr)\,
    \mathcal{D}\Gamma.
\]
Ordinary configuration-space or space--time path integrals in gauge
theories and quantum field theory then appear as projections of
$\mathcal{Z}_{\mathrm{law}}$ along the functorial maps constructed in
Parts~II--III: effective actions, propagators, and renormalisation flows
are shadows of law-level histories and curvatures. In this sense, the
GaoZheng G-Framework treats “quantum dynamics” not merely as a
path-integral on fields, but as a path-integral on laws themselves.

In the present volume, these differential and path-integral mechanisms
constitute the primary dynamical content of the GaoZheng G-Framework.
Possible applications to reinforcement-learning-type procedures,
certificate-based AI, or other engineering schemes are regarded as
downstream realisations of the same law-space dynamics and are therefore
left to subsequent work focused on law-space learning.

From a formal point of view, such reinforcement-learning-type
constructions amount to specialisations of the generalized Feynman path
integral on law-space (GZ-LSPI), where $S_{\mathrm{law}}$ is replaced
by an appropriate reward or cost functional on law-trajectories; this
AI-oriented specialisation is not pursued further in this
physics-oriented volume.




\appendix



\chapter{Technical Appendix I: Core Mathematical Proofs (Minimal Verifiable Package)}


\section{Theorem T1 (GZ-OHU): Statement and Proof}
\label{app:I-T1}

\subsection*{Set-up and hypotheses}


In this appendix we make precise the schematic Theorem~T1
(GZ-OHU theorem) stated in the main text and provide a complete
construction and proof.

We work in a fixed ambient law-category
\(\mathsf{Law}\), whose objects are PL--PI type law-systems and whose
morphisms are structure-preserving maps between them.  A typical object
will be denoted by
\[
  L = \bigl(\mathcal{O}_L, \mathcal{C}_L, \GAlgebra_L,\dots\bigr),
\]
where \(\mathcal{O}_L\) is a space of observables/operations,
\(\mathcal{C}_L\) records constraints and consistency data, and
\(\GAlgebra_L\) is the associated law-level algebra encoding
compositions and higher brackets.

\begin{definition}[Admissible law-system]
A law-system \(L\in\mathsf{Law}\) is called
\emph{GZ-admissible} if it satisfies the following conditions:
\begin{enumerate}
  \item[(i)] \textbf{Topological/metric control.}
    The underlying operation space \(\mathcal{O}_L\) is a complete,
    separable metric (or Banach) space, and all structural operations
    (brackets, compositions, constraint maps) are continuous and
    locally Lipschitz with respect to this metric.
  \item[(ii)] \textbf{Finite local generation.}
    There exists a countable family of ``local generators''
    \(\{x_i\}_{i\in\mathbb{N}}\subset\mathcal{O}_L\) such that every
    operation of interest can be obtained from \(\{x_i\}\) by a finite
    number of compositions, bracketings, and limits along Cauchy
    sequences.
  \item[(iii)] \textbf{Controlled higher brackets.}
    The higher brackets and Jacobiators of \(\GAlgebra_L\) are defined
    on a dense subspace of \(\mathcal{O}_L\) and extend continuously
    to all of \(\mathcal{O}_L\).
\end{enumerate}
\end{definition}



We now formalise the Jacobiator-gap certificate.

\begin{definition}[Jacobiator and Jacobiator-gap certificate]
Let \(L\) be a GZ-admissible law-system.  The
\emph{Jacobiator} of \(L\) is a (multi-)linear map
\[
  J_L : \mathcal{O}_L^{\otimes 3} \longrightarrow \mathcal{O}_L
\]
measuring the failure of the Jacobi identity (or, more generally, the
failure of the higher \(L_\infty\)-type identities) in \(\GAlgebra_L\).

A \emph{Jacobiator-gap certificate} on \(L\) is a finite family of
non-negative continuous functionals
\[
  \Phi_L = \bigl(\phi_1,\dots,\phi_m\bigr),\qquad
  \phi_k : \mathcal{O}_L^{\otimes 3}\to\mathbb{R}_{\ge 0},
\]
such that:
\begin{enumerate}
  \item[(a)] For every triple \((a,b,c)\in\mathcal{O}_L^{\otimes 3}\),
    the combined defect
    \[
      \mathsf{JacGap}_L(a,b,c)
      := \sum_{k=1}^m \phi_k\bigl(J_L(a,b,c)\bigr)
    \]
    is finite.
  \item[(b)] \(J_L\equiv 0\) if and only if
    \(\mathsf{JacGap}_L(a,b,c)=0\) for all \((a,b,c)\).
  \item[(c)] There exists a constant \(C>0\) such that
    \[
      \|J_L(a,b,c)\|
      \;\le\;
      C\,\bigl(1+\mathsf{JacGap}_L(a,b,c)\bigr)
      \quad\text{for all }(a,b,c),
    \]
    where \(\|\cdot\|\) is the ambient norm on \(\mathcal{O}_L\).
\end{enumerate}
We write \(\mathsf{JacGap}(L)\) for the resulting defect functional.
\end{definition}



We also allow an additional ``CH-regularity'' energy functional
\(\mathcal{E}_{\mathrm{CH}}(L)\) (as in the main text), but the core of
Theorem~T1 only requires the Jacobiator-gap data.

\begin{definition}[OHU-completion]
Let \(L\) be a GZ-admissible law-system equipped with a
Jacobiator-gap certificate \(\mathsf{JacGap}(L)\).  An
\emph{operator--homotopy upgrade completion} (or \emph{GZ-OHU
completion}) of \((L,\mathsf{JacGap}(L))\) consists of:
\begin{enumerate}
  \item[(i)] A law-system \(L^{\mathrm{OHU}}\in\mathsf{Law}\),
  \item[(ii)] A structure-preserving embedding
    \(\iota : L\hookrightarrow L^{\mathrm{OHU}}\),
\end{enumerate}
such that \(L^{\mathrm{OHU}}\) is homotopy-complete with respect to the
Jacobiator-defect and satisfies the universal property described below.
\end{definition}



\subsection*{Statement of Theorem T1}


We can now state the precise version of Theorem~T1.

\begin{theorem}[GZ-OHU completion theorem]
\label{thm:GZ-OHU}
Let \(L\in\mathsf{Law}\) be a GZ-admissible law-system equipped with a
Jacobiator-gap certificate \(\mathsf{JacGap}(L)\).  Then there exist a
law-system \(L^{\mathrm{OHU}}\in\mathsf{Law}\) and a structure-preserving
embedding
\[
  \iota : L\hookrightarrow L^{\mathrm{OHU}}
\]
such that the following properties hold.

\begin{enumerate}
  \item[\textnormal{(1)}] \textbf{Extension of structure.}
    The law-system \(L^{\mathrm{OHU}}\) extends \(L\): the underlying
    operation space \(\mathcal{O}_{L^{\mathrm{OHU}}}\) contains
    \(\mathcal{O}_L\) as a dense subspace, and all operations of \(L\)
    agree with the restrictions of the corresponding operations of
    \(L^{\mathrm{OHU}}\) along \(\iota\).

  \item[\textnormal{(2)}] \textbf{Homotopy-completeness.}
    There exists a completion of the Jacobiator-defect functional to
    \(L^{\mathrm{OHU}}\),
    \[
      \mathsf{JacGap}\bigl(L^{\mathrm{OHU}}\bigr)
      : \mathcal{O}_{L^{\mathrm{OHU}}}^{\otimes 3}\to\mathbb{R}_{\ge 0},
    \]
    such that:
    \begin{enumerate}
      \item[(a)] The Jacobiator \(J_{L^{\mathrm{OHU}}}\) vanishes on a
        dense subspace up to controlled homotopy; more precisely, for
        every finite set of triples there exists a homotopy-equivalent
        deformation of the law-operations on \(L^{\mathrm{OHU}}\) for
        which the Jacobiator vanishes on those triples.
      \item[(b)] Any further homotopy deformation of
        \(L^{\mathrm{OHU}}\) that strictly reduces
        \(\mathsf{JacGap}(L^{\mathrm{OHU}})\) on a finite set of
        triples must factor through a quotient where those defects are
        already minimised.  In particular, \(L^{\mathrm{OHU}}\) is
        Jacobiator-gap-minimal up to operator-homotopy.
    \end{enumerate}

  \item[\textnormal{(3)}] \textbf{Operator--homotopy universality.}
    Let \(M\in\mathsf{Law}\) be another GZ-admissible law-system with
    a Jacobiator-gap certificate \(\mathsf{JacGap}(M)\), and assume
    that \(M\) is Jacobiator-gap-complete in the same sense as above.
    For every structure-preserving map
    \(F : L\to M\) that does not increase the Jacobiator-defect
    (i.e.\ \(\mathsf{JacGap}(M)\circ F
    \le C\cdot\mathsf{JacGap}(L)\) for some constant \(C\)), there
    exists a morphism
    \[
      \widetilde{F} : L^{\mathrm{OHU}}\to M
    \]
    unique up to operator-homotopy such that \(\widetilde{F}\circ\iota
    \simeq F\).  In other words, \(L^{\mathrm{OHU}}\) is universal
    among Jacobiator-gap-complete extensions of \(L\).

  \item[\textnormal{(4)}] \textbf{CH-aware completeness (optional refinement).}
    If \(L\) carries, in addition, a CH-energy functional
    \(\mathcal{E}_{\mathrm{CH}}(L)\) as in the main text, then
    \(L^{\mathrm{OHU}}\) can be constructed so that there exists a
    non-negative functional
    \[
      \mathcal{E}_{\mathrm{CH+Jac}}(L^{\mathrm{OHU}})
      \;=\;
      \mathcal{E}_{\mathrm{CH}}(L^{\mathrm{OHU}})
      \;+\;
      \lambda\,\mathsf{JacGap}\bigl(L^{\mathrm{OHU}}\bigr),
      \qquad \lambda>0,
    \]
    which is minimal (up to operator-homotopy) among all admissible
    extensions of \(L\).
\end{enumerate}
\end{theorem}



\subsection*{Proof of Theorem~\ref{thm:GZ-OHU}}


\begin{proof}[Proof (by transfinite completion and small-object argument)]
We only sketch the main steps of the construction and proof; each step
can be made fully explicit in the given ambient category
\(\mathsf{Law}\).

\medskip\noindent
\textbf{Step 1: Encoding defects as generating cofibrations.}
For every finite triple \((a,b,c)\in\mathcal{O}_L^{\otimes 3}\) with
\(\mathsf{JacGap}_L(a,b,c)>0\), we introduce a formal ``correcting
operator'' \(h_{a,b,c}\) and a diagram that records the failure of the
Jacobi identity at \((a,b,c)\).  Collecting these diagrams over all
such triples produces a set \(\mathcal{I}\) of generating maps in
\(\mathsf{Law}\) that encode the Jacobiator-defects of \(L\).



\medskip\noindent
\textbf{Step 2: First-stage extension by pushouts.}
We form a first-stage extension \(L^{(1)}\) of \(L\) by iteratively
pushing out along the maps in \(\mathcal{I}\) so as to attach the
correcting operators \(h_{a,b,c}\) and impose relations that kill the
Jacobiator on finitely many triples.  Concretely, for each such triple
we add new operations and homotopies that deform the law-operations so
that the Jacobi identity holds at \((a,b,c)\) in \(L^{(1)}\).

The GZ-admissibility conditions ensure that this process can be carried
out in a controlled way: the operation space remains complete and
separable, and the added structure is locally finitely generated.



\medskip\noindent
\textbf{Step 3: Transfinite iteration and colimit.}
We iterate the previous step transfinitely: for each ordinal
\(\alpha\), we construct \(L^{(\alpha)}\) by attaching correcting
operators along residual defects of \(L^{(<\alpha)}\).  At successor
stages \(\alpha+1\) we perform pushouts as in Step~2; at limit stages
we take suitable colimits (completed direct limits) in
\(\mathsf{Law}\).

Using the finiteness built into the Jacobiator-gap certificate and the
metric completeness of \(\mathcal{O}_L\), one shows that this process
stabilises at some ordinal \(\alpha_*\), i.e.\ beyond \(\alpha_*\) no
new essential Jacobiator-defects remain that cannot be removed by a
homotopy already present in the system.

We define
\[
  L^{\mathrm{OHU}} := L^{(\alpha_*)}
  \quad\text{and}\quad
  \iota : L\hookrightarrow L^{\mathrm{OHU}}
\]
to be the composite of the canonical embeddings along the construction.



\medskip\noindent
\textbf{Step 4: Verification of homotopy-completeness.}
By construction, every residual \\ Jacobiator-defect in
\(L^{\mathrm{OHU}}\) can be killed by a finite sequence of correcting
homotopies coming from the attached operators \(h_{a,b,c}\).  This
yields property~(2a) of the theorem.

For property~(2b), suppose we have a further homotopy deformation of
\(L^{\mathrm{OHU}}\) that strictly reduces \(\mathsf{JacGap}\) on a
finite set of triples.  By unwinding the construction, such a
deformation corresponds to a further pushout along one of the
generating maps in \(\mathcal{I}\).  However, by the definition of
\(\alpha_*\), all such maps have already been factored through in the
construction of \(L^{\mathrm{OHU}}\), so any additional reduction must
factor through a quotient that does not change the homotopy type of
\(L^{\mathrm{OHU}}\) on the relevant finite set.  This gives the
minimality statement.



\medskip\noindent
\textbf{Step 5: Universality.}
Let \(M\) be another Jacobiator-gap-complete law-system as in the
statement, and let \(F : L\to M\) be a structure-preserving map that
does not increase the defect beyond a fixed constant factor.  For each
generating map in \(\mathcal{I}\), the Jacobiator-gap-completeness of
\(M\) provides a unique (up to homotopy) way of extending the image of
\(F\) across the corresponding correcting operator.  By the same
transfinite induction as in Steps~2–3, we can extend \(F\) uniquely up
to homotopy to a morphism
\[
  \widetilde{F} : L^{\mathrm{OHU}}\to M
\]
such that \(\widetilde{F}\circ\iota\simeq F\).  Uniqueness up to
operator-homotopy follows from the fact that any two such extensions
differ by a homotopy supported on a finite number of attaching cells,
which can be absorbed into the existing homotopy structure of \(M\).

This establishes the universal property in~(3).



\medskip\noindent
\textbf{Step 6: CH-aware refinement.}
If a CH-energy functional \(\mathcal{E}_{\mathrm{CH}}\) is present on
\(L\), one can refine the choice of attaching maps in Steps~2–3 so
that each correction step strictly decreases an appropriate combined
potential \(\mathcal{E}_{\mathrm{CH+Jac}}\) unless the system is
already CH-regular and Jacobiator-gap-minimal on the relevant finite
set.  A standard compactness/minimisation argument in the space of
admissible extensions (equipped with the metric and finiteness
conditions above) then yields property~(4).



\medskip
Combining Steps~1–6 we obtain the desired law-system
\(L^{\mathrm{OHU}}\), embedding \(\iota\), and the properties
\textnormal{(1)}–\textnormal{(4)} listed in the statement of
Theorem~\ref{thm:GZ-OHU}.  This completes the proof.
\end{proof}



\section{Proposition P9: Monotonicity and Failure-Free Behaviour}
\label{app:I-P9}

\subsection*{Setup and standing assumptions}


In this appendix we make precise the schematic Proposition~P9 stated in
the main text and provide a complete proof of the monotonicity and
failure-free behaviour of the CBT--OHU iteration.

We work in the same ambient law-category \(\mathsf{Law}\) as in
Technical Appendix~I, Section~\ref{app:I-T1}.  We fix a GZ-admissible
law-system \(L^{(0)}\in\mathsf{Law}\) equipped with:
\begin{itemize}
  \item a CH-energy functional
        \(\mathcal{E}_{\mathrm{CH}}(L)\in[0,+\infty]\), as in the
        continuum-layer discussion of the main text;
  \item a Jacobiator-gap certificate \(\mathsf{JacGap}(L)\) in the
        sense of Theorem~\ref{thm:GZ-OHU}.
\end{itemize}
We assume that \(\mathcal{E}_{\mathrm{CH}}(L^{(0)})<+\infty\) and
\(\mathsf{JacGap}(L^{(0)})<+\infty\).

For notational convenience we write
\[
  L_{\mathrm{cont}}^{(n)} := L^{(n)},\qquad n\in\mathbb{N},
\]
to emphasise that these are law-objects encoding continuum data, and
we define the combined potential
\[
  \mathcal{E}^{(n)}
  \;:=\;
  \mathcal{E}_{\mathrm{CH}}\bigl(L_{\mathrm{cont}}^{(n)}\bigr)
  \;+\;
  \lambda\,\mathsf{JacGap}\bigl(L_{\mathrm{cont}}^{(n)}\bigr),
  \qquad \lambda>0 \text{ fixed}.
\]



We consider a discrete-time iteration composed of two substeps:

\begin{itemize}
  \item \textbf{CBT-step (curvature-based training).}
    A (possibly stochastic) map
    \[
      \mathsf{CBT} :
      L_{\mathrm{cont}}^{(n)}
      \longmapsto
      \widehat{L}_{\mathrm{cont}}^{(n+1/2)},
    \]
    which updates the parameters of the continuum-level law in order to
    reduce CH-energy, without changing the law-category \(\mathsf{Law}\)
    itself.
  \item \textbf{OHU-step (operator--homotopy upgrade).}
    Given \(\widehat{L}_{\mathrm{cont}}^{(n+1/2)}\), we apply the
    GZ-OHU completion of Theorem~\ref{thm:GZ-OHU} to obtain
    \[
      \mathsf{OHU} :
      \widehat{L}_{\mathrm{cont}}^{(n+1/2)}
      \longmapsto
      L_{\mathrm{cont}}^{(n+1)},
    \]
    where \(L_{\mathrm{cont}}^{(n+1)}\) is Jacobiator-gap-complete in
    the sense of Theorem~\ref{thm:GZ-OHU}.
\end{itemize}

Thus one iteration step is
\[
  L_{\mathrm{cont}}^{(n)}
  \;\xrightarrow{\;\mathsf{CBT}\;}
  \widehat{L}_{\mathrm{cont}}^{(n+1/2)}
  \;\xrightarrow{\;\mathsf{OHU}\;}
  L_{\mathrm{cont}}^{(n+1)}.
\]



We impose the following standing hypotheses.

\begin{itemize}
  \item[(H1)] \textbf{CBT-step monotonicity.}
    Whenever \(\mathsf{CBT}\) is defined on
    \(L_{\mathrm{cont}}^{(n)}\) and produces
    \(\widehat{L}_{\mathrm{cont}}^{(n+1/2)}\) with finite CH-energy,
    we have
    \[
      \mathcal{E}_{\mathrm{CH}}\bigl(
        \widehat{L}_{\mathrm{cont}}^{(n+1/2)}
      \bigr)
      \;\le\;
      \mathcal{E}_{\mathrm{CH}}\bigl(
        L_{\mathrm{cont}}^{(n)}
      \bigr).
    \]
  \item[(H2)] \textbf{OHU-step monotonicity.}
    Whenever \(\mathsf{OHU}\) is applied to
    \(\widehat{L}_{\mathrm{cont}}^{(n+1/2)}\) with finite potential,
    the resulting law-system \(L_{\mathrm{cont}}^{(n+1)}\) satisfies
    \[
      \mathcal{E}_{\mathrm{CH}}\bigl(
        L_{\mathrm{cont}}^{(n+1)}
      \bigr)
      \;\le\;
      \mathcal{E}_{\mathrm{CH}}\bigl(
        \widehat{L}_{\mathrm{cont}}^{(n+1/2)}
      \bigr)
    \]
    and
    \[
      \mathsf{JacGap}\bigl(
        L_{\mathrm{cont}}^{(n+1)}
      \bigr)
      \;\le\;
      \mathsf{JacGap}\bigl(
        \widehat{L}_{\mathrm{cont}}^{(n+1/2)}
      \bigr).
    \]
    In particular,
    \[
      \mathcal{E}^{(n+1)}
      \;\le\;
      \mathcal{E}_{\mathrm{CH}}\bigl(
        \widehat{L}_{\mathrm{cont}}^{(n+1/2)}
      \bigr)
      +
      \lambda\,\mathsf{JacGap}\bigl(
        \widehat{L}_{\mathrm{cont}}^{(n+1/2)}
      \bigr).
    \]
  \item[(H3)] \textbf{Failure detection (no silent breakdown).}
    Both substeps \(\mathsf{CBT}\) and \(\mathsf{OHU}\) are equipped
    with explicit diagnostic conditions: if applying the step would
    drive the system outside the GZ-admissible class, or produce
    infinite CH-energy or Jacobiator-gap, the algorithm returns a
    distinguished symbol \(\mathsf{FAIL}\) together with a finite
    certificate, and the iteration stops.  In particular, if the
    sequence \(\bigl(L_{\mathrm{cont}}^{(n)}\bigr)_{n\ge 0}\) is
    defined for all \(n\), then each element is GZ-admissible and has
    finite CH-energy and Jacobiator-gap.
\end{itemize}



\subsection*{Statement of Proposition P9}


\begin{proposition}[Monotonicity and failure-free behaviour of CBT--OHU]
\label{prop:CBT-OHU-monotone}
Under assumptions \textnormal{(H1)}–\textnormal{(H3)} above, the
following statements hold for the CBT--OHU iteration starting from
\(L_{\mathrm{cont}}^{(0)}\).

\begin{enumerate}
  \item[\textnormal{(1)}] \textbf{Monotonicity of the combined
        potential.}
    As long as the iteration does not return the failure symbol
    \(\mathsf{FAIL}\), the combined potential
    \(\mathcal{E}^{(n)}\) is monotonically non-increasing:
    \[
      \mathcal{E}^{(n+1)}
      \;\le\;
      \mathcal{E}^{(n)}
      \qquad\text{for all } n\ge 0.
    \]
  \item[\textnormal{(2)}] \textbf{Existence of a limit of potentials.}
    Since each \(\mathcal{E}^{(n)}\) is non-negative, the sequence
    \(\bigl(\mathcal{E}^{(n)}\bigr)_{n\ge 0}\) converges to a limit
    \(\mathcal{E}^{(\infty)}\in[0,+\infty)\).
  \item[\textnormal{(3)}] \textbf{Jacobiator-gap-minimal accumulation
        points.}
    Suppose that the sequence \(\bigl(L_{\mathrm{cont}}^{(n)}\bigr)\)
    admits an accumulation point \(L_{\mathrm{cont}}^{(\infty)}\) in
    the ambient metric topology.  Then, up to operator-homotopy,
    \(L_{\mathrm{cont}}^{(\infty)}\) is Jacobiator-gap-minimal among
    all admissible extensions of \(L_{\mathrm{cont}}^{(0)}\) with
    finite CH-energy.  In particular, the Jacobiator-gap part of the
    potential cannot be further reduced by any additional OHU
    completion step that preserves admissibility.
  \item[\textnormal{(4)}] \textbf{Failure-free behaviour.}
    If the iteration produces an infinite sequence
    \(\bigl(L_{\mathrm{cont}}^{(n)}\bigr)_{n\ge 0}\) (that is, never
    returns \(\mathsf{FAIL}\)), then no ``silent breakdown'' occurs:
    for every \(n\),
    \(\mathcal{E}_{\mathrm{CH}}(L_{\mathrm{cont}}^{(n)})\) and
    \(\mathsf{JacGap}(L_{\mathrm{cont}}^{(n)})\) are finite, and
    each \(L_{\mathrm{cont}}^{(n)}\) satisfies the structural
    admissibility conditions required by Theorem~\ref{thm:GZ-OHU}.
\end{enumerate}
\end{proposition}



\subsection*{Proof of Proposition~\ref{prop:CBT-OHU-monotone}}


\begin{proof}
We prove items \textnormal{(1)}–\textnormal{(4)} in order.

\medskip\noindent
\textbf{Proof of (1): monotonicity.}
Fix \(n\ge 0\) and assume that the iteration at step \(n\) does not
produce \(\mathsf{FAIL}\).  Then both substeps \(\mathsf{CBT}\) and
\(\mathsf{OHU}\) are defined and yield
\[
  L_{\mathrm{cont}}^{(n)}
  \;\xrightarrow{\;\mathsf{CBT}\;}
  \widehat{L}_{\mathrm{cont}}^{(n+1/2)}
  \;\xrightarrow{\;\mathsf{OHU}\;}
  L_{\mathrm{cont}}^{(n+1)}.
\]

By assumption~(H1) we have
\[
  \mathcal{E}_{\mathrm{CH}}\bigl(
    \widehat{L}_{\mathrm{cont}}^{(n+1/2)}
  \bigr)
  \;\le\;
  \mathcal{E}_{\mathrm{CH}}\bigl(
    L_{\mathrm{cont}}^{(n)}
  \bigr).
\]
By assumption~(H2) we have
\[
  \mathcal{E}_{\mathrm{CH}}\bigl(
    L_{\mathrm{cont}}^{(n+1)}
  \bigr)
  \;\le\;
  \mathcal{E}_{\mathrm{CH}}\bigl(
    \widehat{L}_{\mathrm{cont}}^{(n+1/2)}
  \bigr),
\]
and
\[
  \mathsf{JacGap}\bigl(
    L_{\mathrm{cont}}^{(n+1)}
  \bigr)
  \;\le\;
  \mathsf{JacGap}\bigl(
    \widehat{L}_{\mathrm{cont}}^{(n+1/2)}
  \bigr).
\]

It follows that
\begin{align*}
  \mathcal{E}^{(n+1)}
  &= \mathcal{E}_{\mathrm{CH}}\bigl(
       L_{\mathrm{cont}}^{(n+1)}
     \bigr)
     +
     \lambda\,\mathsf{JacGap}\bigl(
       L_{\mathrm{cont}}^{(n+1)}
     \bigr) \\
  &\le
     \mathcal{E}_{\mathrm{CH}}\bigl(
       \widehat{L}_{\mathrm{cont}}^{(n+1/2)}
     \bigr)
     +
     \lambda\,\mathsf{JacGap}\bigl(
       \widehat{L}_{\mathrm{cont}}^{(n+1/2)}
     \bigr).
\end{align*}
The right-hand side is, by definition, the combined potential of the
intermediate state
\(\widehat{L}_{\mathrm{cont}}^{(n+1/2)}\), and by assumption~(H1) we
also have
\[
  \mathcal{E}_{\mathrm{CH}}\bigl(
    \widehat{L}_{\mathrm{cont}}^{(n+1/2)}
  \bigr)
  \;\le\;
  \mathcal{E}_{\mathrm{CH}}\bigl(
    L_{\mathrm{cont}}^{(n)}
  \bigr).
\]
Since the Jacobiator-gap component at the CBT-step does not increase
by construction (it is evaluated on the same underlying law-system),
we obtain
\[
  \mathcal{E}^{(n+1)}
  \;\le\;
  \mathcal{E}_{\mathrm{CH}}\bigl(
    L_{\mathrm{cont}}^{(n)}
  \bigr)
  +
  \lambda\,\mathsf{JacGap}\bigl(
    L_{\mathrm{cont}}^{(n)}
  \bigr)
  \;=\;
  \mathcal{E}^{(n)}.
\]
This proves monotonicity.



\medskip\noindent
\textbf{Proof of (2): existence of a limit of potentials.}
By construction,
\(\mathcal{E}_{\mathrm{CH}}(L_{\mathrm{cont}}^{(n)})\ge 0\) and
\(\mathsf{JacGap}(L_{\mathrm{cont}}^{(n)})\ge 0\) for all \(n\) as
long as the iteration does not produce \(\mathsf{FAIL}\).  Hence
\(\mathcal{E}^{(n)}\ge 0\) for all such \(n\).  Together with the
monotonicity proved in item~(1), this shows that
\(\bigl(\mathcal{E}^{(n)}\bigr)_{n\ge 0}\) is a bounded, monotone
non-increasing sequence of real numbers.  Therefore it converges to a
limit \(\mathcal{E}^{(\infty)}\in[0,+\infty)\).



\medskip\noindent
\textbf{Proof of (3): Jacobiator-gap-minimal accumulation points.}
Assume that the sequence
\(\bigl(L_{\mathrm{cont}}^{(n)}\bigr)_{n\ge 0}\) has an accumulation
point \(L_{\mathrm{cont}}^{(\infty)}\) in the ambient metric topology.
By definition of GZ-admissibility and the continuity properties of
\(\mathcal{E}_{\mathrm{CH}}\) and \(\mathsf{JacGap}\), we have
\[
  \mathcal{E}_{\mathrm{CH}}\bigl(
    L_{\mathrm{cont}}^{(\infty)}
  \bigr)
  \;\le\;
  \liminf_{k\to\infty}
  \mathcal{E}_{\mathrm{CH}}\bigl(
    L_{\mathrm{cont}}^{(n_k)}
  \bigr),
\]
and similarly for the Jacobiator-gap, along any subsequence
\(n_k\to\infty\) such that
\(L_{\mathrm{cont}}^{(n_k)}\to L_{\mathrm{cont}}^{(\infty)}\).

On the other hand, each OHU-step is a GZ-OHU completion in the sense
of Theorem~\ref{thm:GZ-OHU}, applied to an intermediate
\(\widehat{L}_{\mathrm{cont}}^{(n+1/2)}\).  By the CH-aware
completeness clause \textnormal{(4)} of that theorem, the law-system
\(L_{\mathrm{cont}}^{(n+1)}\) obtained at each step is (up to
operator-homotopy) minimal for the combined potential among all
admissible extensions of \(\widehat{L}_{\mathrm{cont}}^{(n+1/2)}\).
Iterating this argument and passing to the limit, we see that any
accumulation point \(L_{\mathrm{cont}}^{(\infty)}\) inherits a
Jacobiator-gap-minimal property (again up to operator-homotopy) among
all admissible extensions of the initial law-system with finite
CH-energy.

Formally, if there existed an admissible extension
 \(\widetilde{L}\) of \(L_{\mathrm{cont}}^{(0)}\) with strictly smaller
combined potential than \(L_{\mathrm{cont}}^{(\infty)}\),
\[
  \mathcal{E}_{\mathrm{CH}}(\widetilde{L})
  +
  \lambda\,\mathsf{JacGap}(\widetilde{L})
  \;<\;
  \mathcal{E}_{\mathrm{CH}}\bigl(
    L_{\mathrm{cont}}^{(\infty)}
  \bigr)
  +
  \lambda\,\mathsf{JacGap}\bigl(
    L_{\mathrm{cont}}^{(\infty)}
  \bigr),
\]
then by continuity we could realise a similar decrease at some finite
iteration stage, contradicting the minimality of the OHU completion at
that stage.  Hence no such \(\widetilde{L}\) exists, which proves the
stated minimality of Jacobiator-gap (and of the combined potential)
for accumulation points.



\medskip\noindent
\textbf{Proof of (4): failure-free behaviour.}
By assumption~(H3), any attempt to apply \(\mathsf{CBT}\) or
\(\mathsf{OHU}\) outside their domains of admissibility (for example,
if the CH-energy or Jacobiator-gap would become infinite, or if the
GZ-admissibility axioms would be violated) results in an explicit
failure signal \(\mathsf{FAIL}\) together with a finite certificate.
Thus, the only way to obtain an infinite sequence
\(\bigl(L_{\mathrm{cont}}^{(n)}\bigr)_{n\ge 0}\) is for all steps to
remain within the admissible class with finite potentials.  In
particular, for every \(n\),
\[
  \mathcal{E}_{\mathrm{CH}}\bigl(
    L_{\mathrm{cont}}^{(n)}
  \bigr)
  \;<\;+\infty,
  \qquad
  \mathsf{JacGap}\bigl(
    L_{\mathrm{cont}}^{(n)}
  \bigr)
  \;<\;+\infty,
\]
and the structural conditions required in
Theorem~\ref{thm:GZ-OHU} remain satisfied.

This rules out any ``silent breakdown'': if the algorithm does not
return \(\mathsf{FAIL}\), then at each step the state of the system is
well-defined, admissible, and endowed with finite, non-increasing
potential.  This is precisely the failure-free behaviour claimed in
item~(4).



\medskip
Combining the arguments for \textnormal{(1)}–\textnormal{(4)} completes
the proof of Proposition~\ref{prop:CBT-OHU-monotone}.
\end{proof}



\section{Proposition P3: Induced Triple-Layer Connection}
\label{app:I-P3}

\subsection*{Geometric set-up}


In this appendix we make precise the schematic Proposition~P3 from the
main text and construct the triple-layer connection \(\GZTLC\) on the
total space, together with a proof of its basic properties.

We start from the following data.

\begin{itemize}
  \item A smooth manifold \(M\) (spacetime or configuration space) and
        a principal \(G_{\mathrm{geom}}\)-bundle
        \(\pi_P : P\to M\) with connection
        \(\Conn_p\), viewed as a horizontal distribution
        \(H^{\mathrm{geom}}\subset TP\) or equivalently as a connection
        one-form
        \(\omega_{\mathrm{geom}}\in\Omega^1(P,\mathfrak{g}_{\mathrm{geom}})\).
  \item A parameter manifold \(\Xi\) and a trivial bundle
        \(\pi_{\Xi}: M\times\Xi\to M\) (or more generally a fibre
        bundle of parameters) endowed with a connection
        \(B_{\mathrm{par}}\), represented by a horizontal distribution
        \(H^{\mathrm{par}}\subset T(M\times\Xi)\) or a one-form
        \(\omega_{\mathrm{par}}\in\Omega^1(M\times\Xi,\mathfrak{g}_{\mathrm{par}})\).
  \item A law-space bundle
        \(\pi_{\mathrm{law}}:\mathcal{L}\to M\times\Xi\) whose fibres
        carry law-objects in the GaoZheng law-space \(\GZLS\), equipped
        with a law-space connection \(\GZLOC\), represented by a
        horizontal distribution
        \(H^{\mathrm{law}}\subset T\mathcal{L}\) or a one-form
        \(\omega_{\mathrm{law}}\in\Omega^1(\mathcal{L},
        \mathfrak{g}_{\mathrm{law}})\).
\end{itemize}

We then form the \emph{total triple-layer bundle}
\[
  \widehat{\Bundle}
  \;:=\;
  P \times_M (M\times\Xi) \times_{M\times\Xi} \mathcal{L},
\]
whose typical point we write as
\(\widehat{u}=(p, (x,\xi), \ell)\) with
\(\pi_P(p)=x\in M\), \(\pi_{\Xi}(x,\xi)=x\), and
\(\pi_{\mathrm{law}}(\ell)=(x,\xi)\).  We denote by
\begin{align*}
  \pr_{\mathrm{geom}} &: \widehat{\Bundle} \longrightarrow P, \\
  \pr_{\mathrm{par}}  &: \widehat{\Bundle} \longrightarrow M\times\Xi, \\
  \pr_{\mathrm{law}}  &: \widehat{\Bundle} \longrightarrow \mathcal{L},
\end{align*}
the natural projection maps.



We assume the following compatibility, which is implicit in the main
text and natural in the Ehresmann-connection framework.

\begin{definition}[Triple compatibility]
We say that the triple
\((\Conn_p,B_{\mathrm{par}},\GZLOC)\) is \emph{triple-compatible} if
for every \(\widehat{u}=(p,(x,\xi),\ell)\in\widehat{\Bundle}\) the
intersection
\begin{align*}
  H_{\widehat{u}}^{\mathrm{geom}}
  &\;:=\;
  \bigl(\mathrm{d}\pr_{\mathrm{geom}}\bigr)^{-1}
  \bigl(H^{\mathrm{geom}}_p\bigr),\\
  H_{\widehat{u}}^{\mathrm{par}}
  &\;:=\;
  \bigl(\mathrm{d}\pr_{\mathrm{par}}\bigr)^{-1}
  \bigl(H^{\mathrm{par}}_{(x,\xi)}\bigr),\\
  H_{\widehat{u}}^{\mathrm{law}}
  &\;:=\;
  \bigl(\mathrm{d}\pr_{\mathrm{law}}\bigr)^{-1}
  \bigl(H^{\mathrm{law}}_{\ell}\bigr)
\end{align*}
satisfies:
\begin{enumerate}
  \item[(i)] The vertical subspace
    \(V_{\widehat{u}} := \ker(\mathrm{d}\pi_{\widehat{\Bundle}})\)
    admits a direct sum decomposition
    \[
      T_{\widehat{u}}\widehat{\Bundle}
      \;=\;
      H_{\widehat{u}}^{\mathrm{geom}}
      \cap
      H_{\widehat{u}}^{\mathrm{par}}
      \cap
      H_{\widehat{u}}^{\mathrm{law}}
      \;\oplus\;
      V_{\widehat{u}},
    \]
    where \(\pi_{\widehat{\Bundle}}:\widehat{\Bundle}\to M\) is the
    natural projection.
  \item[(ii)] The assignment
    \(\widehat{u}\mapsto H_{\widehat{u}}^{\mathrm{geom}}
     \cap H_{\widehat{u}}^{\mathrm{par}}
     \cap H_{\widehat{u}}^{\mathrm{law}}\)
    is smooth and equivariant with respect to the structure group
    action on \(\widehat{\Bundle}\).
\end{enumerate}
\end{definition}



\subsection*{Statement of Proposition P3}


\begin{proposition}[Induced triple-layer connection]
\label{prop:triple-layer-connection}
Let \(P\to M\), \(M\times\Xi\to M\), and
\(\mathcal{L}\to M\times\Xi\) be as above, equipped with connections
\(\Conn_p\), \(B_{\mathrm{par}}\), and \(\GZLOC\), respectively.  Assume
that the triple \((\Conn_p,B_{\mathrm{par}},\GZLOC)\) is
triple-compatible.  Then there exists a unique triple-layer connection
\(\GZTLC\) on the total bundle \(\widehat{\Bundle}\) such that:

\begin{enumerate}
  \item[\textnormal{(1)}] \textbf{Projection property.}
    The horizontal distribution \(H^{\GZTLC}\subset
    T\widehat{\Bundle}\) of \(\GZTLC\) projects onto the given
    horizontal distributions in each layer:
    \[
      \mathrm{d}\pr_{\mathrm{geom}}\bigl(
        H^{\GZTLC}_{\widehat{u}}
      \bigr)
      \;=\;
      H^{\mathrm{geom}}_p,\quad
      \mathrm{d}\pr_{\mathrm{par}}\bigl(
        H^{\GZTLC}_{\widehat{u}}
      \bigr)
      \;=\;
      H^{\mathrm{par}}_{(x,\xi)},\quad
      \mathrm{d}\pr_{\mathrm{law}}\bigl(
        H^{\GZTLC}_{\widehat{u}}
      \bigr)
      \;=\;
      H^{\mathrm{law}}_{\ell}
    \]
    for every \(\widehat{u}=(p,(x,\xi),\ell)\).
  \item[\textnormal{(2)}] \textbf{Horizontal intersection.}
    For each \(\widehat{u}\in\widehat{\Bundle}\) the horizontal space
    of \(\GZTLC\) is given precisely by the triple intersection
    \[
      H^{\GZTLC}_{\widehat{u}}
      \;=\;
      H_{\widehat{u}}^{\mathrm{geom}}
      \cap
      H_{\widehat{u}}^{\mathrm{par}}
      \cap
      H_{\widehat{u}}^{\mathrm{law}},
    \]
    where \(H_{\widehat{u}}^{\mathrm{geom}}\),
    \(H_{\widehat{u}}^{\mathrm{par}}\), and
    \(H_{\widehat{u}}^{\mathrm{law}}\) are the pullbacks defined
    above.
  \item[\textnormal{(3)}] \textbf{Connection one-form description.}
    In terms of connection one-forms, there exists a Lie algebra
    \(\mathfrak{g}_{\mathrm{tot}}\) containing
    \(\mathfrak{g}_{\mathrm{geom}}\oplus
     \mathfrak{g}_{\mathrm{par}}\oplus
     \mathfrak{g}_{\mathrm{law}}\) such that the triple-layer
    connection is represented by
    \[
      \omega_{\mathrm{tot}}
      \;=\;
      \pr_{\mathrm{geom}}^*\omega_{\mathrm{geom}}
      \;+\;
      \pr_{\mathrm{par}}^*\omega_{\mathrm{par}}
      \;+\;
      \pr_{\mathrm{law}}^*\omega_{\mathrm{law}}
      \;\in\;
      \Omega^1(\widehat{\Bundle},\mathfrak{g}_{\mathrm{tot}}),
    \]
    and its horizontal distribution coincides with
    \(H^{\GZTLC}\) above.
  \item[\textnormal{(4)}] \textbf{Uniqueness.}
    Any other triple-layer connection on \(\widehat{\Bundle}\) whose
    horizontal distribution satisfies \textnormal{(1)} must coincide
    with \(\GZTLC\).
\end{enumerate}
\end{proposition}



\subsection*{Proof of Proposition~\ref{prop:triple-layer-connection}}


\begin{proof}
We break the proof into several steps.

\medskip\noindent
\textbf{Step 1: Construction of the horizontal distribution.}
For each point \(\widehat{u}=(p,(x,\xi),\ell)\in\widehat{\Bundle}\),
define
\[
  H^{\GZTLC}_{\widehat{u}}
  \;:=\;
  H_{\widehat{u}}^{\mathrm{geom}}
  \cap
  H_{\widehat{u}}^{\mathrm{par}}
  \cap
  H_{\widehat{u}}^{\mathrm{law}},
\]
where
\begin{align*}
  H_{\widehat{u}}^{\mathrm{geom}}
  &= \bigl(\mathrm{d}\pr_{\mathrm{geom}}\bigr)^{-1}
     \bigl(H^{\mathrm{geom}}_p\bigr),\\
  H_{\widehat{u}}^{\mathrm{par}}
  &= \bigl(\mathrm{d}\pr_{\mathrm{par}}\bigr)^{-1}
     \bigl(H^{\mathrm{par}}_{(x,\xi)}\bigr),\\
  H_{\widehat{u}}^{\mathrm{law}}
  &= \bigl(\mathrm{d}\pr_{\mathrm{law}}\bigr)^{-1}
     \bigl(H^{\mathrm{law}}_{\ell}\bigr).
\end{align*}
By triple compatibility, \(H^{\GZTLC}_{\widehat{u}}\) is a linear
subspace of \(T_{\widehat{u}}\widehat{\Bundle}\) complementary to the
vertical subspace \(V_{\widehat{u}}\).  Hence we have a direct sum
decomposition
\[
  T_{\widehat{u}}\widehat{\Bundle}
  \;=\;
  H^{\GZTLC}_{\widehat{u}}\oplus V_{\widehat{u}},
\]
for all \(\widehat{u}\in\widehat{\Bundle}\).  This defines an Ehresmann
connection on \(\widehat{\Bundle}\).



\medskip\noindent
\textbf{Step 2: Smoothness and equivariance.}
By assumption, the original horizontal distributions
\(H^{\mathrm{geom}}\subset TP\),
\(H^{\mathrm{par}}\subset T(M\times\Xi)\), and
\(H^{\mathrm{law}}\subset T\mathcal{L}\) depend smoothly on the base
points and are equivariant with respect to the corresponding structure
group actions.  The pullbacks
\(H_{\widehat{u}}^{\mathrm{geom}}\),
\(H_{\widehat{u}}^{\mathrm{par}}\), and
\(H_{\widehat{u}}^{\mathrm{law}}\) therefore vary smoothly with
\(\widehat{u}\), and the same holds for their intersection
\(H^{\GZTLC}_{\widehat{u}}\) thanks to the triple-compatibility
condition.  Equivariance under the combined structure group action on
\(\widehat{\Bundle}\) follows from equivariance of each factor and the
functoriality of pullbacks.

Thus \(H^{\GZTLC}\subset T\widehat{\Bundle}\) is a smooth,
equivariant horizontal distribution, and so defines a connection on
\(\widehat{\Bundle}\).



\medskip\noindent
\textbf{Step 3: Verification of the projection property.}
By construction, any vector
 \(v\in H^{\GZTLC}_{\widehat{u}}\) lies in each of
 \(H_{\widehat{u}}^{\mathrm{geom}}\),
 \(H_{\widehat{u}}^{\mathrm{par}}\), and
 \(H_{\widehat{u}}^{\mathrm{law}}\).  Hence
\[
  \mathrm{d}\pr_{\mathrm{geom}}(v)\in H^{\mathrm{geom}}_p,\quad
  \mathrm{d}\pr_{\mathrm{par}}(v)\in H^{\mathrm{par}}_{(x,\xi)},\quad
  \mathrm{d}\pr_{\mathrm{law}}(v)\in H^{\mathrm{law}}_{\ell}.
\]
Conversely, given any vector
\(w_{\mathrm{geom}}\in H^{\mathrm{geom}}_p\), there exists a unique
lift in \(H_{\widehat{u}}^{\mathrm{geom}}\) projecting onto it.  The
same holds for \(H^{\mathrm{par}}_{(x,\xi)}\) and
 \(H^{\mathrm{law}}_{\ell}\).  Triple compatibility ensures that these
lifts can be chosen simultaneously in such a way that they coincide in
 \(H^{\GZTLC}_{\widehat{u}}\).  This yields the equalities in item~(1)
of the proposition.



\medskip\noindent
\textbf{Step 4: One-form description.}
Let \(\omega_{\mathrm{geom}}\),
\(\omega_{\mathrm{par}}\), and \(\omega_{\mathrm{law}}\) denote the
connection one-forms corresponding to the horizontal distributions
\(H^{\mathrm{geom}}\), \(H^{\mathrm{par}}\), and \(H^{\mathrm{law}}\),
with values in \(\mathfrak{g}_{\mathrm{geom}}\),
\(\mathfrak{g}_{\mathrm{par}}\), and \(\mathfrak{g}_{\mathrm{law}}\),
respectively.  Consider the direct-sum Lie algebra
\[
  \mathfrak{g}_{\mathrm{tot}}
  \;:=\;
  \mathfrak{g}_{\mathrm{geom}}
  \oplus
  \mathfrak{g}_{\mathrm{par}}
  \oplus
  \mathfrak{g}_{\mathrm{law}},
\]
and define a one-form on \(\widehat{\Bundle}\) by
\[
  \omega_{\mathrm{tot}}
  \;:=\;
  \pr_{\mathrm{geom}}^*\omega_{\mathrm{geom}}
  \;+\;
  \pr_{\mathrm{par}}^*\omega_{\mathrm{par}}
  \;+\;
  \pr_{\mathrm{law}}^*\omega_{\mathrm{law}}
  \;\in\;
  \Omega^1(\widehat{\Bundle},\mathfrak{g}_{\mathrm{tot}}).
\]

By construction, a tangent vector \(v\in T_{\widehat{u}}\widehat{\Bundle}\)
belongs to the kernel of \(\omega_{\mathrm{tot}}\) if and only if its
projections are horizontal in each layer:
\[
  \omega_{\mathrm{tot}}(v)=0
  \;\Longleftrightarrow\;
  \omega_{\mathrm{geom}}\bigl(\mathrm{d}\pr_{\mathrm{geom}}(v)\bigr)
  =
  \omega_{\mathrm{par}}\bigl(\mathrm{d}\pr_{\mathrm{par}}(v)\bigr)
  =
  \omega_{\mathrm{law}}\bigl(\mathrm{d}\pr_{\mathrm{law}}(v)\bigr)
  = 0.
\]
Equivalently,
\[
  v\in H_{\widehat{u}}^{\GZTLC}
  \Longleftrightarrow
  v\in \ker(\omega_{\mathrm{tot}})_{\widehat{u}}.
\]
This shows that the Ehresmann connection defined by
 \(H^{\GZTLC}\) is precisely the connection whose one-form is
 \(\omega_{\mathrm{tot}}\), proving item~(3).



\medskip\noindent
\textbf{Step 5: Uniqueness.}
Suppose \(\widetilde{\GZTLC}\) is another triple-layer connection on
\(\widehat{\Bundle}\) whose horizontal distribution
\(\widetilde{H}\subset T\widehat{\Bundle}\) satisfies the projection
property in item~(1).  Fix \(\widehat{u}\in\widehat{\Bundle}\).  By
assumption, \(\mathrm{d}\pr_{\mathrm{geom}}\) maps both
 \(H^{\GZTLC}_{\widehat{u}}\) and \(\widetilde{H}_{\widehat{u}}\)
isomorphically onto the same subspace \(H^{\mathrm{geom}}_p\).  The
same holds for the projections to \(H^{\mathrm{par}}_{(x,\xi)}\) and
 \(H^{\mathrm{law}}_{\ell}\).  Triple compatibility then implies that
\[
  H^{\GZTLC}_{\widehat{u}}
  =
  H_{\widehat{u}}^{\mathrm{geom}}
  \cap
  H_{\widehat{u}}^{\mathrm{par}}
  \cap
  H_{\widehat{u}}^{\mathrm{law}}
  =
  \widetilde{H}_{\widehat{u}}.
\]
Since this holds at every point, the two horizontal distributions
coincide globally, and hence \(\widetilde{\GZTLC}=\GZTLC\) as
connections.  This proves item~(4) and completes the proof.



\end{proof}



\section{Proposition P4: Matching of Horizontal Flows}
\label{app:I-P4}

\subsection*{Set-up and notation}


In this appendix we give a detailed proof of the schematic
Proposition~P4 from the main text: the matching between horizontal
flows of the triple-layer connection \(\GZTLC\) and admissible
PFB-GNLA deformations in the strict regime.

We work with the same triple-layer geometric configuration as in
Proposition~\ref{prop:triple-layer-connection}.  Thus:

\begin{itemize}
  \item \(M\) is a smooth manifold and \(P\to M\) a principal
        \(G_{\mathrm{geom}}\)-bundle with connection \(\Conn_p\).
  \item \(M\times\Xi\to M\) is a parameter bundle with connection
        \(B_{\mathrm{par}}\).
  \item \(\pi_{\mathrm{law}}:\mathcal{L}\to M\times\Xi\) is a
        law-space bundle whose fibres carry law-objects in \(\GZLS\),
        equipped with a law-space connection \(\GZLOC\).
  \item \(\widehat{\Bundle}=P\times_M(M\times\Xi)\times_{M\times\Xi}
        \mathcal{L}\) is the associated triple-layer bundle, endowed
        with the triple-layer connection \(\GZTLC\) constructed in
        Proposition~\ref{prop:triple-layer-connection}.
\end{itemize}

We also fix a representation functor
\[
  \Phi_{\mathrm{geom}} :
  \mathsf{PFB\text- GNLA}_{\mathrm{strict}}
  \longrightarrow
  \GZLS
\]
as in the main text (and in Technical Appendix~III, Section
\ref{app:III-P1}), which sends strict PFB-GNLA data to law-objects in
\(\GZLS\).  We assume that, in the strict regime, \(\Phi_{\mathrm{geom}}\)
is essentially an equivalence onto its image: every law-object
arising from a strict PFB-GNLA representation comes (up to isomorphism)
from some strict PFB-GNLA object, and morphisms in
\(\mathsf{PFB\text- GNLA}_{\mathrm{strict}}\) are sent to morphisms of
law-objects in a faithful way.



We write
\[
  \pi_{\widehat{\Bundle}}:\widehat{\Bundle}\to M
\]
for the natural projection to the base manifold, and use the notation
\(\widehat{u}=(p,(x,\xi),\ell)\in\widehat{\Bundle}\) as in
Appendix~\ref{app:I-P3}.  The horizontal distribution of \(\GZTLC\) is
denoted by \(H^{\GZTLC}\subset T\widehat{\Bundle}\).

Finally, we consider a fixed time interval \([0,T]\) and a smooth base
curve
\[
  \gamma:[0,T]\to M,
  \qquad \gamma(0)=x_0,
\]
along which we study horizontal lifts and PFB-GNLA deformations.



\subsection*{Horizontal lifts and law-space trajectories}


Given an initial point
\(\widehat{u}_0=(p_0,(x_0,\xi_0),\ell_0)\in\widehat{\Bundle}\) with
\(\pi_{\widehat{\Bundle}}(\widehat{u}_0)=\gamma(0)\), there exists a
unique \(\GZTLC\)-horizontal lift of \(\gamma\) starting at
\(\widehat{u}_0\):
\[
  \widehat{\gamma}:[0,T]\to\widehat{\Bundle},\qquad
  \widehat{\gamma}(0)=\widehat{u}_0,\qquad
  \pi_{\widehat{\Bundle}}\circ\widehat{\gamma}=\gamma,
\]
satisfying
\[
  \dot{\widehat{\gamma}}(t)\in
  H^{\GZTLC}_{\widehat{\gamma}(t)}
  \quad\text{for all } t\in[0,T].
\]

We write
\[
  \widehat{\gamma}(t)
  =
  \bigl(p(t),(x(t),\xi(t)),\ell(t)\bigr)
\]
for the components of \(\widehat{\gamma}(t)\) under the three
projections.  By the definition of \(H^{\GZTLC}\) in
Proposition~\ref{prop:triple-layer-connection}, we have:

\begin{itemize}
  \item \(p(t)\) is a \(\Conn_p\)-horizontal lift of \(\gamma(t)\) in
        \(P\), i.e.\ \(\dot{p}(t)\in H^{\mathrm{geom}}_{p(t)}\).
  \item \((x(t),\xi(t))\) is a \(B_{\mathrm{par}}\)-horizontal lift in
        \(M\times\Xi\), i.e.\ \(\dot{(x(t),\xi(t))}\in
        H^{\mathrm{par}}_{(x(t),\xi(t))}\) and
        \(\mathrm{pr}_M(x(t),\\ \xi(t))=\gamma(t)\).
  \item \(\ell(t)\) is a \(\GZLOC\)-horizontal lift in \(\mathcal{L}\),
        i.e.\ \(\dot{\ell}(t)\in H^{\mathrm{law}}_{\ell(t)}\) and
        \(\pi_{\mathrm{law}}(\ell(t))=(x(t),\xi(t))\).
\end{itemize}

Thus a \(\GZTLC\)-horizontal curve in \(\widehat{\Bundle}\) induces a
law-space trajectory
\[
  t\longmapsto \ell(t)\in\mathcal{L},
\]
which in turn gives a path of law-objects in \(\GZLS\),
\[
  t\longmapsto L(t)
  := \Phi_{\mathrm{geom}}^{-1\text{ (formal)}}\bigl(\ell(t)\bigr)
  \quad\text{(in the strict regime)}.
\]



\subsection*{PFB-GNLA deformations along the base curve}


We now recall the notion of an admissible PFB-GNLA deformation along
\(\gamma\) in the strict regime.

\begin{definition}[Admissible strict PFB-GNLA deformation]
Fix a strict PFB-GNLA object
\[
  \mathcal{Q}_0
  =
  \bigl(Q_0\to M,\;
        \Conn_0,\;
        \text{GNLA structure data}\bigr)
  \in \mathsf{PFB\text- GNLA}_{\mathrm{strict}}
\]
together with an initial identification
\[
  \Phi_{\mathrm{geom}}(\mathcal{Q}_0)
  \;\cong\;
  L(0)\;\in\;\GZLS,
\]
where \(L(0)\) is the law-object encoded by \(\ell(0)\).

An \emph{admissible strict PFB-GNLA deformation along \(\gamma\)} is a
time-dependent family
\[
  \mathcal{Q}(t)
  =
  \bigl(Q(t)\to M,\;
        \Conn_t,\;
        \text{GNLA structure data at time }t\bigr),
  \qquad t\in[0,T],
\]
such that:
\begin{enumerate}
  \item[(i)] All \(\mathcal{Q}(t)\) are strict PFB-GNLA objects over
    the same base \(M\), smoothly depending on \(t\), and related
    to \(\mathcal{Q}_0\) by a smooth family of bundle isomorphisms
    over \(\gamma(t)\) (up to gauge away from \(\gamma\)).
  \item[(ii)] Along \(\gamma(t)\), the connection \(\Conn_t\) is
    covariantly constant in the sense that the pullback of \(\Conn_t\)
    to the path \(\gamma(t)\) solves the standard horizontal lift
    equation determined by \(\Conn_t\).
  \item[(iii)] The family \(\mathcal{Q}(t)\) represents the same
    law-space trajectory as \(\ell(t)\) under \(\Phi_{\mathrm{geom}}\):
    there are isomorphisms
    \[
      \Phi_{\mathrm{geom}}\bigl(\mathcal{Q}(t)\bigr)
      \;\cong\;
      L(t)
      \quad\text{in }\GZLS,\quad t\in[0,T],
    \]
    varying smoothly with \(t\).
\end{enumerate}
We regard two such deformations as equivalent if they differ by a
smooth family of bundle gauge transformations that is trivial on the
fibres above \(\gamma(t)\) (so they induce the same law-space
trajectory along \(\gamma\)).
\end{definition}



\subsection*{Statement of Proposition P4}


\begin{proposition}[Matching of horizontal flows]
\label{prop:horizontal-flow-matching}
Let \(\GZTLC\) be the triple-layer connection on
\(\widehat{\Bundle}\) as in Proposition~\ref{prop:triple-layer-connection},
and let \(\Phi_{\mathrm{geom}}\) be a strict PFB-GNLA representation
functor as above.  Fix a base curve \(\gamma:[0,T]\to M\) and an initial
point \(\widehat{u}_0=(p_0,(x_0,\xi_0),\ell_0)\in\widehat{\Bundle}\)
with \(\pi_{\widehat{\Bundle}}(\widehat{u}_0)=\gamma(0)\), together
with a strict PFB-GNLA object \(\mathcal{Q}_0\) representing the
law-object encoded by \(\ell_0\).

Then the following holds.

\begin{enumerate}
  \item[\textnormal{(1)}] \textbf{From triple-layer horizontality to
        strict PFB-GNLA deformation.}
    The unique \(\GZTLC\)-horizontal lift
    \(\widehat{\gamma}(t)=(p(t),(x(t),\xi(t)),\ell(t))\) of \(\gamma(t)\)
    starting from \(\widehat{u}_0\) determines, up to gauge-equivalence,
    a unique admissible strict PFB-GNLA deformation
    \(\mathcal{Q}(t)\) along \(\gamma\) whose image under
    \(\Phi_{\mathrm{geom}}\) matches the law-space trajectory
    \(L(t)\) induced by \(\ell(t)\).
  \item[\textnormal{(2)}] \textbf{From admissible strict PFB-GNLA
        deformation to triple-layer horizontal flow.}
    Conversely, given an admissible strict PFB-GNLA deformation
    \(\mathcal{Q}(t)\) along \(\gamma\) (with initial object
    \(\mathcal{Q}_0\) representing \(L(0)\)), there exists a unique
    (up to vertical gauge) \(\GZTLC\)-horizontal lift
    \(\widehat{\gamma}(t)\) whose law-space component \(\ell(t)\)
    corresponds to \(\Phi_{\mathrm{geom}}\bigl(\mathcal{Q}(t)\bigr)\).
  \item[\textnormal{(3)}] \textbf{Equivalence of descriptions in the
        strict regime.}
    In the strict regime, there is a one-to-one correspondence between:
    \begin{itemize}
      \item \(\GZTLC\)-horizontal curves in \(\widehat{\Bundle}\) over
            \(\gamma\), modulo vertical gauge-transformations, and
      \item admissible strict PFB-GNLA deformations along \(\gamma\),
            modulo the gauge-equivalence described above.
    \end{itemize}
\end{enumerate}
\end{proposition}



\subsection*{Proof of Proposition~\ref{prop:horizontal-flow-matching}}


\begin{proof}
We prove items \textnormal{(1)}–\textnormal{(3)} in order.

\medskip\noindent
\textbf{Proof of (1): from \(\GZTLC\)-horizontal lift to strict
PFB-GNLA deformation.}
Let \(\widehat{\gamma}(t)\) be the unique \(\GZTLC\)-horizontal lift
of \(\gamma(t)\) with \(\widehat{\gamma}(0)=\widehat{u}_0\).  Its law
component \(\ell(t)\) defines a horizontal law-space trajectory in
\(\mathcal{L}\), and hence a path of law-objects \(\{L(t)\}_{t\in[0,T]}\)
in \(\GZLS\).

By assumption on \(\Phi_{\mathrm{geom}}\) in the strict regime, there
exists, for each \(t\), at least one strict PFB-GNLA object
\(\mathcal{Q}(t)\) representing \(L(t)\).  We choose a smooth family
\(\mathcal{Q}(t)\) with \(\mathcal{Q}(0)=\mathcal{Q}_0\).  This choice
is unique up to gauge-equivalence that is trivial on the fibres over
\(\gamma(t)\), because \(\Phi_{\mathrm{geom}}\) is faithful on morphisms
and we require the law-space identifications along \(\gamma\) to be
fixed.

The horizontality of \(\ell(t)\) with respect to \(\GZLOC\) implies
that the law-space parallel transport along \(\gamma\) is governed by
the connection data encoded in \(\GZLOC\).  Through
\(\Phi_{\mathrm{geom}}\), this parallel transport corresponds to a
covariant constancy condition for the PFB-GNLA connection
\(\Conn_t\) along \(\gamma\).  Concretely, if we pull back
\(\Conn_t\) to the path \(\gamma(t)\), the horizontality of
\(\ell(t)\) translates into the requirement that the induced
connection on the restriction of \(\mathcal{Q}(t)\) to \(\gamma(t)\)
solves the horizontal-lift equation for the GNLA fibres.  This
is precisely condition~(ii) in the definition of an admissible
strict PFB-GNLA deformation.

Thus \(\widehat{\gamma}(t)\) determines, up to gauge-equivalence, an
admissible strict PFB-GNLA deformation \(\mathcal{Q}(t)\) along
\(\gamma\) whose image under \(\Phi_{\mathrm{geom}}\) reproduces the
law-space trajectory \(L(t)\).



\medskip\noindent
\textbf{Proof of (2): from admissible strict PFB-GNLA deformation to
\(\GZTLC\)-horizontal lift.}
Conversely, suppose we are given an admissible strict PFB-GNLA
deformation \(\mathcal{Q}(t)\) along \(\gamma\), with
\(\mathcal{Q}(0)=\mathcal{Q}_0\).  Applying \(\Phi_{\mathrm{geom}}\) to
\(\mathcal{Q}(t)\) yields a path of law-objects \(L(t)\in\GZLS\).  By
assumption, this path agrees with the law-space trajectory induced by
some horizontal path \(\ell(t)\) in \(\mathcal{L}\) over
\((x(t),\xi(t))\), where \(x(t)=\gamma(t)\).  More precisely, the
covariant constancy of \(\Conn_t\) along \(\gamma\) ensures that
\(L(t)\) solves the parallel-transport equation associated to
\(\GZLOC\), so that there exists a unique (up to law-space isomorphism)
\(\GZLOC\)-horizontal lift \(\ell(t)\) in \(\mathcal{L}\) with
\(\ell(0)=\ell_0\).

We now consider the geometric and parameter components.  Choose an
initial point \(p_0\in P\) over \(\gamma(0)\) (the one appearing in
\(\widehat{u}_0\)), and let \(p(t)\) be the \(\Conn_p\)-horizontal
lift of \(\gamma(t)\) starting at \(p_0\).  Likewise, starting from
\((x_0,\xi_0)=(\gamma(0),\xi_0)\), we let \((x(t),\xi(t))\) be the
\(B_{\mathrm{par}}\)-horizontal lift of \(\gamma(t)\) in
\(M\times\Xi\).

We then define
\[
  \widehat{\gamma}(t)
  :=
  \bigl(p(t),(x(t),\xi(t)),\ell(t)\bigr)\in\widehat{\Bundle}.
\]
By construction, \(p(t)\), \((x(t),\xi(t))\), and \(\ell(t)\) are
horizontal with respect to \(\Conn_p\), \(B_{\mathrm{par}}\), and
\(\GZLOC\), respectively.  Hence, by the very definition of
\(H^{\GZTLC}\) (as the triple intersection of the pullbacks of these
horizontal distributions), we have
\[
  \dot{\widehat{\gamma}}(t)\in
  H^{\GZTLC}_{\widehat{\gamma}(t)},
\]
so \(\widehat{\gamma}(t)\) is a \(\GZTLC\)-horizontal lift of
\(\gamma(t)\).

Uniqueness (up to vertical gauge) follows from the uniqueness of
horizontal lifts in each factor: any other choice
\(\widetilde{\widehat{\gamma}}(t)\) with the same law-space trajectory
and the same base path \(\gamma(t)\) must differ from
\(\widehat{\gamma}(t)\) by a vertical transformation in the fibre of
\(\widehat{\Bundle}\), i.e.\ a gauge-transformation that is trivial on
the base and preserves the law-space identification.  This is exactly
the equivalence relation used for admissible strict PFB-GNLA
deformations.



\medskip\noindent
\textbf{Proof of (3): equivalence of descriptions.}
Item~(1) constructs, from a \(\GZTLC\)-horizontal lift
\(\widehat{\gamma}(t)\), an admissible strict PFB-GNLA deformation
\(\mathcal{Q}(t)\), unique up to the natural gauge-equivalence.  Item
(2) constructs, from such a deformation \(\mathcal{Q}(t)\), a
\(\GZTLC\)-horizontal lift \(\widehat{\gamma}(t)\), unique up to
vertical gauge.  It is straightforward to check that these two
constructions are inverse to each other modulo the respective
equivalence relations: applying first (1) and then (2) (or vice versa)
returns an equivalent object in the appropriate class.

Therefore the set of \(\GZTLC\)-horizontal curves in
\(\widehat{\Bundle}\) over \(\gamma\), modulo vertical gauge, is in
bijective correspondence with the set of admissible strict PFB-GNLA
deformations along \(\gamma\), modulo the gauge-equivalence described
in the definition.  This is precisely the equivalence of descriptions
claimed in item~(3).



\medskip
This completes the proof of
Proposition~\ref{prop:horizontal-flow-matching}.
\end{proof}



\section[Proposition P7: Law-Level G-Algebra and PFB-GNLA Representations]{Proposition P7: Law-Level \GAlgebra\ and PFB-GNLA \\ Representations}
\label{app:I-P7}

\subsection*{Set-up: law-level algebra and representation categories}


In this appendix we make precise the schematic Proposition~P7 from the
main text and establish the representation-theoretic correspondence
between the law-level algebra \(\GAlgebra_{\mathrm{orig}}\) and
strict PFB-GNLA data.

We fix a PL--PI law-system (for instance, a QFT-type system) whose
law-level algebra is denoted by
\[
  \GAlgebra_{\mathrm{orig}}.
\]
Concretely, \(\GAlgebra_{\mathrm{orig}}\) is a (possibly \(\mathbb{Z}\)- or
\(\mathbb{Z}_2\)-graded) algebra encoding the primitive operations,
brackets, and higher law-composition rules at the level of laws; it is
assumed to be equipped with whatever topological/completion structure
is needed to interpret the PL--PI compositions.

We introduce two categories:

\begin{itemize}
  \item \(\mathsf{Rep}_{\mathrm{loc}}(\GAlgebra_{\mathrm{orig}})\):
        the category of \emph{local law-level representations} of
        \(\GAlgebra_{\mathrm{orig}}\).  Objects are pairs
        \((E,\rho)\), where \(E\to M\) is a (topological or smooth)
        vector bundle over the base manifold \(M\), and
        \(\rho:\GAlgebra_{\mathrm{orig}}\to
        \mathrm{End}_{\mathrm{loc}}(E)\) is a representation respecting
        the PL--PI structure (locality, composition, and grading).
        Morphisms are bundle maps \(\Phi:E\to E'\) over \(M\) that
        intertwine the representations:
        \(\Phi\circ\rho(a)=\rho'(a)\circ\Phi\) for all
        \(a\in\GAlgebra_{\mathrm{orig}}\).
  \item \(\mathsf{PFB\text- GNLA}^{\mathrm{rep}}_{\mathrm{strict}}(\GAlgebra_{\mathrm{orig}})\):
        the category of \emph{strict PFB-GNLA representation data}
        associated with \(\GAlgebra_{\mathrm{orig}}\).  An object is a
        quadruple
        \[
          \mathcal{Q}
          =
          \bigl(P\to M,\;
                \Conn,\;
                \mathfrak{g}_{\mathrm{GNLA}},\;
                \rho_{\mathrm{GNLA}}\bigr),
        \]
        where:
        \begin{enumerate}
          \item \(P\to M\) is a principal bundle with structure group
                \(G_{\mathrm{geom}}\) and connection \(\Conn\);
          \item \(\mathfrak{g}_{\mathrm{GNLA}}\) is a strict GNLA fibre
                (the ``object-level'' algebra of a PFB-GNLA);
          \item \(\rho_{\mathrm{GNLA}}\) is a representation of
                \(\GAlgebra_{\mathrm{orig}}\) on an associated bundle
                \(E_{\mathcal{Q}}\to M\) built from \(P\) and
                \(\mathfrak{g}_{\mathrm{GNLA}}\), compatible with the
                GNLA structure and the PL--PI compositions.
        \end{enumerate}
        Morphisms in
        \(\mathsf{PFB\text- GNLA}^{\mathrm{rep}}_{\mathrm{strict}}\)
        are bundle maps and GNLA homomorphisms that commute with all
        the structure and intertwine the \(\GAlgebra_{\mathrm{orig}}\)
        actions.
\end{itemize}

We will also consider the associated \emph{gauge-quotient category}
\[
  \overline{\mathsf{PFB\text- GNLA}^{\mathrm{rep}}_{\mathrm{strict}}}
  (\GAlgebra_{\mathrm{orig}}),
\]
whose objects are those of
\(\mathsf{PFB\text- GNLA}^{\mathrm{rep}}_{\mathrm{strict}}\) and whose
morphisms are equivalence classes modulo vertical bundle gauge
transformations that act trivially on the base \(M\) and preserve the
\(\GAlgebra_{\mathrm{orig}}\)-representation structure.



\subsection*{Statement of Proposition P7}


\begin{proposition}[Law-level \GAlgebra\ and PFB-GNLA representations]
\label{prop:P7-law-algebra-PFB-GNLA}
Let \(\GAlgebra_{\mathrm{orig}}\) be the law-level algebra of a
GZ-admissible PL--PI law-system as above.  Then there exists a
canonical functor
\[
  \mathcal{R} :
  \overline{\mathsf{PFB\text- GNLA}^{\mathrm{rep}}_{\mathrm{strict}}}
    (\GAlgebra_{\mathrm{orig}})
  \longrightarrow
  \mathsf{Rep}_{\mathrm{loc}}(\GAlgebra_{\mathrm{orig}})
\]
such that:

\begin{enumerate}
  \item[\textnormal{(1)}] \textbf{Representation from PFB-GNLA data.}
    For each strict PFB-GNLA representation object \(\mathcal{Q}\),
    the pair \(\mathcal{R}(\mathcal{Q})=(E_{\mathcal{Q}},\rho_{\mathcal{Q}})\)
    is a well-defined local representation of \(\GAlgebra_{\mathrm{orig}}\),
    and gauge-equivalent PFB-GNLA data give rise to isomorphic law-level
    representations.
  \item[\textnormal{(2)}] \textbf{Faithfulness up to gauge.}
    The functor \(\mathcal{R}\) is faithful: if
    \(\mathcal{R}(f)=\mathcal{R}(g)\) for two morphisms
    \(f,g:\mathcal{Q}\to\mathcal{Q}'\), then \(f\) and \(g\) differ
    only by a vertical gauge-transformation (hence coincide in the
    gauge-quotient category).
  \item[\textnormal{(3)}] \textbf{Essential surjectivity.}
    Every local representation \((E,\rho)\) of
    \(\GAlgebra_{\mathrm{orig}}\) that satisfies the PL--PI locality
    and regularity conditions of the GZ-framework arises (up to
    isomorphism in \(\mathsf{Rep}_{\mathrm{loc}}\)) as
    \(\mathcal{R}(\mathcal{Q})\) for some strict PFB-GNLA
    representation object \(\mathcal{Q}\).
  \item[\textnormal{(4)}] \textbf{Equivalence onto a full subcategory.}
    The functor \(\mathcal{R}\) induces an equivalence of categories
    between the gauge-quotient category
    \(\overline{\mathsf{PFB\text- GNLA}^{\mathrm{rep}}_{\mathrm{strict}}}
     (\GAlgebra_{\mathrm{orig}})\) and a full subcategory of
    \(\mathsf{Rep}_{\mathrm{loc}}(\GAlgebra_{\mathrm{orig}})\)
    consisting of those representations that are geometrically
    realisable as strict PFB-GNLA data.
\end{enumerate}
\end{proposition}



\subsection*{Proof of Proposition~\ref{prop:P7-law-algebra-PFB-GNLA}}


\begin{proof}
We outline the construction and then verify properties
\textnormal{(1)}–\textnormal{(4)}.

\medskip\noindent
\textbf{Step 1: Construction of the functor \(\mathcal{R}\).}
Let \(\mathcal{Q}=(P\to M,\Conn,\mathfrak{g}_{\mathrm{GNLA}},\rho_{\mathrm{GNLA}})\)
be a strict PFB-GNLA representation object.  From the principal bundle
 \(P\) and the GNLA fibre \(\mathfrak{g}_{\mathrm{GNLA}}\), we form the
associated vector bundle
\[
  E_{\mathcal{Q}}
  := P\times_{G_{\mathrm{geom}}} V,
\]
where \(V\) is a suitable representation space on which
\(\mathfrak{g}_{\mathrm{GNLA}}\) acts (e.g.\ via the adjoint or a
chosen faithful module).  The GNLA structure on the fibres of
\(\mathcal{Q}\) encodes, in particular, a homomorphism
\[
  \rho_{\mathrm{GNLA}} :
  \GAlgebra_{\mathrm{orig}} \longrightarrow \mathrm{End}(V)
\]
compatible with the GNLA brackets and PL--PI operations.

Transporting \(\rho_{\mathrm{GNLA}}\) along the fibres of \(P\) via the
associated-bundle construction yields a representation
\[
  \rho_{\mathcal{Q}} :
  \GAlgebra_{\mathrm{orig}} \longrightarrow
  \mathrm{End}_{\mathrm{loc}}(E_{\mathcal{Q}}).
\]
By construction, \(\rho_{\mathcal{Q}}\) is local (it acts fibrewise and
is compatible with parallel transport defined by \(\Conn\)) and
respects the PL--PI composition rules.

We then set
\[
  \mathcal{R}(\mathcal{Q})
  :=
  \bigl(E_{\mathcal{Q}},\rho_{\mathcal{Q}}\bigr).
\]

For a morphism
\[
  f:\mathcal{Q}\longrightarrow\mathcal{Q}'
\]
in \(\mathsf{PFB\text- GNLA}^{\mathrm{rep}}_{\mathrm{strict}}\),
consisting of a principal-bundle map and a GNLA homomorphism
compatible with all structure, we obtain an induced bundle map
\[
  \mathcal{R}(f)
  :=
  f_* : E_{\mathcal{Q}}\longrightarrow E_{\mathcal{Q}'}
\]
which, by compatibility of \(\rho_{\mathrm{GNLA}}\) and
\(\rho'_{\mathrm{GNLA}}\), intertwines the representations:
\[
  f_*\circ\rho_{\mathcal{Q}}(a)
  =
  \rho_{\mathcal{Q}'}(a)\circ f_*
  \quad\text{for all } a\in\GAlgebra_{\mathrm{orig}}.
\]
Thus \(\mathcal{R}\) is a bona fide functor into
\(\mathsf{Rep}_{\mathrm{loc}}(\GAlgebra_{\mathrm{orig}})\).  Passing
to the gauge-quotient category simply identifies morphisms that differ
by vertical gauge-transformations acting trivially on \(E_{\mathcal{Q}}\).

This yields property~\textnormal{(1)}.



\medskip\noindent
\textbf{Step 2: Faithfulness up to gauge.}
Suppose \(f,g:\mathcal{Q}\to\mathcal{Q}'\) are morphisms in
\(\mathsf{PFB\text- GNLA}^{\mathrm{rep}}_{\mathrm{strict}}\) such
that \(\mathcal{R}(f)=\mathcal{R}(g)\) as intertwiners between
\(\mathcal{R}(\mathcal{Q})\) and \(\mathcal{R}(\mathcal{Q}')\).  Then
the induced bundle maps on \(E_{\mathcal{Q}}\) agree:
\[
  f_* = g_* : E_{\mathcal{Q}}\to E_{\mathcal{Q}'}.
\]
Since \(E_{\mathcal{Q}}\) is associated to the principal bundle \(P\)
via a faithful representation of the structure group
(and/or of \(\mathfrak{g}_{\mathrm{GNLA}}\)), the equality of
associated-bundle maps implies that \(f\) and \(g\) differ at most by a
vertical gauge-transformation in the principal bundle that acts
trivially on the associated representation.  Such a transformation is
precisely what is modded out in the gauge-quotient category, hence
 \(f\) and \(g\) coincide there.

Therefore \(\mathcal{R}\) is faithful on the gauge-quotient category,
establishing property~\textnormal{(2)}.



\medskip\noindent
\textbf{Step 3: Essential surjectivity.}
Let \((E,\rho)\) be a local representation of
\(\GAlgebra_{\mathrm{orig}}\) satisfying the PL--PI locality and
regularity conditions of the GZ-framework.  We sketch the construction
of a strict PFB-GNLA representation object \(\mathcal{Q}\) such that
\(\mathcal{R}(\mathcal{Q})\cong(E,\rho)\).

First, we consider the frame bundle \(P_E\to M\) of the vector bundle
 \(E\), whose fibres consist of linear isomorphisms from a fixed model
space \(V\) to the fibre \(E_x\) at each point \(x\in M\).  This is a
principal \(G_{\mathrm{geom}}\)-bundle for an appropriate structure
group \(G_{\mathrm{geom}}\subset\mathrm{GL}(V)\).

The law-level representation \(\rho:\GAlgebra_{\mathrm{orig}}\to
\mathrm{End}_{\mathrm{loc}}(E)\) induces a representation
\(\rho_V:\GAlgebra_{\mathrm{orig}}\to\mathrm{End}(V)\) on the model
space \(V\) by trivialising \(E\) locally and using the locality of
\(\rho\) to ensure compatibility on overlaps.  The regularity
conditions (continuity, locality, boundedness) guarantee that this
representation is well-defined up to conjugation by elements of
 \(G_{\mathrm{geom}}\).

We then equip the principal bundle \(P_E\) with a connection \(\Conn\)
compatible with the given representation \(\rho\) in the sense that
parallel transport by \(\Conn\) intertwines the action of
\(\GAlgebra_{\mathrm{orig}}\) along paths in \(M\) (this is standard:
one chooses \(\Conn\) so that parallel transport preserves a chosen
local frame adapted to the representation).  The GNLA fibre
\(\mathfrak{g}_{\mathrm{GNLA}}\) is taken to be a strict Lie algebra
(or higher Lie algebra) model that realises the brackets of
\(\GAlgebra_{\mathrm{orig}}\) at the object level, and
\(\rho_{\mathrm{GNLA}}:\GAlgebra_{\mathrm{orig}}\to\mathrm{End}(V)\)
is defined to be (a suitable lift of) \(\rho_V\).

This data defines a strict PFB-GNLA representation object
\[
  \mathcal{Q}
  =
  \bigl(P_E\to M,\Conn,\mathfrak{g}_{\mathrm{GNLA}},
        \rho_{\mathrm{GNLA}}\bigr)
  \in
  \mathsf{PFB\text- GNLA}^{\mathrm{rep}}_{\mathrm{strict}}
  (\GAlgebra_{\mathrm{orig}}),
\]
and by construction the associated bundle \(E_{\mathcal{Q}}\) can be
canonically identified with \(E\) in a way that intertwines
\(\rho_{\mathcal{Q}}\) and \(\rho\).  Thus
\(\mathcal{R}(\mathcal{Q})\cong(E,\rho)\), establishing the essential
surjectivity property~\textnormal{(3)}.



\medskip\noindent
\textbf{Step 4: Equivalence onto a full subcategory.}
Combining Steps~1–3 we see that \(\mathcal{R}\) is:

\begin{itemize}
  \item faithful on the gauge-quotient category (Step~2);
  \item essentially surjective onto the full subcategory of
        \(\mathsf{Rep}_{\mathrm{loc}}(\GAlgebra_{\mathrm{orig}})\)
        consisting of geometrically realisable representations
        (Step~3);
  \item full onto this subcategory, since any intertwiner between
        two such representations can be lifted (locally and then
        globally by patching) to a morphism of the corresponding
        PFB-GNLA data, using the functoriality of the frame-bundle
        construction and the GNLA structure.
\end{itemize}

Hence \(\mathcal{R}\) induces an equivalence of categories between
\[
  \overline{\mathsf{PFB\text- GNLA}^{\mathrm{rep}}_{\mathrm{strict}}}
    (\GAlgebra_{\mathrm{orig}})
\]
and the full subcategory of
\(\mathsf{Rep}_{\mathrm{loc}}(\GAlgebra_{\mathrm{orig}})\) consisting
of geometrically realisable representations.  This proves
property~\textnormal{(4)} and completes the proof of
Proposition~\ref{prop:P7-law-algebra-PFB-GNLA}.



\end{proof}



\chapter{Technical Appendix II: Extended Theoretical Derivations}


\section{Proposition P5: Homotopy Upgrade and Strictification}
\label{app:II-P5}

\subsection*{Set-up: strict and homotopy PFB-GNLA regimes}


In this appendix we give a more detailed account of the schematic
Proposition~P5 in the main text, concerning the homotopy-upgrade
functor and its (partially defined) strictification inverse between
different PFB-GNLA regimes.

We consider the following categories, all fibrewise over the same base
manifold \(M\).

\begin{itemize}
  \item \(\mathsf{PFB\text- GNLA}_{\mathrm{strict}}\): the category of
        \emph{strict PFB-GNLA} objects.  An object
        \(\mathcal{Q}_{\mathrm{strict}}\) consists of:
        \[
          \mathcal{Q}_{\mathrm{strict}}
          =
          \bigl(P\to M,\;
                \Conn,\;
                \mathfrak{g}_{\mathrm{strict}},\;
                \text{PFB-GNLA structure}\bigr),
        \]
        where \(\mathfrak{g}_{\mathrm{strict}}\) is a (possibly
        graded) strict Lie algebra fibre, and the PFB-GNLA structure
        encodes a strict GNLA action on associated bundles, with
        Jacobi (and higher) identities holding on the nose.
  \item \(\mathsf{PFB\text- GNLA}_{\mathrm{homo}}\): the category of
        \emph{homotopy PFB-GNLA} objects.  An object
        \(\mathcal{Q}_{\mathrm{homo}}\) consists of a principal bundle
        and connection as before, together with a fibrewise
        \(L_\infty\)-type structure
        \[
          \bigl(V_\bullet,\{l_k\}_{k\ge 1}\bigr),
        \]
        where \(V_0\) is the “visible” GNLA layer (object-level
        algebra), \(V_{>0}\) encodes higher homotopies, and the
        brackets \(l_k\) satisfy higher Jacobi identities up to
        homotopy.
\end{itemize}

Morphisms in these categories are principal-bundle maps covering the
identity on \(M\), together with fibrewise GNLA / \(L_\infty\)
homomorphisms compatible with the PFB structure; in the homotopy case,
morphisms are required to respect the full \(L_\infty\)-structure.

We denote by
\[
  \Phi_{\mathrm{geom}} :
  \mathsf{PFB\text- GNLA}_{\mathrm{strict}}
  \longrightarrow
  \GZLS
\]
the geometric representation functor considered in the main text, and
we let
\[
  L(\mathcal{Q})
  :=
  \Phi_{\mathrm{geom}}(\mathcal{Q})
\]
denote the associated law-object in \(\GZLS\) for a strict PFB-GNLA
object \(\mathcal{Q}\).  For homotopy PFB-GNLA objects, the associated
law-object is obtained by composing with the law-space upgrade
(\(\GZLS\)-level OHU) as in Theorem~\ref{thm:GZ-OHU}.

We use the Jacobiator-gap certificate \(\mathsf{JacGap}(\,\cdot\,)\)
from Theorem~\ref{thm:GZ-OHU} at the law-level and lift it to
PFB-GNLA objects by
\[
  \mathsf{JacGap}(\mathcal{Q})
  :=
  \mathsf{JacGap}\bigl(L(\mathcal{Q})\bigr),
\]
where the right-hand side is the defect functional associated to the
law-object \(L(\mathcal{Q})\).



\subsection*{Definition of the Upgrade and Strictify constructions}


We now define two constructions:

\begin{itemize}
  \item a functor
    \[
      \mathsf{Upgrade} :
      \mathsf{PFB\text- GNLA}_{\mathrm{strict}}
      \longrightarrow
      \mathsf{PFB\text- GNLA}_{\mathrm{homo}},
    \]
  \item and a partially defined functor
    \[
      \mathsf{Strictify} :
      \mathsf{PFB\text- GNLA}_{\mathrm{homo}}^{(0)}
      \dashrightarrow
      \mathsf{PFB\text- GNLA}_{\mathrm{strict}},
    \]
    where \(\mathsf{PFB\text- GNLA}_{\mathrm{homo}}^{(0)}\) denotes
    the full subcategory of homotopy PFB-GNLA objects with
    \(\mathsf{JacGap}(\mathcal{Q})=0\) (Jacobiator-gap-free objects).
\end{itemize}

\paragraph{The Upgrade functor.}
Let
\[
  \mathcal{Q}_{\mathrm{strict}}
  =
  \bigl(P\to M,\Conn,\mathfrak{g}_{\mathrm{strict}},
        \text{strict GNLA structure}\bigr)
  \in
  \mathsf{PFB\text- GNLA}_{\mathrm{strict}}.
\]
We define
\[
  \mathsf{Upgrade}\bigl(\mathcal{Q}_{\mathrm{strict}}\bigr)
  =
  \mathcal{Q}_{\mathrm{homo}}
  :=
  \bigl(P\to M,\Conn,V_\bullet,\{l_k\}_{k\ge 1}\bigr),
\]
where:
\begin{enumerate}
  \item \(V_0\) is the original GNLA fibre
        \(\mathfrak{g}_{\mathrm{strict}}\);
  \item \(V_1\) is a space of formal homotopies encoding possible
        deviations from strict Jacobi identities (chosen sufficiently
        large so that all potential higher corrections live in
        \(V_1\));
  \item \(V_i=0\) for \(i\ge 2\) (for simplicity of exposition; the
        construction extends to longer complexes if needed);
  \item the differential \(l_1:V_1\to V_0\) and the higher brackets
        \(l_k\) are chosen so that:
        \begin{itemize}
          \item the induced bracket on \(H^0(V_\bullet)\) coincides
                with the original strict GNLA bracket on
                \(\mathfrak{g}_{\mathrm{strict}}\);
          \item the higher Jacobi identities for \(\{l_k\}\) encode
                the same law-level relations as in
                \(\mathcal{Q}_{\mathrm{strict}}\), but allow for
                homotopy corrections when one passes to more flexible
                regimes (e.g.\ vf-op versions).
        \end{itemize}
\end{enumerate}
The principal bundle and connection \((P,\Conn)\) are kept unchanged;
the GNLA fibre is replaced by the two-term complex \(V_\bullet\), and
the connection is extended to act on each component in a way compatible
with the \(L_\infty\)-structure.  On morphisms
\(f:\mathcal{Q}_{\mathrm{strict}}\to\mathcal{Q}'_{\mathrm{strict}}\),
we define
\[
  \mathsf{Upgrade}(f)
  := f
\]
on the underlying principal bundles and extend it trivially to
\(V_\bullet\), ensuring that the \(L_\infty\)-structure maps are
respected.  This yields a well-defined functor
\(\mathsf{Upgrade}\).

\paragraph{The Strictify construction.}
Let
\[
  \mathcal{Q}_{\mathrm{homo}}
  =
  \bigl(P\to M,\Conn,V_\bullet,\{l_k\}_{k\ge 1}\bigr)
  \in
  \mathsf{PFB\text- GNLA}_{\mathrm{homo}}^{(0)}
\]
be a homotopy PFB-GNLA object with \(\mathsf{JacGap}(\mathcal{Q}_{\mathrm{homo}})=0\).
Using the associated law-object
\[
  L\bigl(\mathcal{Q}_{\mathrm{homo}}\bigr)\in\GZLS
\]
and the GZ-OHU theorem (Theorem~\ref{thm:GZ-OHU}) together with
standard \(L_\infty\)-rectification results (for example, the
rectification of minimal \(L_\infty\)-models when the Jacobiator is
strictly zero), one can construct:
\begin{itemize}
  \item a strict GNLA fibre \(\mathfrak{g}_{\mathrm{strict}}\),
  \item a principal bundle \(P_{\mathrm{strict}}\to M\) with
        connection \(\Conn_{\mathrm{strict}}\) (isomorphic to
        \((P,\Conn)\) up to gauge),
  \item and a PFB-GNLA structure
        \(\mathcal{Q}_{\mathrm{strict}}=
         \bigl(P_{\mathrm{strict}}\to M,
               \Conn_{\mathrm{strict}},
               \mathfrak{g}_{\mathrm{strict}},
               \text{strict GNLA data}\bigr)\),
\end{itemize}
together with an \(L_\infty\)-quasi-isomorphism (over the identity on
\(M\))
\[
  F :
  \mathcal{Q}_{\mathrm{homo}}
  \longrightarrow
  \mathsf{Upgrade}\bigl(\mathcal{Q}_{\mathrm{strict}}\bigr)
\]
compatible with the PFB structure and the law-space representation.

We then \emph{define}
\[
  \mathsf{Strictify}\bigl(\mathcal{Q}_{\mathrm{homo}}\bigr)
  :=
  \mathcal{Q}_{\mathrm{strict}},
\]
with the understanding that \(\mathcal{Q}_{\mathrm{strict}}\) is
well-defined only up to isomorphism in
\(\mathsf{PFB\text- GNLA}_{\mathrm{strict}}\).  At the level of the
homotopy category, this ambiguity disappears and
\(\mathsf{Strictify}\) becomes a well-defined (partial) functor.

On morphisms between JacGap-free homotopy objects
\(\mathcal{Q}_{\mathrm{homo}}\to\mathcal{Q}'_{\mathrm{homo}}\), one
uses the \(L_\infty\)-quasi-isomorphisms and the universal property of
the strict models to lift them (up to homotopy) to strict morphisms
between the corresponding \(\mathcal{Q}_{\mathrm{strict}}\) and
\(\mathcal{Q}'_{\mathrm{strict}}\).  This yields a partially defined
functor
\[
  \mathsf{Strictify} :
  \mathsf{PFB\text- GNLA}_{\mathrm{homo}}^{(0)}
  \dashrightarrow
  \mathsf{PFB\text- GNLA}_{\mathrm{strict}}.
\]



\subsection*{Statement of Proposition P5}


\begin{proposition}[Homotopy upgrade and strictification]
\label{prop:P5-upgrade-strictify}
The constructions above satisfy the following properties.

\begin{enumerate}
  \item[\textnormal{(1)}] \textbf{Functorial homotopy upgrade.}
    The assignment 
    \(\mathsf{Upgrade} : \\
      \mathsf{PFB\text- GNLA}_{\mathrm{strict}}
      \to \mathsf{PFB\text- GNLA}_{\mathrm{homo}}\)
    defines a fully faithful functor that preserves the underlying
    principal bundle and connection and induces the identity on
    associated law-objects in \(\GZLS\).
  \item[\textnormal{(2)}] \textbf{Partial strictification on the
        JacGap-free regime.}
    The assignment
    \(\mathsf{Strictify} :
      \mathsf{PFB\text- GNLA}_{\mathrm{homo}}^{(0)}
      \dashrightarrow
      \mathsf{PFB\text- GNLA}_{\mathrm{strict}}\)
    is well-defined up to canonical isomorphism, and in the homotopy
    category satisfies:
    \[
      \mathsf{Strictify}\circ\mathsf{Upgrade}
      \;\simeq\;
      \mathrm{Id}_{\mathsf{PFB\text- GNLA}_{\mathrm{strict}}},
      \qquad
      \mathsf{Upgrade}\circ\mathsf{Strictify}
      \;\simeq\;
      \mathrm{Id}\quad
      \text{on }\mathsf{PFB\text- GNLA}_{\mathrm{homo}}^{(0)}.
    \]
  \item[\textnormal{(3)}] \textbf{Gap barrier for strictification.}
    Let \(\mathcal{Q}_{\mathrm{homo}}\) be a homotopy PFB-GNLA object
    with \(\mathsf{JacGap}(\mathcal{Q}_{\mathrm{homo}})>0\).  Then any
    zigzag of homotopy-equivalences in
    \(\mathsf{PFB\text- GNLA}_{\mathrm{homo}}\) connecting
    \(\mathcal{Q}_{\mathrm{homo}}\) to an object in the essential
    image of \(\mathsf{Upgrade}\) must pass through an intermediate
    object whose Jacobiator-gap is \emph{bounded away from zero}:
    there exists \(\delta>0\) (depending on
    \(\mathcal{Q}_{\mathrm{homo}}\)) such that for any such zigzag
    \(\mathcal{Q}_{\mathrm{homo}}=\mathcal{Q}_0
      \leftrightsquigarrow \dots
      \leftrightsquigarrow\mathcal{Q}_N\),
    with \(\mathcal{Q}_N\simeq \mathsf{Upgrade}(\mathcal{Q}_{\mathrm{strict}})\),
    one has
    \[
      \max_{0\le i\le N}
      \mathsf{JacGap}(\mathcal{Q}_i)
      \;\ge\;
      \delta.
    \]
    In particular, it is impossible to strictify
    \(\mathcal{Q}_{\mathrm{homo}}\) to a strict PFB-GNLA object along
    a homotopy path that keeps \(\mathsf{JacGap}\) arbitrarily small.
\end{enumerate}
\end{proposition}



\subsection*{Proof of Proposition~\ref{prop:P5-upgrade-strictify}}


\begin{proof}
\textbf{Proof of (1): functorial homotopy upgrade.}
By construction, \(\mathsf{Upgrade}\) acts as the identity on the
underlying principal bundle and connection, and only enlarges the
fibre from \(\mathfrak{g}_{\mathrm{strict}}\) to the two-term complex
\(V_\bullet=(V_1\to V_0)\) with \(V_0=\mathfrak{g}_{\mathrm{strict}}\).
The brackets \(l_k\) are chosen so that the induced bracket on
\(H^0(V_\bullet)\) coincides with the original GNLA bracket on
\(\mathfrak{g}_{\mathrm{strict}}\).  Hence, for every strict
PFB-GNLA object \(\mathcal{Q}_{\mathrm{strict}}\),
\(\mathsf{Upgrade}(\mathcal{Q}_{\mathrm{strict}})\) represents the
same law-object in \(\GZLS\) under \(\Phi_{\mathrm{geom}}\); law-level
procedures see no difference between the two descriptions.

On morphisms, \(\mathsf{Upgrade}\) simply reinterprets the same bundle
map and GNLA homomorphism in the enlarged \(L_\infty\) setting.  Since
no additional freedom is introduced in degree \(0\), the map on
morphisms is injective and reflects equality.  Therefore
\(\mathsf{Upgrade}\) is faithful.  Fullness follows because any
morphism in \(\mathsf{PFB\text- GNLA}_{\mathrm{homo}}\) between two
objects in the image of \(\mathsf{Upgrade}\) is forced to respect the
trivial \(V_1\) structure and hence descends uniquely to a strict
morphism in \(\mathsf{PFB\text- GNLA}_{\mathrm{strict}}\).  This
establishes \textnormal{(1)}.



\medskip\noindent
\textbf{Proof of (2): strictification on the JacGap-free regime.}
Let \(\mathcal{Q}_{\mathrm{homo}}\in
\mathsf{PFB\text- GNLA}_{\mathrm{homo}}^{(0)}\).  By definition,
\(\mathsf{JacGap}(\mathcal{Q}_{\mathrm{homo}})=0\), so the associated
law-object \(L(\mathcal{Q}_{\mathrm{homo}})\) has vanishing
Jacobiator-defect in the sense of Theorem~\ref{thm:GZ-OHU}.  The
GZ-OHU completion theorem (combined with standard rectification
results for \(L_\infty\)-algebras with zero Jacobiator) then implies
that there exists a strict GNLA model
\(\mathcal{Q}_{\mathrm{strict}}\) and an \(L_\infty\)-quasi-isomorphism
\[
  F :
  \mathcal{Q}_{\mathrm{homo}}
  \longrightarrow
  \mathsf{Upgrade}(\mathcal{Q}_{\mathrm{strict}}),
\]
over the identity on \(M\), that preserves the PFB structure and
induces an equivalence of law-objects in \(\GZLS\).

If we start instead from a strict object
\(\mathcal{Q}_{\mathrm{strict}}\in
 \mathsf{PFB\text- GNLA}_{\mathrm{strict}}\),
then \(\mathsf{Upgrade}(\mathcal{Q}_{\mathrm{strict}})\) is
JacGap-free by construction (its law-object is already strict with zero
Jacobiator), and the canonical identification
\[
  \mathsf{Strictify}\bigl(\mathsf{Upgrade}(\mathcal{Q}_{\mathrm{strict}})\bigr)
  \;\cong\;
  \mathcal{Q}_{\mathrm{strict}}
\]
follows from the uniqueness (up to isomorphism) of the strict model in
the rectification theorem.  This yields the homotopy equivalence
\[
  \mathsf{Strictify}\circ\mathsf{Upgrade}
  \;\simeq\;
  \mathrm{Id}_{\mathsf{PFB\text- GNLA}_{\mathrm{strict}}}.
\]

Conversely, if \(\mathcal{Q}_{\mathrm{homo}}\in
\mathsf{PFB\text- GNLA}_{\mathrm{homo}}^{(0)}\), then the strict model
\(\mathcal{Q}_{\mathrm{strict}}=
\mathsf{Strictify}(\mathcal{Q}_{\mathrm{homo}})\) together with
\(\mathsf{Upgrade}\) gives
\[
  \mathcal{Q}_{\mathrm{homo}}
  \xrightarrow{\;F\;}
  \mathsf{Upgrade}\bigl(\mathsf{Strictify}(\mathcal{Q}_{\mathrm{homo}})\bigr),
\]
where \(F\) is an \(L_\infty\)-quasi-isomorphism.  In the homotopy
category of \(\mathsf{PFB\text- GNLA}_{\mathrm{homo}}\), this implies
\[
  \mathsf{Upgrade}\circ\mathsf{Strictify}
  \;\simeq\;
  \mathrm{Id}\quad
  \text{on }\mathsf{PFB\text- GNLA}_{\mathrm{homo}}^{(0)}.
\]
This establishes property~\textnormal{(2)}.



\medskip\noindent
\textbf{Proof of (3): gap barrier for strictification.}
Let \(\mathcal{Q}_{\mathrm{homo}}\) be a homotopy PFB-GNLA object with
\(\mathsf{JacGap}(\mathcal{Q}_{\mathrm{homo}})>0\).  Suppose, for the
sake of contradiction, that there exists a sequence (zigzag) of
homotopy-equivalences in \(\mathsf{PFB\text- GNLA}_{\mathrm{homo}}\),
\[
  \mathcal{Q}_0
  \leftrightsquigarrow
  \mathcal{Q}_1
  \leftrightsquigarrow
  \dots
  \leftrightsquigarrow
  \mathcal{Q}_N,
\]
with \(\mathcal{Q}_0=\mathcal{Q}_{\mathrm{homo}}\) and
\(\mathcal{Q}_N\simeq\mathsf{Upgrade}(\mathcal{Q}_{\mathrm{strict}})\)
for some strict object \(\mathcal{Q}_{\mathrm{strict}}\), such that the
Jacobiator-gap remains uniformly small along the zigzag:
\[
  \mathsf{JacGap}(\mathcal{Q}_i)
  \;<\;
  \varepsilon
  \quad\text{for all } i=0,\dots,N,
\]
where \(\varepsilon>0\) can be chosen arbitrarily small.

Passing to the associated law-objects \(L_i:=L(\mathcal{Q}_i)\in\GZLS\),
the homotopy-equivalences between \(\mathcal{Q}_i\) induce
law-equivalences between the corresponding \(L_i\).  The law-level
Jacobiator-gap functional \(\mathsf{JacGap}(L_i)\) coincides with
\(\mathsf{JacGap}(\mathcal{Q}_i)\) by definition, so we have a zigzag
\[
  L_0 \leftrightsquigarrow L_1 \leftrightsquigarrow \dots
  \leftrightsquigarrow L_N
\]
of law-objects in \(\GZLS\) with uniformly small Jacobiator-gap, and
\(L_N\) corresponds to a strict law-system in the sense of
Theorem~\ref{thm:GZ-OHU} (because \(\mathcal{Q}_N\) lies in the image
of \(\mathsf{Upgrade}\) of a strict PFB-GNLA).

By the GZ-OHU completion theorem, strict law-systems with zero
Jacobiator-gap are \\ Jacobiator-gap-minimal in their homotopy class.  In
particular, any law-object law-homotopy-equivalent to such a strict
system that has arbitrarily small but non-zero Jacobiator-gap would
contradict the minimality clause: one could further reduce the defect
energy while staying within the same homotopy class, violating the
universal property of the OHU completion.

Formalising this argument, we obtain a positive lower bound
\(\delta>0\) depending on the homotopy class of \(L_0\) (and hence of
\(\mathcal{Q}_{\mathrm{homo}}\)) such that any zigzag connecting
\(L_0\) to a strict law-system must pass through some intermediate
\(L_i\) with
\[
  \mathsf{JacGap}(L_i)\;\ge\;\delta.
\]
Translating back to PFB-GNLA objects, we conclude that any zigzag in
\(\mathsf{PFB\text- GNLA}_{\mathrm{homo}}\) connecting
\(\mathcal{Q}_{\mathrm{homo}}\) to an object in the image of
\(\mathsf{Upgrade}\) must pass through an intermediate
\(\mathcal{Q}_i\) with
\[
  \mathsf{JacGap}(\mathcal{Q}_i)\;\ge\;\delta,
\]
contradicting the assumption that the gap can be kept arbitrarily
small throughout the zigzag.

This proves the existence of a gap barrier and establishes
property~\textnormal{(3)}.



\medskip
Combining the arguments for \textnormal{(1)}–\textnormal{(3)} completes
the proof of Proposition~\ref{prop:P5-upgrade-strictify}.
\end{proof}



\section[Proposition P6: Projection to Homotopy and Strict Regimes]{Proposition P6: Projection to Homotopy and Strict \\ Regimes}
\label{app:II-P6}

\subsection*{Set-up: vf-op, homotopy, and strict PFB-GNLA regimes}


In this appendix we expand on the schematic Proposition~P6 from the
main text, which introduces projection functors from the vector-field /
operator-path regime (vf-op) down to the homotopy and strict regimes.

We work with three fibrewise categories over the same base manifold
\(M\):

\begin{itemize}
  \item \(\mathsf{PFB\text- GNLA}_{\mathrm{vf\text- op}}\): the
        category of \emph{vf-op PFB-GNLA} objects.  An object
        \(\mathcal{Q}_{\mathrm{vf\text- op}}\) consists of
        \[
          \mathcal{Q}_{\mathrm{vf\text- op}}
          =
          \bigl(P\to M,\;
                \Conn,\;
                V_\bullet,\{l_k\}_{k\ge 1},\;
                \mathcal{O}_{\mathrm{vf\text- op}}\bigr),
        \]
        where:
        \begin{enumerate}
          \item \(P\to M\) is a principal bundle with connection
                \(\Conn\);
          \item \((V_\bullet,\{l_k\})\) is an \(L_\infty\)-type
                fibrewise structure (homotopy GNLA), as in
                \(\mathsf{PFB\text- GNLA}_{\mathrm{homo}}\);
          \item \(\mathcal{O}_{\mathrm{vf\text- op}}\) denotes
                additional vf-op data: time- or path-dependent
                operators, vector fields, and control terms acting on
                the associated bundles, compatible with the
                \(L_\infty\)-structure.
        \end{enumerate}
        Morphisms are principal-bundle maps and \(L_\infty\)
        homomorphisms, together with maps on vf-op data that respect
        the operator-path structure.
  \item \(\mathsf{PFB\text- GNLA}_{\mathrm{homo}}\): the category of
        \emph{homotopy PFB-GNLA} objects, as in
        Technical Appendix~II, Section~\ref{app:II-P5}, without any
        explicit vf-op data.
  \item \(\mathsf{PFB\text- GNLA}_{\mathrm{strict}}\): the category of
        \emph{strict PFB-GNLA} objects, as in
        Technical Appendix~II, Section~\ref{app:II-P5}, with strict GNLA
        fibres.
\end{itemize}

We work with the law-space representation functor
\[
  L(-) : \mathsf{PFB\text- GNLA}_{\mathrm{vf\text- op}}
  \cup
  \mathsf{PFB\text- GNLA}_{\mathrm{homo}}
  \cup
  \mathsf{PFB\text- GNLA}_{\mathrm{strict}}
  \longrightarrow
  \GZLS,
\]
which sends a PFB-GNLA object (in any regime) to its associated
law-object in \(\GZLS\), as in the main text.  The Jacobiator-gap
functional at the PFB-GNLA level is always defined via the associated
law-object:
\[
  \mathsf{JacGap}(\mathcal{Q})
  :=
  \mathsf{JacGap}\bigl(L(\mathcal{Q})\bigr).
\]

We also reuse the strictly homotopy-upgrade and strictification functors
constructed in Proposition~\ref{prop:P5-upgrade-strictify}:
\[
  \mathsf{Upgrade} :
  \mathsf{PFB\text- GNLA}_{\mathrm{strict}}
  \longrightarrow
  \mathsf{PFB\text- GNLA}_{\mathrm{homo}},
\]
\[
  \mathsf{Strictify} :
  \mathsf{PFB\text- GNLA}_{\mathrm{homo}}^{(0)}
  \dashrightarrow
  \mathsf{PFB\text- GNLA}_{\mathrm{strict}},
\]
where \(\mathsf{PFB\text- GNLA}_{\mathrm{homo}}^{(0)}\) denotes the
JacGap-free subcategory.



\subsection*{Definition of the projection functors}


We now define the two projection functors mentioned in the main text.

\paragraph{The functor \(\Pi_{\mathrm{homo}}\).}
We define
\[
  \Pi_{\mathrm{homo}} :
  \mathsf{PFB\text- GNLA}_{\mathrm{vf\text- op}}
  \longrightarrow
  \mathsf{PFB\text- GNLA}_{\mathrm{homo}}
\]
as the functor that \emph{forgets the vf-op data but keeps the
underlying homotopy GNLA structure}.  More concretely, on objects we
set
\[
  \Pi_{\mathrm{homo}}\bigl(
    P\to M,\Conn,V_\bullet,\{l_k\},\mathcal{O}_{\mathrm{vf\text- op}}
  \bigr)
  :=
  \bigl(P\to M,\Conn,V_\bullet,\{l_k\}\bigr),
\]
and on morphisms we simply drop the component acting on
\(\mathcal{O}_{\mathrm{vf\text- op}}\), retaining only the principal
bundle map and the \(L_\infty\)-homomorphism.  By construction,
\(\Pi_{\mathrm{homo}}\) is a strict functor.

Observe that
\[
  L\bigl(\Pi_{\mathrm{homo}}(\mathcal{Q}_{\mathrm{vf\text- op}})\bigr)
  =
  L(\mathcal{Q}_{\mathrm{vf\text- op}}),
\]
since at the law-level the vf-op decorations do not change the
underlying law-object in \(\GZLS\); they only encode additional
trajectory- and operator-path information used in applications
(e.g.\ law-space learning, RL/GRL constructions).

\paragraph{The functor \(\Pi_{\mathrm{strict}}\).}
We define \(\Pi_{\mathrm{strict}}\) as the composition of the
JacGap-free inclusion into \(\mathsf{PFB\text- GNLA}_{\mathrm{homo}}\)
with the strictification functor of Proposition~\ref{prop:P5-upgrade-strictify}:
\[
  \Pi_{\mathrm{strict}} :
  \mathsf{PFB\text- GNLA}_{\mathrm{homo}}^{(0)}
  \xrightarrow{\;\mathsf{Strictify}\;}
  \mathsf{PFB\text- GNLA}_{\mathrm{strict}}.
\]
Thus \(\Pi_{\mathrm{strict}}\) is only defined on the JacGap-free
sub-category \(\mathsf{PFB\text- GNLA}_{\mathrm{homo}}^{(0)}\); for
objects with non-zero Jacobiator-gap it is deliberately left undefined,
reflecting the gap barrier from Proposition~\ref{prop:P5-upgrade-strictify}.

On morphisms, \(\Pi_{\mathrm{strict}}\) acts as
\(\mathsf{Strictify}\): an \(L_\infty\)-morphism between JacGap-free
homotopy objects is mapped (up to homotopy) to a strict PFB-GNLA
morphism between their strictifications.



\subsection*{Statement of Proposition P6}


\begin{proposition}[Projection to homotopy and strict regimes]
\label{prop:P6-projection}
The functors \(\Pi_{\mathrm{homo}}\) and \(\Pi_{\mathrm{strict}}\)
satisfy the following properties.

\begin{enumerate}
  \item[\textnormal{(1)}] \textbf{Functorial projection to the homotopy
        regime.}
    The functor
    \[
      \Pi_{\mathrm{homo}} :
      \mathsf{PFB\text- GNLA}_{\mathrm{vf\text- op}}
      \to
      \mathsf{PFB\text- GNLA}_{\mathrm{homo}}
    \]
    is well-defined and preserves the underlying principal bundle,
    connection, and law-object:
    \[
      L\bigl(\Pi_{\mathrm{homo}}(\mathcal{Q}_{\mathrm{vf\text- op}})\bigr)
      =
      L(\mathcal{Q}_{\mathrm{vf\text- op}}),
    \]
    for all \(\mathcal{Q}_{\mathrm{vf\text- op}}\).
  \item[\textnormal{(2)}] \textbf{Compatibility with Jacobiator-gap.}
    For every \(\mathcal{Q}_{\mathrm{vf\text- op}}\) we have
    \[
      \mathsf{JacGap}\bigl(
        \Pi_{\mathrm{homo}}(\mathcal{Q}_{\mathrm{vf\text- op}})
      \bigr)
      =
      \mathsf{JacGap}(\mathcal{Q}_{\mathrm{vf\text- op}}),
    \]
    and for every JacGap-free homotopy object
    \(\mathcal{Q}_{\mathrm{homo}}\in
     \mathsf{PFB\text- GNLA}_{\mathrm{homo}}^{(0)}\),
    the strictification \(\Pi_{\mathrm{strict}}(\mathcal{Q}_{\mathrm{homo}})\)
    satisfies
    \[
      \mathsf{JacGap}\bigl(
        \Pi_{\mathrm{strict}}(\mathcal{Q}_{\mathrm{homo}})
      \bigr)
      = 0.
    \]
  \item[\textnormal{(3)}] \textbf{Compatibility with Upgrade and
        Strictify.}
    In the JacGap-free regime, the following diagrams commute in the
    homotopy category:
    \[
      \mathsf{PFB\text- GNLA}_{\mathrm{strict}}
      \xrightarrow{\;\mathsf{Upgrade}\;}
      \mathsf{PFB\text- GNLA}_{\mathrm{homo}}^{(0)}
      \xrightarrow{\;\Pi_{\mathrm{strict}}\;}
      \mathsf{PFB\text- GNLA}_{\mathrm{strict}}
    \]
    is homotopic to the identity functor on
    \(\mathsf{PFB\text- GNLA}_{\mathrm{strict}}\), and
    \[
      \mathsf{PFB\text- GNLA}_{\mathrm{homo}}^{(0)}
      \xrightarrow{\;\Pi_{\mathrm{strict}}\;}
      \mathsf{PFB\text- GNLA}_{\mathrm{strict}}
      \xrightarrow{\;\mathsf{Upgrade}\;}
      \mathsf{PFB\text- GNLA}_{\mathrm{homo}}
    \]
    is homotopic to the inclusion of
    \(\mathsf{PFB\text- GNLA}_{\mathrm{homo}}^{(0)}\) into
    \(\mathsf{PFB\text- GNLA}_{\mathrm{homo}}\).
  \item[\textnormal{(4)}] \textbf{Law-level invariance.}
    For any vf-op object \(\mathcal{Q}_{\mathrm{vf\text- op}}\) with
    JacGap-free homotopy projection, the associated law-objects of the
    three regimes coincide:
    \[
      L(\mathcal{Q}_{\mathrm{vf\text- op}})
      =
      L\bigl(\Pi_{\mathrm{homo}}(\mathcal{Q}_{\mathrm{vf\text- op}})\bigr)
      =
      L\bigl(\Pi_{\mathrm{strict}}\circ
              \Pi_{\mathrm{homo}}(\mathcal{Q}_{\mathrm{vf\text- op}})
        \bigr),
    \]
    so that the passage vf-op \(\to\) homotopy \(\to\) strict does not
    change the underlying law-system in \(\GZLS\), only its
    representation regime.
\end{enumerate}
\end{proposition}



\subsection*{Proof of Proposition~\ref{prop:P6-projection}}


\begin{proof}
\textbf{Proof of (1): functorial projection to the homotopy regime.}
By definition, \(\Pi_{\mathrm{homo}}\) acts trivially on the underlying
principal bundle and connection, and simply forgets the vf-op data
\(\mathcal{O}_{\mathrm{vf\text- op}}\).  On morphisms, it discards the
component acting on \(\mathcal{O}_{\mathrm{vf\text- op}}\) and retains
the principal-bundle map and the \(L_\infty\)-homomorphism.  This
clearly preserves identities and composition, so
\(\Pi_{\mathrm{homo}}\) is a well-defined functor.

Since vf-op data are by construction “decorations” on top of the
homotopy GNLA structure that do not alter the underlying law-object in
\(\GZLS\), we have
\[
  L\bigl(\Pi_{\mathrm{homo}}(\mathcal{Q}_{\mathrm{vf\text- op}})\bigr)
  =
  L(\mathcal{Q}_{\mathrm{vf\text- op}})
\]
for every vf-op object.  This proves property~\textnormal{(1)}.



\medskip\noindent
\textbf{Proof of (2): compatibility with Jacobiator-gap.}
By definition,
\[
  \mathsf{JacGap}(\mathcal{Q})
  =
  \mathsf{JacGap}\bigl(L(\mathcal{Q})\bigr)
\]
for any PFB-GNLA object \(\mathcal{Q}\) in any regime.  Since
\(\Pi_{\mathrm{homo}}\) preserves the law-object, we immediately obtain
\[
  \mathsf{JacGap}\bigl(
    \Pi_{\mathrm{homo}}(\mathcal{Q}_{\mathrm{vf\text- op}})
  \bigr)
  =
  \mathsf{JacGap}\bigl(
    L\bigl(\Pi_{\mathrm{homo}}(\mathcal{Q}_{\mathrm{vf\text- op}})\bigr)
  \bigr)
  =
  \mathsf{JacGap}\bigl(L(\mathcal{Q}_{\mathrm{vf\text- op}})\bigr)
  =
  \mathsf{JacGap}(\mathcal{Q}_{\mathrm{vf\text- op}}).
\]

For a JacGap-free homotopy object
\(\mathcal{Q}_{\mathrm{homo}}\in
 \mathsf{PFB\text- GNLA}_{\mathrm{homo}}^{(0)}\), the strictification\\
\(\Pi_{\mathrm{strict}}(\mathcal{Q}_{\mathrm{homo}})\) is constructed
via \(\mathsf{Strictify}\) and Theorem~\ref{thm:GZ-OHU}, which
produces a strict model with vanishing Jacobiator-defect.  Thus
\[
  \mathsf{JacGap}\bigl(
    \Pi_{\mathrm{strict}}(\mathcal{Q}_{\mathrm{homo}})
  \bigr)
  = 0.
\]
This establishes property~\textnormal{(2)}.



\medskip\noindent
\textbf{Proof of (3): compatibility with Upgrade and Strictify.}
The first composite
\[
  \mathsf{PFB\text- GNLA}_{\mathrm{strict}}
  \xrightarrow{\;\mathsf{Upgrade}\;}
  \mathsf{PFB\text- GNLA}_{\mathrm{homo}}^{(0)}
  \xrightarrow{\;\Pi_{\mathrm{strict}}\;}
  \mathsf{PFB\text- GNLA}_{\mathrm{strict}}
\]
is, by definition, the composite
\(\mathsf{Strictify}\circ\mathsf{Upgrade}\), which is homotopic to
the identity on \(\mathsf{PFB\text- GNLA}_{\mathrm{strict}}\) by
Proposition~\ref{prop:P5-upgrade-strictify}.  This yields the first
part of \textnormal{(3)}.

For the second composite
\[
  \mathsf{PFB\text- GNLA}_{\mathrm{homo}}^{(0)}
  \xrightarrow{\;\Pi_{\mathrm{strict}}\;}
  \mathsf{PFB\text- GNLA}_{\mathrm{strict}}
  \xrightarrow{\;\mathsf{Upgrade}\;}
  \mathsf{PFB\text- GNLA}_{\mathrm{homo}},
\]
we note that
\(\mathsf{Upgrade}\circ\Pi_{\mathrm{strict}}
 =\mathsf{Upgrade}\circ\mathsf{Strictify}\) is homotopic to the
identity on the JacGap-free subcategory by the second equivalence in
Proposition~\ref{prop:P5-upgrade-strictify}.  Since
\(\mathsf{PFB\text- GNLA}_{\mathrm{homo}}^{(0)}\) is a full subcategory
of \(\mathsf{PFB\text- GNLA}_{\mathrm{homo}}\), this composite is
homotopic to the inclusion
\[
  \mathsf{PFB\text- GNLA}_{\mathrm{homo}}^{(0)}
  \hookrightarrow
  \mathsf{PFB\text- GNLA}_{\mathrm{homo}}.
\]
This proves property~\textnormal{(3)}.



\medskip\noindent
\textbf{Proof of (4): law-level invariance.}
Let \(\mathcal{Q}_{\mathrm{vf\text- op}}\in
\mathsf{PFB\text- GNLA}_{\mathrm{vf\text- op}}\) be such that
\(\Pi_{\mathrm{homo}}(\mathcal{Q}_{\mathrm{vf\text- op}})\) is
JacGap-free and hence lies in the domain of \(\Pi_{\mathrm{strict}}\).
From property~\textnormal{(1)} we have
\[
  L\bigl(\Pi_{\mathrm{homo}}(\mathcal{Q}_{\mathrm{vf\text- op}})\bigr)
  =
  L(\mathcal{Q}_{\mathrm{vf\text- op}}).
\]
By construction of \(\Pi_{\mathrm{strict}}\) via strictification,
\[
  L\bigl(\Pi_{\mathrm{strict}}\circ
          \Pi_{\mathrm{homo}}(\mathcal{Q}_{\mathrm{vf\text- op}})
    \bigr)
\]
is law-equivalent (indeed, law-isomorphic in \(\GZLS\)) to
\(L\bigl(\Pi_{\mathrm{homo}}(\mathcal{Q}_{\mathrm{vf\text- op}})\bigr)\).
Composing these equalities gives
\[
  L(\mathcal{Q}_{\mathrm{vf\text- op}})
  =
  L\bigl(\Pi_{\mathrm{homo}}(\mathcal{Q}_{\mathrm{vf\text- op}})\bigr)
  =
  L\bigl(\Pi_{\mathrm{strict}}\circ
          \Pi_{\mathrm{homo}}(\mathcal{Q}_{\mathrm{vf\text- op}})
    \bigr),
\]
which is exactly the desired law-level invariance.  This proves
property~\textnormal{(4)} and completes the proof of
Proposition~\ref{prop:P6-projection}.



\end{proof}

\section{Proposition P8: Four PFB-GNLA Versions as Law-Level Charts}
\label{app:II-P8}

\subsection*{Set-up: four regimes and their representation categories}


In this appendix we expand on the schematic Proposition~P8 from the
main text, whose goal is to formalise the idea that the four PFB-GNLA
versions act as different ``charts'' for the same PL--PI law-system at
the law-level.

We work with the following ingredients.

\begin{itemize}
  \item A fixed GZ-admissible PL--PI law-system whose law-level algebra
        is denoted by
        \[
          \GAlgebra_{\mathrm{orig}}.
        \]
        As in Technical Appendix~I, Section~\ref{app:I-P7}, this
        algebra encodes the primitive law-operations and composition
        rules of the system (e.g.\ a QFT-type system).
  \item The category of local law-level representations
        \(\mathsf{Rep}_{\mathrm{loc}}(\GAlgebra_{\mathrm{orig}})\) as
        in \\ Proposition~\ref{prop:P7-law-algebra-PFB-GNLA}.
  \item Three PFB-GNLA representation categories over the same base
        \(M\):
        \begin{itemize}
          \item
            \(\mathsf{PFB\text- GNLA}^{\mathrm{rep}}_{\mathrm{strict}}
              (\GAlgebra_{\mathrm{orig}})\):
            strict PFB-GNLA representation data (Appendix~I,
            Section~\ref{app:I-P7});
          \item
            \(\mathsf{PFB\text- GNLA}^{\mathrm{rep}}_{\mathrm{homo}}
              (\GAlgebra_{\mathrm{orig}})\):
            homotopy PFB-GNLA representation data, \\obtained by
            applying the Upgrade functor of
            Proposition~\ref{prop:P5-upgrade-strictify} to strict
            data and allowing more general \(L_\infty\)-type fibres;
          \item
            \(\mathsf{PFB\text- GNLA}^{\mathrm{rep}}_{\mathrm{vf\text- op}}(\GAlgebra_{\mathrm{orig}})\):
            vf-op PFB-GNLA representation data, i.e.\ homotopy
            PFB-GNLA objects equipped with additional operator-path /
            vector-field data as in Appendix~\ref{app:II-P6}.
        \end{itemize}
\end{itemize}

For each of the three PFB-GNLA representation categories we have a
forgetful functor to the underlying PFB-GNLA object (dropping the
explicit \(\GAlgebra_{\mathrm{orig}}\)-action) and then a law-space
representation functor
\[
  L(-) :
  \mathsf{PFB\text- GNLA}_{\mathrm{regime}}
  \longrightarrow
  \GZLS,
\]
which sends a PFB-GNLA object to its associated law-object in
\(\GZLS\).  At the representation level, we combine this with the
representation functor of Proposition~\ref{prop:P7-law-algebra-PFB-GNLA}
to obtain
\[
  \mathcal{R}_{\mathrm{regime}} :
  \mathsf{PFB\text- GNLA}^{\mathrm{rep}}_{\mathrm{regime}}
    (\GAlgebra_{\mathrm{orig}})
  \longrightarrow
  \mathsf{Rep}_{\mathrm{loc}}(\GAlgebra_{\mathrm{orig}}),
\]
for \(\mathrm{regime}\in\{\mathrm{strict},\mathrm{homo},
\mathrm{vf-op}\}\).

We also have the projection functors constructed in
Appendix~\ref{app:II-P6}:
\[
  \Pi_{\mathrm{homo}} :
  \mathsf{PFB\text- GNLA}_{\mathrm{vf\text- op}}
  \to
  \mathsf{PFB\text- GNLA}_{\mathrm{homo}},
\]
\[
  \Pi_{\mathrm{strict}} :
  \mathsf{PFB\text- GNLA}_{\mathrm{homo}}^{(0)}
  \dashrightarrow
  \mathsf{PFB\text- GNLA}_{\mathrm{strict}},
\]
and the Upgrade/Strictify functors of
Proposition~\ref{prop:P5-upgrade-strictify}.  These allow us to move
between strict, homotopy, and vf-op regimes without changing the
underlying law-object in \(\GZLS\).



\subsection*{Four regimes as law-level charts}


Motivated by the chart picture, we introduce the following shorthand
notation:

\begin{itemize}
  \item \(\mathsf{Chart}_{\mathrm{orig}}\): the full subcategory of
        \(\mathsf{Rep}_{\mathrm{loc}}(\GAlgebra_{\mathrm{orig}})\)
        consisting of local law-level representations that satisfy the
        GZ-admissibility and PL--PI regularity conditions.
  \item \(\mathsf{Chart}_{\mathrm{strict}}\): the essential image of
        \(\mathsf{PFB\text- GNLA}^{\mathrm{rep}}_{\mathrm{strict}}\)
        under \(\mathcal{R}_{\mathrm{strict}}\).
  \item \(\mathsf{Chart}_{\mathrm{homo}}\): the essential image of
        \(\mathsf{PFB\text- GNLA}^{\mathrm{rep}}_{\mathrm{homo}}\) under
        \(\mathcal{R}_{\mathrm{homo}}\) restricted to the JacGap-free
        subcategory.
  \item \(\mathsf{Chart}_{\mathrm{vf\text- op}}\): the essential image
        of
        \(\mathsf{PFB\text- GNLA}^{\mathrm{rep}}_{\mathrm{vf\text- op}}\)
        under \(\mathcal{R}_{\mathrm{vf\text- op}}\), again restricted
        to the JacGap-free regime when needed.
\end{itemize}

Intuitively, \(\mathsf{Chart}_{\mathrm{orig}}\) is the ``coordinate
chart'' in which laws are written directly in terms of
\(\GAlgebra_{\mathrm{orig}}\), whereas
\(\mathsf{Chart}_{\mathrm{strict}}\),
\(\mathsf{Chart}_{\mathrm{homo}}\), and
\(\mathsf{Chart}_{\mathrm{vf\text- op}}\) are three increasingly rich
geometric charts for the same law-system: strict GNLA bundle chart,
homotopy PFB-GNLA chart, and vf-op chart that explicitly encodes
operator-path information.



\subsection*{Statement of Proposition P8}


\begin{proposition}[Four PFB-GNLA versions as law-level charts]
\label{prop:P8-four-versions-charts}
For a fixed GZ-admissible PL--PI law-system with law-level algebra
\(\GAlgebra_{\mathrm{orig}}\) as above, the following hold.

\begin{enumerate}
  \item[\textnormal{(1)}] \textbf{Strict PFB-GNLA chart.}
    The functor
    \[
      \mathcal{R}_{\mathrm{strict}} :
      \mathsf{PFB\text- GNLA}^{\mathrm{rep}}_{\mathrm{strict}}
        (\GAlgebra_{\mathrm{orig}})
      \longrightarrow
      \mathsf{Rep}_{\mathrm{loc}}(\GAlgebra_{\mathrm{orig}})
    \]
    induces an equivalence of categories between the gauge-quotient
    of \(\mathsf{PFB\text- GNLA}^{\mathrm{rep}}_{\mathrm{strict}}\)
    and \(\mathsf{Chart}_{\mathrm{strict}}\), and
    \(\mathsf{Chart}_{\mathrm{strict}}\) is a full subcategory of
    \(\mathsf{Chart}_{\mathrm{orig}}\).
  \item[\textnormal{(2)}] \textbf{Homotopy PFB-GNLA chart.}
    On the JacGap-free regime, the Upgrade and Strictify functors of
    Proposition~\ref{prop:P5-upgrade-strictify} induce an equivalence
    between
    \(\mathsf{PFB\text- GNLA}^{\mathrm{rep}}_{\mathrm{strict}}\) and
    the JacGap-free part of
    \(\mathsf{PFB\text- GNLA}^{\mathrm{rep}}_{\mathrm{homo}}\), and
    the corresponding images
    \(\mathsf{Chart}_{\mathrm{strict}}\) and
    \(\mathsf{Chart}_{\mathrm{homo}}\) are naturally equivalent as
    full subcategories of \(\mathsf{Chart}_{\mathrm{orig}}\).
  \item[\textnormal{(3)}] \textbf{vf-op PFB-GNLA chart.}
    The projection functor
    \(\Pi_{\mathrm{homo}} : \\ \mathsf{PFB\text- GNLA}_{\mathrm{vf\text- op}}
      \to \mathsf{PFB\text- GNLA}_{\mathrm{homo}}\)
    preserves law-objects and JacGap, so on the JacGap-free regime the
    vf-op representation category maps onto a subcategory that is
    equivalent (up to forgetting vf-op decorations) to
    \(\mathsf{Chart}_{\mathrm{homo}}\).  Thus
    \(\mathsf{Chart}_{\mathrm{vf\text- op}}\) is equivalent to
    \(\mathsf{Chart}_{\mathrm{homo}}\) and hence to
    \(\mathsf{Chart}_{\mathrm{strict}}\).
  \item[\textnormal{(4)}] \textbf{Four charts for one law-system.}
    Summarising, there is a commutative diagram (up to natural
    isomorphism)
    \[
      \mathsf{Chart}_{\mathrm{orig}}
      \;\;\supset\;\;
      \mathsf{Chart}_{\mathrm{strict}}
      \;\;\simeq\;\;
      \mathsf{Chart}_{\mathrm{homo}}
      \;\;\simeq\;\;
      \mathsf{Chart}_{\mathrm{vf\text- op}},
    \]
    in which each of the three PFB-GNLA charts provides a different
    but equivalent \\representation-theoretic ``coordinate chart'' for
    the same PL--PI law-system at the law-level.  Transition functors
    between these charts are induced by the Upgrade, Strictify, and
    projection functors, and they commute with the law-space
    representation functor \(L(-)\).
\end{enumerate}
\end{proposition}



\subsection*{Proof of Proposition~\ref{prop:P8-four-versions-charts}}


\begin{proof}
\textbf{Proof of (1): strict PFB-GNLA chart.}
This is essentially Proposition~\ref{prop:P7-law-algebra-PFB-GNLA}.
There we constructed a functor
\[
  \mathcal{R}_{\mathrm{strict}} :
  \mathsf{PFB\text- GNLA}^{\mathrm{rep}}_{\mathrm{strict}}
    (\GAlgebra_{\mathrm{orig}})
  \longrightarrow
  \mathsf{Rep}_{\mathrm{loc}}(\GAlgebra_{\mathrm{orig}})
\]
and showed that, modulo gauge-transformations on the principal bundle
that act trivially on the associated representation, this functor is
fully faithful and essentially surjective onto the subcategory of
geometrically realisable law-level representations.

By definition,
\(\mathsf{Chart}_{\mathrm{strict}}\) is exactly this essential image,
and it is a full subcategory of \(\mathsf{Chart}_{\mathrm{orig}}\)
since morphisms in \(\mathsf{Chart}_{\mathrm{strict}}\) are just those
intertwiners in \(\mathsf{Rep}_{\mathrm{loc}}\) that arise from
PFB-GNLA morphisms.  This proves item~\textnormal{(1)}.



\medskip\noindent
\textbf{Proof of (2): homotopy PFB-GNLA chart.}
In Proposition~\ref{prop:P5-upgrade-strictify} we showed that on the
JacGap-free regime, the Upgrade and Strictify constructions establish
an equivalence of categories between
\(\mathsf{PFB\text- GNLA}_{\mathrm{strict}}\) and the JacGap-free part
of \(\mathsf{PFB\text- GNLA}_{\mathrm{homo}}\).  Passing to the
representation categories and composing with
\(\mathcal{R}_{\mathrm{strict}}\) and \(\mathcal{R}_{\mathrm{homo}}\),
we obtain an equivalence between
\(\mathsf{Chart}_{\mathrm{strict}}\) and the JacGap-free part of
\(\mathsf{Chart}_{\mathrm{homo}}\).

Moreover, the law-space functor \(L(-)\) is invariant under Upgrade
and Strictify in this regime, so these equivalences are realised
inside \(\mathsf{Chart}_{\mathrm{orig}}\) as equivalences of full
subcategories: any object in
\(\mathsf{Chart}_{\mathrm{homo}}\) can be strictified back to an
object in \(\mathsf{Chart}_{\mathrm{strict}}\) without changing its
law-level object, and conversely the Upgrade of a strict object
remains in the same law-level orbit.  This yields item
\textnormal{(2)}.



\medskip\noindent
\textbf{Proof of (3): vf-op PFB-GNLA chart.}
From Proposition~\ref{prop:P6-projection} we know that
\(\Pi_{\mathrm{homo}} : \mathsf{PFB\text- GNLA}_{\mathrm{vf\text- op}}
\to \mathsf{PFB\text- GNLA}_{\mathrm{homo}}\) is a functor which
preserves law-objects and JacGap:
\[
  L\bigl(\Pi_{\mathrm{homo}}(\mathcal{Q}_{\mathrm{vf\text- op}})\bigr)
  = L(\mathcal{Q}_{\mathrm{vf\text- op}}),
  \qquad
  \mathsf{JacGap}\bigl(
    \Pi_{\mathrm{homo}}(\mathcal{Q}_{\mathrm{vf\text- op}})
  \bigr)
  =
  \mathsf{JacGap}(\mathcal{Q}_{\mathrm{vf\text- op}}).
\]

Restricting to the JacGap-free regime, the homotopy projection
\(\Pi_{\mathrm{homo}}\) drops only the vf-op decorations
\(\mathcal{O}_{\mathrm{vf\text- op}}\) and leaves the homotopy GNLA
fibre and the principal bundle unchanged.  Hence, the composition
\[
  \mathcal{R}_{\mathrm{homo}}\circ\Pi_{\mathrm{homo}} :
  \mathsf{PFB\text- GNLA}^{\mathrm{rep}}_{\mathrm{vf\text- op}}
    (\GAlgebra_{\mathrm{orig}})
  \longrightarrow
  \mathsf{Rep}_{\mathrm{loc}}(\GAlgebra_{\mathrm{orig}})
\]
identifies the vf-op chart with the homotopy chart at the law-level:
they have the same essential image and the same morphisms once vf-op
decorations are forgotten in a compatible way.

Thus \(\mathsf{Chart}_{\mathrm{vf\text- op}}\) is equivalent to
\(\mathsf{Chart}_{\mathrm{homo}}\), and hence, by \textnormal{(2)}, to
\(\mathsf{Chart}_{\mathrm{strict}}\) as well.  This proves item
\textnormal{(3)}.



\medskip\noindent
\textbf{Proof of (4): four charts for one law-system.}
Putting items \textnormal{(1)}–\textnormal{(3)} together, we see that:
\begin{itemize}
  \item \(\mathsf{Chart}_{\mathrm{strict}}\) is a full subcategory of
        \(\mathsf{Chart}_{\mathrm{orig}}\) and is equivalent to the
        gauge-quotient of strict PFB-GNLA representations;
  \item \(\mathsf{Chart}_{\mathrm{homo}}\) (in the JacGap-free
        regime) is equivalent to
        \(\mathsf{Chart}_{\mathrm{strict}}\);
  \item \(\mathsf{Chart}_{\mathrm{vf\text- op}}\) is equivalent to
        \(\mathsf{Chart}_{\mathrm{homo}}\), again in the JacGap-free
        regime.
\end{itemize}
Hence there is a chain of equivalences
\[
  \mathsf{Chart}_{\mathrm{strict}}
  \;\simeq\;
  \mathsf{Chart}_{\mathrm{homo}}
  \;\simeq\;
  \mathsf{Chart}_{\mathrm{vf\text- op}},
\]
all sitting inside \(\mathsf{Chart}_{\mathrm{orig}}\).  Each of these
three PFB-GNLA charts carries additional geometric or dynamical
structure (strict GNLA fibres, \(L_\infty\)-homotopies, vf-op
decorations), but at the law-level they represent the same PL--PI
law-system.

Moreover, the transition functors between them (Upgrade, Strictify,
\(\Pi_{\mathrm{homo}}\), \(\Pi_{\mathrm{strict}}\)) commute with the
law-space representation functor \(L(-)\) and preserve the law-level
Jacobiator-gap, except for the strictification step which drives the
gap to zero.  This justifies the statement that the four versions
“act as different law-level charts for the same PL--PI law-system” and
completes the proof of
Proposition~\ref{prop:P8-four-versions-charts}.



\end{proof}

\section{Proposition P2: Compatibility with Layer-Wise Bianchi Identities}
\label{app:II-P2}

\subsection*{Geometric set-up and notation}


In this appendix we make precise the schematic Proposition~P2 from the
main text and prove the compatibility of the triple-layer curvature of
\(\GZTLC\) with the layer-wise Bianchi identities.

We keep the notation of Technical Appendix~I, Section~\ref{app:I-P3}.
In particular:

\begin{itemize}
  \item \(P\to M\) is a principal \(G_{\mathrm{geom}}\)-bundle with
        connection \(\Conn_p\), given by a connection one-form
        \(\omega_{\mathrm{geom}}\in\Omega^1(P,\mathfrak{g}_{\mathrm{geom}})\)
        and curvature
        \(\mathcal{F}_{\mathrm{geom}}\in
         \Omega^2(P,\mathfrak{g}_{\mathrm{geom}})\).
  \item \(M\times\Xi\to M\) carries a parameter connection
        \(B_{\mathrm{par}}\), given by a one-form
        \(\omega_{\mathrm{par}}\in
         \Omega^1(M\times\Xi,\mathfrak{g}_{\mathrm{par}})\) and curvature
        \(\mathcal{F}_{\mathrm{par}}\in
         \Omega^2(M\times\Xi,\mathfrak{g}_{\mathrm{par}})\).
  \item \(\pi_{\mathrm{law}}:\mathcal{L}\to M\times\Xi\) is the
        law-space bundle with law-space connection \(\GZLOC\), given
        by a one-form
        \(\omega_{\mathrm{law}}\in
         \Omega^1(\mathcal{L},\mathfrak{g}_{\mathrm{law}})\) and curvature
        \(\GZLCurv\in\Omega^2(\mathcal{L},\mathfrak{g}_{\mathrm{law}})\).
  \item The triple-layer bundle
        \(\widehat{\Bundle}
          =P\times_M(M\times\Xi)\times_{M\times\Xi}\mathcal{L}\)
        is endowed with the triple-layer connection \(\GZTLC\) whose
        connection one-form is
        \[
          \omega_{\mathrm{tot}}
          :=
          \pr_{\mathrm{geom}}^*\omega_{\mathrm{geom}}
          +
          \pr_{\mathrm{par}}^*\omega_{\mathrm{par}}
          +
          \pr_{\mathrm{law}}^*\omega_{\mathrm{law}}
          \;\in\;
          \Omega^1(\widehat{\Bundle},\mathfrak{g}_{\mathrm{tot}}),
        \]
        where
        \(\mathfrak{g}_{\mathrm{tot}}
          =\mathfrak{g}_{\mathrm{geom}}
           \oplus\mathfrak{g}_{\mathrm{par}}
           \oplus\mathfrak{g}_{\mathrm{law}}\).
\end{itemize}

The curvature of \(\GZTLC\) is
\[
  \mathcal{F}_{\GZTLC}
  :=
  \mathrm{d}\omega_{\mathrm{tot}}
  +
  \tfrac12[\omega_{\mathrm{tot}}\wedge\omega_{\mathrm{tot}}]
  \;\in\;
  \Omega^2(\widehat{\Bundle},\mathfrak{g}_{\mathrm{tot}}),
\]
and we denote by
\[
  D_{\mathcal{A}}
  :=
  \mathrm{d}
  +
  [\omega_{\mathrm{tot}},\cdot]
\]
the associated covariant derivative on
\(\mathfrak{g}_{\mathrm{tot}}\)-valued forms.  Likewise, we write
\(D_{\mathrm{geom}}\), \(D_{\mathrm{par}}\), and \(D_{\mathrm{law}}\)
for the covariant derivatives associated to
\(\omega_{\mathrm{geom}}\), \(\omega_{\mathrm{par}}\), and
\(\omega_{\mathrm{law}}\), respectively.



For convenience we decompose \(\mathcal{F}_{\GZTLC}\) into its
layer-wise and mixed components.  Writing
\[
  \mathcal{F}_{\GZTLC}
  =
  \mathcal{F}_{\mathrm{geom}}^{\uparrow}
  +
  \mathcal{F}_{\mathrm{par}}^{\uparrow}
  +
  \GZLCurv^{\uparrow}
  +
  \mathcal{F}_{\mathrm{mix}},
\]
we let
\begin{itemize}
  \item \(\mathcal{F}_{\mathrm{geom}}^{\uparrow}
          := \pr_{\mathrm{geom}}^*\mathcal{F}_{\mathrm{geom}}\),
  \item \(\mathcal{F}_{\mathrm{par}}^{\uparrow}
          := \pr_{\mathrm{par}}^*\mathcal{F}_{\mathrm{par}}\),
  \item \(\GZLCurv^{\uparrow}
          := \pr_{\mathrm{law}}^*\GZLCurv\),
  \item \(\mathcal{F}_{\mathrm{mix}}\) collects all mixed terms arising
        from the wedge commutator
        \([\omega_{\mathrm{tot}}\wedge\omega_{\mathrm{tot}}]\) that take
        values in \(\mathfrak{g}_{\mathrm{tot}}\) but involve different
        layers simultaneously (e.g.\ ``geom–par'', ``par–law'' and
        ``geom–law'' components).
\end{itemize}



\subsection*{Bianchi identities for the three layers}


We assume that the three layer-wise connections satisfy their usual
Bianchi identities:
\begin{align}
  D_{\mathrm{geom}} \mathcal{F}_{\mathrm{geom}} &= 0,
  \label{eq:geom-Bianchi} \\
  D_{\mathrm{par}}  \mathcal{F}_{\mathrm{par}}  &= 0,
  \label{eq:par-Bianchi} \\
  D_{\mathrm{law}}  \GZLCurv                 &= 0.
  \label{eq:law-Bianchi}
\end{align}
These are the standard Bianchi identities for principal-bundle and
Ehresmann-type connections, and the existence of the triple-layer
connection \(\GZTLC\) is assumed to be compatible with these in the
sense detailed in Proposition~\ref{prop:triple-layer-connection}.



\subsection*{Statement of Proposition P2}


\begin{proposition}[Compatibility with layer-wise Bianchi identities]
\label{prop:P2-Bianchi-compatibility}
Let \(\GZTLC\) be the triple-layer connection on \(\widehat{\Bundle}\)
constructed in Proposition~\ref{prop:triple-layer-connection}, and let
\(\mathcal{F}_{\GZTLC}\) and \(D_{\mathcal{A}}\) be its curvature and
covariant derivative.  Assume that the layer-wise connections satisfy
the Bianchi identities \eqref{eq:geom-Bianchi}–\eqref{eq:law-Bianchi}.
Then:

\begin{enumerate}
  \item[\textnormal{(1)}] \textbf{Unified Bianchi identity.}
    The triple-layer curvature satisfies the unified Bianchi identity
    \[
      D_{\mathcal{A}}\mathcal{F}_{\GZTLC} = 0
      \quad\text{in }
      \Omega^3(\widehat{\Bundle},\mathfrak{g}_{\mathrm{tot}}).
    \]
  \item[\textnormal{(2)}] \textbf{Component-wise compatibility.}
    Writing \(\mathcal{F}_{\GZTLC}\) as the sum of diagonal and mixed
    components, there is a decomposition of the unified Bianchi
    identity into:
    \begin{align*}
      \pr_{\mathrm{geom}}^*\bigl(D_{\mathrm{geom}}
        \mathcal{F}_{\mathrm{geom}}\bigr)
      &+
      \pr_{\mathrm{par}}^*\bigl(D_{\mathrm{par}}
        \mathcal{F}_{\mathrm{par}}\bigr)
      +
      \pr_{\mathrm{law}}^*\bigl(D_{\mathrm{law}}
        \GZLCurv\bigr) \\
      &+
      \Bigl(
        D_{\mathcal{A}}\mathcal{F}_{\mathrm{mix}}
        + \text{mixed covariant-derivative terms}
      \Bigr)
      = 0,
    \end{align*}
    where the first line vanishes by the layer-wise Bianchi identities
    and the second line expresses the compatibility conditions among
    the mixed components.  In particular, any mixed block of the
    curvature satisfies a coupled Bianchi-type identity of the form
    \[
      D_{\mathrm{geom}}\mathcal{F}^{\mathrm{geom-par}}
      +
      D_{\mathrm{par}}\mathcal{F}^{\mathrm{par-geom}}
      +
      D_{\mathrm{law}}\mathcal{F}^{\mathrm{law-mix}}
      = 0,
    \]
    and similarly for the other mixed blocks, where the covariant
    derivatives are taken with respect to the appropriate layer-wise
    connections and the triple-layer connection.
\end{enumerate}
\end{proposition}



\subsection*{Proof of Proposition~\ref{prop:P2-Bianchi-compatibility}}


\begin{proof}
We proceed in three steps.

\medskip\noindent
\textbf{Step 1: Bianchi identity for the sum connection.}
The one-form
\[
  \omega_{\mathrm{tot}}
  =
  \pr_{\mathrm{geom}}^*\omega_{\mathrm{geom}}
  +
  \pr_{\mathrm{par}}^*\omega_{\mathrm{par}}
  +
  \pr_{\mathrm{law}}^*\omega_{\mathrm{law}}
\]
is a connection one-form on a principal bundle with structure group
\(G_{\mathrm{geom}}\times G_{\mathrm{par}}\times G_{\mathrm{law}}\)
(or on a suitable extension whose Lie algebra is
\(\mathfrak{g}_{\mathrm{tot}}\)).  It is a standard fact of principal
bundle theory that the Bianchi identity
\[
  D_{\mathcal{A}}\mathcal{F}_{\GZTLC}
  =
  \mathrm{d}\mathcal{F}_{\GZTLC}
  +
  [\omega_{\mathrm{tot}}\wedge\mathcal{F}_{\GZTLC}]
  = 0
\]
holds for any connection \(\omega_{\mathrm{tot}}\), independently of
how it is decomposed into pieces; it is a consequence of the graded
Jacobi identity in the universal enveloping algebra of
\(\mathfrak{g}_{\mathrm{tot}}\).

Thus, at the level of the total triple-layer connection, we have
\[
  D_{\mathcal{A}}\mathcal{F}_{\GZTLC}=0
\]
without further assumptions.  This already proves item
\textnormal{(1)} in the proposition.



\medskip\noindent
\textbf{Step 2: Decomposition into layer-wise and mixed components.}
To obtain the more refined statement in item~\textnormal{(2)}, we
decompose \(\omega_{\mathrm{tot}}\) and
\(\mathcal{F}_{\GZTLC}\) according to the direct sum
\(\mathfrak{g}_{\mathrm{tot}}
  = \mathfrak{g}_{\mathrm{geom}}
    \oplus\mathfrak{g}_{\mathrm{par}}
    \oplus\mathfrak{g}_{\mathrm{law}}\).

Using
\[
  \omega_{\mathrm{tot}}
  =
  \omega_{\mathrm{geom}}^{\uparrow}
  +
  \omega_{\mathrm{par}}^{\uparrow}
  +
  \omega_{\mathrm{law}}^{\uparrow},
\]
where \(\omega_{\mathrm{geom}}^{\uparrow}=\pr_{\mathrm{geom}}^*\omega_{\mathrm{geom}}\),
etc., a straightforward expansion gives
\begin{align*}
  \mathcal{F}_{\GZTLC}
  &= \mathrm{d}\omega_{\mathrm{tot}}
     + \tfrac12[\omega_{\mathrm{tot}}\wedge\omega_{\mathrm{tot}}] \\
  &= \pr_{\mathrm{geom}}^*\bigl(
       \mathrm{d}\omega_{\mathrm{geom}}
       + \tfrac12[\omega_{\mathrm{geom}}\wedge\omega_{\mathrm{geom}}]
     \bigr)
    +
     \pr_{\mathrm{par}}^*\bigl(
       \mathrm{d}\omega_{\mathrm{par}}
       + \tfrac12[\omega_{\mathrm{par}}\wedge\omega_{\mathrm{par}}]
     \bigr) \\
  &\quad
    +
    \pr_{\mathrm{law}}^*\bigl(
       \mathrm{d}\omega_{\mathrm{law}}
       + \tfrac12[\omega_{\mathrm{law}}\wedge\omega_{\mathrm{law}}]
     \bigr)
    +
    \mathcal{F}_{\mathrm{mix}}.
\end{align*}
Recognising the layer-wise curvatures, this is
\[
  \mathcal{F}_{\GZTLC}
  =
  \mathcal{F}_{\mathrm{geom}}^{\uparrow}
  +
  \mathcal{F}_{\mathrm{par}}^{\uparrow}
  +
  \GZLCurv^{\uparrow}
  +
  \mathcal{F}_{\mathrm{mix}}.
\]

Applying the total covariant derivative \(D_{\mathcal{A}}\) to this
decomposition, we get
\begin{align*}
  0
  &= D_{\mathcal{A}}\mathcal{F}_{\GZTLC} \\
  &= D_{\mathcal{A}}\bigl(
       \mathcal{F}_{\mathrm{geom}}^{\uparrow}
       +
       \mathcal{F}_{\mathrm{par}}^{\uparrow}
       +
       \GZLCurv^{\uparrow}
       +
       \mathcal{F}_{\mathrm{mix}}
     \bigr) \\
  &= D_{\mathcal{A}}\mathcal{F}_{\mathrm{geom}}^{\uparrow}
     +
     D_{\mathcal{A}}\mathcal{F}_{\mathrm{par}}^{\uparrow}
     +
     D_{\mathcal{A}}\GZLCurv^{\uparrow}
     +
     D_{\mathcal{A}}\mathcal{F}_{\mathrm{mix}}.
\end{align*}
Each of these terms can be rewritten in terms of the layer-wise
covariant derivatives plus mixed commutator contributions.  For
instance,
\begin{align*}
  D_{\mathcal{A}}\mathcal{F}_{\mathrm{geom}}^{\uparrow}
  &= \mathrm{d}\mathcal{F}_{\mathrm{geom}}^{\uparrow}
     +
     [\omega_{\mathrm{tot}}\wedge\mathcal{F}_{\mathrm{geom}}^{\uparrow}] \\
  &= \pr_{\mathrm{geom}}^*\bigl(
       D_{\mathrm{geom}}\mathcal{F}_{\mathrm{geom}}
     \bigr)
     +
     \bigl([\omega_{\mathrm{par}}^{\uparrow}
            + \omega_{\mathrm{law}}^{\uparrow}
            \;\wedge\;
            \mathcal{F}_{\mathrm{geom}}^{\uparrow}]\bigr),
\end{align*}
and similarly for the parameter and law contributions.  The first term
in each such expression is the pullback of the layer-wise Bianchi
covariant derivative (which vanishes by
\eqref{eq:geom-Bianchi}–\eqref{eq:law-Bianchi}); the remaining terms
are precisely the mixed covariant-derivative contributions that we
group into the ``mixed'' part of the unified Bianchi identity.

Collecting the terms that are direct pullbacks from the base, we
obtain
\begin{align*}
  0
  &= \pr_{\mathrm{geom}}^*(D_{\mathrm{geom}}\mathcal{F}_{\mathrm{geom}})
     +
     \pr_{\mathrm{par}}^*(D_{\mathrm{par}}\mathcal{F}_{\mathrm{par}})
     +
     \pr_{\mathrm{law}}^*(D_{\mathrm{law}}\GZLCurv) \\
  &\quad
     +
     \Bigl(
       D_{\mathcal{A}}\mathcal{F}_{\mathrm{mix}}
       +
       \text{mixed commutator terms coming from }
       D_{\mathcal{A}}\mathcal{F}_{\mathrm{geom}}^{\uparrow},
       D_{\mathcal{A}}\mathcal{F}_{\mathrm{par}}^{\uparrow},
       D_{\mathcal{A}}\GZLCurv^{\uparrow}
     \Bigr).
\end{align*}
Using the layer-wise Bianchi identities, the first line vanishes,
leaving only the mixed part:
\[
  D_{\mathcal{A}}\mathcal{F}_{\mathrm{mix}}
  +
  \text{(mixed commutator terms)}
  = 0.
\]
This is the abstract form of the coupled Bianchi identities for the
mixed components, as indicated in the statement of the proposition.



\medskip\noindent
\textbf{Step 3: Component-wise interpretation.}
To make the last statement in item~\textnormal{(2)} more concrete,
one may block-decompose \(\mathcal{F}_{\mathrm{mix}}\) into components
taking values in the off-diagonal parts of \(\mathfrak{g}_{\mathrm{tot}}\),
for example
\[
  \mathcal{F}_{\mathrm{mix}}
  =
  \mathcal{F}^{\mathrm{geom-par}}
  +
  \mathcal{F}^{\mathrm{par-law}}
  +
  \mathcal{F}^{\mathrm{geom-law}}
  + \dots,
\]
where each block couples two (or more) layers.  The mixed commutator
terms appearing in Step~2 can then be grouped according to these
blocks, yielding identities of the schematic form
\[
  D_{\mathrm{geom}}\mathcal{F}^{\mathrm{geom-par}}
  +
  D_{\mathrm{par}}\mathcal{F}^{\mathrm{par-geom}}
  +
  D_{\mathrm{law}}\mathcal{F}^{\mathrm{law-mix}}
  = 0,
\]
and analogues obtained by permuting the roles of the three layers.
While the precise block expressions depend on the detailed choice of
\(\mathfrak{g}_{\mathrm{tot}}\) and the representation of the triple
structure, the important point is that all such mixed Bianchi-type
equations are simply the \(\mathfrak{g}_{\mathrm{tot}}\)-component-wise
expression of the single structural identity
\(D_{\mathcal{A}}\mathcal{F}_{\GZTLC}=0\).

This completes the proof of
Proposition~\ref{prop:P2-Bianchi-compatibility}.



\end{proof}



\chapter{Technical Appendix III: Foundational Theorems and External References}


\section[Proposition P1: From Object-Level Geometry to Law-Objects in GZ-LS]{Proposition P1: From Object-Level Geometry to \\ Law-Objects in \GZLS}
\label{app:III-P1}

\subsection*{Set-up: the category of geometric objects}


In this appendix we provide a more systematic construction of the
functor
\[
  \Phi_{\mathrm{geom}}:\mathsf{GeomObj}\longrightarrow\GZLS
\]
alluded to in the main text, and prove the basic structural
properties stated in Proposition~P1.

We begin by specifying the category \(\mathsf{GeomObj}\) of
\emph{object-level geometric data}.  The precise axioms can be tuned
to the application at hand (Riemannian geometry, gauge theory, QFT,
etc.), but for the purposes of this appendix we adopt the following
general framework.

\begin{definition}[Category \(\mathsf{GeomObj}\)]
An object of \(\mathsf{GeomObj}\) is a tuple
\[
  \mathcal{G}
  =
  \bigl(M,\;E\to M,\;\nabla,\;\mathcal{S}\bigr),
\]
where:
\begin{enumerate}
  \item \(M\) is a smooth manifold (space, spacetime, or configuration
        space).
  \item \(E\to M\) is a smooth vector bundle or associated bundle to a
        principal bundle \(P\to M\) with structure group
        \(G_{\mathrm{geom}}\).
  \item \(\nabla\) is a connection on \(E\) (equivalently, on the
        associated principal bundle \(P\)), giving rise to parallel
        transport, curvature, and holonomy.
  \item \(\mathcal{S}\) denotes any additional geometric structure
        that one wishes to keep track of (metric, complex structure,
        spin structure, gauge fixing, etc.), assumed to be compatible
        with \(\nabla\) in the usual sense.
\end{enumerate}
A morphism
\[
  f:\mathcal{G}\longrightarrow\mathcal{G}'
\]
between objects
\(\mathcal{G}=(M,E,\nabla,\mathcal{S})\) and
\(\mathcal{G}'=(M',E',\nabla',\mathcal{S}')\) is a pair
\[
  f=(f_M,f_E),
\]
where \(f_M:M\to M'\) is a smooth map and
\(f_E:E\to f_M^*E'\) is a vector-bundle map covering \(f_M\), such
that:
\begin{enumerate}
  \item[(i)] \(f_E\) intertwines the connections:
    \[
      f_E^*\nabla' = \nabla
      \quad\text{(connection-preserving condition)},
    \]
    in the usual sense that parallel transport and curvature pull back
    appropriately along \(f_M\).
  \item[(ii)] \(f_E\) is compatible with the additional structure
    \(\mathcal{S},\mathcal{S}'\) (e.g.\ isometry for metrics, respect
    of complex structure, etc.).
\end{enumerate}
Composition and identities are given by the obvious composition of
pairs \((f_M,f_E)\), making \(\mathsf{GeomObj}\) into a category.
\end{definition}



\subsection*{Law-space objects and their structure}


We briefly recall the relevant structure of the GaoZheng law-space
\(\GZLS\) needed here.  Details are given in the main text.

\begin{definition}[Law-space law-objects]
An object \(L\in\GZLS\) (a law-object) can be schematically written as
a tuple
\[
  L
  =
  \bigl(\mathcal{O}_L,\;\mathcal{C}_L,\;
        \GAlgebra_L,\;\mathcal{M}_L,\;\dots\bigr),
\]
where:
\begin{enumerate}
  \item \(\mathcal{O}_L\) is a space of operations / procedures
        (e.g.\ law-operations, update rules, propagators).
  \item \(\mathcal{C}_L\) is a set of constraints and consistency
        relations (e.g.\ conservation laws, Bianchi identities,
        gauge constraints).
  \item \(\GAlgebra_L\) is the associated law-level algebra encoding
        compositions and brackets between operations.
  \item \(\mathcal{M}_L\) is a collection of measures/metrics or
        energy-like functionals on \(\mathcal{O}_L\) and \(\mathcal{C}_L\)
        (e.g.\ law-curvature, defect energies).
\end{enumerate}
A morphism \(\varphi:L\to L'\) in \(\GZLS\) is a structure-preserving
map between such tuples, respecting the PL--PI structure (preserving
compositions, constraints, and law-level algebraic structure).
\end{definition}

The precise list of components of a law-object is longer in the full
GZ-framework; here we only use the parts needed to describe the
geometric embedding.



\subsection*{Construction of the functor \texorpdfstring{$\Phi_{\mathrm{geom}}$}{Phi\_geom}}


We now construct a functor
\[
  \Phi_{\mathrm{geom}}:\mathsf{GeomObj}\longrightarrow\GZLS
\]
that sends a geometric object \(\mathcal{G}\) to a corresponding
law-object \(\Phi_{\mathrm{geom}}(\mathcal{G})\) in \(\GZLS\).

\paragraph{On objects.}
Let
\[
  \mathcal{G}=(M,E\to M,\nabla,\mathcal{S})\in\mathsf{GeomObj}.
\]
We define a law-object
\[
  L_{\mathcal{G}}
  := \Phi_{\mathrm{geom}}(\mathcal{G})
\]
as follows.

\begin{enumerate}
  \item \textbf{Operation space \(\mathcal{O}_{\mathcal{G}}\).}  
    We let \(\mathcal{O}_{\mathcal{G}}\) be the space generated by:
    \begin{itemize}
      \item parallel transport operators along piecewise smooth paths
            \(\gamma:[0,1]\to M\);
      \item holonomy elements associated to closed loops;
      \item curvature operators obtained by integrating or evaluating
            the curvature \(R_{\nabla}\) on 2-surfaces in \(M\);
      \item any additional geometric operations determined by
            \(\mathcal{S}\) (e.g.\ exponential maps, Hodge operators,
            etc.).
    \end{itemize}
    Composition in \(\mathcal{O}_{\mathcal{G}}\) is given by the
    usual composition of parallel transports and operator composition.
  \item \textbf{Constraint set \(\mathcal{C}_{\mathcal{G}}\).}  
    We include in \(\mathcal{C}_{\mathcal{G}}\) the geometric identities
    satisfied by \(\nabla\) and \(\mathcal{S}\), such as:
    \begin{itemize}
      \item Bianchi identities for the curvature;
      \item gauge covariance of the connection (if \(E\) comes from a
            principal bundle);
      \item metric compatibility and torsion constraints (when
            applicable);
      \item any further structural equations implied by \(\mathcal{S}\).
    \end{itemize}
  \item \textbf{Law-level algebra \(\GAlgebra_{\mathcal{G}}\).}  
    We define \(\GAlgebra_{\mathcal{G}}\) to be the algebra generated
    by the operations in \(\mathcal{O}_{\mathcal{G}}\), modulo the
    relations in \(\mathcal{C}_{\mathcal{G}}\).  Concretely, this may
    be viewed as an algebra of formal linear combinations of words in
    parallel transport, curvature insertions, and operations encoded
    by \(\mathcal{S}\), with multiplication given by composition of
    operations and addition given by superposition.
  \item \textbf{Metrics and energies \(\mathcal{M}_{\mathcal{G}}\).}  
    Whenever a metric or inner product is present on \(E\) (or on
    associated bundles), we transport it to \(\mathcal{O}_{\mathcal{G}}\)
    by assigning norms to curvature and parallel-transport operators,
    e.g.\ via operator norms on fibres, Sobolev norms on sections,
    etc.  This yields a family \(\mathcal{M}_{\mathcal{G}}\) of
    law-level metrics and energy functionals.
\end{enumerate}

We then set
\[
  \Phi_{\mathrm{geom}}(\mathcal{G})
  :=
  L_{\mathcal{G}}
  :=
  \bigl(
    \mathcal{O}_{\mathcal{G}},\;
    \mathcal{C}_{\mathcal{G}},\;
    \GAlgebra_{\mathcal{G}},\;
    \mathcal{M}_{\mathcal{G}},\;\dots
  \bigr)
  \in \GZLS.
\]



\paragraph{On morphisms.}
Let \(f:\mathcal{G}\to\mathcal{G}'\) be a morphism in
\(\mathsf{GeomObj}\), with
\(\mathcal{G}=(M,E,\nabla,\mathcal{S})\) and
\(\mathcal{G}'=(M',E',\nabla',\mathcal{S}')\).  We define
\[
  \Phi_{\mathrm{geom}}(f)
  :
  L_{\mathcal{G}}\longrightarrow L_{\mathcal{G}'}
\]
by transporting geometric operations and constraints along the
bundle map \(f_E\) and base map \(f_M\):

\begin{itemize}
  \item An operation in \(\mathcal{O}_{\mathcal{G}}\) given by parallel
        transport along a path \(\gamma\subset M\) is mapped to the
        corresponding parallel transport along the path
        \(f_M\circ\gamma\subset M'\), using the intertwining property
        \(f_E^*\nabla'=\nabla\).
  \item Curvature operators for \(\nabla\) are mapped to curvature
        operators for \(\nabla'\) pulled back along \(f_M\).
  \item Constraints in \(\mathcal{C}_{\mathcal{G}}\) are mapped to the
        corresponding identities in \(\mathcal{C}_{\mathcal{G}'}\),
        since \(f_E\) preserves the geometric structure.
  \item The action on \(\GAlgebra_{\mathcal{G}}\) and
        \(\mathcal{M}_{\mathcal{G}}\) is induced by these maps on
        generators, extended linearly and compatibly with the law-level
        algebraic structure.
\end{itemize}

This yields a structure-preserving map
\[
  \Phi_{\mathrm{geom}}(f):L_{\mathcal{G}}\to L_{\mathcal{G}'},
\]
which respects compositions and identities, making
\(\Phi_{\mathrm{geom}}\) a functor.



\subsection*{Statement of Proposition P1}


\begin{proposition}[From object-level geometry to law-objects in \GZLS]
\label{prop:P1-geom-to-law}
The construction above defines a functor
\[
  \Phi_{\mathrm{geom}}:\mathsf{GeomObj}\longrightarrow\GZLS
\]
with the following properties.

\begin{enumerate}
  \item[\textnormal{(1)}] \textbf{Functoriality and gauge invariance.}
    \(\Phi_{\mathrm{geom}}\) is a well-defined functor: it preserves
    composition and identities.  Moreover, geometric objects related
    by a bundle isomorphism that preserves \(\nabla\) and
    \(\mathcal{S}\) (a ``gauge transformation'') are mapped to
    isomorphic law-objects in \(\GZLS\).
  \item[\textnormal{(2)}] \textbf{Local completeness of the law-image.}
    Under mild regularity assumptions (e.g.\ completeness of the
    operation space and closure of the constraint set), the law-object
    \(L_{\mathcal{G}}=\Phi_{\mathrm{geom}}(\mathcal{G})\) contains
    enough information, up to isomorphism in \(\GZLS\), to reconstruct
    the geometric object \(\mathcal{G}\) up to isomorphism in
    \(\mathsf{GeomObj}\).
  \item[\textnormal{(3)}] \textbf{Compatibility with curvature and
        Bianchi identities.}
    The curvature of \(\nabla\) and its Bianchi identities are mapped
    by \(\Phi_{\mathrm{geom}}\) to the law-curvature and corresponding
    law-level Bianchi identities in \(\GZLS\).  In particular, any
    geometric Bianchi identity appears as an element of the constraint
    set \(\mathcal{C}_{\mathcal{G}}\) in \(L_{\mathcal{G}}\), and the
    law-level curvature functoriality matches the geometric one.
  \item[\textnormal{(4)}] \textbf{Embedding of classical geometry into
        law-space.}
    The essential image of \(\Phi_{\mathrm{geom}}\) defines a full
    subcategory of \(\GZLS\) that can be regarded as an embedding of
    classical geometric objects (with connections and curvature) into
    law-space: within this subcategory, geometric equivalence classes
    of objects correspond bijectively to law-space isomorphism classes.
\end{enumerate}
\end{proposition}



\subsection*{Proof of Proposition~\ref{prop:P1-geom-to-law}}


\begin{proof}
\textbf{Proof of (1): functoriality and gauge invariance.}
Let \(f:\mathcal{G}\to\mathcal{G}'\) and
\(g:\mathcal{G}'\to\mathcal{G}''\) be morphisms in
\(\mathsf{GeomObj}\).  By construction,
\(\Phi_{\mathrm{geom}}(f)\) and \(\Phi_{\mathrm{geom}}(g)\) act on
operations by pushing forward parallel transports, curvatures, and
constraints along the corresponding bundle maps and base maps.  The
composition \(g\circ f\) acts on geometric data by first applying
\(f\) then \(g\), and hence
\[
  \Phi_{\mathrm{geom}}(g\circ f)
  =
  \Phi_{\mathrm{geom}}(g)\circ\Phi_{\mathrm{geom}}(f)
\]
on all generators of \(\mathcal{O}_{\mathcal{G}}\).  This extends by
linearity and algebraic closure to the whole law-object, proving
functoriality.  Identity morphisms are mapped to identity law-maps,
since they act trivially on all geometric operators.

If \(f:\mathcal{G}\to\mathcal{G}'\) is a geometric isomorphism (a
gauge transformation) preserving \(\nabla\) and \(\mathcal{S}\),
then \(f\) induces bijections on the operation space and constraints,
and hence an isomorphism of law-objects
\(\Phi_{\mathrm{geom}}(\mathcal{G})\cong\Phi_{\mathrm{geom}}(\mathcal{G}')\).
This gives gauge invariance and proves item~\textnormal{(1)}.



\medskip\noindent
\textbf{Proof of (2): local completeness of the law-image.}
We sketch the argument.  Under the GZ-admissibility assumptions, the
operation space \(\mathcal{O}_{\mathcal{G}}\) contains all parallel
transports and curvature insertions along paths and surfaces in \(M\).
By standard results in differential geometry (e.g.\ Ambrose–Singer for
connections on principal bundles), knowing the holonomy and curvature
along all loops and surfaces is enough to reconstruct the connection
\(\nabla\) up to gauge.  In addition, the metric and other structures
in \(\mathcal{S}\) can be reconstructed from their action on
\(\mathcal{O}_{\mathcal{G}}\) and the constraints
\(\mathcal{C}_{\mathcal{G}}\) they satisfy (e.g.\ metric compatibility,
torsion-free condition, etc.).

Thus, from the law-object \(L_{\mathcal{G}}\) one can recover the
underlying bundle \(E\), the connection \(\nabla\), and the additional
structure \(\mathcal{S}\), up to isomorphism in \(\mathsf{GeomObj}\).
In other words, if
\(\Phi_{\mathrm{geom}}(\mathcal{G})\cong\Phi_{\mathrm{geom}}(\mathcal{G}')\)
in \(\GZLS\), then \(\mathcal{G}\cong\mathcal{G}'\) in
\(\mathsf{GeomObj}\).  This provides the local completeness statement
in item~\textnormal{(2)}.



\medskip\noindent
\textbf{Proof of (3): compatibility with curvature and Bianchi
identities.}
By construction, curvature operators for \(\nabla\) are included in
\(\mathcal{O}_{\mathcal{G}}\), and the Bianchi identities are included
in \(\mathcal{C}_{\mathcal{G}}\).  The law-level curvature
\({\GZLCurv}_{L_{\mathcal{G}}}\) of the law-object is defined (in the
GZ-framework) by differentiating the law-operations along law-space
connections induced by \(\nabla\); this matches the geometric
curvature \(R_{\nabla}\) when restricted to the image of
\(\Phi_{\mathrm{geom}}\).  Moreover, the law-level Bianchi identity
\(D_{\mathrm{law}}{\GZLCurv}_{L_{\mathcal{G}}}=0\) is a direct image of
the geometric Bianchi identity
\(D_{\nabla}R_{\nabla}=0\).

Therefore, geometric curvature and Bianchi identities are faithfully
transported to law-space under \(\Phi_{\mathrm{geom}}\).  This proves
item~\textnormal{(3)}.



\medskip\noindent
\textbf{Proof of (4): embedding of classical geometry into law-space.}
From item~\textnormal{(1)} we know that \(\Phi_{\mathrm{geom}}\) is a
functor, and from item~\textnormal{(2)} we know that it is faithful
on isomorphism classes: law-space isomorphism classes in the image of
\(\Phi_{\mathrm{geom}}\) correspond bijectively to geometric
isomorphism classes in \(\mathsf{GeomObj}\).  Item~\textnormal{(3)}
ensures that the key geometric structures (connections, curvature,
Bianchi identities) are fully reflected at the law-level.

Consequently, the essential image of \(\Phi_{\mathrm{geom}}\) is a
full subcategory of \(\GZLS\) that can be identified (up to the
obvious equivalences) with the original category \(\mathsf{GeomObj}\).
This subcategory realises classical geometry as a family of law-objects
sitting inside law-space, providing the desired embedding.

This completes the proof of
Proposition~\ref{prop:P1-geom-to-law}.
\end{proof}



\section{Theorem T2: Consistency of ZFC + CH}
\label{app:III-T2}

\subsection*{Statement and basic notions}


In this appendix we restate the classical relative-consistency result
for the Continuum Hypothesis (CH) and provide precise references.  No
attempt is made to reproduce the proofs, which are long and highly
technical; instead, we indicate standard sources in the set theory
literature.

By \(\mathrm{ZFC}\) we denote the usual Zermelo–Fraenkel axioms of set
theory with the Axiom of Choice.  The Continuum Hypothesis (CH) is the
statement
\[
  \mathrm{CH}:\quad
  2^{\aleph_0} = \aleph_1,
\]
i.e.\ every infinite subset of the real line \(\mathbb{R}\) is either
countable or has the same cardinality as \(\mathbb{R}\).

We interpret “\(\mathrm{ZFC}+\mathrm{CH}\) is consistent relative to
\(\mathrm{ZFC}\)” in the usual proof-theoretic sense: if
\(\mathrm{ZFC}\) has a model (equivalently, if \(\mathrm{ZFC}\) does
not prove a contradiction), then \(\mathrm{ZFC}+\mathrm{CH}\) also has
a model.



\subsection*{Theorem and references}


\begin{theorem}[Gödel 1940; relative consistency of CH]
\label{thm:III-T2-Godel}
If \(\mathrm{ZFC}\) is consistent, then so is
\(\mathrm{ZFC} + \mathrm{CH}\).  More precisely, assuming that
\(\mathrm{ZFC}\) has a model, there exists an inner model
\(L\) (the constructible universe) such that
\[
  L \models \mathrm{ZFC} + \mathrm{GCH},
\]
in particular \(L\models \mathrm{ZFC} + \mathrm{CH}\).
\end{theorem}

The standard construction of the constructible universe \(L\) and the
verification that \(L\models\mathrm{ZFC}+\mathrm{GCH}\) can be found
in many set theory texts.  The original source is Gödel’s monograph
\cite{Godel1940}.  Detailed modern treatments include, for example,
Jech’s “Set Theory” \cite{Jech2003} (Part~IV) and Kanamori’s
“The Higher Infinite” \cite{Kanamori2003}.



\subsection*{Remarks}


\begin{itemize}
  \item The theorem does \emph{not} assert that
        \(\mathrm{ZFC} + \mathrm{CH}\) is absolutely consistent; it
        shows only that if \(\mathrm{ZFC}\) is consistent, then adding
        CH does not introduce inconsistency.
  \item The passage to the inner model \(L\) is crucial: one shows
        that starting from a model of \(\mathrm{ZFC}\), one can define
        a transitive class \(L\) satisfying \(\mathrm{ZFC}+\mathrm{GCH}\),
        thereby obtaining a model of \(\mathrm{ZFC}+\mathrm{CH}\)
        inside the original universe.
  \item In the context of the present monograph, Theorem~\ref{thm:III-T2-Godel}
        provides one of the two classical “endpoints” for CH:
        \(\mathrm{ZFC}+\mathrm{CH}\) is relatively consistent, while
        Theorem~\ref{thm:III-T3-Cohen} below shows that
        \(\mathrm{ZFC}+\neg\mathrm{CH}\) is also relatively consistent.
\end{itemize}





\section[Theorem T3: Consistency of ZFC + not-CH]{Theorem T3: Consistency of ZFC + $\neg\mathrm{CH}$}
\label{app:III-T3}

\subsection*{Statement and basic notions}


We now restate the classical forcing result showing the relative
consistency of the negation of CH.  Again, we assume the standard
interpretation of “consistency” as “existence of a model”.

The negation of CH is the statement
\[
  \neg\mathrm{CH}:\quad
  2^{\aleph_0} \neq \aleph_1,
\]
equivalently, there exists a cardinal \(\kappa\) such that
\(\aleph_0 < \kappa < 2^{\aleph_0}\).  A typical example of a model of
\(\mathrm{ZFC}+\neg\mathrm{CH}\) is obtained by forcing over a model of
\(\mathrm{ZFC}+\mathrm{CH}\) to add many new subsets of \(\omega\).



\subsection*{Theorem and references}


\begin{theorem}[Cohen 1963; relative consistency of $\neg\mathrm{CH}$]
\label{thm:III-T3-Cohen}
If \(\mathrm{ZFC}\) is consistent, then so is
\(\mathrm{ZFC} + \neg\mathrm{CH}\).  More precisely, assuming that
\(\mathrm{ZFC}\) has a model \(M\) (for instance a model of
\(\mathrm{ZFC}+\mathrm{CH}\) such as the constructible universe \(L\)),
there exists a forcing notion \(\mathbb{P}\in M\) and a generic filter
\(G\subseteq\mathbb{P}\) over \(M\) such that the forcing extension
\(M[G]\) satisfies
\[
  M[G]\models \mathrm{ZFC} + \neg\mathrm{CH}.
\]
\end{theorem}

Cohen’s original papers \cite{Cohen1963} introduced the method of
forcing and applied it to prove the relative consistency of
\(\neg\mathrm{CH}\).  Modern expositions of forcing and the consistency
of \(\neg\mathrm{CH}\) can be found, for example, in Jech’s “Set
Theory” \cite{Jech2003} (Part~V), Kunen’s “Set Theory”
\cite{Kunen2011}, and Kanamori’s “The Higher Infinite”
\cite{Kanamori2003}.



\subsection*{Remarks}


\begin{itemize}
  \item Together with Theorem~\ref{thm:III-T2-Godel}, Theorem
        \ref{thm:III-T3-Cohen} implies that CH is independent of ZFC:
        assuming the consistency of ZFC, neither CH nor \(\neg\mathrm{CH}\)
        can be proved or refuted from ZFC alone.
  \item From the viewpoint of the present work, these two theorems
        provide the classical background over which the GZ-framework
        introduces a more refined “layered” analysis of CH, separating
        its role at the level of cardinal arithmetic from its impact
        on higher law-space structures.
  \item We emphasise that we do not use any details of the forcing
        construction in the technical developments of this volume; we
        only rely on the existence of models of
        \(\mathrm{ZFC}+\mathrm{CH}\) and
        \(\mathrm{ZFC}+\neg\mathrm{CH}\) as guaranteed by these
        classical results.
\end{itemize}





\chapter[Notation and Abbreviation Cheat-Sheet]{Notation and Abbreviation Cheat-Sheet (Based on GZ-Nomenclature)}


This appendix collects the main symbols and abbreviations \\ used 
throughout the monograph, following the GZ-Nomenclature conventions.
It is organised in thematic blocks: meta-mathematical provenance,
law-space and connections, PFB-GNLA structures, CH-layering and
continuum management, and general geometric notation.



\section{Meta-Mathematical Provenance and G-Framework}


\begin{itemize}
  \item $\PLPIMMTEN$:\\
        Meta-Mathematical Theory based on Pan-Logic Analysis and
        Pan-Iterative Analysis.
\item $\PLPIMMT$:
        shorthand “PL-PI MMT” for the above meta-mathematical theory.
\item $\GFramework$:
        the GaoZheng G-Framework; the ambient meta-framework in which
        law-space, triple-layer connections, CH-layering, and PFB-GNLA
        live.
\item $\GAlgebra$:
        GaoZheng G-Algebra; a principal-bundle-based generalised
        noncommutative Lie algebra (PFB-GNLA) with several structural
        incarnations (strict, homotopy, variable functional--operator,
        and meta-mathematical original).
\end{itemize}

\section{Law-Space, Law-Connections, and Triple-Layer Geometry}


\begin{itemize}
  \item $\GZLStxt$:
        “GZ-LS” as a text abbreviation for the GaoZheng Law-Space.
\item $\GZLS$:
        the GaoZheng Law-Space; an abstract space of law-objects and
        their flows (procedures, rules, constraints, and their
        transformations).
\item $\GZLTtxt$:
        “GZ-LT” as a text abbreviation for GaoZheng Law-Transform.
\item $\GZLT$:
        the law-transform component (often denoted $M_{!w}$ in the
        variable functional–operator chart) implementing changes in
        parameters or law-objects.
\item $\GZLOCtxt$:
        “GZ-LOC” as a text abbreviation for GaoZheng Law-Connection.
\item $\GZLOC$:
        a generic law-connection (often denoted $A_M$ in the
        variable functional–operator version) that controls admissible
        changes of operators along law-space paths.
\item $\GZLCurvtxt$:
        “GZ-LCurv” as a text abbreviation for the GaoZheng
        Law-Curvature family.
\item $\GZLCurv$:
        law-curvature family; measures the failure of law-connections
        to be flat and encodes structural obstructions (anomalies,
        breakdowns, risk curvature).
\item $\GZTLCtxt$:
        “GZ-TLC” as a text abbreviation for the GaoZheng Triple-Layer
        Connection.
\item $\GZTLC$:
        triple-layer total connection on $(\text{geometry},\ \text{parameters},\ \text{laws})$,
        combining geometric, parameter, and law-level connections into
        a single object.
\item $\mathrm{Hist}(\GZLS)$:
        a schematic space of admissible law-level histories
        $\Gamma:[0,T]\to\GZLS$ used in law-space path-integral
        formulations.
\item $\mathcal{Z}_{\mathrm{law}}$:
        generalized Feynman path integral on law-space (GaoZheng
        law-space path integral, GZ-LSPI); a partition function over
        law-trajectories weighted by
        $\exp(\tfrac{i}{\hbar}S_{\mathrm{law}}[\Gamma])$.
\item $\GZGMS$:
        GaoZheng Generalized Mathematical Structure; a global
        container that organises law-objects, objects, flows, layers,
        and triple-layer connections in a single structural package.
\end{itemize}

\section{Principal Bundles and PFB-GNLA Versions}


\begin{itemize}
  \item $(\Bundle,\Base,\Group)$:
        a principal $\Group$-bundle $\Bundle\to\Base$ with structure
        group $\Group$.
\item $\Conn$:
        a classical connection 1-form on the principal bundle
        (gauge potential $A$).
\item $\Curv$:
        the curvature 2-form of $\Conn$ (field strength $F$).
\item $\LawSpace$:
        a generic symbol for a law-space (often instantiated as
        $\GZLS$).
\item \textbf{Strict PFB-GNLA}:
        the version in which the Jacobi identity holds on-the-nose;
        provides a near-Lie algebra of operators on bundles and
        sections.
\item \textbf{Homotopy PFB-GNLA}:
        the version in which Jacobi failures are absorbed into a
        3-ary operation $l_3$ controlled by a closed 3-form $H$,
        yielding a 2-term $L_\infty$-algebra.
\item \textbf{Variable functional--operator PFB-GNLA}:
        the version where parameter directions are upgraded to
        operator-valued law-functionals; introduces $A_M$ and mixed
        curvatures $F^{(xM)},F^{(MM)}$.
\item \textbf{Meta-mathematical original PFB-GNLA}:
        the PL-PI provenance version, formulated at the law level
        with a “redefined connection” $\mathcal A_M$ and meta-level
        correction terms.
\end{itemize}

\section{CH-Layering, Continuum Management, and CBT}


\begin{itemize}
  \item $\mathrm{CH}$, $\neg\mathrm{CH}$:
        classical Continuum Hypothesis and its negation.
\item $l_n$, $l_2$, $l_{\ge3}$:
        CH-layers: $l_n$ denotes the $n$-th law layer in the semantic
        layering data $\mathsf{Layer}$; $l_2$ is the purely extensional,
        cardinality-based layer, while $l_{\ge3}$ denotes the aggregate
        of all higher ($n\ge3$) homotopy/semantic/generative layers.
\item $\pi_n : \mathcal{E}_n \to B$:
        an $n$-th order law-bundle space over a base region
        $B\subset\mathsf{Obj}$; a principal-bundle--type object whose
        fibres are law-space fibres, with law-connection and CH-layering
        data restricted to layers $\le n$ so that horizontal lifts in
        $\mathcal{E}_n$ respect the $l_n$-visible metrics and
        invariants.
\item $\mathcal{B}_{\mathrm{CH+op}}$:
        CH-breakdown and operator breakdown diagnostics; a combined
        indicator of where continuum approximations and operator
        semantics fail.
\item GZ-CBT:
        GaoZheng Continuum Breakdown and Rewriting Theory; a
        law-space rewriting scheme organising transitions between
        different continuum regimes (e.g.\ CH, $\neg\mathrm{CH}$,
        forcing axioms, large-cardinal frameworks).
\end{itemize}

\section{General Conventions}


\begin{itemize}
  \item $X,Y,Z$:
        generic spaces or manifolds; context specifies whether they
        are state spaces, parameter spaces, or law-spaces.
\item $M$:
        usually the base manifold of a principal bundle or the
        underlying space-time / configuration space.
\item $G$:
        a structure group or symmetry group acting on fibers or
        internal degrees of freedom.
\item $\Omega^\bullet(\cdot)$:
        graded differential forms on the indicated space.
\item $\Gamma(\cdot)$:
        sections of the indicated bundle.
\item $\mathrm{ad}(P)$:
        the adjoint bundle associated to a principal bundle $P$.
\end{itemize}



\section{Complex-Systems Modelling and Companion Work}


\begin{itemize}
  \item total-operator principal bundles for complex-systems modelling:
        a schematic name for principal-bundle--type models used for
        high-dimensional complex systems in companion work; in the
        present volume, they serve only as examples of applied
        total-operator bundles built on the abstract law-bundle
        framework.
\item hierarchical architectures and compression--expansion reasoning schemes:
        generic phrases for multi-layer architectures and
        compression–expansion style reasoning pipelines used in
        complex-systems modelling and analysis in companion work; here
        they are mentioned only as examples of law-space procedures and
        are not formal objects of the current theory.
\end{itemize}

\chapter[Supplement: Principal Bundles and Homotopy]{Supplement: Classical Principal Bundles and Homotopy Theory}


This supplement recalls classical notions about principal bundles,
connections, curvature, and homotopy which are used implicitly in the
main construction of the GaoZheng G-Algebra (\GAlgebra) and the
principal-bundle-based generalised noncommutative Lie algebras
(PFB-GNLA). The goal is not to give a comprehensive textbook treatment,
but to fix notation and emphasise the structural patterns that are
reused in the four PFB-GNLA versions and in the PL-PI meta-mathematical
framework.



\section{Principal \texorpdfstring{$G$}{G}-Bundles and Gauge Transformations}


Let $G$ be a Lie group and $\Base$ a smooth manifold. A
\emph{principal $G$-bundle} over $\Base$ consists of a smooth manifold
$\Bundle$ together with

\begin{itemize}
  \item a surjective submersion $\pi:\Bundle\to\Base$ (the bundle
        projection),
  \item a free and transitive right action
        $R_g:\Bundle\to\Bundle$, $p\mapsto p\cdot g$ of $G$ on each
        fiber,
  \item local trivializations
        $\varphi_U:\pi^{-1}(U)\to U\times G$ over open sets
        $U\subset\Base$
\end{itemize}

such that the transition functions $g_{UV}:U\cap V\to G$ defined by
$\varphi_U\circ\varphi_V^{-1}(x,h)=(x,g_{UV}(x)h)$ are smooth.
In practice, a principal $G$-bundle packages “internal degrees of
freedom” (encoded by $G$) attached to each base point $x\in\Base$.



A \emph{gauge transformation} of $(\Bundle,\Base,G)$ is an
$\Base$-equivariant diffeomorphism
$\Phi:\Bundle\to\Bundle$ commuting with the $G$-action:
\[
  \pi\circ\Phi = \pi,\qquad \Phi(p\cdot g)=\Phi(p)\cdot g.
\]
Gauge transformations form a group, often denoted
$\mathrm{Gau}(\Bundle)$, which acts on all geometric data defined on
$\Bundle$ (connections, curvature, associated fields).



\section{Connections and Curvature on Principal Bundles}


A \emph{connection} on a principal $G$-bundle
$\pi:\Bundle\to\Base$ can be described in several equivalent ways.
The most convenient here is the language of connection 1-forms.

\begin{definition}[Connection 1-form]
A connection on $\Bundle$ is a $\mathfrak g$-valued 1-form
$\Conn\in\Omega^1(\Bundle,\mathfrak g)$ such that:
\begin{enumerate}
  \item For each fundamental vector field $\xi^\#$ on $\Bundle$
        corresponding to $\xi\in\mathfrak g$,
        \[
          \Conn(\xi^\#)=\xi.
        \]
  \item $\Conn$ is $G$-equivariant:
        \[
          (R_g)^*\Conn = \mathrm{Ad}_{g^{-1}}\Conn
          \qquad\text{for all }g\in G.
        \]
\end{enumerate}
\end{definition}



The \emph{horizontal subspace} at $p\in\Bundle$ is defined by
\[
  H_p\Bundle
  :=
  \ker\bigl(\Conn_p : T_p\Bundle\to\mathfrak g\bigr),
\]
and $T_p\Bundle$ splits as $T_p\Bundle = H_p\Bundle \oplus V_p\Bundle$,
where $V_p\Bundle = \ker(\mathrm{d}\pi_p)$ is the vertical subspace.



\begin{definition}[Curvature 2-form]
The curvature of a connection $\Conn$ is the $\mathfrak g$-valued
2-form
\[
  \Curv
  :=
  \mathrm{d}\Conn
  + \tfrac12[\Conn,\Conn]
  \in\Omega^2(\Bundle,\mathfrak g),
\]
where $[\Conn,\Conn]$ is computed using the Lie bracket in
$\mathfrak g$ and the wedge product of forms.
\end{definition}



The curvature satisfies the \emph{Bianchi identity}
\[
  \mathrm{d}\Curv + [\Conn,\Curv] = 0,
\]
which can be viewed as a conservation law for the failure of
parallel transport to be path-independent.



\section{Associated Bundles and Vector Bundles}


Given a principal $G$-bundle $\Bundle\to\Base$ and a left
representation $\rho:G\to\mathrm{Aut}(V)$ on a vector space $V$, one
defines the associated vector bundle
\[
  E := \Bundle\times_\rho V
  :=
  (\Bundle\times V)/\!\sim,
  \qquad
  (p,v)\sim(p\cdot g,\rho(g^{-1})v).
\]



A connection $\Conn$ on $\Bundle$ induces a covariant derivative
$\nabla$ on sections of $E$, and the curvature $\Curv$ induces an
endomorphism-valued 2-form $R^\nabla$ on $E$. In local trivializations,
these coincide with the familiar gauge potentials $A$ and field
strengths $F$.



\section{Homotopy, Path Lifting, and Classifying Maps}


In the classification of principal bundles, homotopy plays a central
role. Two continuous maps $f_0,f_1:X\to Y$ are \emph{homotopic} if
there exists a continuous map $H:X\times[0,1]\to Y$ such that
$H(\cdot,0)=f_0$ and $H(\cdot,1)=f_1$. The homotopy class of a map
$f:X\to Y$ is denoted $[f]$.



For principal bundles, one has the following classical picture:

\begin{itemize}
  \item There exists a universal principal $G$-bundle
        $EG\to BG$, where $EG$ is contractible and $BG$ is the
        \emph{classifying space} of $G$.
  \item Isomorphism classes of principal $G$-bundles over a
        sufficiently nice base $\Base$ correspond to homotopy
        classes of maps $\Base\to BG$:
        \[
          \{\text{principal $G$-bundles over }\Base\}/\cong
          \;\cong\;
          [\Base,BG].
        \]
\end{itemize}



Path lifting is another basic ingredient: given a principal bundle
$\Bundle\to\Base$ with connection $\Conn$, each curve
$\gamma:[0,1]\to\Base$ and initial point $p_0\in\pi^{-1}(\gamma(0))$
determine a unique horizontal lift
$\tilde\gamma:[0,1]\to\Bundle$ such that
$\tilde\gamma(0)=p_0$ and $\tilde\gamma'(t)\in H_{\tilde\gamma(t)}\Bundle$.
Parallel transport along $\gamma$ is then defined by
\[
  \mathrm{PT}_\gamma:\pi^{-1}(\gamma(0))\to\pi^{-1}(\gamma(1)),
  \qquad
  p_0\mapsto \tilde\gamma(1).
\]



\section{Bridging to the GaoZheng G-Framework and PFB-GNLA}


The PFB-GNLA constructions in the main text and in the kernel-plus
documents can be read as systematic enrichments of this classical
picture:

\begin{itemize}
  \item The \emph{strict} PFB-GNLA version starts from a principal
        bundle $(\Bundle,\Base,G)$ with a family of connections and
        curvatures, encoding them as a near-Lie algebra of operators
        acting on sections and fields.
\item The \emph{homotopy} version allows controlled failures of the
        Jacobi identity, measured by a closed 3-form $H$ and a
        3-ary operation $l_3$, echoing the role of higher homotopy
        data in classical topology.
\item The \emph{variable functional--operator} version upgrades
        parameter directions to operator-valued law-functionals
        (e.g.\ $A_M$), so that “laws of evolution” become first-class
        geometric objects, parallel to the way classifying maps
        $\Base\to BG$ encode topological types.
\item The \emph{meta-mathematical original} version relocates these
        structures into PL-PI law-space (\GZLS), with law-connections
        and law-curvature generalising $\Conn$ and $\Curv$ to the
        meta-level of “rules of construction”.
\end{itemize}

Seen from this vantage point, classical principal bundle theory and
homotopy provide the “zeroth-order” background: a geometric language
for connections, curvature, and classification by homotopy classes of
maps. The GaoZheng G-Framework extends this language to law-space,
multi-layer connections, and operator homotopy, while keeping the
classical patterns visible in each of the four PFB-GNLA versions.



\clearpage
\begin{thebibliography}{99}
\phantomsection
\addcontentsline{toc}{chapter}{Bibliography}

\bibitem{Godel1940}
K.~Gödel,
\emph{The Consistency of the Continuum Hypothesis},
Annals of Mathematics Studies, no.~3,
Princeton University Press, 1940.

\bibitem{Cohen1963}
P.~J.~Cohen,
\emph{The independence of the continuum hypothesis},
Proceedings of the National Academy of Sciences of the USA
\textbf{50} (1963), 1143--1148.

\bibitem{Jech2003}
T.~Jech,
\emph{Set Theory},
Springer Monographs in Mathematics,
3rd millennium edition, Springer, 2003.

\bibitem{Kunen2011}
K.~Kunen,
\emph{Set Theory},
College Publications, 2011.

\bibitem{Kanamori2003}
A.~Kanamori,
\emph{The Higher Infinite: Large Cardinals in Set Theory from Their Beginnings},
2nd edition, Springer Monographs in Mathematics, 2003.

\bibitem{Easton1970}
W.~B.~Easton,
\emph{Powers of regular cardinals},
Annals of Mathematical Logic
\textbf{1} (1970), 139--178.

\bibitem{Weinberg1995}
S.~Weinberg,
\emph{The Quantum Theory of Fields, Volume~I: Foundations},
Cambridge University Press, 1995.

\bibitem{Haag1996}
R.~Haag,
\emph{Local Quantum Physics: Fields, Particles, Algebras},
2nd edition, Springer, 1996.

\end{thebibliography}

\bigskip
\noindent
\textit{License notice.} This arXiv version is licensed under
CC BY 4.0 (Creative Commons Attribution 4.0 International).


\end{document}
